\documentclass[11pt]{article}

% OT1 (default) encoding -- no T1 fontenc/lmodern, per working convention
\usepackage[margin=1.05in]{geometry}
\usepackage{amsmath,amssymb,amsthm,mathtools,mathrsfs}
\usepackage{enumitem}
\usepackage{booktabs,array}
\usepackage{tabularx}
\usepackage[colorlinks=true,linkcolor=blue,citecolor=blue,urlcolor=blue]{hyperref}
\usepackage[nameinlink,capitalise]{cleveref}
\setlength{\emergencystretch}{3em}
\hbadness=10000
\raggedbottom

% ---------------------------------------------------------------------
% theorem environments
% ---------------------------------------------------------------------
\newtheorem{theorem}{Theorem}[section]
\newtheorem{proposition}[theorem]{Proposition}
\newtheorem{lemma}[theorem]{Lemma}
\newtheorem{corollary}[theorem]{Corollary}
\newtheorem{openproblem}[theorem]{Open Problem}
\theoremstyle{definition}
\newtheorem{definition}[theorem]{Definition}
\newtheorem{assumption}{Assumption}[section]
\newtheorem{principle}[theorem]{Principle}
\newtheorem{construction}[theorem]{Construction}
\newtheorem{example}[theorem]{Example}
\theoremstyle{remark}
\newtheorem{remark}[theorem]{Remark}
\crefname{assumption}{Assumption}{Assumptions}
\Crefname{assumption}{Assumption}{Assumptions}
\crefname{principle}{Principle}{Principles}
\Crefname{principle}{Principle}{Principles}
\crefname{remark}{Remark}{Remarks}
\Crefname{remark}{Remark}{Remarks}
\crefname{example}{Example}{Examples}
\Crefname{example}{Example}{Examples}
\crefname{proposition}{Proposition}{Propositions}
\Crefname{proposition}{Proposition}{Propositions}
\crefname{corollary}{Corollary}{Corollaries}
\Crefname{corollary}{Corollary}{Corollaries}
\crefname{definition}{Definition}{Definitions}
\Crefname{definition}{Definition}{Definitions}
\crefname{lemma}{Lemma}{Lemmas}
\Crefname{lemma}{Lemma}{Lemmas}
\crefname{theorem}{Theorem}{Theorems}
\Crefname{theorem}{Theorem}{Theorems}

% ---------------------------------------------------------------------
% macros
% ---------------------------------------------------------------------
\newcommand{\N}{\mathbb{N}}
\newcommand{\Z}{\mathbb{Z}}
\newcommand{\R}{\mathbb{R}}
\newcommand{\C}{\mathbb{C}}
\newcommand{\F}{\mathbb{F}}
\newcommand{\Ztwo}{\mathbb{Z}/2}
\newcommand{\Zfour}{\mathbb{Z}/4}
\newcommand{\Zsix}{\mathbb{Z}/6}
\newcommand{\Arev}{\mathcal{A}_{\mathrm{rev}}}
\newcommand{\Hf}{H_{\mathrm{full}}}
\newcommand{\Acc}{\mathcal{R}}
\newcommand{\Cl}{\mathrm{Cl}}
\newcommand{\rad}{\operatorname{rad}}
\newcommand{\argmax}{\operatorname*{arg\,max}}
\newcommand{\T}{\mathbb{T}}
\newcommand{\A}{\mathcal{A}}
\newcommand{\B}{\mathcal{B}}
\newcommand{\Hh}{\mathcal{H}}
\newcommand{\Kh}{\mathcal{K}}
\newcommand{\W}{\mathcal{W}}
\newcommand{\Pp}{\mathfrak{P}}
\newcommand{\M}{\mathcal{M}}
\newcommand{\CP}{\mathrm{CP}}
\newcommand{\UCP}{\mathrm{UCP}}
\newcommand{\qalg}{q_{\mathrm{alg}}}
\newcommand{\qmet}{q_{\mathrm{met}}}
\newcommand{\rmon}{r_{\mathrm{mon}}}
\newcommand{\dCl}{d_{\mathrm{Cl}}}
\newcommand{\ds}{d_s}
\newcommand{\beff}{b_1^{\mathrm{eff}}}
\newcommand{\Ncp}{N_{\mathrm{CP}}}
\newcommand{\Ncb}{N_{\mathrm{cb}}}
\newcommand{\Go}{G_0}
\newcommand{\Gt}{\widetilde G_0}
\newcommand{\Lat}{L}
\newcommand{\qcp}{q_{\CP}}
\newcommand{\Lam}{\Lambda}
\newcommand{\eps}{\varepsilon}
\newcommand{\id}{\mathrm{id}}
\newcommand{\Tr}{\operatorname{Tr}}
\newcommand{\spec}{\operatorname{spec}}
\newcommand{\GKdim}{\operatorname{GKdim}}
\newcommand{\rank}{\operatorname{rank}}
\newcommand{\Gr}{\operatorname{Gr}}
\newcommand{\diam}{\operatorname{diam}}
\newcommand{\Res}{\operatorname{Res}}
\newcommand{\Leb}{\operatorname{Leb}}
\newcommand{\sgn}{\operatorname{sgn}}
\newcommand{\cb}{\mathrm{cb}}
\newcommand{\op}{\mathrm{op}}
\newcommand{\taug}{\boldsymbol{\tau}}
\newcommand{\Sphere}{\mathbb{S}}
\newcommand{\tr}{\operatorname{tr}}
\newcommand{\sig}{\operatorname{sig}}
\newcommand{\spn}{\operatorname{span}}
\newcommand{\cone}{\operatorname{cone}}
\newcommand{\supp}{\operatorname{supp}}
\newcommand{\sym}{\boldsymbol{\sigma}}
\newcommand{\one}{\mathbf{1}}

\newcommand{\Dd}{\mathcal{D}}
\newcommand{\Wgrp}{\mathcal{W}}
\newcommand{\Jfam}{\mathfrak{J}}
\newcommand{\diag}{\operatorname{diag}}
\newcommand{\Sch}{\mathcal{S}}
\newcommand{\Dirac}{\mathsf{D}}
\newcommand{\Ham}{\mathsf{H}}
\newcommand{\norm}[1]{\lVert#1\rVert}
\newcommand{\Sheets}{\mathsf{S}}
\newcommand{\Alt}{\operatorname{Alt}}
\newcommand{\Hom}{\operatorname{Hom}}
\newcommand{\Ann}{\operatorname{Ann}}
\newcommand{\Dt}{\widetilde{\mathcal{D}}}
\newcommand{\Arf}{\operatorname{Arf}}
\newcommand{\Kproj}{K_{\mathrm{proj}}}
\newcommand{\Grp}{\Gamma}
\newcommand{\Vsp}{V_{\mathrm{sp}}}
\newcommand{\Iint}{\mathcal{I}}
\newcommand{\Hol}{\mathrm{Hol}}
\newcommand{\Holp}{\mathrm{Hol}^{+}}
\newcommand{\Volsp}{\Omega_{\mathrm{sp}}}
\newcommand{\detl}{\operatorname{det}}

\newcommand{\Lie}{\operatorname{Lie}}
\newcommand{\Img}{\operatorname{Im}}
\newcommand{\Volr}{\Omega_{\mathrm{ren}}}
\newcommand{\Hren}{\mathcal{H}_{\mathrm{ren}}}
\newcommand{\bren}{\beta_{\mathrm{ren}}}
\newcommand{\Tt}{\mathcal{T}}




% ---------------------------------------------------------------------
% intrinsic-canonicity additions
% ---------------------------------------------------------------------
\newtheorem{axiom}{Condition}[section]
\crefname{axiom}{Condition}{Conditions}
\Crefname{axiom}{Condition}{Conditions}
\newcommand{\Hprim}{H_{\mathrm{prim}}}
\newcommand{\Ad}{\operatorname{Ad}}
\newcommand{\Asym}{\operatorname{Asym}}
\newcommand{\Sym}{\operatorname{Sym}}
\newcommand{\Homm}{\operatorname{Hom}}
\newcommand{\so}{\mathfrak{so}}
\newcommand{\gl}{\mathfrak{gl}}
\newcommand{\HS}{\mathrm{HS}}
\newcommand{\Jchi}{J_\chi}
\newcommand{\Stilde}{\widetilde{S}}
\newcommand{\Pitr}{\Pi_{\mathrm{trans}}}
\newcommand{\Pisp}{\Pi_{\mathrm{spin}}}
\newcommand{\mone}{\mathfrak{m}_1}
\newcommand{\Ocal}{\mathcal{O}}
\newcommand{\Rb}{\mathscr{R}_b}
\newcommand{\Pcal}{\mathscr{P}}
\newcommand{\Fcal}{F}
\newcommand{\Eucl}{\mathrm{E}(\dCl)}
\newcommand{\SO}{\mathrm{SO}}

\title{\textbf{Renewal Spectral Geometry and the Emergence of Lorentzian Spacetime}}
\author{A.~P\'elissier}
\date{}

\begin{document}
\maketitle




\begin{abstract}
We show that Lorentzian spacetime can emerge from predictive memory rather than being imposed as prior geometry.  Starting from a finite renewal process of completely positive channels, with no initial metric, causal order, time direction, or spacetime dimension, we reconstruct a positive spectral geometry and an intrinsic causal-order skeleton.  We prove that positivity alone cannot produce a Lorentzian signature.  The necessary Lorentzian datum is carried by a minimal signed double cover of the renewal geometry, with a minimal finite ($\mathbb Z/4$) lift when the signed holonomy is retained at amplitude level.  Reversible predictive transformations then generate the projective Clifford sector and distinguish a single temporal direction, while a Perron--Frobenius pressure condition uniquely selects the modular scale.
%
Under nondegenerate coarse-graining, the signed geometry converges to a Lorentzian Dirac continuum.  The associated discrete Hamiltonians converge globally in strong resolvent sense and in norm on fixed momentum bands.  For slowly varying renewal statistics, covariance of transported second moments yields metric-compatible frame transport, while a primitive exchange-symmetric stationary diamond law forces translational plaquette closure.  Channel composition therefore reconstructs torsion-free Levi--Civita spin transport while leaving rotational curvature unconstrained.  The primitive revision algebra forces an odd spatial rank and makes (3+1)-dimensional spacetime the minimal nondegenerate Lorentzian endpoint.
%
Thus a single operator-algebraic renewal object contains the kinematic, causal, spectral, and spinorial architecture required for Lorentzian spacetime.  The construction uses no gravitational field equations, matter content, or phenomenological assumptions.
%
A \texttt{Lean} formalization of the results is available at
\url{https://github.com/AI-MathPhys/renewal-geometry}.
\end{abstract}



\setcounter{tocdepth}{2}
\tableofcontents

 
\section{Introduction}
\label{sec:intro}

A persistent ambition of mathematical physics is to \emph{derive} spacetime
geometry rather than to postulate it.  Ideally, dimension, causal order,
Lorentzian signature, and relativistic propagation should arise as different
aspects of a common underlying structure, rather than being introduced as
independent kinematic assumptions.  This is difficult because the mathematical
frameworks that most naturally encode geometry tend to presuppose precisely
the structures one would like to explain.  A manifold already carries a
topology and a differentiable structure; a causal set begins with an order;
and a Lorentzian spectral geometry is ordinarily supplied with an indefinite
or temporal datum from the outset.

The tension is particularly sharp in noncommutative geometry.  In Connes'
spectral formulation, the metric information of a Riemannian spin space is
encoded by a spectral triple ($(\A,\Hh,D)$), consisting of an algebra of
observables, a Hilbert-space representation, and a Dirac operator
\cite{connes}.  This operator-theoretic formulation is well adapted to a
setting in which geometry and quantum observables should be described in a
common language.  Length-function spectral triples further show that
nontrivial metric and dimensional information can be carried by discrete or
totally disconnected structures
\cite{connes1989,christensenivan2006,pearson-bellissard,
kellendonksavinien2012}.  Yet the ordinary Hilbert-space framework is
positive-definite and therefore naturally Riemannian.  Physical spacetime,
by contrast, is Lorentzian.

Krein, pseudo-Riemannian, temporal, and twisted spectral triples address this
signature problem by enriching the spectral data with an indefinite inner
product, a fundamental symmetry, a temporal operator, or a modular twist
\cite{bognar,strohmaier,PaschkeSitarz2006,
VanDenDungenPaschkeRennie2013,franco,franco-eckstein,
vandendungen2016,bizibrouderbesnard2018,connes-moscovici}.
Temporal Lorentzian triples isolate a distinguished time direction through an
exact time operator and lapse, while spectral-spacetime formulations promote
time orientation to structural data
\cite{Franco2014Temporal,EcksteinFranco2015,
Besnard2017Spacetimes}.  Twisted almost-commutative constructions provide
another route from Euclidean to pseudo-Riemannian spectral geometry
\cite{DevastatoFarnsworthLizziMartinetti2018,
DevastatoFilaciMartinettiSingh2020,nieuviarts2026a,
nieuviarts2026b}.  These theories explain how Lorentzian geometry is encoded
once suitable signed or temporal data are supplied.  The question pursued
here is more primitive: can the need for such data, their minimal form, and
the associated Dirac structure be recovered from a process that is not
geometric at the outset?

We propose predictive renewal memory as such a primitive object.  Renewal
theory studies processes that repeatedly forget and restart, and its
resolvents control systems with memory throughout probability, statistical
mechanics, reliability, and information theory \cite{feller}.  Here renewal
is not placed inside a pre-existing spacetime.  Instead, a finite alphabet of
reset operations acts on an internal observable algebra through unital
completely positive maps.  A history is merely an ordered sequence of such
operations.  It is not yet a trajectory in space or time.

The basic reduction is predictive rather than geometric.  Two histories are
identified when they induce the same accumulated channel and therefore have
the same future predictive action.  This resembles causal-state reduction in
computational mechanics, where histories are grouped according to their
predictive equivalence
\cite{shalizicrutchfield,TraversCrutchfield2014}.  Continuous-time
causal-state models and their labelled transition operators provide the
closest probabilistic comparison
\cite{MarzenCrutchfield2017,Marzen2018MWC,
RiechersCrutchfield2017a,RiechersCrutchfield2017b,
LoomisCrutchfield2019}.  The distinction here is that predictive labels are
completely positive operators rather than scalar transition probabilities.

This quotient already has geometric content.  Its minimal history length
defines a positive spectral geometry under explicit regularity assumptions,
and the growth of distinguishable predictive classes becomes a spectral
dimension.  At the same time, continuation of one accumulated channel by
another defines predictive divisibility.  After quotienting mutual
reachability, this becomes a partial order; under the stated monotonicity
condition, the quotient length becomes a time function and grading.  Thus the
order is not inserted as causal-set kinematics.  It is the order induced by
the possibility of future predictive continuation.  In this sense, geometry
is what prediction remembers.

This positive construction reaches a sharp boundary.  Complete positivity
produces a stable Hilbert-positive geometry, but that same positivity cannot
supply a nontrivial Lorentzian fundamental symmetry in the natural
irreducible representation.  More generally, any compatible nonconstant
diagonal involution is equivalent to introducing additional
($\Z/2$)-valued cocycle data on the predictive graph.  The signed sector is
therefore not an optional decoration of an already Lorentzian geometry.  It
is the minimal structure required to pass beyond positivity.

The form of this enrichment is rigid even though its cohomology class is not
selected by the graph-general classification alone.  Preservation of the
distinguished renewal cone forces local unitary transports to permute its
extremal rays.  Fibrewise indefiniteness requires at least two local directions,
and at minimal rank the only possible transports are the identity and the
sheet exchange.  Consequently every objectwise-minimal renewal-positive
indefinite enrichment is a principal double cover, classified by a class in
$H^1(G,\Z/2)$.  In record-regular operational realizations, this abstract class
can be selected by the determinant parity of stable future-distinct record
transport; the present paper classifies the geometry resulting from the
selected class.  The corresponding sheet-sign operator supplies the Krein
fundamental symmetry.  If the same signed holonomy is retained at the level
of unitary amplitudes, a nontrivial sign must admit a square root; the minimal
finite amplitude completion is then a $\Z/4$ lift.  The double cover carries
the real/Krein signature data, while the lift retains the additional
fourth-root phase information.

A second structural ingredient comes from reversibility.  Reversible
predictive transformations are not assumed to be unitary at the microscopic
level.  Rather, a unital completely positive map with a unital completely
positive inverse is automatically a (*)-automorphism.  On the primitive
finite factor these automorphisms are inner, and their unitary implementers
therefore compose projectively.  The resulting commutator pairing supplies
the symplectic and Clifford structure needed for a Dirac geometry.  Thus the
Clifford sector is not attached independently to the signed cover: it arises
from the reversible algebra of predictions.

The renewal grading also provides a natural modular clock.  When elementary
resets transform homogeneously under a faithful modular flow, the modular
Hamiltonian is forced to be affine in the renewal length.  On a recurrent
weighted component, the remaining exponent is selected by a Perron--Frobenius
pressure equation: under supercriticality, the spectral radius of the
length-weighted transfer matrix decreases through one at a unique positive
value $\beta$.  Together with deck oddness, this determines the unbounded
signed Dirac magnitude up to its overall physical-unit conversion and a
bounded Krein-self-adjoint perturbation.  This gives a renewal realization of
the modular-time perspective of Connes and Rovelli
\cite{connesrovelli1994}, within the broader theory of twisted and modular
spectral triples
\cite{connes-moscovici,careyphillipsrennie2010,
careyneshveyevnestrennie2011}.

The continuum problem is then not to manufacture an indefinite metric from
positive averages, which is impossible, but to understand the scaling limit
of the signed Clifford geometry.  A finite deterministic reset frame
generally produces a polyhedral order cone.  At the level of the signed
principal symbol, however, a nondegenerate spatial second moment yields a
Lorentzian quadratic cone with a single possible timelike direction.  No
microscopic rotational invariance is required: a positive-definite limiting
second moment defines a spatial vielbein, and the corresponding flat
Hamiltonian is unitarily equivalent to the standard massless Dirac
Hamiltonian.  Stationary ergodic reset statistics therefore produce the flat
Minkowski Dirac isometry class almost surely whenever the spatial reset
directions span.

The curved regime reveals why channel composition is geometrically
significant.  For slowly varying renewal statistics, the defect between the
two orders of traversing an elementary plaquette separates into a
translational and a rotational component.  The former is discrete torsion;
the latter is curvature.  Metric compatibility need not be imposed
independently: covariance of neighbouring spatial second moments makes the
normalised comparison transport orthogonal, and preservation of the selected
orientation places it in $SO(\dCl)$.  Torsion freedom is likewise derived
within a distinguished phase.  For a primitive exchange-covariant stationary
microscopic diamond process, the unique stationary law is exchange invariant,
so every exchange-odd translational plaquette defect has vanishing expectation,
while the rotational record remains unconstrained.  The resulting
translational closure, together with metric compatibility, allows the discrete
fundamental theorem to select the Levi--Civita spin connection without forcing
rotational holonomy to vanish.  The corresponding covariant Hamiltonians
converge to the variable-coefficient Lorentzian Dirac Hamiltonian through
subprincipal order.  Curvature survives the construction, although no equation
governing that curvature is imposed.

Dimension enters through the primitive revision algebra rather than through
a chosen number of coordinates.  Its reduced commutator pairing has even
rank, and one of the corresponding Clifford directions is distinguished as
temporal.  The spatial rank is therefore odd.  The spatially
one-dimensional branch is degenerate for the nontrivial continuum mechanism
considered here, so three spatial directions give the first nondegenerate
Lorentzian Dirac endpoint.  The conclusion is an unconditional
\emph{minimality} result: (3+1) is the lowest nondegenerate dimension in the
primitive modular class.  Parity alone does not exclude higher odd spatial
ranks, and their elimination requires the additional closure condition
stated explicitly later.  Related parity and single-time constraints for
first-order relativistic equations have classical precedents
\cite{mankocborstniknielsen2000,vandamng2001,nielsenrugh1994};
the role of renewal memory is to generate the relevant projective Clifford
factor rather than to postulate it.

The construction is related to, but distinct from, other programmes of
emergent spacetime.  Thermodynamic approaches seek gravitational dynamics
from entropy and horizon laws \cite{jacobson1995,verlinde2011}, while
holographic approaches relate spacetime connectivity to entanglement
\cite{vanraamsdonk2010}.  Causal-set theory begins with a locally finite order
\cite{bombelli-sorkin,surya}, and graph, simplicial, and group-field models
begin with combinatorial structures whose collective phases approximate
geometry
\cite{konopkamarkopouloseverini2008,
ambjornjurkiewiczloll2009,oriti2009}.  Here the primitive datum is neither a
metric nor an order, but predictive memory.  Order arises from divisibility
of accumulated channels, the necessity of a signed extension is exposed by
a positivity obstruction, and the continuum endpoint is a controlled
Krein--Clifford Dirac evolution.

The result is kinematic and operator-theoretic.  We construct a positive
predictive spectral geometry, an intrinsic causal-order skeleton, a minimal
signed/Krein enrichment, a projective Clifford revision sector, and
controlled flat and adiabatic Lorentzian Dirac limits.  We do not derive
Einstein's equations, a gravitational variational principle, matter
multiplets, gauge interactions, conformal completion, or phenomenological
predictions.  The claim is instead that predictive renewal memory already
contains enough structure to support the causal, spectral, temporal, and
spinorial architecture of Lorentzian spacetime.

The hypotheses used below do not all have the same logical status.  The finite UCP
renewal memory is the mathematical starting object, while finite fibres, bounded
increments, and continuum regularity are analytic domain conditions for particular
spectral or convergence results.  Nontrivial signed holonomy, a realised reversible
projective sector, pressure supercriticality, elliptic spatial scaling, and
exchange-symmetric stationary transport are phase conditions selecting the
Lorentzian regime.  Regular variation, exact isotropy, metric coding, and the
higher-rank uniqueness conditions are optional strengthenings rather than premises
of the core emergence theorem.  A detailed condition audit and the canonical
reductions of global connectivity, faithfulness, irreducibility, and projective
degeneracy are recorded in Appendix \Cref{app:hypothesis-status}.

The paper is organized as follows.  \Cref{sec:object} constructs the positive
predictive spectral geometry and order and establishes the positivity
boundary.  \Cref{sec:signed} classifies the minimal signed enrichment,
constructs the associated Krein and modular Dirac structures, and proves the
signed reconstruction theorem.  \Cref{sec:primitive-dimension-selection}
derives the projective revision algebra from reversible predictive dynamics,
performs its canonical symplectic reduction to the primitive Clifford factor,
and proves the (3+1)-dimensional minimality theorem.  With this Clifford
structure in place, \Cref{sec:lorentzian} establishes the flat and slowly
varying Lorentzian Dirac limits, including operator convergence,
self-averaging, and the reconstruction of Levi--Civita spin transport from
channel plaquettes.  \Cref{sec:conclusion} discusses the
interpretation, precise scope, and limitations of the construction.  The
appendices collect regularity criteria, auxiliary dimension results,
renewal-tail singularities, deterministic product-order limits, conditional
dimension-selection mechanisms, propagation and global-topology results, and
worked examples.




 



\section{Renewal memories and the positive predictive geometry}
\label{sec:object}


We recall briefly the part of the spectral-triple formalism that will be used in
this section.  The reader familiar with noncommutative geometry may pass directly
to \Cref{subsec:renewal-collapse}; the purpose of the present subsection is only
to fix conventions and to explain why the quotient constructed below should be
viewed as a positive noncommutative geometry rather than merely as a labelled
counting problem.

\subsection{Spectral triples and predictive noncommutative spaces}
\label{subsec:spectral-background}

In noncommutative geometry a compact Riemannian spin geometry is encoded by a
spectral triple.  In the form used here this consists of a unital involutive
algebra $\A$ represented by bounded operators on a separable Hilbert space
$\Hh$, together with a self-adjoint, usually unbounded, operator $D$ such that
$(1+D^2)^{-1/2}$ is compact and $[D,a]$ is bounded for all $a$ in a dense
$*$-subalgebra of $\A$ \cite{connes}.  For a compact spin manifold $X$ the model
example is
\[
   \bigl(C^\infty(X),L^2(X,S),D_X\bigr),
\]
where $S$ is the spinor bundle and $D_X$ is the classical Dirac operator.  In
this commutative case the metric can be recovered from the spectral data; in the
noncommutative case the same conditions are taken as the definition of a
``space'' whose metric information is carried by the operator $D$.

Two analytic features of this definition are important below.  First, compact
resolvent is the replacement for compactness of the underlying space: it means
that the spectrum of $D$ escapes to infinity with finite multiplicities.  Second,
the boundedness of commutators is the Lipschitz condition: the algebra elements
are precisely the observables whose variation is controlled by the metric
measured by $D$.  Thus, once an algebra of creation operators and a length
operator are specified, the question of whether they form a spectral triple is
not formal.  It is an operator-theoretic regularity condition, and it is exactly
this condition that will force finite fibres and bounded length increments in
\Cref{thm:fibre-dichotomy}.

The dimension attached to a spectral triple is also spectral.  One considers
zeta functions of the form
\[
   \zeta_D(s)=\Tr(|D|^{-s})
\]
(or, more generally, localized zeta functions $\Tr(a|D|^{-s})$), and the
abscissa of convergence gives the metric growth exponent of the geometry.  In
regular cases the poles form the dimension spectrum and enter the spectral
action expansion \cite{chamseddineconnes1997}.  We shall not use the spectral
action in this paper.  What we use is the same principle at the kinematic level:
a length Dirac on the predictive quotient turns the growth of distinguishable
predictive states into a spectral dimension.  Length-function spectral triples of
this discrete kind go back to Connes \cite{connes1989} and were developed for AF
algebras and Cantor sets by Christensen--Ivan \cite{christensenivan2006} and, for
ultrametric and aperiodic media, by Pearson--Bellissard \cite{pearson-bellissard}
and Kellendonk--Savinien \cite{kellendonksavinien2012}; the predictive-length
triple constructed below is the renewal instance of this lineage.

The adjective \emph{positive} in this section is meant literally.  The Hilbert
space is an ordinary Hilbert space, the length Dirac is self-adjoint for a
positive definite scalar product, and the channels used to define the quotient
are completely positive maps.  Consequently the geometry constructed in this
section is Riemannian in the sense relevant to spectral triples.  It will supply
length, dimension, volume residues, and an order skeleton.  It will not supply a
Lorentzian fundamental symmetry.  The latter appears only in the signed sector,
where the Hilbert-positive construction is replaced by a Krein datum.

\subsection{Renewal memories and predictive quotient}
\label{subsec:renewal-collapse}

We now introduce the primitive object.  It is deliberately smaller than a
manifold, a metric graph, or a causal set: it is only a finite family of
completely positive reset operations and the histories obtained by composing
them.  The predictive quotient identifies histories that have the same accumulated
channel.  It is the space on which the positive spectral triple will be built.

\begin{definition}[Renewal memory]
\label{def:renewal-memory}
A finite completely positive renewal memory consists of a finite alphabet $E$, a finite-dimensional unital $C^*$-algebra $\A_{\rm int}$, and a family of unital completely positive maps
\[
 \Phi_\sigma:\A_{\rm int}\longrightarrow \A_{\rm int},\qquad \sigma\in E.
\]
A history is a word $w=\sigma_1\cdots\sigma_k\in E^*$.  Its accumulated channel is
\[
 \Phi_w=\Phi_{\sigma_k}\circ\cdots\circ\Phi_{\sigma_1}.
\]
\end{definition}

\begin{definition}[Predictive quotient]
\label{def:predictive-quotient}
Two histories are predictively equivalent, written $w\sim_{\CP}w'$, if $\Phi_w=\Phi_{w'}$.  The predictive quotient is
\[
 \W_{\CP}=E^*/{\sim_{\CP}}.
\]
The quotient length is
\[
 \Lam_{\min}([w])=1+\min\{|v|:\Phi_v=\Phi_w\}.
\]
\end{definition}



\begin{definition}[Predictive channel monoid and its units]
\label{def:predictive-channel-monoid}
The accumulated channels form the composition monoid
\[
  \M_{\rm pred}:=\{\Phi_w:w\in E^*\}.
\]
Its unit group is
\[
  \mathsf U_{\rm pred}
  :=\left\{\Psi\in\M_{\rm pred}:
  \exists\Theta\in\M_{\rm pred},\quad
  \Theta\Psi=\Psi\Theta=\id\right\}.
\]
This definition depends only on the predictive channel quotient and not on a
choice of Kraus representation.
\end{definition}

\begin{theorem}[Predictive-unit theorem]
\label{thm:predictive-unit}
Let \(\Psi\) be a unital completely positive map on a finite-dimensional
\(C^*\)-algebra.  If \(\Psi\) has a unital completely positive inverse, then
\(\Psi\) is a \(*\)-automorphism.  Consequently
\[
  \boxed{\ \mathsf U_{\rm pred}
  \subseteq \operatorname{Aut}(\A_{\rm int}).\ }
\]
\end{theorem}

\begin{proof}
Let \(\Theta\) be the UCP inverse of \(\Psi\).  Kadison--Schwarz gives
\[
  \Psi(x)^*\Psi(x)\leq\Psi(x^*x).
\]
Applying \(\Theta\) yields
\[
  \Theta\!\left(\Psi(x)^*\Psi(x)\right)\leq x^*x.
\]
Kadison--Schwarz for \(\Theta\), applied to \(\Psi(x)\), gives the reverse
inequality,
\[
  x^*x
  =\Theta(\Psi(x))^*\Theta(\Psi(x))
  \leq\Theta\!\left(\Psi(x)^*\Psi(x)\right).
\]
Thus equality holds.  Applying the same argument to \(xx^*\) shows that every
\(x\) belongs to the multiplicative domain of \(\Psi\).  Hence \(\Psi\) is a
unital \(*\)-homomorphism, and bijectivity makes it a \(*\)-automorphism.
\end{proof}

\begin{remark}[Sharp scope]
\label{rem:predictive-unit-scope}
A generic UCP renewal tuple need not possess nontrivial units.  For example, a
completely depolarising channel annihilates every traceless observable and is
not invertible.  Reversible revisions are therefore a property of the signed
recurrent phase, not a consequence of bare complete positivity.
\end{remark}

\begin{remark}[The channel family controls the predictive phase]
\label{rem:channel-family-controls-phase}
The maps \(\{\Phi_\sigma\}_{\sigma\in E}\) are not decorative input data: they
determine the phase of the predictive quotient.  If the channels erase all
distinctions, for instance if every non-empty history induces the same
accumulated channel, then \(N_{\CP}(R)=O(1)\) and the algebraic predictive
dimension is \(q_{\mathrm{alg}}=0\).  If, at the opposite extreme, the channels
preserve essentially all symbolic distinctions, then \(N_{\CP}(R)\) grows
exponentially and the algebraic predictive dimension is infinite.  The
finite-dimensional geometric regime lies between these extremes: enough memory
must survive to produce growth, but enough predictive identification must occur
to reduce the free symbolic explosion to polynomial order.
\end{remark}




\begin{lemma}[Shift congruence]
\label{lem:shift-congruence}
The equivalence relation $\sim_{\CP}$ is a left shift congruence: if $w\sim_{\CP}w'$, then $\sigma w\sim_{\CP}\sigma w'$ for all $\sigma\in E$.  Consequently the creation operators descend to
\[
 \widehat S_\sigma e_{[w]}=e_{[\sigma w]}
\]
on $\ell^2(\W_{\CP})$.
\end{lemma}

\begin{proof}
With the convention of Definition~\ref{def:renewal-memory}, $\Phi_{\sigma w}=\Phi_w\circ\Phi_\sigma$.  Therefore $\Phi_w=\Phi_{w'}$ implies $\Phi_{\sigma w}=\Phi_{\sigma w'}$.
\end{proof}

\begin{assumption}[Finite fibres and bounded increments]
\label{ass:fibres}
The descended shifts have uniformly finite fibres.  Moreover, for every $\sigma\in E$ there exists $C_\sigma<\infty$ such that
\[
 |\Lam_{\min}([\sigma w])-\Lam_{\min}([w])|\le C_\sigma
\]
for all $[w]\in\W_{\CP}$.
\end{assumption}

Let $\B_{\CP,0}$ denote the $*$-algebra generated on $\ell^2(\W_{\CP})$ by the descended shifts, their adjoints, and finite-support diagonal multipliers.  Let $\Hh_{\rm int}$ be a faithful Hilbert-space representation of $\A_{\rm int}$.



\subsection{The predictive-length regularity criterion}
\label{subsec:predictive-length-criterion}

For the quotient length to define a spectral geometry, the descended renewal
shifts must be bounded and their commutators with the length operator must be
bounded.  In the present setting these conditions are equivalent to uniformly
finite shift fibres and uniformly bounded length increments.  Several useful
finite-combinatorial criteria verifying these hypotheses are collected in
\Cref{app:regularity-positive}; the main text records the sharp operator-theoretic
dichotomy.

\begin{theorem}[Finite-fibre dichotomy for the predictive triple]
\label{thm:fibre-dichotomy}
Let $\W_{\CP}$ be the predictive quotient of a finite CP renewal memory, with
descended shifts $\widehat S_\sigma$ \textup{(}\Cref{lem:shift-congruence}\textup{)}
and canonical length Dirac $D_{\CP}e_{[w]}=\Lam_{\min}([w])\,e_{[w]}$.  Exactly
one of the following holds.
\begin{enumerate}[label=\textup{(\arabic*)},leftmargin=2.4em]
\item \textup{(regular)} every $\widehat S_\sigma$ has uniformly finite fibres and
      every $\Lam_{\min}$-increment is bounded \textup{(}\Cref{ass:fibres}\textup{)}.
      Then each $\widehat S_\sigma$ is bounded with
      $\norm{\widehat S_\sigma}=(\sup_{[u]}\#\text{fibre}_\sigma[u])^{1/2}$, each
      $[D_{\CP},\widehat S_\sigma]$ is bounded, and
      $(\B_{\CP,0},\ell^2(\W_{\CP})\otimes\Hh_{\rm int},D_{\CP})$ is a spectral
      triple \textup{(}\Cref{thm:triple}\textup{)}; it is $\theta$-summable, and
      finitely summable if and only if $N_{\CP}(R)$ grows polynomially.
\item \textup{(singular)} some $\widehat S_\sigma$ has an infinite fibre, or some
      $\Lam_{\min}$-increment is unbounded.  Then the canonical data fail to be a
      spectral triple, and the failure is not reparable within the
      predictive-length geometry:
      \begin{enumerate}[label=\textup{(2\alph*)},leftmargin=2.4em]
      \item an infinite fibre makes $\widehat S_\sigma$ unbounded
            ($\norm{\widehat S_\sigma}=\infty$), so no choice of diagonal Dirac
            can make $\widehat S_\sigma$ a bounded algebra element;
      \item an unbounded increment makes $[D_{\CP},\widehat S_\sigma]$ unbounded,
            and more generally $[f(D_{\CP}),\widehat S_\sigma]$ is unbounded for
            every $f$ that is bi-Lipschitz on the length spectrum
            \textup{(}every length-comparable reparametrisation\textup{)}.
      \end{enumerate}
\end{enumerate}
\end{theorem}

\begin{proof}
The two alternatives are the negation of one another, so exactly one holds.

\emph{(1).}  On basis vectors $\widehat S_\sigma^{*}\widehat S_\sigma e_{[u]}=
(\#\text{fibre}_\sigma[u])\,e_{[u]}$, so $\norm{\widehat S_\sigma}^2=
\sup_{[u]}\#\text{fibre}_\sigma[u]$, finite by hypothesis.  By
\Cref{thm:triple}, $[D_{\CP},\widehat S_\sigma]e_{[w]}=(\Lam_{\min}([\sigma w])-
\Lam_{\min}([w]))e_{[\sigma w]}$ has norm $\le\sup$-increment $<\infty$, and
$D_{\CP}$ has compact resolvent because each length shell is finite (finite
alphabet, finite fibres, bounded increments) and $\Lam_{\min}\to\infty$.  Thus
the triple exists; $\Tr e^{-tD_{\CP}^2}=\dim(\Hh_{\rm int})\sum_R h(R)e^{-tR^2}<
\infty$ gives $\theta$-summability, and finite summability
$\sum_R h(R)R^{-s}<\infty$ for some $s$ holds iff the abscissa
$\qalg=\limsup\log N_{\CP}(R)/\log R$ is finite, i.e.\ iff $N_{\CP}$ is
polynomial.

\emph{(2a).}  If $\#\text{fibre}_\sigma[u]=\infty$ for some $[u]$, then
$\widehat S_\sigma^{*}\widehat S_\sigma$ has the eigenvalue $+\infty$ on
$e_{[u]}$, so $\widehat S_\sigma$ is unbounded.  Boundedness of an algebra
element is a property of $\widehat S_\sigma$ alone and does not depend on the
choice of $D$; hence no diagonal Dirac repairs it.

\emph{(2b).}  If $\Lam_{\min}([\sigma w_n])-\Lam_{\min}([w_n])\to\infty$ along a
sequence, then for $f$ bi-Lipschitz on the length spectrum with lower slope
$c>0$, $|f(\Lam_{\min}([\sigma w_n]))-f(\Lam_{\min}([w_n]))|\ge c\,
|\Lam_{\min}([\sigma w_n])-\Lam_{\min}([w_n])|\to\infty$, so
$[f(D_{\CP}),\widehat S_\sigma]$ is unbounded.
\end{proof}

\begin{remark}[Why the dichotomy is sharp, and its one genuine limitation]
\label{rem:dichotomy-sharp}
The two cases are exhaustive and mutually exclusive at the level of hypotheses;
the content is that each forces a definite operator-theoretic verdict for the
\emph{predictive-length} geometry --- the object whose growth exponent is the
predictive dimension $\qalg$.  The infinite-fibre obstruction (2a) is absolute:
no diagonal Dirac can rescue an unbounded shift.  The unbounded-increment
obstruction (2b) is sharp \emph{within length-comparable Diracs}, and this is the precise limitation of the statement: it does \emph{not} assert that no diagonal Dirac
whatsoever can be summable in case (2b).  The graph-distance function
$\lambda_{[w]}=d_G(x_0,[w])$ has unit increments and tends to infinity whenever
balls are finite, so $D'=\mathrm{diag}(\lambda)$ gives bounded commutators even
when $\Lam_{\min}$-increments are unbounded.  But $D'$ is no longer the quotient
length: it measures the graph metric, a different dimension, not the predictive
growth exponent.  The dichotomy is therefore exactly a statement about the
predictive-length Dirac; swapping in an unrelated metric changes the geometry
being measured.
\end{remark}

\begin{corollary}[Sharp existence criterion for the predictive-length triple]
\label{cor:sharp-existence}
For the canonical length Dirac $D_{\CP}e_{[w]}=\Lam_{\min}([w])e_{[w]}$,
\[
  \boxed{\
  \begin{aligned}
  &\bigl(\B_{\CP,0},\ell^2(\W_{\CP})\otimes\Hh_{\rm int},D_{\CP}\bigr)
   \ \text{is a spectral triple}\\
  &\qquad\iff\ \text{bounded shift fibres}\ \textbf{and}\
   \text{bounded }\Lam_{\min}\text{-increments.}
  \end{aligned}\ }
\]
When the right-hand side holds the triple is $\theta$-summable, and finitely
summable if and only if $N_{\CP}(R)$ grows polynomially.  When it fails,
\[
  \boxed{\
  \text{either some }\widehat S_\sigma\text{ is unbounded, or some }
  [D_{\CP},\widehat S_\sigma]\text{ is unbounded.}
  \ }
\]
Both equivalences are statements about the canonical predictive-length geometry;
the qualification of \Cref{rem:dichotomy-sharp} applies to the second box (an
unrelated diagonal metric may still be summable, but it no longer measures the
predictive growth dimension $\qalg$).
\end{corollary}

\begin{proof}
The forward implication of the first box and the failure dichotomy in the second
box are, respectively, alternatives (1) and (2) of \Cref{thm:fibre-dichotomy};
the converse of the first box is immediate, since a spectral triple requires the
$\widehat S_\sigma$ to be bounded operators (whence finite fibres,
$\norm{\widehat S_\sigma}^2=\sup_{[u]}\#\text{fibre}_\sigma[u]<\infty$) and the
commutators $[D_{\CP},\widehat S_\sigma]$ to be bounded (whence bounded
increments).  The summability statements are \Cref{thm:fibre-dichotomy}\,(1).
\end{proof}






% =====================================================================
\subsection{The positive predictive spectral triple}
\label{sec:triple}

We can now pass from the regularity criterion to the actual spectral triple.
The algebra is generated by the predictive shifts and finite-rank diagonal
observables; the Hilbert space is the square-summable space of predictive
classes, with the internal representation carried along as a finite fibre; and
the Dirac operator is the quotient length.  The theorem identifies the spectral
dimension with the polynomial growth exponent of the predictive quotient.

\begin{theorem}[Predictive spectral triple]
\label{thm:triple}
Under \Cref{ass:fibres}, define
\[
 D_{\CP}e_{[w]}\otimes\xi=\Lam_{\min}([w])e_{[w]}\otimes\xi.
\]
Then
\[
 (\B_{\CP,0},\ell^2(\W_{\CP})\otimes\Hh_{\rm int},D_{\CP})
\]
is a spectral triple.  Its zeta function
\[
 \zeta_{\CP}(s)=\dim(\Hh_{\rm int})\sum_{[w]\in\W_{\CP}}\Lam_{\min}([w])^{-s}
\]
has abscissa of convergence
\[
 \qalg=\limsup_{R\to\infty}\frac{\log N_{\CP}(R)}{\log R},\qquad
 N_{\CP}(R)=\#\{[w]:\Lam_{\min}([w])\le R\}.
\]
\end{theorem}

\begin{proof}
Uniform finite fibres give boundedness of the descended shifts.  Bounded increments give
\[
 [D_{\CP},\widehat S_\sigma]e_{[w]}
 =\bigl(\Lam_{\min}([\sigma w])-\Lam_{\min}([w])\bigr)e_{[\sigma w]},
\]
so $[D_{\CP},\widehat S_\sigma]$ is bounded.  The adjoints are handled similarly, and finite-support diagonal multipliers commute with $D_{\CP}$ up to finite rank operators.  Compactness of the resolvent follows from finiteness of rank balls.  The statement about the zeta abscissa is the standard counting-function criterion for Dirichlet series with non-negative coefficients.
\end{proof}

\begin{remark}[Compact resolvent uses finite rank balls]
\label{rem:compact-resolvent-finite-balls}
The compact-resolvent assertion in \Cref{thm:triple} uses the quotient-length
definition in an essential way: for each $R$ there are at most finitely many
classes with $\Lam_{\min}\le R$, because each such class has a representative
word of length at most $R-1$ over the finite alphabet $E$.  Thus the rank balls
are finite even before imposing polynomial growth.  Polynomial growth is needed
for finite summability, not for compactness of the resolvent.
\end{remark}

\begin{definition}[Algebraic predictive dimension]
\label{def:qalg}
The number $\qalg$ of \Cref{thm:triple} is the algebraic predictive dimension.  It measures the polynomial growth rate of distinguishable accumulated channels as a function of quotient length.
\end{definition}

\begin{theorem}[Cancellative commuting regime]
\label{thm:commuting}
Suppose the channels $\Phi_\sigma$ commute and the resulting predictive monoid is
cancellative and embeds as a finitely generated affine semigroup in its
Grothendieck group.  Then its growth exponent is an integer
\[
 \qalg=\rmon\in\{0,1,\ldots,|E|\},
\]
and $N_{\CP}(R)\asymp R^{\rmon}$.  More precisely, $\rmon$ is the rank of the
abelian group generated by the image monoid of exponent vectors.  Without the
cancellative/affine-semigroup hypothesis, commutativity alone gives a finitely
generated commutative monoid but not this lattice-cone counting theorem.
\end{theorem}

\begin{proof}
When the channels commute, $\Phi_w$ depends only on the exponent vector
$n(w)\in\N^{|E|}$.  The map $n\mapsto\Phi_n$ has image a finitely generated
commutative monoid.  Under the additional cancellative affine-semigroup
hypothesis, this image embeds in its Grothendieck group, a finitely generated
abelian group of rank $\rmon$.  For $\rmon\ge1$, counting elements of the
image monoid of length at most $R$ is then equivalent, up to constants, to
counting lattice points in a rational cone of dimension $\rmon$, which gives
$N_{\CP}(R)\asymp R^{\rmon}$.  In the degenerate case $\rmon=0$ the image
monoid is a finite group, so $N_{\CP}$ is eventually constant and
$N_{\CP}(R)\asymp R^0$ holds trivially.
\end{proof}

\begin{assumption}[Regularly varying predictive multiplicities]
\label{ass:regular-variation}
Let $h(R)=\#\{[w]\in\W_{\CP}:\Lam_{\min}([w])=R\}$ be the number of predictive
classes of quotient length $R$.  Assume $h$ is regularly varying of index
$q-1$ with $q=\qalg>0$: there is a constant $a>0$ and a slowly varying
$L(R)\to1$ with
\[
  h(R)=a\,R^{\,q-1}L(R),\qquad R\to\infty .
\]
In the free commutative case on $\rmon$ generators this holds with
$L\equiv1$ and $h(R)=\binom{R+\rmon-1}{\rmon-1}\sim
\tfrac{1}{(\rmon-1)!}R^{\rmon-1}$, so $q=\rmon$ and $a=1/(\rmon-1)!$.  More
generally, affine-semigroup regimes with a positive rational degree have
Ehrhart-type cumulative asymptotics; the shell regular-variation hypothesis is
automatic only in the standard aperiodic/unimodular cases, or should be checked
as an explicit hypothesis.
\end{assumption}

\begin{theorem}[Renewal Weyl law and residue volume]
\label{thm:weyl-law}
Under \Cref{ass:fibres,ass:regular-variation}, the positive predictive triple has
a genuine Weyl asymptotic governed by the single constant $a$ of
\Cref{ass:regular-variation}, with $\qalg=q$:
\begin{enumerate}[label=\textup{(\roman*)},leftmargin=2.4em]
\item \textup{(counting / Weyl asymptotics)}
      \[
        N_{\CP}(R)\ \sim\ C_W\,R^{\,q},\qquad C_W=\frac{a}{q};
      \]
\item \textup{(zeta boundary residue)} the predictive zeta function
      $\zeta_{\CP}(s)=\dim(\Hh_{\rm int})\sum_{[w]}\Lam_{\min}([w])^{-s}$ has
      abscissa of convergence $q$, and
      \[
        \lim_{s\downarrow q}(s-q)\zeta_{\CP}(s)
        =\dim(\Hh_{\rm int})\,a
        =\dim(\Hh_{\rm int})\,q\,C_W;
      \]
\item \textup{(heat trace)}
      \[
        \Tr\bigl(e^{-tD_{\CP}^2}\bigr)\ \sim\
        \dim(\Hh_{\rm int})\,\frac{a\,\Gamma(q/2)}{2}\;t^{-q/2},
        \qquad t\downarrow0;
      \]
\item \textup{(residue volume)} the operator $|D_{\CP}|^{-q}$ lies in the weak
      trace ideal $\mathcal L^{1,\infty}$, is measurable, and its Dixmier trace
      equals the Weyl constant,
      \[
        \Tr_\omega\bigl(|D_{\CP}|^{-q}\bigr)
        =\dim(\Hh_{\rm int})\,C_W .
      \]
\end{enumerate}
Thus the renewal volume $C_{\rm ren}:=\Tr_\omega(|D_{\CP}|^{-q})$ is the leading
Weyl coefficient $\dim(\Hh_{\rm int})\,C_W$; the predictive triple obeys a
genuine Weyl law rather than merely possessing a finite zeta abscissa.
\end{theorem}

\begin{proof}
Write $m=\dim(\Hh_{\rm int})$; each predictive class contributes the eigenvalue
$\Lam_{\min}([w])=R$ with internal multiplicity $m$, so the eigenvalue $R$ of
$D_{\CP}$ has total multiplicity $m\,h(R)$.

\emph{(i)} By Karamata's theorem for regularly varying sequences
\cite{feller}, partial sums of a regularly varying sequence of index $q-1$
satisfy $\sum_{r\le R}h(r)\sim a R^{q}/q$.  Hence $N_{\CP}(R)\sim(a/q)R^{q}$,
which also identifies the abscissa of \Cref{thm:triple} as $\qalg=q$.

\emph{(ii)} For $s>q$,
$\zeta_{\CP}(s)=m\sum_{R\ge1}h(R)R^{-s}$.  The Abelian/Karamata theorem for
Dirichlet series with regularly varying coefficients gives
\[
  \sum_{R\ge1}h(R)R^{-s}\sim \frac{a}{s-q},\qquad s\downarrow q,
\]
so $(s-q)\zeta_{\CP}(s)\to m a$.  This is a boundary-residue statement; by
itself regular variation does not imply meromorphic continuation of
$\zeta_{\CP}$ across $s=q$.

\emph{(iii)} $\Tr(e^{-tD_{\CP}^2})=m\sum_{R\ge1}h(R)e^{-tR^2}$.  By the
Tauberian/Karamata correspondence applied to the Laplace--Stieltjes transform of
the regularly varying $h$,
\[
  \sum_{R\ge1}h(R)e^{-tR^2}\sim a\int_0^\infty r^{q-1}e^{-tr^2}\,dr
  =\frac{a}{2}\,\Gamma\!\Bigl(\frac q2\Bigr)\,t^{-q/2},
\]
giving the stated constant after multiplying by $m$.

\emph{(iv)} The counting asymptotic in (i) says that the singular-value counting
function of $|D_{\CP}|^{-q}$ has logarithmic Ces\`aro mean $mC_W$.  Equivalently,
$|D_{\CP}|^{-q}\in\mathcal L^{1,\infty}$ is measurable and
$\Tr_\omega(|D_{\CP}|^{-q})=mC_W$ by the standard Weyl-asymptotic form of
Connes' trace theorem \cite{connes}.
\end{proof}

\begin{remark}[Meromorphic strengthening]
\label{rem:weyl-meromorphic}
If the regular variation hypothesis is strengthened to an error estimate such as
$h(R)=aR^{q-1}+O(R^{q-1-\delta})$ for some $\delta>0$, or to the exact polynomial
shell formula available in the free commutative case, then the zeta function has
a genuine simple pole at $s=q$ and
\[
  \Res_{s=q}\zeta_{\CP}(s)=\dim(\Hh_{\rm int})\,a,
  \qquad
  \Tr_\omega(|D_{\CP}|^{-q})=\frac1q\Res_{s=q}\zeta_{\CP}(s).
\]
Without such additional analytic control, \Cref{thm:weyl-law} asserts the
Tauberian boundary residue and Dixmier-trace volume, not meromorphic
continuation.
\end{remark}

\begin{remark}[One constant, four readings]
\label{rem:weyl-one-constant}
\Cref{thm:weyl-law} ties the counting measure, the zeta boundary residue, the
small-time heat kernel, and the noncommutative volume to the \emph{same} leading
coefficient.  This is the spectral-geometric content that a bare length triple
lacks: the predictive dimension $\qalg$ is not only the zeta abscissa but a
genuine Weyl exponent, and the renewal volume
$C_{\rm ren}=\Tr_\omega(|D_{\CP}|^{-q})$ is its Dixmier-trace volume.  The
regular-variation hypothesis \Cref{ass:regular-variation} is the precise
condition converting the $\limsup$ of \Cref{thm:triple} into a genuine asymptotic law; meromorphic continuation is the separate strengthening recorded in
\Cref{rem:weyl-meromorphic}.
\end{remark}

% =====================================================================


\subsection{Predictive order}

\label{sec:order}

The same quotient that carries the positive spectral triple also carries a
natural order.  This order is not imposed externally.  It is the divisibility
order of accumulated predictive channels: one predictive state precedes another
when the latter can be obtained from the former by further renewal evolution.
Thus the positive sector supplies both a spectral dimension and the order
skeleton that will later be used as the causal substrate.

Let $\M$ be the monoid generated by the channels $\Phi_\sigma$.

\begin{definition}[Predictive divisibility order]
\label{def:predictive-order}
For $a,b\in\M$ write $a\preceq b$ if $b=\Psi a$ for some $\Psi\in\M$.  Write $a\approx b$ if $a\preceq b$ and $b\preceq a$.  The predictive poset is
\[
 \Pp=\M/{\approx}
\]
with the induced order.
\end{definition}

\begin{lemma}[Poset reflection]
\label{lem:poset}
The relation on $\Pp$ induced by Definition~\ref{def:predictive-order} is a partial order.
\end{lemma}

\begin{proof}
Left divisibility is a preorder on $\M$.  Quotienting a preorder by its symmetric part gives a partial order.
\end{proof}

\begin{assumption}[Aperiodic elementary grading]
\label{ass:aperiodic}
On the condensation \(\Pp=\M/{\approx}\), the quotient length satisfies:
\begin{enumerate}[label=\textup{(\roman*)},leftmargin=2.3em]
\item it is constant on every mutual-reachability class, and therefore descends
      to a function \(\Lambda_{\Pp}:\Pp\to\N\);
\item it is strictly increasing on strict divisibility classes;
\item every Hasse cover \(x\lessdot y\) is elementary in the sense that
      \[
        \Lambda_{\Pp}(y)=\Lambda_{\Pp}(x)+1;
      \]
\item every length ball \(\{x\in\Pp:\Lambda_{\Pp}(x)\le R\}\) is finite.
\end{enumerate}
The first two items supply a time function; the third is the additional
hypothesis that promotes that time function to a genuine rank function.
\end{assumption}

\begin{proposition}[Causal-order skeleton]
\label{prop:causal-skeleton}
Under \Cref{ass:aperiodic}, \((\Pp,\preceq)\) is a locally finite graded poset
with rank function
\[
  T(x)=\Lambda_{\Pp}(x).
\]
In particular it is a causal-set-like order skeleton.  If only
\Cref{ass:aperiodic}\textup{(i),(ii),(iv)} are assumed, the same conclusion holds
with ``graded rank'' replaced by ``strictly increasing integer-valued time
function.''
\end{proposition}

\begin{proof}
The quotient by mutual reachability is a partial order by \Cref{lem:poset}.
Items \textup{(i)}--\textup{(ii)} make \(T\) well defined and strictly increasing
on strict comparable pairs.  Item \textup{(iii)} says exactly that every Hasse
cover raises \(T\) by one, so \(T\) is a rank function and the poset is graded.
Finally, an interval \([x,y]\) lies in the finite length ball of radius
\(T(y)\), hence is finite by \textup{(iv)}.  This proves local finiteness.
\end{proof}

\begin{proposition}[Chain-counting renewal resolvent]
\label{prop:incidence-zeta}
Let $\Phi$ denote the cover transfer operator of the Hasse diagram of $\Pp$.  In
the completed path-counting incidence algebra \cite{stanley} graded by rank,
\[
 \mathcal R_{\Pp}:=\sum_{k\ge0}\Phi^k=(1-\Phi)^{-1}.
\]
The kernel $\mathcal R_{\Pp}(x,y)$ counts saturated chains from $x$ to $y$.  It
coincides with the ordinary incidence zeta $\zeta_{\Pp}(x,y)=\one_{x\preceq y}$
only when each interval has a unique saturated chain, or after passing to the
Boolean incidence algebra in which all positive multiplicities are identified.
\end{proposition}

\begin{proof}
The $k$th power $\Phi^k$ counts saturated chains of length $k$.  Summing over all
$k$ therefore gives the chain-counting renewal resolvent.  In a general poset
this is not the ordinary incidence zeta: for example, in $\N^2$ there are two
saturated chains from $(0,0)$ to $(1,1)$, while the incidence zeta takes the
value $1$.  The two notions agree precisely in the unique-chain case, or after
Booleanising multiplicities.
\end{proof}

\begin{remark}[Markov and renewal singularities]
\label{rem:markov-renewal}
This identification makes the analytic type of memory visible in the order.
Markov holding times give rational renewal resolvents and pole singularities.
Heavy-tailed renewal laws give branch-point singularities and fractional
exponents.  Thus genuine renewal memory can appear as a fractional singularity
of the chain-counting renewal resolvent, not merely as a probabilistic feature of
the source.
\end{remark}




% =====================================================================
\subsection{The positivity boundary}
\label{sec:positive-obstruction}

The completely positive construction supplies a Hilbert-positive spectral
geometry, residue data, and a predictive order.  It does not, by itself,
distinguish a negative subspace.  The following two statements isolate the
precise obstruction and identify the extra information carried by any compatible
nontrivial fundamental symmetry.

\begin{assumption}[Irreducible natural positive sector]
\label{ass:irreducible-positive-sector}
Let \(\mathfrak B_+\subset\mathcal B(\ell^2(\W_{\CP})\otimes\Hh_{\rm int})\)
be the unital \(*\)-algebra generated by the represented renewal shifts,
pointwise internal observables, and admissible relabellings.  Assume
\[
  \mathfrak B_+'=\C\one.
\]
A fundamental symmetry extracted functorially from the positive sector is
required to commute with \(\mathfrak B_+\) and may not use an added sign, phase,
orientation, or sheet label.
\end{assumption}

\begin{theorem}[Irreducible positive-sector no-go]
\label{thm:irreducible-positive-no-go}
Under \Cref{ass:irreducible-positive-sector}, every functorially extracted
self-adjoint involution is scalar,
\[
  J=+\one\qquad\text{or}\qquad J=-\one.
\]
Its associated form is definite.  Hence the natural irreducible positive
representation does not determine a Lorentzian splitting.
\end{theorem}

\begin{proof}
Naturality places \(J\) in \(\mathfrak B_+'=\C\one\), so \(J=c\one\).
Self-adjointness and \(J^2=\one\) force \(c=\pm1\).
\end{proof}

\begin{theorem}[Compatible fundamental symmetries are signed data]
\label{thm:functorial-positive-obstruction}
Let a recurrent positive predictive component be fixed.  Suppose a self-adjoint
involution \(J\) commutes with the atomic diagonal algebra and, for every
descended generator \(\widehat S_\sigma\), satisfies
\[
  J\widehat S_\sigma=c(\sigma)\widehat S_\sigma J,
  \qquad c(\sigma)\in\{\pm1\}.
\]
Then the edgewise ratios of the diagonal signs of \(J\) define a
\(\Z/2\)-valued cocycle on the predictive graph.  Conversely every such cocycle
produces the corresponding exchange relations on its signed double cover.
Thus every compatible nonconstant diagonal fundamental symmetry contains
precisely the additional signed information classified in \Cref{sec:signed}.
\end{theorem}

\begin{proof}
Write \(Je_x=\eta(x)e_x\), with \(\eta(x)\in\{\pm1\}\).  Along a generator
edge \(x\to\sigma x\), the exchange relation gives
\(\eta(\sigma x)=c(\sigma)\eta(x)\).  These edge ratios form a
\(\Z/2\)-valued one-cochain; on a graph every one-cochain is a cocycle, and a
change of vertex signs changes it by a coboundary.  The converse is the signed
cover construction of \Cref{constr:signed-cover}.
\end{proof}

\begin{corollary}[CP positivity boundary]
\label{cor:positivity-obstruction}
The completely positive sector canonically supplies positive spectral geometry,
dimension, volume, and order.  A nontrivial Lorentzian/Krein splitting requires
additional signed data; in the irreducible natural representation it cannot be
extracted from the positive sector alone.
\end{corollary}


% =====================================================================
\section{The signed Krein sector}
\label{sec:signed}

The positivity boundary identifies what must be added without determining a
particular nonzero class.  We first characterize the objectwise-minimal
renewal-positive indefinite enrichment, then give its concrete double-cover
realization, distinguish the real signed layer from its amplitude lift, and
finally construct and characterize the modular signed Dirac operator.  The
section closes with a reconstruction theorem for the graded signed predictive
geometry.

\subsection{Minimal indefinite enrichments over a fixed component}
\label{sec:setup}
\label{sec:signed-reconstruction}

We now strengthen the signed-cover construction from existence to categorical
minimality.  Throughout this subsection \(G=(V,E_G,\Lambda)\) is one fixed
finite connected recurrent predictive component.  The edge count
\(|E_{\mathrm{geom}}|\) below counts unoriented geometric edges, equivalently one
chosen orientation from each pair \(e,\bar e\).  Thus
\[
  H^1(G,\Z/2)\cong(\Z/2)^{b_1(G)},\qquad
  b_1(G)=|E_{\mathrm{geom}}|-|V|+1.
\]

The categorical level is important.  We classify enrichments \emph{over this
fixed \(G\)} and then show that the classification itself is contravariantly
natural under renewal isomorphisms.  We do not require one global enrichment
functor to choose a class on every object simultaneously; such a requirement
would restrict the value at \(G\) to automorphism-invariant classes.

\begin{definition}[Renewal-positive indefinite enrichment over \(G\)]
\label{def:fixed-enrichment}
An enrichment over the fixed positive predictive component \(G\) consists of:
\begin{enumerate}[label=\textup{(\roman*)},leftmargin=2.3em]
\item a finite Hermitian local system \(F=(F_x)_{x\in V}\) with a distinguished
      orthonormal renewal basis and its simplicial cone \(\mathcal C_x\);
\item unitary edge transports \(T_e:F_x\to F_y\), with
      \(T_{\bar e}=T_e^{-1}\), augmenting the renewal shifts;
\item fibrewise self-adjoint involutions \(J_x\), commuting with the full atomic
      diagonal algebra generated by the distinguished renewal-basis projections,
      as well as with the length grading;
\item a homogeneous degree \(\epsilon(e)\in\Z/2\) for each edge,
      \[
        J_yT_e=(-1)^{\epsilon(e)}T_eJ_x;
      \]
\item structural recovery of the untwisted positive base by the deck quotient.
\end{enumerate}
It is \emph{renewal-positive} when
\(T_e(\mathcal C_x)=\mathcal C_y\) for every edge, and
\emph{fibrewise indefinite} when both eigenspaces of every \(J_x\) are nonzero.
Its fibre rank is \(\rank F:=\dim_{\C}F_x\), independent of \(x\) because \(G\)
is connected and the transports are unitary.  It is \emph{objectwise minimal}
when this constant rank is least among all renewal-positive
fibrewise-indefinite enrichments over \(G\).
\end{definition}

\begin{definition}[Gauge equivalence and pullback]
\label{def:fixed-gauge}
An isomorphism over the identity of \(G\) is a vertexwise cone-preserving unitary
which intertwines the transports and transforms the local fundamental
symmetries covariantly.  On a rank-two fibre a sheet trivialisation writes
\[
  J_x=(-1)^{j(x)}J_0,\qquad
  J_0=\diag(1,-1),\qquad
  T_e=S^{c(e)},\qquad
  S=\begin{pmatrix}0&1\\1&0\end{pmatrix}.
\]
A change of sheet orientation \(a:V\to\Z/2\) acts by
\[
  j\longmapsto j+a,\qquad c\longmapsto c+\delta a.
\]
The homogeneity degree \(\epsilon\) is invariant and satisfies
\[
  \epsilon=c+\delta j.
\]
For a renewal isomorphism \(\varphi:G\to G'\), pullback transports, cones, and
local symmetries define
\[
  \varphi^*:\operatorname{Enr}_{\min}(G')
  \longrightarrow\operatorname{Enr}_{\min}(G).
\]
\end{definition}

\begin{lemma}[Cone-positive unitaries are permutations]
\label{lem:cone-positive-permutation}
A unitary \(T:F_x\to F_y\) satisfying
\(T(\mathcal C_x)=\mathcal C_y\) permutes the distinguished renewal basis.
\end{lemma}

\begin{proof}
A linear isomorphism carrying a simplicial cone onto a simplicial cone maps
extremal rays bijectively to extremal rays.  Hence \(T\) sends each renewal basis
vector to a positive scalar multiple of a renewal basis vector.  Unitarity makes
that scalar equal to one.
\end{proof}

\begin{lemma}[Minimal rank and reduction of the structure group]
\label{lem:minimal-rank-two}
Every fibrewise-indefinite enrichment has \(\rank F\ge2\).  The signed cover
realises rank two.  On an objectwise-minimal enrichment,
\[
  \rank F=2,\qquad T_e\in\{I,S\}.
\]
In particular all \(U(1)\) phases and residual signs are excluded and the
structure group is \(\{I,S\}\cong\Z/2\).
\end{lemma}

\begin{proof}
Fibrewise indefiniteness gives
\(\dim F_x^+\ge1\) and \(\dim F_x^-\ge1\), hence \(\rank F\ge2\).
The two-sheeted signed cover has one positive and one negative line in each
fibre, so rank two is attainable.  Commutation with the full atomic diagonal
algebra makes every \(J_x\) diagonal in the distinguished renewal basis.  At
rank two, fibrewise indefiniteness therefore forces
\(J_x=\pm\diag(1,-1)\).  Moreover,
\Cref{lem:cone-positive-permutation} leaves exactly the two permutation
matrices \(I\) and \(S\) as edge transports.  They are respectively even and
odd relative to \(J_0\).  Matrices \(-I\), \(-S\), and arbitrary phases fail
cone positivity.
\end{proof}

\begin{lemma}[Two cochains, one cohomology class]
\label{lem:two-cochains-one-class}
For an objectwise-minimal enrichment, choose a sheet trivialisation and write
\(T_e=S^{c(e)}\), \(J_x=(-1)^{j(x)}J_0\).  Then:
\begin{enumerate}[label=\textup{(\roman*)},leftmargin=2.3em]
\item \(c\in C^1(G;\Z/2)\) is the transition cochain of a principal
      \(\Z/2\)-cover;
\item because \(G\) is a \(1\)-complex, \(C^2(G;\Z/2)=0\), so every
      \(1\)-cochain is a cocycle;
\item a sheet change \(a\in C^0(G;\Z/2)\) sends
      \(c\mapsto c+\delta a\);
\item the gauge-invariant homogeneity degree is
      \[
        \epsilon(e)=c(e)+j(x)+j(y)=c(e)+(\delta j)(e),
      \]
      and therefore \([\epsilon]=[c]\in H^1(G,\Z/2)\).
\end{enumerate}
\end{lemma}

\begin{proof}
The edge transport either preserves or exchanges the two sheets, producing
\(c(e)\).  Since the graph has no \(2\)-cells, \(\delta^1=0\), so
\(Z^1=C^1\).  Conjugating by \(S^{a(x)}\) at the vertices gives the standard
coboundary action.  Finally,
\[
 (-1)^{j(y)}J_0S^{c(e)}
 =(-1)^{j(y)+c(e)}S^{c(e)}J_0
 =(-1)^{c(e)+j(x)+j(y)}T_eJ_x,
\]
which proves the formula for \(\epsilon\) and the equality of classes.
\end{proof}

\begin{remark}[The invariant is the cover]
\label{rem:cover-is-invariant}
The principal cover \(P\to G\), equivalently \([c]\), is the intrinsic invariant.
The representatives \(c\), \(j\), and \(J\) are gauge-dependent realization
data.  For a fixed cover all sheet-sign choices form one gauge orbit, so
\([J]\) contributes no independent invariant.  The deck involution and the
fundamental symmetry are distinct: in a sheet trivialisation the deck swap
anticommutes with \(J_0\).
\end{remark}

\begin{theorem}[Categorical minimality of the signed enrichment]
\label{thm:classification}
Let \(G\) be a finite connected recurrent predictive component.  Let
\(\operatorname{Enr}_{\min}(G)\) denote the groupoid of objectwise-minimal
renewal-positive fibrewise-indefinite enrichments over \(G\), with morphisms
covering the identity.  Then
\[
  \boxed{\qquad
  \pi_0\operatorname{Enr}_{\min}(G)
  \ \cong\ H^1(G,\Z/2)
  \ \cong\ (\Z/2)^{b_1(G)}.
  \qquad}
\]
Every such enrichment is the signed cover \(G^\chi\), its constant fibre rank is
two, and its fundamental symmetry is gauge-covariant realization data of the
cover.  These bijections are contravariantly natural under renewal
isomorphisms.
\end{theorem}

\begin{proof}
By \Cref{lem:minimal-rank-two}, every minimal enrichment has rank two and every
edge transport is \(S^{c(e)}\).  By
\Cref{lem:two-cochains-one-class}, these transitions define a principal
\(\Z/2\)-cover and gauge equivalence is precisely quotienting \(C^1\) by
coboundaries.  Thus
\[
 \pi_0\operatorname{Enr}_{\min}(G)
 \cong C^1(G;\Z/2)/\delta C^0(G;\Z/2)
 =H^1(G,\Z/2).
\]
Conversely every principal double cover gives the rank-two signed enrichment,
so the map is bijective.  For a connected finite graph, the cohomology dimension
is \(|E_{\mathrm{geom}}|-|V|+1\).  Pullback along a renewal isomorphism pulls back
the transition cochain and all associated local data, so the class of the
pulled-back enrichment is \(\varphi^*[c]\).
\end{proof}

\begin{corollary}[Removable and connected sectors]
\label{cor:sector-count}
There are \(2^{b_1(G)}\) gauge classes.  The zero class gives the split cover
\(G\sqcup G\).  The \(2^{b_1(G)}-1\) nonzero classes are exactly the
non-removable sectors, and each corresponding double cover is connected.
\end{corollary}

\begin{proof}
A principal double cover is classified by the homomorphism
\(\pi_1(G)\to\Z/2\) represented by its cohomology class.  The zero map gives the
split cover.  Every nonzero map has image all of \(\Z/2\), hence its cover is
connected.  Triviality of the class is exactly removability by a sheet gauge.
\end{proof}

\begin{remark}[Edge-locality, label-locality, and multiplicativity]
\label{rem:locality-multiplicativity}
The theorem uses independent edge-local transitions.  If signs are required to
be constant on reset-label classes, the realised classes form the subgroup
\[
 H^1_{\mathrm{loc}}(G,\Z/2)
 :=\{[c]:c\text{ has a label-local representative}\}.
\]
Literal multiplicativity is not a hypothesis of the classification.  Tensoring
two rank-two minimal enrichments gives rank four on the product component,
although that component itself admits a rank-two signed cover.  Thus global
objectwise minimality and literal tensor-product multiplicativity are
incompatible.  Multiplicativity may be imposed only as optional downstream
coherence on chosen monoidal generators, producing rank \(2^r\) on an
\(r\)-fold product and the usual Koszul/super structure.
\end{remark}

\begin{remark}[Role of cone preservation]
\label{rem:positivity-load-bearing}
Without preservation of the distinguished renewal cone, rank-two homogeneous
unitaries carry arbitrary phase holonomy and the classification enlarges beyond
\(H^1(G,\Z/2)\).  Without fibrewise indefiniteness, one-dimensional vertex-sign
assignments survive but realise only the trivial class.  Without structural base
recovery, the quotient need not return the positive predictive component.
Thus positivity, locality, base recovery, and fibrewise indefiniteness are the
essential hypotheses; naturality belongs to the family of classifications,
not to the definition of one enrichment over a fixed \(G\).
\end{remark}

\begin{remark}[Holonomy, not deficiency]
\label{rem:not-deficiency}
The classifying invariant is the holonomy class \([c]\), not a spectral-radius
deficiency.  Deficiency depends on a chosen transfer operator and may fail to
detect half-period or otherwise spectrally silent classes.
\end{remark}


\subsection{A record-theoretic origin of the signed class}
\label{subsec:record-origin-signed-class}

The classification above treats \([\chi]\) as the intrinsic datum of the
minimal indefinite enrichment.  In a record-regular operational realisation,
the class can itself be generated by the orientation of the stable
future-distinct record bundle.

\begin{proposition}[Stable-record determinant orientation]
\label{prop:record-determinant-orientation}
Let \(G\) be a connected recurrent predictive graph.  Assign to every vertex
\(x\) a finite-dimensional real vector space \(\mathcal R_x\) of stable
future-distinct records, of constant rank, and to every oriented edge
\(e:x\to y\) an invertible record transport
\[
   L_e:\mathcal R_x\longrightarrow\mathcal R_y,
   \qquad L_{\bar e}=L_e^{-1}.
\]
After choosing local orientations, define
\[
   \chi_{\rm rec}(e):=\sgn\det L_e\in\{\pm1\}.
\]
Then a change of local record orientation changes \(\chi_{\rm rec}\) by a
vertex coboundary, while the product around every loop is invariant.  Hence
\[
   [\chi_{\rm rec}]
   =w_1(\det\mathcal R)
   \in H^1(G,\Ztwo)
\]
is well defined, and the principal double cover classified by this class is
exactly the orientation cover of the stable-record determinant line.
\end{proposition}

\begin{proof}
Choose an oriented basis in each \(\mathcal R_x\).  Reversing the chosen
orientation at a vertex \(x\) multiplies the determinant sign of every incident
transport by \(-1\); if \(g(x)\in\{\pm1\}\) records that choice, then
\[
   \chi_{\rm rec}(e:x\to y)
   \longmapsto g(y)\chi_{\rm rec}(e)g(x),
\]
which is the standard coboundary action.  The vertex factors cancel in the
product around a closed loop, so the loop sign depends only on the cohomology
class.  A lift of a path to the orientation torsor records whether its
transport preserves or reverses the determinant orientation; its monodromy is
therefore precisely the loop product of \(\chi_{\rm rec}\).  This is the
principal cover of \Cref{thm:cover}.
\end{proof}

\begin{remark}[Affinity and orientation are different classes]
\label{rem:affinity-versus-orientation}
The same transports also define the real one-cochain
\[
   A_{\rm rec}(e):=\log|\det L_e|.
\]
Under a positive rescaling of local determinant frames it changes by a real
coboundary and therefore defines an additive class
\([A_{\rm rec}]\in H^1(G;\R)\).  This class measures determinant magnitude,
whereas \(w_1(\det\mathcal R)\) measures determinant parity.  Neither determines
the other.  The present paper classifies the signed geometry produced by the
second class; a companion upstream construction derives stable record spaces
and their transports from operational prediction.
\end{remark}

\subsection{Concrete signed covers and Krein symmetry}
\label{subsec:signed-cover}

Fix a recurrent predictive component, represented by a connected directed graph
$G=(V,E_G)$.  Its vertices are predictive classes, and an edge
$e:x\to y$ records a generator transition $y=\sigma\!\cdot x$.  The quotient
length is denoted by $\Lambda:V\to\N$.

\begin{definition}[Sign cocycle]
\label{def:sign-cocycle}
A sign cocycle on $G$ is a map
\[
  \chi:E_G\longrightarrow \{\pm1\}.
\]
For an oriented path $p=e_k\cdots e_1$ we write
\[
  \chi(p)=\prod_{j=1}^k\chi(e_j).
\]
Since a graph has no $2$-cells, every $\{\pm1\}$-valued $1$-cochain is a
$1$-cocycle.  A flat gauge is a vertex function $g:V\to\{\pm1\}$, acting by
\[
  \chi^g(e:x\to y)=g(y)\,\chi(e)\,g(x).
\]
The gauge-invariant content of $\chi$ is its holonomy class
\[
  [\chi]\in H^1(G,\Z/2).
\]
Equivalently, $[\chi]$ is the homomorphism from the cycle space of $G$ to
$\{\pm1\}$ given by the product of signs around loops.
\end{definition}

\begin{remark}[Relation with history signs]
\label{rem:history-signs}
If the original renewal memory carries a multiplicative history sign
$\eps:E^*\to\{\pm1\}$, then it defines a sign cocycle on the predictive graph
provided it is compatible with the predictive quotient: whenever two generator
transitions define the same edge of $G$, they have the same sign.  In that case
$\chi(e)$ is the sign of any history segment representing the edge $e$.  This
compatibility is the precise form of ``internality'' needed for the sign to
descend to predictive geometry.
\end{remark}

\begin{construction}[Signed predictive cover]
\label{constr:signed-cover}
The signed predictive cover associated with $\chi$ is the graph
$G^\chi=(V^\chi,E_G^\chi)$ with
\[
  V^\chi=V\times\{\pm1\}.
\]
For every edge $e:x\to y$ in $G$ and every sheet label $\eta\in\{\pm1\}$,
there is an edge
\[
  (e,\eta):(x,\eta)\longrightarrow (y,\eta\,\chi(e)).
\]
The projection
\[
  \pi:G^\chi\longrightarrow G,\qquad \pi(x,\eta)=x,
\]
is a two-sheeted graph cover, and the deck transformation is
\[
  \tau(x,\eta)=(x,-\eta).
\]
We write $\W_{\CP}^{\chi}$ for the corresponding signed predictive quotient.
\end{construction}

\begin{theorem}[Principal $\Z/2$-cover and holonomy]
\label{thm:cover}
The projection $\pi:G^\chi\to G$ is a principal $\Z/2$-cover.  Its monodromy
around a loop $\ell$ is multiplication by $\chi(\ell)$.  The cover is
gauge-equivalent to the trivial double cover if and only if
\[
  [\chi]=0\in H^1(G,\Z/2).
\]
Thus the non-removable content of the signed cover is precisely the holonomy
class $[\chi]$.
\end{theorem}

\begin{proof}
The deck transformation $\tau$ acts freely and transitively on each fibre
$\pi^{-1}(x)=\{(x,+1),(x,-1)\}$, so $\pi$ is a principal $\Z/2$-cover.  Lifting
a loop $\ell=e_k\cdots e_1$ from $(x,\eta)$ gives the endpoint
$(x,\eta\chi(\ell))$, hence the monodromy is the cycle-sign homomorphism.

If $[\chi]=0$, then $\chi(e:x\to y)=g(y)g(x)$ for some
$g:V\to\{\pm1\}$, and the map $(x,\eta)\mapsto (x,\eta g(x))$ identifies
$G^\chi$ with the trivial cover.  Conversely, a trivialisation chooses one
sheet over each vertex; comparing the chosen sheet across each edge gives a
vertex gauge whose coboundary is $\chi$.  Hence $\chi$ is cohomologically
trivial exactly when the cover is removable.
\end{proof}

\begin{definition}[Canonical fundamental symmetry on the signed cover]
\label{def:J-eps}
On
\[
  \Kh_\chi=\ell^2(V^\chi)\otimes\Hh_{\rm int}
\]
define
\[
  J_\chi e_{(x,\eta)}=\eta\,e_{(x,\eta)}.
\]
The quotient length lifts to the cover by
\[
  N e_{(x,\eta)}=\Lambda(x)\,e_{(x,\eta)}.
\]
For an edge $e:x\to y$, the lifted shift is
\[
  \widetilde S_e e_{(x,\eta)}=e_{(y,\eta\chi(e))},
\]
on its natural domain, and zero away from the source fibre.
\end{definition}

\begin{proposition}[Krein datum]
\label{prop:krein-datum}
The operator $J_\chi$ is a self-adjoint unitary.  The form
\[
  [u,v]_\chi=\langle u,J_\chi v\rangle
\]
is a Krein form whenever both sheets occur.  Moreover the lifted shifts satisfy
the exchange relation
\[
  J_\chi\,\widetilde S_e=\chi(e)\,\widetilde S_e\,J_\chi .
\]
A flat gauge changes the representative cocycle but leaves the loop holonomies,
and hence the class $[\chi]$, unchanged.
\end{proposition}

\begin{proof}
Since $\eta=\pm1$, $J_\chi^2=1$ and $J_\chi^*=J_\chi$.  The form is positive
on the $+1$ sheet and negative on the $-1$ sheet, giving a fundamental
decomposition.  Finally,
\[
  J_\chi\widetilde S_e e_{(x,\eta)}
  =J_\chi e_{(y,\eta\chi(e))}
  =\eta\chi(e)e_{(y,\eta\chi(e))}
  =\chi(e)\widetilde S_eJ_\chi e_{(x,\eta)} .
\]
Gauge transformations multiply edge signs by coboundaries, which do not change
products around closed loops.
\end{proof}

\begin{corollary}[Removability is cohomology]
\label{cor:removability}
The signed sector is removable by a flat gauge if and only if
\[
  [\chi]=0\in H^1(G,\Z/2).
\]
Equivalently, nonzero classes in $H^1(G,\Z/2)$ are precisely the
non-removable signed/Krein sectors.
\end{corollary}

\begin{proof}
This is the triviality criterion in \Cref{thm:cover}, translated from covers
to the associated fundamental symmetries.
\end{proof}

\subsection{Real signed quotients and amplitude lifts}
\label{subsec:real-amplitude-lifts}

The classification above is intentionally a real/Krein classification.  It
answers the question: what is the objectwise-minimal indefinite enrichment of a
positive renewal component that still preserves the renewal cone and recovers the
positive base?  The answer is a principal \(\Z/2\)-cover.  A different question
appears only when the same holonomy is asked to act on amplitudes or charged
local systems: a sign is then not the transported quantity itself, but the square
of a unitary phase.  The next elementary facts record this distinction without
changing the signed-cover theorem.

\begin{definition}[Amplitude lift of a signed cohomology class]
\label{def:amplitude-lift}
Let \([c]\in H^1(G,\Ztwo)\) be the signed-cover class of
\Cref{thm:classification}.  A finite amplitude lift of \([c]\) is a class
\([\widetilde c]\in H^1(G,\Zfour)\) whose reduction modulo two is \([c]\):
\[
  \rho([\widetilde c])=[c],
  \qquad \rho:\Zfour\longrightarrow\Ztwo .
\]
Equivalently, if a cycle has lifted exponent \(k\in\Zfour\), the associated
unitary phase is \(i^k\) and its real signed shadow is
\[
  (i^k)^2=(-1)^k .
\]
Thus the \(\Z/2\) class records the square of the amplitude transport, while the
\(\Z/4\) lift records a choice of square root.
\end{definition}

\begin{lemma}[Surjectivity of the amplitude-lift map on graphs]
\label{lem:z4-lift-surjects}
For every finite connected graph \(G\), reduction modulo two gives a surjection
\[
  \rho:H^1(G,\Zfour)\twoheadrightarrow H^1(G,\Ztwo).
\]
If \(b_1(G)=r\), then every signed class has exactly \(2^r\) amplitude lifts.
\end{lemma}

\begin{proof}
A connected graph is homotopy equivalent to a wedge of \(r=b_1(G)\) circles.
Hence
\[
  H^1(G,\Zfour)\cong(\Zfour)^r,
  \qquad
  H^1(G,\Ztwo)\cong(\Ztwo)^r .
\]
Under these identifications \(\rho\) is coordinatewise reduction
\(\Zfour\to\Ztwo\), which is surjective and has a two-point fibre in each
coordinate.  Therefore the fibre over every element of \((\Ztwo)^r\) has size
\(2^r\).
\end{proof}

\begin{theorem}[Minimal finite square-root completion]
\label{thm:minimal-z4-amplitude-lift}
Let a signed renewal sector contain a loop whose real holonomy is \(-1\).  Any
finite phase group acting on amplitudes and containing a unitary square root of
that sign contains a cyclic subgroup of order four.  Consequently the minimal
finite amplitude completion of a nontrivial signed sector is \(\Zfour\).  Its
reduction modulo two is exactly the principal \(\Ztwo\)-cover of
\Cref{thm:classification}.
\end{theorem}

\begin{proof}
Let \(u\) be the amplitude phase assigned to a loop with real signed holonomy
\(-1\).  The square-root condition is \(u^2=-1\).  Then \(u^4=1\), while
\(u^2\ne1\); hence \(u\) has order exactly four.  Therefore the subgroup
\(\langle u\rangle\) is cyclic of order four, so any finite phase group carrying
such a square root contains a copy of \(\Zfour\).  Conversely \(\Zfour\) works:
its generator may be represented by \(i\), whose square is \(-1\).  Reducing the
exponent of \(i^k\) modulo two records precisely whether the square is \(+1\) or
\(-1\), which is the original \(\Ztwo\) signed-cover class.
\end{proof}

\begin{corollary}[Conservativity over the real/Krein layer]
\label{cor:z4-conservative-real-layer}
Replacing a signed class \([c]\in H^1(G,\Ztwo)\) by any amplitude lift
\([\widetilde c]\in H^1(G,\Zfour)\) with \(\rho([\widetilde c])=[c]\) does not
alter any construction depending only on the principal double cover, the Krein
fundamental symmetry, or the real commutator polar form.  The \(\Zfour\) datum is
extra amplitude information above the real quotient, not a replacement of the
\(\Ztwo\) Lorentzian-signature mechanism.
\end{corollary}

\begin{proof}
The principal double cover, sheet swap, and fundamental symmetry are obtained by
applying \(\rho\) to the amplitude lift.  If two amplitude lifts have the same
reduction, they define the same \(\Ztwo\)-cover and therefore the same signed
sheet Hilbert space, the same diagonal fundamental symmetry \(J_\chi\), and the
same exchange signs used in \Cref{lem:cycle-krein-exchange} and
\Cref{thm:canonical-temporal-row}.  Likewise the polar commutator form is the
antisymmetric real quotient of the projective multiplier; scalar rephasings and
choices of square root may change the lifted fourth-root phases, but not their
reduction to the \(\{\pm1\}\)-valued commutator.  Thus the real/Krein theorems
are unchanged.
\end{proof}

\begin{remark}[Where the lift is used]
\label{rem:z4-lift-where-used}
The \(\Ztwo\) layer is sufficient for Lorentzian signature and for the temporal
Krein exchange.  The \(\Zfour\) layer becomes relevant only when one retains
amplitude-level local systems: charged holonomy, square-root phases of
projective transports, or finite matter-sector refinements.  In particular, a
proof carried out purely in the real category should be stated as a real-category
statement; it need not control all amplitude-level splittings after a
\(\Zfour\) lift is retained.
\end{remark}

% =====================================================================
The cover and the fundamental symmetry provide the indefinite inner product.
To obtain the signed Dirac operator we combine the symmetry with the renewal
clock.  The clock is the same quotient length as in the positive sector, but it
is inserted through an exponential modular weight.

\begin{definition}[Signed modular twist]
\label{def:signed-modular-twist}
Let
\[
  \Delta=e^{\beta N},\qquad \beta>0.
\]
For an operator $a$ for which $\Delta a\Delta^{-1}$ extends boundedly, set
\[
  \sigma_\beta(a)=\Delta a\Delta^{-1},
  \qquad
  \rho_\chi(a)=\operatorname{Ad}(J_\chi)\sigma_\beta(a)
  =J_\chi\Delta a\Delta^{-1}J_\chi .
\]
The signed modular Dirac operator is
\[
  D_\chi=J_\chi\Delta+B,
\]
where $B$ is a bounded $J_\chi$-self-adjoint perturbation, i.e.
$J_\chi B^*J_\chi=B$.
\end{definition}

\begin{theorem}[Canonical signed modular Dirac]
\label{thm:signed-dirac}
Assume that the length shells of $N$ are finite and that $N\to\infty$ along the
signed quotient.  Let $\mathcal A_\chi$ be a $*$-algebra of bounded operators on
$\Kh_\chi$ such that, for every $a\in\mathcal A_\chi$,
\[
  \rho_\chi(a)\ \text{is bounded},
  \qquad
  [B,a]_{\rho_\chi}:=Ba-\rho_\chi(a)B\ \text{is bounded}.
\]
Then
\[
  D_\chi=J_\chi e^{\beta N}+B
\]
is Krein-self-adjoint with respect to $J_\chi$ and defines a twisted Krein
spectral triple
\[
  (\mathcal A_\chi,\Kh_\chi,D_\chi;\rho_\chi).
\]
Its resolvent is compact.  If the shell multiplicities have subexponential
growth, and in particular if they have the polynomial growth of
\Cref{ass:regular-variation}, then
\[
  \operatorname{Tr}\bigl(1+D_\chi^*D_\chi\bigr)^{-s/2}<\infty
  \qquad\text{for every }s>0 .
\]
Equivalently, the regularised singular-value summability abscissa is zero.  If
\(D_\chi\) is invertible, the unregularised statement
\[
  \operatorname{Tr}|D_\chi|^{-s}<\infty
  \qquad\text{for every }s>0
\]
also holds.  The summability is signature-blind and depends only on the growth
of length shells, not on the sign holonomy.
\end{theorem}

\begin{proof}
Put \(A:=J_\chi\Delta\), where \(\Delta=e^{\beta N}\).  Since \(J_\chi\)
commutes with \(N\), the operator \(A\) is Hilbert-self-adjoint on
\(\operatorname{Dom}\Delta\), satisfies \(|A|=\Delta\), and has compact
resolvent.  Moreover
\[
  J_\chi A^*J_\chi=A.
\]
Together with \(J_\chi B^*J_\chi=B\), this gives
\[
  J_\chi D_\chi^*J_\chi=D_\chi,
\]
so \(D_\chi\) is Krein-self-adjoint.

Compactness of the resolvent can be verified without invoking an unproved
operator inequality.  Choose \(\eta>\|B\|\).  Since \(A\) is self-adjoint,
\(\|(A-i\eta)^{-1}\|\le\eta^{-1}\), and
\[
  D_\chi-i\eta
  =\bigl(\one+B(A-i\eta)^{-1}\bigr)(A-i\eta).
\]
The first factor is invertible by a Neumann series.  Hence
\[
  (D_\chi-i\eta)^{-1}
  =(A-i\eta)^{-1}
   \bigl(\one+B(A-i\eta)^{-1}\bigr)^{-1},
\]
which is compact because \((A-i\eta)^{-1}\) is compact.

For \(a\in\mathcal A_\chi\), compute on the common core of finitely supported
vectors:
\[
\begin{aligned}
[J_\chi\Delta,a]_{\rho_\chi}
&=J_\chi\Delta a-\rho_\chi(a)J_\chi\Delta  \\
&=J_\chi\Delta a-(J_\chi\Delta a\Delta^{-1}J_\chi)J_\chi\Delta=0 .
\end{aligned}
\]
Therefore
\[
  [D_\chi,a]_{\rho_\chi}=[B,a]_{\rho_\chi},
\]
which is bounded by hypothesis.

It remains to justify the summability statement.  The bounded perturbation
preserves the operator domain and gives equivalent graph norms:
\[
\begin{aligned}
  \|D_\chi u\|+\|u\|
  &\le \|\Delta u\|+(1+\|B\|)\|u\|,\\
  \|\Delta u\|+\|u\|
  &\le \|D_\chi u\|+(1+\|B\|)\|u\|,
  \qquad u\in\operatorname{Dom}\Delta .
\end{aligned}
\]
Consequently the compact embeddings
\[
  \bigl(\operatorname{Dom}D_\chi,\|\cdot\|_{D_\chi}\bigr)
  \hookrightarrow\Kh_\chi,
  \qquad
  \bigl(\operatorname{Dom}\Delta,\|\cdot\|_{\Delta}\bigr)
  \hookrightarrow\Kh_\chi
\]
have comparable approximation numbers.  Equivalently, the singular values of
\((\one+D_\chi^*D_\chi)^{-1/2}\) are comparable, up to fixed constants, with
those of \((\one+\Delta^2)^{-1/2}\).  The latter have shell multiplicity
\(h(R)\) and size asymptotic to \(e^{-\beta R}\).  Thus
\[
  \operatorname{Tr}\bigl(1+D_\chi^*D_\chi\bigr)^{-s/2}
  <\infty
\]
is governed by
\[
  \sum_R h(R)e^{-s\beta R},
\]
which converges for every \(s>0\) under subexponential shell growth.  If
\(D_\chi\) is invertible, replacing the regularised modulus by
\(|D_\chi|^{-1}\) changes only finitely many low singular values and gives the
last assertion.
\end{proof}

\begin{remark}[Position relative to modular spectral triples]
\label{rem:modular-spectral-triple-lineage}
The exponential weight $\Delta=e^{\beta N}$ is a Gibbs weight, and the modular
automorphism $\sigma_\beta=\operatorname{Ad}\Delta$ makes $D_\chi$ an instance of
a \emph{modular} (twisted, type III) spectral triple in the sense of
Connes--Moscovici \cite{connes-moscovici} and of Carey--Phillips--Rennie and
Carey--Neshveyev--Nest--Rennie
\cite{careyphillipsrennie2010,careyneshveyevnestrennie2011}; its all-orders finite summability under subexponential shell growth lies in an entire/JLO-type
regime \cite{jlo1988}.  The descended renewal shifts generate a
Cuntz--Krieger/graph algebra \cite{cuntz,paskrennie2006}, and signed graph weights
of the kind used here---weights allowed to carry a sign rather than being
positive---already appear in the modular index theory of Carey--Marcolli--Rennie
\cite{careymarcollirennie2011}.  None of this machinery is claimed as new.  What
is new is the \emph{role} assigned to the sign: it is the carrier of Lorentzian
signature, represented by the Krein fundamental symmetry $J_\chi$ and classified
by $[\chi]\in H^1(G,\Z/2)$, rather than a device for computing a modular index.
\end{remark}


% =====================================================================
\subsection{Checkable conditions for bounded twisted commutators}
\label{subsec:checkable-modular}

The previous theorem isolates the analytic condition needed for the signed
modular triple.  We now give a concrete sufficient condition in terms of the
renewal graph.  The following graph-theoretic condition gives an explicit sufficient
criterion for bounded twisted commutators.

\begin{assumption}[Bounded geometry and bounded clock increments]
\label{ass:signed-bounded-geometry}
The recurrent predictive graph has bounded geometry: for each $r$ there is a
uniform bound on the number of vertices in a graph ball of radius $r$.  Moreover
there is $L<\infty$ such that every generator edge $e:x\to y$ satisfies
\[
  |\Lambda(y)-\Lambda(x)|\le L .
\]
\end{assumption}

\begin{definition}[Finite-propagation renewal algebra]
\label{def:finite-propagation}
An operator $a$ on $\Kh_\chi$ has propagation $\le r$ if its matrix coefficient
$\langle e_{(y,\eta')},a e_{(x,\eta)}\rangle$ vanishes whenever the graph
distance between $x$ and $y$ is greater than $r$.  Let
$\mathcal A_{\mathrm{fp}}^\chi$ be the algebra generated by lifted shifts,
their adjoints, and finite-propagation diagonal multipliers for which the
matrix entries are uniformly bounded.  Under bounded geometry these operators
are bounded, and products remain finite-propagation.
\end{definition}

\begin{lemma}[Modular conjugation of finite-propagation operators]
\label{lem:bounded-twisted}
Assume \Cref{ass:signed-bounded-geometry}.  If $a\in\mathcal A_{\mathrm{fp}}^\chi$
has propagation $\le r$, then $\sigma_\beta(a)=\Delta a\Delta^{-1}$ extends to
a bounded operator.  More explicitly, for a lifted generator
$\widetilde S_e$ one has
\[
  \|\sigma_\beta(\widetilde S_e)\|\le e^{\beta L}\|\widetilde S_e\|.
\]
For general propagation-$r$ elements, boundedness follows with a constant
depending only on $r$, $\beta L$, and the bounded-geometry constants of $G$.
Consequently, if $B$ is bounded and $B\in\mathcal A_{\mathrm{fp}}^\chi$, then
\[
  [D_\chi,a]_{\rho_\chi}=[B,a]_{\rho_\chi}
\]
is bounded for every $a\in\mathcal A_{\mathrm{fp}}^\chi$.
\end{lemma}

\begin{proof}
For a lifted generator $e:x\to y$,
\[
  \sigma_\beta(\widetilde S_e)e_{(x,\eta)}
  =e^{\beta(\Lambda(y)-\Lambda(x))}\,
    e_{(y,\eta\chi(e))},
\]
and the multiplier is bounded in modulus by $e^{\beta L}$.  This gives the
stated bound.

If $a$ has propagation $\le r$, then every nonzero matrix coefficient connects
vertices whose lengths differ by at most $rL$.  Thus conjugation by $\Delta$
multiplies each nonzero coefficient by a scalar bounded by $e^{\beta rL}$.
Bounded geometry gives a uniform row and column bound for the propagation-$r$
matrix pattern, so the Schur test yields boundedness of $\sigma_\beta(a)$, with
a constant depending only on the propagation radius and the bounded-geometry
constants.

By \Cref{thm:signed-dirac},
\[
  [D_\chi,a]_{\rho_\chi}=[B,a]_{\rho_\chi}
  =Ba-\rho_\chi(a)B .
\]
The two terms on the right are bounded because $B$, $a$, and $\rho_\chi(a)$ are
bounded.
\end{proof}

\begin{corollary}[Checkable signed modular triple]
\label{cor:checkable-signed-triple}
Under \Cref{ass:signed-bounded-geometry}, if $B\in\mathcal A_{\mathrm{fp}}^\chi$
is bounded and $J_\chi$-self-adjoint, then
\[
  \bigl(\mathcal A_{\mathrm{fp}}^\chi,\Kh_\chi,D_\chi;\rho_\chi\bigr),
  \qquad
  D_\chi=J_\chi e^{\beta N}+B,
\]
is a twisted Krein spectral triple.  If the length-shell multiplicities are
subexponential, it is $s$-summable for every $s>0$.
\end{corollary}

\begin{proof}
This is \Cref{thm:signed-dirac} with the boundedness of the modular conjugates
and twisted commutators supplied by \Cref{lem:bounded-twisted}.
\end{proof}

\begin{remark}[The degenerate unperturbed case]
\label{rem:checkable-honest}
For $B=0$ the matched modular twist makes the twisted commutator of
$J_\chi\Delta$ vanish.  Thus the unperturbed signed modular triple is
analytically valid but differentially degenerate.  A nontrivial signed
differential geometry requires a genuine finite-propagation Dirac-type
perturbation $B$; the role of the exponential clock is to supply modular time
and summability, not by itself to create a nonzero twisted differential.
\end{remark}

\begin{remark}[Holonomy and deficiency]
\label{rem:holonomy-deficiency}
The holonomy class and spectral-radius deficiency are distinct invariants.  The
class $[\chi]$ controls removability of the signed cover.  A deficiency ratio
such as
\[
  r_\chi=\frac{\rho(A_\chi)}{\rho(A)}
\]
measures cancellation in a chosen weighted transfer operator.  Nontrivial
holonomy need not imply a strict spectral-radius drop; periodic or half-period
classes can be cohomologically nontrivial but spectrally silent.  The present
paper uses the holonomy class, the fundamental symmetry, and the signed Dirac
operator, not a strict deficiency hypothesis.
\end{remark}

\subsection{Characterisation of the modular signed Dirac}
\label{sec:cviii}

The preceding construction introduces $D_-=\Jchi e^{\beta N}+B$ and verifies it is
Krein-self-adjoint.  We show the \emph{form} $e^{\beta N}$ is not a choice: it is
the unique positive operator implementing the renewal modular flow on one-reset
operations.  The only residual input is the scalar $\beta$, which we then tie to
a faithful stationary weight.

\subsubsection{The deck-odd diagonal normal form}

Recall $\Jchi e_{(x,\eta)}=\eta\,e_{(x,\eta)}$, $N e_{(x,\eta)}=\Lambda(x)\,
e_{(x,\eta)}$, $\tau e_{(x,\eta)}=e_{(x,-\eta)}$.  The atomic (diagonal) algebra
$\Dd_\chi$ is generated by the rank-one projections onto the basis vectors
$e_{(x,\eta)}$.

\begin{lemma}[Deck-odd diagonal operators are signed magnitudes]
\label{lem:deck-odd-normal-form}
Let $D_0$ be affiliated to $\Dd_\chi$ \textup{(}i.e.\ diagonal in the signed-cover
basis\textup{)} and deck-odd, $\tau D_0\tau^{-1}=-D_0$.  Then there is a
sheet-independent diagonal operator $g$, with $g\,e_{(x,\eta)}=h(x)\,e_{(x,\eta)}$
for some real $h:V\to\R$, such that
\[
  D_0=\Jchi\,g .
\]
If in addition $D_0$ has trivial kernel and $|D_0|$ is its magnitude, then
$h>0$ can be chosen and $g=|D_0|>0$.
\end{lemma}

\begin{proof}
Diagonality gives $D_0 e_{(x,\eta)}=d(x,\eta)e_{(x,\eta)}$ with $d(x,\eta)\in\C$.
Deck-oddness:
\[
  \tau D_0\tau^{-1}e_{(x,\eta)}
  =\tau D_0 e_{(x,-\eta)}
  =d(x,-\eta)\,e_{(x,\eta)}
  \stackrel{!}{=}-d(x,\eta)e_{(x,\eta)},
\]
so $d(x,-\eta)=-d(x,\eta)$.  Writing $h(x):=d(x,+1)$ gives $d(x,-1)=-h(x)$, hence
$d(x,\eta)=\eta\,h(x)$ and
\[
  D_0 e_{(x,\eta)}=\eta\,h(x)\,e_{(x,\eta)}=\Jchi\bigl(h(x)e_{(x,\eta)}\bigr),
\]
which is $\Jchi g$ with $g e_{(x,\eta)}=h(x)e_{(x,\eta)}$, manifestly
sheet-independent.  Self-adjointness of $D_0$ forces $h$ real; trivial kernel
forces $h(x)\neq0$, and $g=\Jchi D_0=|D_0|>0$ after fixing the sign $h>0$.
\end{proof}

So a diagonal, deck-odd unbounded part is \emph{automatically} of the form
$\Jchi g$ with $g$ a sheet-independent positive diagonal magnitude.  This fixes the factor $\Jchi$ from deck oddness.

\subsubsection{The renewal modular flow scales elementary resets}

\begin{lemma}[Graded clock scaling]
\label{lem:clock-scaling}
For every elementary reset edge $e:x\to y$ \textup{(}so $\Lambda(y)=\Lambda(x)+1$\textup{)},
\[
  \Ad(e^{\beta N})(\Stilde_e)
  =e^{\beta N}\,\Stilde_e\,e^{-\beta N}
  =e^{\beta}\,\Stilde_e .
\]
\end{lemma}

\begin{proof}
On $e_{(x,\eta)}$,
\[
  e^{\beta N}\Stilde_e e^{-\beta N}e_{(x,\eta)}
  =e^{-\beta\Lambda(x)}e^{\beta N}\Stilde_e e_{(x,\eta)}
  =e^{-\beta\Lambda(x)}e^{\beta N}e_{(y,\eta\chi(e))}
  =e^{\beta(\Lambda(y)-\Lambda(x))}e_{(y,\eta\chi(e))}
  =e^{\beta}\,\Stilde_e e_{(x,\eta)},
\]
using $\Lambda(y)-\Lambda(x)=1$.  As this holds on every basis vector and
$\Stilde_e$ annihilates vectors off its source fibre, the operator identity
follows.
\end{proof}

Thus the modular automorphism $\sigma_\beta:=\Ad(e^{\beta N})$ scales every
elementary reset by the \emph{single} constant $e^{\beta}$ -- the renewal
incarnation of the Connes--Rovelli thermal-time flow.

\subsubsection{The characterisation theorem}

\begin{axiom}[Renewal modular Dirac]
\label{ax:modular-dirac}
A closed operator $D$ on $\Kh_\chi$ is a \emph{renewal modular Dirac} with
parameter $\beta>0$ if:
\begin{enumerate}[label=\textup{(M\arabic*)},leftmargin=2.8em]
\item \textup{[decomposition]} $D=D_0+B$ with $B$ bounded and $\Jchi$-self-adjoint
      \textup{(}$\Jchi B^*\Jchi=B$\textup{)}, and $D_0$ unbounded and affiliated to the atomic
      algebra $\Dd_\chi$;
\item \textup{[deck-odd]} $\tau D_0\tau^{-1}=-D_0$;
\item \textup{[modular implementation]} the magnitude $g:=|D_0|$ implements the
      renewal modular flow on one-reset operations:
      $\Ad(g)(\Stilde_e)=e^{\beta}\,\Stilde_e$ for every elementary reset edge $e$;
\item \textup{[recurrence]} $G^\chi$ is connected \textup{(}automatic when
      $[\chi]\neq0$; otherwise applied componentwise\textup{)}.
\end{enumerate}
\end{axiom}

\begin{theorem}[The exponential weight is forced]
\label{thm:modular-characterisation}
Every renewal modular Dirac with parameter $\beta$ has the form
\[
  \boxed{\,D=\Jchi\,e^{\beta N}+B\,}
\]
after normalising the magnitude scale to $g\,e_{(x_0,\eta)}=e_{(x_0,\eta)}$ at a
base vertex $x_0$ with $\Lambda(x_0)=0$.  Conversely $\Jchi e^{\beta N}+B$ is a
renewal modular Dirac and is Krein-self-adjoint with respect to $\Jchi$.
\end{theorem}

\begin{proof}
By (M1)--(M2) and \Cref{lem:deck-odd-normal-form}, $D_0=\Jchi g$ with $g>0$
sheet-independent and diagonal.  Consider $m:=g\,e^{-\beta N}$, positive and
diagonal.  For any elementary reset $e$, \Cref{lem:clock-scaling} and (M3) give,
since $g,e^{\beta N}$ commute,
\[
  \Ad(m)(\Stilde_e)
  =\Ad(g)\Ad(e^{-\beta N})(\Stilde_e)
  =\Ad(g)\bigl(e^{-\beta}\Stilde_e\bigr)
  =e^{-\beta}\,e^{\beta}\,\Stilde_e
  =\Stilde_e .
\]
Hence $m$ commutes with every lifted shift.  Writing $m\,e_{(x,\eta)}=
\mu(x,\eta)e_{(x,\eta)}$, the relation $m\Stilde_e=\Stilde_e m$ on $e_{(x,\eta)}$
reads
\[
  \mu(y,\eta\chi(e))=\mu(x,\eta)
  \qquad\text{for every edge }e:x\to y .
\]
By (M4) the signed cover is connected, so $\mu$ is constant: $\mu\equiv c>0$, i.e.
$m=c\,\one$ and $g=c\,e^{\beta N}$.  Normalising at $x_0$ gives $c=1$, so
$g=e^{\beta N}$ and $D_0=\Jchi e^{\beta N}$.

Conversely $\Jchi$ commutes with $N$, hence with $e^{\beta N}$, so $\Jchi
e^{\beta N}$ is Hilbert-self-adjoint and $\Jchi(\Jchi e^{\beta N})^*\Jchi=
\Jchi e^{\beta N}$; with $\Jchi B^*\Jchi=B$ this gives $\Jchi D^*\Jchi=D$.
Deck-oddness holds since $\tau\Jchi\tau^{-1}=-\Jchi$ and $\tau$ commutes with the
sheet-independent $e^{\beta N}$; modular implementation is
\Cref{lem:clock-scaling}.
\end{proof}

\begin{remark}[Equivalent recurrence formulation]
\label{rem:recurrence-route}
If one keeps shell-scalarity of $g$ as an explicit hypothesis,
$g\,e_{(x,\eta)}=f(\Lambda(x))e_{(x,\eta)}$, then (M3) reduces on each edge to the
linear recurrence $f(n+1)=e^{\beta}f(n)$, hence $f(n)=f(0)e^{\beta n}$.  The proof
above is sharper: it \emph{derives} shell-scalarity from (M3) plus connectivity
rather than assuming it.
\end{remark}

\begin{axiom}[Modular homogeneity of elementary resets]
\label{ax:modular-homogeneity}
There is a faithful normal weight \(\varphi\) whose modular automorphism group
\(\sigma_t^\varphi\) leaves the recurrent renewal algebra invariant.  Every
elementary reset shift \(\Stilde_e\) is analytic and is a modular eigenoperator
with one common eigenvalue:
\[
  \sigma_t^\varphi(\Stilde_e)=e^{it\beta}\Stilde_e
  \qquad(t\in\R,\ \text{all elementary }e),
\]
for some \(\beta\in\R\).  On the represented analytic reset algebra, the
commutant of all elementary reset shifts is assumed to be \(\C\one\).
\end{axiom}

\begin{theorem}[Affine modular Hamiltonian]
\label{thm:affine-modular}
Under \Cref{ax:modular-homogeneity}, write
\(\sigma_t^\varphi=\operatorname{Ad}(\Delta_\varphi^{it})\) and
\(\Delta_\varphi=e^H\) in standard form.  Then, on the analytic one-reset
algebra,
\[
  [H,\Stilde_e]=\beta\Stilde_e=\beta[N,\Stilde_e].
\]
Consequently \(H-\beta N\) belongs to the reset-shift commutant and hence
\[
  \boxed{\ H=\beta N+c\one,\qquad
  \Delta_\varphi=e^{\beta N+c\one}.\ }
\]
After normalising the weight at a base vertex, \(c=0\).  Thus the modular
Hamiltonian is affine in the renewal clock, and \(\beta\) is the common modular
eigenvalue of the elementary resets.
\end{theorem}

\begin{proof}
The eigenoperator relation reads
\(\operatorname{Ad}(e^{itH})(\Stilde_e)=e^{it\beta}\Stilde_e\).  Differentiating
at \(t=0\) gives \([H,\Stilde_e]=\beta\Stilde_e\).  Since every elementary reset
raises the grading by one, \([N,\Stilde_e]=\Stilde_e\), and therefore
\([H-\beta N,\Stilde_e]=0\) for every \(e\).  The scalar-commutant hypothesis in
\Cref{ax:modular-homogeneity} gives \(H-\beta N=c\one\).  Exponentiation and base
normalisation give the asserted formula.
\end{proof}

\begin{remark}[Sign convention]
\label{rem:modular-sign-convention}
We define \(\beta\) by
\(\sigma_t^\varphi(\Stilde_e)=e^{it\beta}\Stilde_e\), so the corresponding
modular operator is \(e^{\beta N}\) after normalisation.  A thermodynamic
convention written with \(e^{-N/T}\) describes the inverse modular operator and
therefore has \(\beta=-1/T\) in this convention.  No imaginary-time continuation
is used in the theorem.
\end{remark}

\begin{remark}[Scope of the modular characterisation]
\label{rem:cviii-status}
The functional form $\Jchi e^{\beta N}+B$ is forced by (M1)--(M4), and the
value of $\beta$ is fixed by the common modular eigenvalue of the elementary
reset shifts under \Cref{ax:modular-homogeneity}.  The bounded
$\Jchi$-self-adjoint term $B$ remains free and carries the differential content
of the geometry; $B=0$ is differentially degenerate.  Thus the signed modular
Dirac is intrinsically characterised up to its bounded differential part.
\end{remark}


\subsection{Pressure selection of the modular scale}
\label{subsec:pressure-selection}

The preceding characterisation forces the functional form \(e^{\beta N}\).
The remaining dimensionless exponent is selected by the recurrent weighted
renewal transfer whenever the selected component is supercritical.

\begin{theorem}[Perron pressure selects the modular exponent]
\label{thm:pressure-selects-beta}
Let \(G_0\) be a finite strongly connected recurrent component and let
\[
   K(s)_{xy}
   :=\sum_{e:x\to y} q_e e^{-s\ell_e},
   \qquad q_e>0,\qquad \ell_e\ge\ell_0>0 .
\]
Write \(r(s)=\rho(K(s))\).  Then \(r\) is continuous and strictly decreasing on
\([0,\infty)\), and \(r(s)\to0\) as \(s\to\infty\).  Consequently:
\begin{enumerate}[label=\textup{(\roman*)},leftmargin=2.4em]
\item if \(r(0)>1\), there is a unique \(\beta>0\) satisfying
      \[
         \boxed{\ \rho(K(\beta))=1\ };
      \]
\item if \(r(0)=1\), the unique nonnegative pressure root is \(\beta=0\);
\item if \(r(0)<1\), no nonnegative pressure root exists.
\end{enumerate}
At the root, the Perron eigenvalue is simple and has strictly positive left and
right eigenvectors.  Thus, in the supercritical recurrent phase, the renewal
transfer selects the parameter appearing in
\(\Delta=e^{\beta N}\); conversion of this dimensionless modular scale into
seconds or energy units is a separate metrological calibration.
\end{theorem}

\begin{proof}
For \(s_2>s_1\), every positive matrix entry associated with an edge is
strictly smaller in \(K(s_2)\) than in \(K(s_1)\), and the two matrices have
the same irreducible zero pattern.  Perron--Frobenius strict monotonicity
therefore gives \(r(s_2)<r(s_1)\).  Continuity follows from finite
dimensionality.  Since \(\ell_e\ge\ell_0>0\) and there are finitely many edges,
all entries of \(K(s)\) tend to zero, hence \(r(s)\to0\).  The three alternatives
follow from the intermediate value theorem and strict monotonicity.
Perron--Frobenius also gives simplicity of the spectral radius and strict
positivity of its left and right eigenvectors.
\end{proof}

\subsection{Signed predictive reconstruction}
\label{sec:reconstruction}

The classification theorem counts the possible signature sectors over a fixed
predictive graph.  We next record the reconstruction statement: in a fixed
gauge, the signed spectral data recover the predictive graph, its length
grading, and the sign cocycle.  Modulo flat gauge, the recovered sign invariant
is the cohomology class $[\chi]$.

A fixed-gauge signed predictive geometry is the tuple
\[
  \mathfrak G_\chi=
  (\mathcal A_\chi,\Kh_\chi,N,J_\chi,\tau,\{\widetilde S_e\}_{e\in E_G}),
\]
where $\Kh_\chi=\ell^2(V^\chi)\otimes\Hh_{\rm int}$, $N$ is the lifted length
operator, $J_\chi$ is the sheet fundamental symmetry, $\tau$ is the deck
transformation, and the $\widetilde S_e$ are the lifted renewal shifts.

\begin{lemma}[Atomic-diagonal rigidity]
\label{lem:atom}
Let $\Dd_\chi$ and $\Dd_{\chi'}$ be the atomic diagonal algebras generated by
the rank-one projections onto the basis vectors of $V^\chi$ and
$V^{\chi'}$.  If a unitary $U:\Kh_\chi\to\Kh_{\chi'}$ satisfies
$U\Dd_\chi U^*=\Dd_{\chi'}$, then there is a bijection
\[
  \psi:V^\chi\longrightarrow V^{\chi'}
\]
and phases $\lambda_z\in\T$ such that
\[
  Ue_z=\lambda_z e'_{\psi(z)} .
\]
\end{lemma}

\begin{proof}
The atomic diagonal algebras are maximal abelian with rank-one minimal
projections.  Conjugation by $U$ carries minimal projections to minimal
projections, giving the bijection $\psi$.  Each corresponding one-dimensional
range is then mapped by a phase.
\end{proof}

\begin{theorem}[Signed predictive reconstruction in a fixed gauge]
\label{thm:reconstruction}
Let $\mathfrak G_\chi$ and $\mathfrak G_{\chi'}'$ be fixed-gauge signed
predictive geometries.  Suppose a unitary $U:\Kh_\chi\to\Kh_{\chi'}'$ satisfies
\begin{enumerate}[label=\textup{(R\arabic*)},leftmargin=2.6em]
\item $U\Dd_\chi U^*=\Dd_{\chi'}'$;
\item $U\tau U^*=\tau'$;
\item $U\widetilde S_eU^*=\widetilde S'_{e'}$ under an induced relabelling of
      lifted generator transitions;
\item $UNU^*=N'$;
\item $UJ_\chi U^*=J_{\chi'}'$.
\end{enumerate}
Then $U$ descends to an isomorphism $\varphi:G\to G'$ of predictive graphs
which preserves the length grading:
\[
  \Lambda'(\varphi(x))=\Lambda(x),
\]
intertwines the generated predictive monoids and orders,
\[
  (\M,\preceq,\Lambda)\cong(\M',\preceq',\Lambda'),
\]
and identifies the fixed-gauge cocycles:
\[
  \chi(e)=\chi'(\varphi(e)).
\]
In particular,
\[
  [\chi]=\varphi^*[\chi']\in H^1(G,\Z/2).
\]
\end{theorem}

\begin{proof}
By \Cref{lem:atom}, $U$ is a phase-weighted permutation of the signed-cover
basis.  The condition $U\tau U^*=\tau'$ implies that this permutation maps deck
orbits to deck orbits, hence descends to a bijection
$\varphi:V\to V'$ of base vertices.  Intertwining the lifted shifts gives
\[
  \varphi(y)=\varphi(e\cdot x)=e'\cdot\varphi(x)
\]
whenever $e:x\to y$, so $\varphi$ is an isomorphism of the predictive transition
graphs, with the stated relabelling of generator transitions.  The equality
$UNU^*=N'$ gives $\Lambda'(\varphi(x))=\Lambda(x)$, so the grading is
preserved.  Since the predictive order is divisibility in the generated
transition monoid, the graph isomorphism gives an isomorphism of graded ordered
predictive monoids.

It remains to identify the sign.  Let $e:x\to y$ and lift it from
$(x,\eta)$ to $(y,\eta\chi(e))$.  The equality $UJ_\chi U^*=J_{\chi'}'$ fixes
the sheet label in the chosen gauge, while the intertwining of lifted shifts
sends this lifted edge to the lifted edge over $\varphi(e)$.  Therefore the
endpoint sheet must be
\[
  \eta\chi(e)=\eta\chi'(\varphi(e)),
\]
for both choices of $\eta$, and hence $\chi(e)=\chi'(\varphi(e))$.  Taking loop
products gives $[\chi]=\varphi^*[\chi']$.
\end{proof}

\begin{corollary}[Gauge-invariant complete invariant]
\label{cor:iff}
Two signed predictive geometries are equivalent up to unitary equivalence and
flat gauge if and only if there is a predictive graph isomorphism
$\varphi:G\to G'$ preserving $\Lambda$ and satisfying
\[
  [\chi]=\varphi^*[\chi']\in H^1(G,\Z/2).
\]
Equivalently, the pair
\[
  \Bigl((\M,\preceq,\Lambda)\ \text{up to graded-ordered-monoid isomorphism},
  \ [\chi]\in H^1(G,\Z/2)\Bigr)
\]
is a complete invariant of signed predictive geometries modulo flat gauge.
\end{corollary}

\begin{proof}
The forward implication follows from \Cref{thm:reconstruction}, with the
observation that composing with a flat gauge changes $\chi$ by a coboundary and
therefore leaves $[\chi]$ unchanged.  Conversely, if
$[\chi]=\varphi^*[\chi']$, choose a gauge representative of $\chi'$ for which
the equality holds as cochains.  Then the map
\[
  e_{(x,\eta)}\longmapsto e'_{(\varphi(x),\eta)}
\]
intertwines the diagonal algebra, deck transformation, lifted shifts, length
operator, and sheet fundamental symmetry in that gauge.
\end{proof}

\begin{remark}[Position relative to Connes reconstruction]
\label{rem:vs-connes}
This is a discrete reconstruction theorem.  Connes' reconstruction theorem
recovers a Riemannian spin manifold from a commutative spectral triple via
orientability, regularity, finiteness, and Poincar\'e-duality hypotheses
\cite{connesreconstruction}.  Here the recovered object is not a manifold but a
graded ordered predictive monoid decorated by a $\Z/2$ holonomy class.  The
mechanism is correspondingly simpler: the atomic diagonal recovers the
predictive points, the lifted shifts recover the transition monoid, the length
operator recovers the grading, and the signed cover recovers the holonomy.
\end{remark}

\begin{remark}[Transition to the continuum endpoint]
\label{rem:signed-transition}
The result of this section is finite and operator-theoretic.  Starting from the
positive predictive geometry of \Cref{sec:object}, a sign cocycle produces a
principal double cover, a Krein fundamental symmetry, and a signed modular
Dirac operator.  The non-removable signature sectors are exactly the classes in
$H^1(G,\Z/2)$, and the signed spectral data reconstruct the predictive monoid
together with this holonomy.  The next section derives the projective Clifford factor from reversible
predictive dynamics before taking any continuum limit.
\end{remark}





% =====================================================================
\section{Reversible revision dynamics and dimensional minimality}
\label{sec:primitive-dimension-selection}

The signed cover supplies a flat holonomy class and a distinguished temporal
exchange relation, but flat graph topology alone cannot generate a
noncommutative Clifford algebra.  The missing spatial commutator form is derived
from reversible predictive dynamics.  Predictive units are automorphisms; on a
finite recurrent factor they are inner, so their implementers compose
projectively.  Factoring the radical of the resulting alternating form gives a
canonical symplectic quotient and its Clifford matrix factor.  Only after this
finite operator-algebraic derivation do we address dimensional minimality.

\subsection{Flat holonomy does not produce the Clifford twist}
\label{subsec:flat-holonomy-no-go}

We begin by separating two roles that are easy to conflate.  The signed cover
of \Cref{sec:signed} provides holonomy classes, hence fundamental symmetries
modulo flat gauge.  A Clifford algebra, however, requires noncommuting
generators, or equivalently a nontrivial commutator bicharacter.  Flat graph
holonomy by itself has only \(1\)-dimensional topology and therefore cannot
supply this twist.

Let \(G\) be a recurrent predictive component and write
\(H=H^1(G,\Z/2)\) for its group of signed holonomy sectors.  The classification
result of \Cref{thm:classification} shows that \(H\) classifies compatible
fundamental symmetries modulo flat gauge.  It is tempting to identify \(H\)
directly with the spatial Clifford grading group.  The following observation
explains why that identification cannot be a consequence of graph topology
alone.

\begin{proposition}[Flat-holonomy obstruction]
\label{prop:flat-holonomy-obstruction}
The signed predictive cover is a flat \(\Z/2\)-cover.  Since a graph is a
\(1\)-complex, \(H^2(G,\Z/2)=0\).  Consequently any twist obtained only from the
topology of the cover, for example by a cup product or a commutator \(2\)-cell,
is cohomologically trivial.  The topological holonomy sector algebra is
therefore the commutative group algebra \(\C[H]\), not a Clifford algebra.
\end{proposition}

\begin{proof}
A connected graph is homotopy equivalent to a wedge of circles, hence has no
cohomology in degree \(2\).  The signed cover is classified by a class in
\(H^1(G,\Z/2)\), so it is flat: it has holonomy around loops but no curvature
\(2\)-class.  A Clifford exchange sign is a projective commutator bicharacter, equivalently a
twist of the \emph{group algebra of the holonomy group}.  Such a twist may exist
after extra dynamical input, but it is not produced by the topology of the graph
cover itself: the graph has no \(2\)-cells on which a curvature class could live.
In the flat topological sector the fundamental symmetries multiply as
\(J_\alpha J_\beta=J_{\alpha+\beta}=J_\beta J_\alpha\), so their algebra is
commutative.  The spatial--spatial noncommutative Clifford twist must therefore be supplied by
modular/internal dynamics rather than by flat topology.  The temporal row is an
exception: \Cref{thm:canonical-temporal-row} shows that it is forced by the
Krein exchange relation of the signed cover.
\end{proof}

Thus topology supplies the signed sectors, but not their Clifford
anticommutation.  The missing sign must be dynamical: it must come from how
revision operations are ordered by the modular renewal clock.


\subsection{The temporal row is forced by the signed cover}
\label{subsec:temporal-row-forced}

The flat-cover obstruction does not mean that every entry of the modular
commutator form is additional.  The temporal row is already fixed by the signed
cover.

\begin{definition}[Realized reversible sign subgroup]
\label{def:realized-reversible-subgroup}
Fix the recurrent algebra on which revisions are to act.  A class
\(\alpha\in H=H^1(G,\Ztwo)\) is \emph{realized reversibly} if it is represented
by a predictive unit preserving that recurrent algebra and admits a
clock-preserving signed lift whose sheet action is the geometric sheet shift
associated with \(\alpha\).  The set of such classes is denoted
\[
  K_{\mathrm{rev}}\subseteq H .
\]
\end{definition}

\begin{lemma}[Realized reversible classes form a subgroup]
\label{lem:realized-reversible-subgroup}
One has
\[
  K_{\mathrm{rev}}\le H .
\]
\end{lemma}

\begin{proof}
The identity channel realizes the zero class.  If revisions realizing
\(\alpha\) and \(\beta\) are composed, their predictive units compose and their
sheet shifts add, so the composite realizes \(\alpha+\beta\).  The inverse of
a predictive unit is again a predictive unit and reverses the same
elementary-two sheet class.  Hence the realized classes are closed under
identity, addition, and inversion.
\end{proof}

Let $\widetilde S_e$ denote the lifted shift associated with an oriented
edge $e:x\to y$ of the recurrent predictive graph.  With
\[
 J_\chi e_{(x,\eta)}=\eta e_{(x,\eta)},\qquad
 \widetilde S_e e_{(x,\eta)}=e_{(y,\eta\chi(e))},
\]
the defining exchange relation is
\begin{equation}
 \label{eq:krein-edge-exchange}
 J_\chi\widetilde S_e=\chi(e)\widetilde S_eJ_\chi.
\end{equation}
For a recurrent cycle representing $\alpha\in H=H^1(G,\Ztwo)$, write
$\widetilde S_\alpha$ for the ordered product of its lifted edge shifts and
$\chi(\alpha)=(-1)^{\langle[\chi],\alpha\rangle}$ for its sign holonomy.

\begin{lemma}[Krein exchange along a recurrent cycle]
\label{lem:cycle-krein-exchange}
For every recurrent cycle class $\alpha$,
\[
 J_\chi\widetilde S_\alpha
 =(-1)^{\langle[\chi],\alpha\rangle}
  \widetilde S_\alpha J_\chi.
\]
The sign depends only on the cohomology--homology pairing and not on the chosen
cycle representative.
\end{lemma}

\begin{proof}
Move $J_\chi$ through the ordered edge product using
\eqref{eq:krein-edge-exchange}.  The accumulated sign is the product of the
edge signs, namely $\chi(\alpha)$.  Representative independence is exactly the
invariance of the pairing of the cocycle $[\chi]$ with the cycle class.
\end{proof}

\begin{theorem}[Canonical temporal row of the modular form]
\label{thm:canonical-temporal-row}
Let $R_t=J_\chi$ and let $R_\alpha$ be any clock-preserving revision operator
whose action on the signed sheet is the geometric sheet shift associated with
\(\alpha\in K_{\mathrm{rev}}\).  Then
\[
 \boxed{\ \omega(\alpha,t)=\langle[\chi],\alpha\rangle
 \qquad(\alpha\in K_{\mathrm{rev}}).\ }
\]
Equivalently,
\[
 J_\chi R_\alpha
 =(-1)^{\langle[\chi],\alpha\rangle}R_\alpha J_\chi.
\]
Thus the temporal row of $\omega$ is forced by the signed cover and is not a
freely chosen part of the modular revision datum.
\end{theorem}

\begin{proof}
The internal and base-preserving factors of $R_\alpha$ commute with the diagonal
sheet operator $J_\chi$, while its sheet-shift factor obeys
\Cref{lem:cycle-krein-exchange}.  Hence the commutator of $R_\alpha$ with
$R_t=J_\chi$ is the scalar
$(-1)^{\langle[\chi],\alpha\rangle}$.  Comparing with the definition of the
modular commutator form gives the identity.
\end{proof}

\begin{corollary}[Canonical hyperbolic core]
\label{cor:canonical-hyperbolic-core}
If the restriction of \([\chi]\) to \(K_{\mathrm{rev}}\) is nonzero, there is
\(\alpha_0\in K_{\mathrm{rev}}\) with \(\omega(\alpha_0,t)=1\).  The restriction
of \(\omega\) to \(\langle t,\alpha_0\rangle\) is the standard rank-two
alternating form, and the
noncentral revision algebra contains a canonical factor
\[
 \Cl_2(\C)\cong M_2(\C).
\]
This rank-two temporal Clifford core is supplied by the signed cover before any
spatial--spatial commutator is specified.
\end{corollary}

\begin{proof}
Nontriviality of the restriction means that the functional
\(\alpha\mapsto\langle[\chi],\alpha\rangle\) is nonzero on
\(K_{\mathrm{rev}}\), so some \(\alpha_0\in K_{\mathrm{rev}}\) pairs
nontrivially with \(t\).  The resulting two-dimensional alternating form is
nondegenerate.  Its twisted group algebra has noncentral factor
$\Cl_2(\C)\cong M_2(\C)$.
\end{proof}

\begin{proposition}[Flat sheet transport has trivial spatial block]
\label{prop:flat-spatial-block-zero}
For the undressed flat realization \(R_\alpha=\widetilde S_\alpha\) on the
abelian signed cover, restricted to realized reversible classes and with
reorderable recurrent base transport,
\[
 \omega(\alpha,\beta)=0
 \qquad(\alpha,\beta\in K_{\mathrm{rev}}).
\]
Consequently the signed cover alone fixes the temporal row but produces no
spatial--spatial Clifford twist.
\end{proposition}

\begin{proof}
The sheet contribution to the group commutator is
$\chi(\alpha)\chi(\beta)\chi(\alpha)^{-1}\chi(\beta)^{-1}=1$, and the
reorderable base transport contributes no scalar phase.  Hence the spatial
commutator is trivial.
\end{proof}

\begin{remark}[The remaining cocycle content]
\label{rem:only-spatial-block-remains}
After \Cref{thm:canonical-temporal-row}, the only portion of the modular form
not fixed by the signed cover is
\[
 \omega\big|_{K_{\mathrm{rev}}\times K_{\mathrm{rev}}}.
\]
The cover supplies the temporal Clifford pair.  Ranks beyond the canonical
rank-two core require noncommutative internal transport.  Thus only the spatial--spatial block remains to be generated dynamically.
\end{remark}

\subsection{The modular revision algebra from reversible predictive dynamics}
\label{subsec:modular-revision-algebra}

The signed cover supplies the topologically available spatial label module
\[
  H=H^1(G,\Z/2),
\]
while the dynamics realizes only the reversible subgroup
\(K_{\mathrm{rev}}\le H\) of
\Cref{def:realized-reversible-subgroup}.  The Krein grading supplies the
temporal label \(t\).  Put
\[
  \Hf=K_{\mathrm{rev}}\oplus\langle t\rangle .
\]
The full-access phase is the special case \(K_{\mathrm{rev}}=H\).  The
projective revision law will not be assumed.  It follows from reversible
predictive channels: units of the predictive channel monoid are automorphisms by
\Cref{thm:predictive-unit}; on a finite irreducible matrix block those
automorphisms are inner; equality of the labelled automorphism products then
forces their unitary implementers to compose up to a scalar cocycle.

\begin{assumption}[Signed recurrent revision phase]
\label{ass:signed-revision-phase}
\label{def:intrinsic-spatial-revisions}
The recurrent signed phase has the following properties.
\begin{enumerate}[label=\textup{(R\arabic*)},leftmargin=2.6em]
\item every \(\alpha\in K_{\mathrm{rev}}\) is represented by a predictive unit
      \(\Phi_\alpha\in\mathsf U_{\rm pred}\) preserving the fixed recurrent
      algebra;
\item the chosen labels are compatible with the subgroup law, so labelled
      concatenation is additive,
      \[
        \Phi_\alpha\Phi_\beta=\Phi_{\alpha+\beta};
      \]
\item the revisions preserve the minimal sufficient recurrent algebra;
\item after quotienting revision labels that act trivially, the preserved
      algebra is a finite matrix factor \(\A_{\rm rev}\), and the induced action
      is faithful and irreducible.
\end{enumerate}
The temporal revision is represented separately by
\(R_t=J_\chi\), with the canonical exchange relation of
\Cref{thm:canonical-temporal-row}.
\end{assumption}

\begin{theorem}[Reversible predictive revision theorem]
\label{thm:reversible-predictive-revision}
Under \Cref{ass:signed-revision-phase}, every spatial revision
\(\alpha\in K_{\mathrm{rev}}\) acts on \(\A_{\rm rev}\) by an involutive
\(*\)-automorphism,
\[
  \Phi_\alpha\in\operatorname{Aut}(\A_{\rm rev}),
  \qquad
  \Phi_\alpha^2=\id,
  \qquad
  \Phi_\alpha\Phi_\beta=\Phi_{\alpha+\beta}.
\]
The action is intrinsic to the predictive channel quotient.
\end{theorem}

\begin{proof}
Each \(\Phi_\alpha\) is a unit of the predictive channel monoid, hence a
\(*\)-automorphism by \Cref{thm:predictive-unit}.  The elementary-two label law
gives \(\Phi_\alpha^2=\id\), and labelled concatenation gives the additive
composition rule.  Preservation of the sufficient algebra and passage to the
faithful finite factor give the stated restricted action.
\end{proof}

\begin{theorem}[Automatic projective lift]
\label{thm:spatial-multiplier-scalar-main}
Choose unitary implementers on the matrix factor,
\[
  \Phi_\alpha=\operatorname{Ad}(R_\alpha),
  \qquad R_0=\one.
\]
Then there is a normalised scalar two-cocycle
\(\sigma_{\rm s}\in Z^2(K_{\mathrm{rev}},U(1))\) such that
\[
  \boxed{\ R_\alpha R_\beta
  =\sigma_{\rm s}(\alpha,\beta)R_{\alpha+\beta}.\ }
\]
Changing the implementing section by phases changes \(\sigma_{\rm s}\) by a
coboundary.  Its antisymmetric ratio is an intrinsic alternating form
\[
  \sigma_{\rm s}(\alpha,\beta)
  \sigma_{\rm s}(\beta,\alpha)^{-1}
  =(-1)^{\omega_{\rm s}(\alpha,\beta)},
  \qquad
  \omega_{\rm s}\in\Alt^2(K_{\mathrm{rev}},\F_2).
\]
No preferred Kraus section or preferred unitary section is required.
\end{theorem}

\begin{proof}
Every automorphism of a complex matrix factor is inner.  The unitaries
\(R_\alpha R_\beta\) and \(R_{\alpha+\beta}\) implement the same automorphism,
so \(R_\alpha R_\beta R_{\alpha+\beta}^{-1}\) lies in the centre of
\(\A_{\rm rev}\) and is a scalar phase.  Associativity is exactly the scalar
cocycle identity.  A scalar change of section produces the usual coboundary.
For a scalar cocycle on an abelian group the antisymmetric ratio is a
bicharacter; since \(2\alpha=0\), it takes values in \(\{\pm1\}\) and therefore
defines the displayed \(\F_2\)-valued alternating form.
\end{proof}

\begin{proposition}[Canonical extension to the full revision group]
\label{prop:full-revision-cocycle-main}
Let \(R_t=J_\chi\) and write
\[
  J_\chi R_\alpha
  =(-1)^{\langle[\chi],\alpha\rangle}R_\alpha J_\chi
  \qquad(\alpha\in K_{\mathrm{rev}})
\]
as in \Cref{thm:canonical-temporal-row}.  Define
\[
  R_{\alpha+\epsilon t}:=R_\alpha J_\chi^\epsilon,
  \qquad \alpha\in K_{\mathrm{rev}},\quad\epsilon\in\F_2.
\]
Then
\[
  R_{\alpha+\epsilon t}R_{\beta+\delta t}
  =\sigma\bigl((\alpha,\epsilon),(\beta,\delta)\bigr)
   R_{\alpha+\beta+(\epsilon+\delta)t},
\]
where
\begin{equation}
\label{eq:full-revision-cocycle}
  \sigma\bigl((\alpha,\epsilon),(\beta,\delta)\bigr)
  =\sigma_{\rm s}(\alpha,\beta)
   (-1)^{\epsilon\langle[\chi],\beta\rangle}.
\end{equation}
This is a normalised scalar cocycle on \(\Hf\), with alternating form
\begin{equation}
\label{eq:full-revision-form}
  \omega\bigl((\alpha,\epsilon),(\beta,\delta)\bigr)
  =\omega_{\rm s}(\alpha,\beta)
   +\epsilon\langle[\chi],\beta\rangle
   +\delta\langle[\chi],\alpha\rangle.
\end{equation}
Thus the spatial block and the canonical temporal row assemble intrinsically.
\end{proposition}

\begin{proof}
Move \(J_\chi^\epsilon\) through \(R_\beta\) using the exchange relation and
then apply \Cref{thm:spatial-multiplier-scalar-main}.  Associativity proves the
cocycle identity, and taking the antisymmetric ratio gives
\eqref{eq:full-revision-form}.
\end{proof}

\begin{definition}[Modular revision datum]
\label{def:revision-operators-main}
The modular revision datum is the derived scalar projective family
\(\{R_\gamma:\gamma\in\Hf\}\) with multiplier
\(\sigma\in Z^2(\Hf,U(1))\).  Its finite-dimensional revision algebra is
\[
  \Arev:=\operatorname{span}_{\C}\{R_\gamma:\gamma\in\Hf\}
  \cong\C_\sigma[\Hf].
\]
Its gauge-invariant commutator bicharacter is
\[
  (-1)^{\omega(\alpha,\beta)}
  =\sigma(\alpha,\beta)\sigma(\beta,\alpha)^{-1}.
\]
\end{definition}

\begin{remark}[The datum is derived, not imposed]
\label{rem:revision-datum-assumption}
On the signed recurrent revision phase, neither the scalar multiplier, the
projective law, nor bilinearity of the commutator is assumed.  They follow from
predictive reversibility, the finite-factor reduction, innerness, and
associativity.  Geometric verticality and exact recurrent covariance may be
used as sufficient criteria for identifying the same finite factor, but are not
additional hypotheses in the main projective-lift theorem.
\end{remark}

\begin{definition}[Modular commutator form]
\label{def:modular-form-main}
The modular commutator form is the alternating form \(\omega\) characterised by
\[
  R_\alpha R_\beta R_\alpha^{-1}R_\beta^{-1}
  =(-1)^{\omega(\alpha,\beta)}\one.
\]
For a length-shell projection \(P_N\), the same identity holds after compression
to every nonempty shell.
\end{definition}

\begin{theorem}[Exact-bicharacter stabilisation]
\label{thm:exact-bicharacter-stabilisation}
For the derived modular revision datum, the compressed commutator form is
shell-independent:
\[
  \omega_N(\alpha,\beta)=\omega(\alpha,\beta)
\]
on every nonempty shell.
\end{theorem}

\begin{proof}
The group commutator is the scalar operator
\((-1)^{\omega(\alpha,\beta)}\one\) before compression, so compression to any
nonzero shell leaves the same scalar.
\end{proof}

\subsection{Quadratic refinement and the finite central extension}
\label{subsec:revision-central-extension}

The scalar cocycle determines more than the alternating form.  Its lift squares
form a quadratic refinement, and the corresponding projective representation
has a canonical finite extraspecial realisation in the primitive case.

\begin{theorem}[Quadratic refinement of the modular form]
\label{thm:revision-quadratic-refinement-main}
After scalar rephasing, the cocycle $\sigma$ may be chosen $\{\pm1\}$-valued:
\[
  \sigma(\alpha,\beta)=(-1)^{B(\alpha,\beta)},
\]
where $B$ is bilinear and $B-B^{\mathsf T}=\omega$.  The lift squares are
\[
  R_\gamma^2=(-1)^{q(\gamma)}\one,
  \qquad q(\gamma):=B(\gamma,\gamma),
\]
and satisfy
\[
  q(\alpha+\beta)
  =q(\alpha)+q(\beta)+\omega(\alpha,\beta).
\]
Thus $q$ is a quadratic refinement of $\omega$.  Basis lifts may be chosen
involutive, but when $\omega\ne0$ not every lift can square to $\one$ while the
cocycle remains $\{\pm1\}$-valued.
\end{theorem}

\begin{proof}
For an elementary abelian $2$-group,
$H^2(\Hf,U(1))$ is $2$-torsion and every class admits a $\{\pm1\}$-valued
bilinear representative.  The square and polarisation identities follow from
$2\gamma=0$ and bilinearity of $B$.  If $q$ vanished identically, its polar form
$\omega$ would vanish.
\end{proof}

\begin{theorem}[Central extension and projective kernel]
\label{thm:revision-central-extension-main}
Choose projectively equivalent representatives $\widetilde R_\gamma$ with the
$\{\pm1\}$-valued cocycle of
\Cref{thm:revision-quadratic-refinement-main}, and let
\[
  \Grp:=\langle\widetilde R_\gamma:\gamma\in\Hf\rangle,
  \qquad
  A:=\Grp\cap U(1)\one.
\]
The canonical map
\[
  \pi_{\rm proj}:\Hf\longrightarrow\Grp/A,
  \qquad \gamma\longmapsto[\widetilde R_\gamma],
\]
is surjective, with kernel
\[
  \Kproj:=\{\gamma:\widetilde R_\gamma\in A\}
  \subseteq\rad(\omega).
\]
Hence
\[
  \Grp/A\cong\Hf/\Kproj,
  \qquad A\le\{\pm1\}.
\]
If $\omega$ is nondegenerate, then $\Kproj=0$, $A=\{\pm1\}$, and $\Grp$ is an
extraspecial $2$-group of order
\[
  |\Grp|=2^{1+\dim_{\F_2}\Hf}.
\]
The projectively equivalent section may be chosen with
$\widetilde R_t=R_t=J$.
\end{theorem}

\begin{proof}
The quotient map is well defined because products of the rephased
representatives differ from the representative of the sum by a sign.  Its
kernel is exactly $\Kproj$, giving the quotient by the first isomorphism theorem.
A projectively scalar lift commutes with every other lift, hence lies in
$\rad(\omega)$.  If $\omega$ is nondegenerate, the projective kernel is trivial
and some commutator equals $-\one$, so $A=\{\pm1\}$.  Every group element has a
unique normal form up to the central sign, giving the stated order and the
extraspecial property.  The rephasing cochain is determined up to a character;
using that freedom one normalises its value at $t$ to one, preserving $R_t=J$.
\end{proof}

\begin{proposition}[Projective revision law]
\label{prop:projective-revision}
The commutator bicharacter
\[
  \chi_{\rm rev}(\alpha,\beta)
  :=\sigma(\alpha,\beta)\sigma(\beta,\alpha)^{-1}
  =(-1)^{\omega(\alpha,\beta)}
\]
is alternating, bilinear over $\F_2$, and invariant under scalar rephasing of
the projective section.  The quadratic refinement changes by an affine linear
term under rephasing, while its polar form $\omega$ is gauge invariant.
\end{proposition}

\begin{proof}
The bicharacter identities follow from the scalar cocycle identity on the
abelian group $\Hf$.  A coboundary is symmetric, so it cancels in the
antisymmetric ratio.  Rephasing changes the lift squares by a character, hence
changes $q$ by a linear function without changing its polar form.
\end{proof}

\begin{corollary}[Amplitude square-root lift of the revision squares]
\label{cor:revision-z4-square-root}
Let \(q:\Hf\to\F_2\) be the lift-square function of
\Cref{thm:revision-quadratic-refinement-main}.  A choice of amplitude square
roots is equivalently a lift \(\widetilde q:\Hf\to\Zfour\) with
\(\widetilde q(\gamma)\bmod2=q(\gamma)\), assigning the phase
\(i^{\widetilde q(\gamma)}\) to the lift of \(R_\gamma\).  If
\(q(\gamma)=1\) for some \(\gamma\), the corresponding amplitude phase has order
four.  Hence the nontrivial quadratic-square sector has minimal finite amplitude
coefficient group \(\Zfour\).  The polar commutator form \(\omega\) is unchanged
by the choice of \(\widetilde q\).
\end{corollary}

\begin{proof}
The identity \(R_\gamma^2=(-1)^{q(\gamma)}\one\) says that the real square of the
amplitude lift is determined by \(q(\gamma)\).  Choosing
\(\widetilde q(\gamma)\in\Zfour\) with parity \(q(\gamma)\) is exactly choosing a
phase whose square is \((-1)^{q(\gamma)}\).  If \(q(\gamma)=1\), this phase
squares to \(-1\), so by \Cref{thm:minimal-z4-amplitude-lift} it generates a
cyclic order-four subgroup and no smaller finite phase group can carry it.  The
commutator form \(\omega\) is the antisymmetric ratio of products in the real
signed quotient; multiplying individual lifts by square-root phases changes lift
squares but cancels from the antisymmetric ratio except through the already fixed
\(\{\pm1\}\)-valued polar form.
\end{proof}



\subsection{Canonical symplectic reduction of the derived revision algebra}
\label{sec:cx}

The alternating form is already determined by the automatic projective lift.
Primitivity is therefore not an assumption on a postulated form: it is the
canonical passage from the derived twisted group algebra to its noncentral
symplectic quotient.

\begin{lemma}[The real polar commutator is \(\Ztwo\)-valued]
\label{lem:commutator-pm1}
For the derived scalar projective family,
\[
  R_aR_bR_a^{-1}R_b^{-1}
  =\sigma(a,b)\sigma(b,a)^{-1}\one
  =(-1)^{\omega(a,b)}\one,
\]
where \(\omega\in\Alt^2(\Hf,\F_2)\).  The form is independent of the
implementing section.
\end{lemma}

\begin{proof}
The cocycle ratio is a bicharacter on the abelian group \(\Hf\).  Since
\(2a=0\), its square is one, so its values lie in \(\{\pm1\}\).  Coboundaries
are symmetric and cancel from the ratio.
\end{proof}

\begin{definition}[Primitive revision sector]
\label{def:primitive-reduction}
Set
\[
  \rad\omega:=\{a\in\Hf:\omega(a,b)=0\ \text{for all }b\in\Hf\},
  \qquad
  \boxed{\ \Hprim:=\Hf/\rad\omega.\ }
\]
The induced form is denoted \(\bar\omega\), and the reduced twisted group
algebra by \(\Arev^{\rm prim}\).
\end{definition}

\begin{theorem}[Radical--centre dictionary]
\label{thm:radical-centre}
The centre of the derived revision algebra is
\[
  Z(\Arev)
  =\operatorname{span}_{\C}\{R_r:r\in\rad\omega\}.
\]
It is a commutative unital finite-dimensional \(C^*\)-algebra and hence
\[
  Z(\Arev)\cong\C^{\,2^{\dim_{\F_2}\rad\omega}}.
\]
In particular, \(\Arev\) is a factor if and only if \(\rad\omega=0\).
\end{theorem}

\begin{proof}
A homogeneous generator \(R_r\) is central exactly when its commutator with
every \(R_b\) is trivial, equivalently when \(r\in\rad\omega\).  Linear
independence of the homogeneous basis gives the displayed span.  It is closed
under multiplication and adjoint, and its homogeneous generators commute
because \(\omega\) vanishes on the radical.  Its complex dimension is
\(|\rad\omega|=2^{\dim\rad\omega}\), so the finite-dimensional commutative
\(C^*\)-algebra has the stated form.  It is one-dimensional exactly when the
radical is zero.
\end{proof}

\begin{theorem}[Block decomposition and factor quotients]
\label{thm:factor-quotients-corrected}
Let
\[
  2m:=\dim_{\F_2}(\Hf/\rad\omega).
\]
Then
\[
  \boxed{\ \Arev
  \cong\bigoplus_{\psi\in\widehat{\rad\omega}}M_{2^m}(\C).\ }
\]
For each character \(\psi\in\widehat{\rad\omega}\), the ideal
\[
  I_\psi
  :=\bigl\langle R_r-\psi(r)\one:r\in\rad\omega\bigr\rangle
\]
gives a factor quotient
\[
  \Arev/I_\psi\cong M_{2^m}(\C)\cong\Cl_{2m}(\C).
\]
The quotients are pairwise isomorphic but their quotient maps are not canonical.
The canonical invariant is the symplectic quotient
\((\Hprim,\bar\omega)\) and the common matrix/Morita type; choosing a specific
factor representation additionally requires a central character \(\psi\).
\end{theorem}

\begin{proof}
By \Cref{thm:radical-centre}, the centre decomposes into minimal central
projections indexed by \(\widehat{\rad\omega}\).  On each central block the
radical acts by the corresponding character, while the induced form on
\(\Hprim\) is nondegenerate.  A symplectic basis identifies the resulting
twisted group algebra with the tensor product of \(m\) rank-two Weyl pairs,
namely \(M_{2^m}(\C)\).  Summing the central blocks gives the decomposition.
Different characters select different central projections, so the abstract
factor type is canonical but a quotient map is not.
\end{proof}

\begin{theorem}[Canonical primitive Clifford factor]
\label{thm:primitivity-canonical}
The induced form \(\bar\omega\) on \(\Hprim\) is nondegenerate and therefore has
even rank \(2m\).  Every factor representation of \(\Arev\) has the canonical
matrix type
\[
  \boxed{\ \Arev^{\rm prim}\cong M_{2^m}(\C)\cong\Cl_{2m}(\C).\ }
\]
It carries the unique faithful irreducible projective representation of the
nondegenerate quotient, of dimension \(2^m\), up to unitary equivalence.
\end{theorem}

\begin{proof}
Well-definedness of \(\bar\omega\) follows from quotienting by the radical.  If
\([a]\) lies in the radical of \(\bar\omega\), then \(a\in\rad\omega\), so
\([a]=0\); thus \(\bar\omega\) is nondegenerate.  A nondegenerate alternating
form has even rank.  Choosing a symplectic basis gives \(m\) finite Weyl pairs,
and finite Stone--von Neumann gives the unique irreducible projective
representation of dimension \(2^m\), with full matrix algebra
\(M_{2^m}(\C)\cong\Cl_{2m}(\C)\).
\end{proof}

\begin{theorem}[Internal revision twirl]
\label{thm:revision-twirl-main}
If the primitive quotient has rank \(2m\) and acts irreducibly on
\(M_{2^m}(\C)\) through projective implementers \(R_\alpha\), then
\[
  \frac1{|\Hprim|}\sum_{\alpha\in\Hprim}R_\alpha X R_\alpha^*
  =\frac{\Tr X}{2^m}\one.
\]
Consequently the internal fixed algebra is \(\C\one\).
\end{theorem}

\begin{proof}
The range of the twirl is the commutant of the irreducible projective
representation, hence is scalar by Schur's lemma.  Trace preservation fixes the
scalar.
\end{proof}

\begin{theorem}[Sheet ergodicity of the maximal recurrent cover]
\label{thm:sheet-ergodicity-main}
For the maximal elementary-two signed cover with tautological surjective
holonomy, recurrent first-return loops act transitively on each sheet fibre.
The only recurrently invariant sheet functions are constants.
\end{theorem}

\begin{proof}
First-return loops generate the recurrent fundamental group, and surjectivity of
the tautological holonomy makes their deck transformations generate the full
sheet group, which acts transitively on its torsor.
\end{proof}

\begin{corollary}[Double ergodicity is derived in the primitive phase]
\label{cor:double-ergodicity-derived}
After passage to an irreducible primitive factor and the maximal recurrent
signed cover, internal and sheet ergodicity follow from the construction.  They
need not be imposed as independent scalar-centralisation assumptions.
\end{corollary}


\begin{corollary}[Primitive data are not an input]
\label{cor:primitive-not-input}
The primitive modular revision sector is the canonical symplectic reduction of
the projective action derived in \Cref{thm:spatial-multiplier-scalar-main,prop:full-revision-cocycle-main}.  ``Primitive'' means passage to the noncentral
factor type, not an assumed nondegeneracy of a separately postulated form.  Its
Clifford rank is determined by \(\dim_{\F_2}\Hprim\), independently of the
central character used to realise one factor block.
\end{corollary}

\begin{theorem}[Structure and rigidity of the revision sector]
\label{thm:revision-structure-main}
Under \Cref{ass:signed-revision-phase}, the spatial projective multiplier and
its alternating form are automatic, the temporal grading extends them
canonically to \(\Hf\), and the commutator form is shell-independent.  After
canonical symplectic reduction, if
\(\dim_{\F_2}\Hprim=1+\dCl\), then
\[
  \Arev^{\rm prim}\cong\Cl_{1+\dCl}(\C)
  \cong M_{2^{(1+\dCl)/2}}(\C),
  \qquad
  \dCl=\dim_{\F_2}\Hprim-1\ \text{is odd}.
\]
The irreducible primitive revision representation is unique up to unitary
equivalence; a realization of one factor block is parametrised only by a
central character of the radical.
\end{theorem}

\begin{proof}
Automatic projectivity is \Cref{thm:spatial-multiplier-scalar-main}; the temporal
extension is \Cref{prop:full-revision-cocycle-main}; shell independence is
\Cref{thm:exact-bicharacter-stabilisation}; and the factor/Clifford statement is
\Cref{thm:factor-quotients-corrected,thm:primitivity-canonical}.  Evenness of the
nondegenerate alternating rank makes \(1+\dCl\) even.
\end{proof}

\begin{corollary}[Full-access revision phase]
\label{cor:full-access-revision-phase}
If
\[
  K_{\mathrm{rev}}=H^1(G,\Z/2),
\]
then
\[
  \Hf=H^1(G,\Z/2)\oplus\langle t\rangle,
\]
and all projective-lift, radical-reduction, and primitive Clifford conclusions
of this section apply to the complete signed cohomology module.  Thus the
earlier full-\(H^1\) formulation is recovered as a corollary, while the main
theorems require only the dynamically realized reversible subgroup.
\end{corollary}

\begin{proof}
This is the specialization of
\Cref{thm:spatial-multiplier-scalar-main,prop:full-revision-cocycle-main,thm:revision-structure-main}
to \(K_{\mathrm{rev}}=H^1(G,\Z/2)\).
\end{proof}

\begin{corollary}[Universality of the Clifford datum]
\label{cor:universality-clifford-main}
On the reversible, faithful finite-factor signed phase, neither the scalar
projective law nor the primitive Clifford representation is an additional
hypothesis.  The graded Clifford renewal memory constructed below is an explicit
representative of the unique primitive matrix type.
\end{corollary}

\begin{remark}[Sharp scope]
\label{rem:cx-status}
The predictive-unit theorem does not assert that every CP renewal tuple has a
nontrivial reversible phase, nor does it derive full access to the
elementary-two label module.  The remaining hypotheses are precisely the
existence of a realized reversible subgroup, preservation of a finite
sufficient factor, and faithfulness after
quotienting neutral labels.  Once those hold, projectivity, the commutator form,
its radical, and the primitive Clifford type are derived.  The \(\Zfour\)
amplitude lift records square-root phases but does not alter the real polar form.
\end{remark}

\subsection{Intrinsic generation by a graded Clifford renewal memory}
\label{subsec:graded-clifford-memory}

The preceding results still allow the projective revision datum to be external
to the renewal memory.  We now exhibit a fully specified completely positive
memory for which the signed reset operators themselves generate the required
projective law and its shell-stable commutator form.

\begin{construction}[Graded Clifford renewal memory]
\label{constr:graded-clifford-memory}
Fix an odd integer \(\dCl\ge1\), put \(n=1+\dCl\), and let
\(\Gamma_0,\ldots,\Gamma_{\dCl}\) be Hermitian generators of
\(\Cl_n(\C)\) on its irreducible spinor module \(S\), so that
\[
 \Gamma_\mu\Gamma_\nu+\Gamma_\nu\Gamma_\mu
 =2\delta_{\mu\nu}\one .
\]
Take
\[
 E=\{a_0,\ldots,a_{\dCl}\},\qquad
 \A_{\rm int}=\Cl_n(\C),\qquad
 \Phi_{a_\mu}=\operatorname{Ad}(\Gamma_\mu).
\]
Each reset channel is a unital \(*\)-automorphism and hence completely
positive.  On the signed lift retain the implementing operator,
\[
 \widehat S_{a_\mu}=\widetilde S_\mu\otimes\Gamma_\mu.
\]
For \(\alpha=(\alpha_0,\ldots,\alpha_{\dCl})\in\F_2^n\), define
\[
 \Gamma^\alpha=
 \Gamma_0^{\alpha_0}\cdots\Gamma_{\dCl}^{\alpha_{\dCl}},
 \qquad
 R_\alpha=\one\otimes\Gamma^\alpha .
\]
We identify the rank-\(n\) noncentral revision block with
\(\Hf=\langle t,e_1,\ldots,e_{\dCl}\rangle\), with \(R_t=\Gamma_0\).
\end{construction}

\begin{remark}[Grading generator versus timelike Clifford generator]
\label{rem:graded-model-temporal-generator}
In the primitive factor, \(R_t=\Gamma_0\) represents the image of the Krein
grading \(J_\chi\).  It is a Hermitian involution and is not itself the
anti-Hermitian Lorentzian generator.  Choosing, for example,
\(\alpha_0=e_1\), one has
\[
  \omega(\alpha_0,t)=1,
  \qquad
  R_{\alpha_0}=\Gamma_1,
\]
and the physical timelike Clifford generator is
\[
  \gamma^0=R_tR_{\alpha_0}=\Gamma_0\Gamma_1,
  \qquad
  (\gamma^0)^*=-\gamma^0,
  \qquad
  (\gamma^0)^2=-\one.
\]
Thus \(R_t\) carries the indefinite grading, while the canonical hyperbolic
pair \((R_t,R_{\alpha_0})\) carries the Lorentzian temporal square.
\end{remark}

The completely positive quotient sees only the conjugation channels.  Those
channels commute because the central sign in
\(\Gamma_\mu\Gamma_\nu=-\Gamma_\nu\Gamma_\mu\) disappears under
\(\operatorname{Ad}\).  The signed lift retains precisely the sign that the
positive quotient forgets.

\begin{lemma}[Clifford commutator form]
\label{lem:graded-clifford-form}
For \(\alpha,\beta\in\F_2^n\),
\[
 \Gamma^\alpha\Gamma^\beta
 =(-1)^{\omega(\alpha,\beta)}\Gamma^\beta\Gamma^\alpha,
 \qquad
 \omega(\alpha,\beta)=
 \Bigl(\sum_\mu\alpha_\mu\Bigr)
 \Bigl(\sum_\nu\beta_\nu\Bigr)
 -\sum_\mu\alpha_\mu\beta_\mu\pmod2 .
\]
In the generator basis, the Gram matrix is
\(\Omega=J_n-I_n\).  It is nonsingular exactly when \(n\) is even.
\end{lemma}

\begin{proof}
Reordering \(\Gamma^\alpha\Gamma^\beta\) requires one sign for each pair of
distinct occupied indices and no sign for coincident indices, giving the
displayed exponent.  It is bilinear and alternating over \(\F_2\).  On basis
vectors it equals \(1-\delta_{\mu\nu}\), hence \(\Omega=J_n-I_n\).  Over
\(\F_2\), \((J_n-I_n)^2=I_n\) when \(n\) is even; when \(n\) is odd the
all-ones vector lies in its kernel.
\end{proof}

\begin{theorem}[Intrinsic primitive and shell-stable revision block]
\label{thm:intrinsic-graded-clifford-datum}
For the graded Clifford renewal memory, the operators \(R_\alpha\) satisfy the
projective law
\[
 R_\alpha R_\beta=\tau(\alpha,\beta)R_{\alpha+\beta},
 \qquad
 \tau(\alpha,\beta)=
 (-1)^{\sum_{\mu>\nu}\alpha_\mu\beta_\nu},
\]
and are linearly independent.  Their commutator is the exact operator identity
\[
 R_\alpha R_\beta R_\alpha^{-1}R_\beta^{-1}
 =(-1)^{\omega(\alpha,\beta)}\one .
\]
Consequently \(\omega_N\equiv\omega\) on every nonempty shell.  For every odd
\(\dCl\), the rank-\((1+\dCl)\) noncentral revision block is primitive and
\[
 \Arev\cong\Cl_{1+\dCl}(\C),\qquad
 \rank_{\F_2}\Omega=1+\dCl .
\]
In particular, for \(\dCl=3\), \(\Omega=J_4-I_4\) has rank four and
\(\Arev\cong\Cl_4(\C)\cong M_4(\C)\).
\end{theorem}

\begin{proof}
The cocycle formula follows by putting each Clifford monomial back into the
fixed ordered basis.  The monomials form a vector-space basis of
\(\Cl_n(\C)\), hence are linearly independent.  They act only on the internal
signed fibre, so they commute with the renewal clock; the temporal grading
generator is the Hermitian involution \(R_t=\Gamma_0\), while the physical
timelike generator is the composite described in
\Cref{rem:graded-model-temporal-generator}.  The exact
commutator identity is \Cref{lem:graded-clifford-form}, and
\Cref{thm:exact-bicharacter-stabilisation} gives shell independence.  Since
\(n=1+\dCl\) is even, \(J_n-I_n\) is nonsingular, so the radical vanishes on
the noncentral block and \Cref{thm:primitivity-canonical} identifies the twisted
group algebra with the full Clifford matrix algebra.
\end{proof}

\begin{remark}[What has and has not been derived]
\label{rem:graded-clifford-scope}
For this named model, existence of the projective operators, their linear
independence, primitivity, and shell stabilisation are direct calculations.  By
\Cref{thm:revision-structure-main}, however, the resulting primitive block is
not special: every local, vertical, spatially exactly covariant, doubly
ergodic primitive revision sector is unitarily equivalent to it.  The full
predictive graph may still contain additional central cohomology directions;
these are factored through the radical before applying the primitive rank
theorem.  Universality therefore applies to the noncentral primitive factor,
while sectors outside the doubly ergodic phase may retain operator-valued or
sheet-dependent commutator data.
\end{remark}

\begin{corollary}[Unconditional minimality in the graded class]
\label{cor:graded-unconditional-minimality}
Within the graded Clifford class, \(2+1\) is impossible, \(1+1\) is spatially
degenerate, and \(3+1\) is the minimal spatially nondegenerate Lorentzian Dirac
endpoint, with no separate existence or stabilisation hypothesis.  Higher odd
spatial ranks remain admissible, so this is a minimality theorem rather than a
uniqueness theorem.
\end{corollary}

\begin{proof}
By \Cref{thm:intrinsic-graded-clifford-datum}, the primitive revision block
exists and stabilises for every odd \(\dCl\).  The parity and spatial
nondegeneracy argument of \Cref{thm:minimal-nondegenerate-3plus1} then gives the
claim.  The same construction for \(\dCl=5,7,\ldots\) proves that no higher odd
rank is excluded without an additional selection principle.
\end{proof}

The rank criterion tells us which Clifford ranks are possible.  It does not yet
distinguish a spatially meaningful endpoint from a one-axis toy model.  We make
that distinction explicit before stating the minimality theorem.

\begin{definition}[Spatially nondegenerate endpoint]
\label{def:spatially-nondegenerate}
A Lorentzian renewal endpoint is \emph{spatially nondegenerate} if its spatial
revision module has a nonzero two-direction sector,
\[
  \bigwedge\nolimits^2 V_{\mathrm{sp}}\neq0.
\]
Equivalently, the endpoint supports a nontrivial spatial plaquette or
antisymmetric two-reset interference observable.  This excludes the one-axis
case \(\dCl=1\) for an algebraic reason intrinsic to the renewal construction.
Since a primitive alternating form forces \(\dCl\) to be odd, the first
spatially nondegenerate primitive rank is \(\dCl=3\).
\end{definition}

\subsection{Unconditional minimality of the \texorpdfstring{$3+1$}{3+1} endpoint}
\label{subsec:minimal-3plus1}

We can now state the unconditional dimension result.  The only inputs are
primitivity of the modular revision algebra and the spatial nondegeneracy
condition just defined.  No access principle is used in this subsection.

The parity consequence should be read in relation to existing
dimension-signature arguments.  Manko\v{c} Bor\v{s}tnik and Nielsen proved that
linear-in-momentum equations for nonzero spin in even spacetime dimension,
subject to Hermiticity and Lorentz irreducibility, allow only signatures with an
odd number of time directions; van Dam and Ng gave a Wigner-representation
argument against multiple time directions; and Nielsen and Rugh observed that
Weyl-matrix counting singles out \(d=4\) through the equality
\(n_{\mathrm{Weyl}}^2=d\) among nontrivial even-dimensional Weyl systems
\cite{mankocborstniknielsen2000,vandamng2001,nielsenrugh1994}.  The theorem
below uses a different microscopic input: the nondegenerate alternating
\(\F_2\)-commutator form is generated by renewal holonomy itself, and the
classical parity conclusion is recovered as a consequence of that generated
symplectic rank.

\begin{theorem}[Minimal primitive endpoint with nontrivial spatial bivector sector]
\label{thm:minimal-nondegenerate-3plus1}
Let \(G\) be a primitive signed renewal component with
\[
  \Hf=K_{\mathrm{rev}}\oplus\langle t\rangle,
  \qquad
  \dCl:=\dim_{\F_2}K_{\mathrm{rev}}.
\]
In the full-access phase, \(\dCl=b_1(G)\).  Then:
\begin{enumerate}[label=\textup{(\roman*)},leftmargin=2.4em]
\item \(1+\dCl\) is even, hence \(\dCl\) is odd;
\item \(2+1\) is impossible in the primitive class, since \(1+\dCl=3\) carries no
      nondegenerate alternating \(\F_2\)-form;
\item \(1+1\) is primitive but spatially degenerate in the sense of
      \Cref{def:spatially-nondegenerate};
\item the minimal spatially nondegenerate primitive endpoint is
      \[
        \boxed{\dCl=3},
      \]
      hence the minimal nondegenerate Lorentzian Dirac endpoint is \(3+1\).
\end{enumerate}
This statement uses only primitivity and the nontrivial spatial-bivector
requirement; no access or cost selection principle is used.
\end{theorem}

\begin{proof}
A nondegenerate alternating bilinear form over \(\F_2\) exists only in even
dimension.  Primitivity is nondegeneracy of \(\omega\) on \(\Hf\), so
\(1+\dCl=\dim_{\F_2}\Hf\) must be even and \(\dCl\) must be odd.  Thus \(\dCl=2\) is
forbidden.  For \(\dCl=1\), \(\Hf\cong(\Z/2)^2\) admits the standard symplectic
form, so primitivity is possible; but \(\bigwedge^2\R=0\), so it has no nonzero spatial plaquette or
two-reset interference sector and is spatially degenerate by
\Cref{def:spatially-nondegenerate}.  The next odd spatial rank is \(\dCl=3\), for
which \(\bigwedge^2\R^3\neq0\).  Hence \(3+1\) is the minimal
spatially nondegenerate primitive endpoint.
\end{proof}

\begin{proposition}[Primitive revision data exist in every odd spatial rank]
\label{prop:primitive-all-odd}
For every odd \(\dCl\) there exists a recurrent signed graph with \(b_1=\dCl\)
and a modular revision datum on
\(\Hf=H^1(G,\Z/2)\oplus\langle t\rangle\) whose commutator form is
nondegenerate.  Consequently the primitive revision-algebra axioms do not select
\(\dCl=3\) over \(\dCl=5,7,\ldots\).
\end{proposition}

\begin{proof}
Take a recurrent component whose graph is the bouquet of \(\dCl\) loops on one
vertex, so \(b_1=\dCl\).  Since \(\dCl\) is odd, \(1+\dCl\) is even, and
\(\Hf\cong(\Z/2)^{1+\dCl}\) carries the standard symplectic form
\[
  \Omega=\begin{pmatrix}0&I\\ I&0\end{pmatrix}
\]
after choosing a symplectic basis.  Define a twisted group algebra
\(\C^\tau[\Hf]\) whose commutator bicharacter is \((-1)^{\Omega}\); equivalently,
take the finite Heisenberg representation associated with this symplectic
\(\F_2\)-space.  This is a modular revision datum in the sense of
\Cref{def:revision-operators-main}, and its commutator form is nondegenerate.
Thus no algebraic obstruction in the primitive revision axioms excludes higher
odd Clifford ranks.
\end{proof}

\begin{remark}[Minimality, not unconditional uniqueness]
\label{rem:minimal-not-unique}
\Cref{thm:minimal-nondegenerate-3plus1} proves that \(3+1\) is the first
nondegenerate primitive endpoint and that \(2+1\) is impossible.  It does not
prove that primitive dynamics alone chooses \(\dCl=3\).  The higher odd primitive
endpoints are algebraically admissible by \Cref{prop:primitive-all-odd}; a
selection statement requires an additional principle.
\end{remark}





\begin{proposition}[Intrinsic timelike square from a hyperbolic pair]
\label{prop:forced-timelike-generator}
Let \(\alpha_0\in K_{\mathrm{rev}}\) satisfy
\(\omega(\alpha_0,t)=1\).  Normalize the order-two projective implementers so
that
\[
  R_t=J_\chi=R_t^*=R_t^{-1},
  \qquad
  R_{\alpha_0}=R_{\alpha_0}^*=R_{\alpha_0}^{-1}.
\]
Such a normalization is unique up to independent signs.  Then
\[
  \boxed{\ \gamma^0:=J_\chi R_{\alpha_0}\ }
\]
is anti-Hermitian, Krein-odd, and has the forced square
\[
  (\gamma^0)^*=-\gamma^0,
  \qquad
  J_\chi\gamma^0J_\chi=-\gamma^0,
  \qquad
  (\gamma^0)^2=-\one.
\]
Thus the Lorentzian temporal square is fixed by the normalized hyperbolic pair
and is independent of every scalar rephasing compatible with the Hermitian
involution normalization.
\end{proposition}

\begin{proof}
The temporal-row theorem gives
\(J_\chi R_{\alpha_0}=-R_{\alpha_0}J_\chi\).  Hence
\[
 (J_\chi R_{\alpha_0})^*=R_{\alpha_0}J_\chi
 =-J_\chi R_{\alpha_0},
 \qquad
 (J_\chi R_{\alpha_0})^2=-J_\chi^2R_{\alpha_0}^2=-\one.
\]
Conjugation by \(J_\chi\) gives Krein oddness.  An order-two inner
automorphism has an implementing unitary whose square is scalar; scalar
normalization makes it a Hermitian involution, after which the residual phase
freedom is only \(\pm1\), which does not alter the displayed square.
\end{proof}


\begin{lemma}[Clifford completion of the normalized temporal generator]
\label{lem:temporal-clifford-completion}
Suppose the primitive quotient has rank
\(1+\dCl=2m\), so that its factor is
\(M_{2^m}(\C)\cong\Cl_{1+\dCl}(\C)\).  The operator
\(\gamma^0=J_\chi R_{\alpha_0}\) of
\Cref{prop:forced-timelike-generator} can be completed by Hermitian involutions
\(\gamma^1,\ldots,\gamma^{\dCl}\) satisfying
\[
  J_\chi\gamma^iJ_\chi=\gamma^i,
  \qquad
  \{\gamma^0,\gamma^i\}=0,
  \qquad
  \{\gamma^i,\gamma^j\}=2\delta^{ij}\one.
\]
Any two such completions determine unitarily equivalent representations of
the full complex Clifford algebra.
\end{lemma}

\begin{proof}
The anticommuting Hermitian involutions
\(J_\chi\) and \(R_{\alpha_0}\) form a hyperbolic pair.  Their \(+1\) and
\(-1\) eigenspaces have equal dimension, and after a unitary change of basis
they take the form
\[
  J_\chi=
  \begin{pmatrix}1&0\\0&-1\end{pmatrix}\otimes\one,
  \qquad
  R_{\alpha_0}=
  \begin{pmatrix}0&1\\1&0\end{pmatrix}\otimes\one.
\]
Consequently
\[
  \gamma^0=J_\chi R_{\alpha_0}
  =
  \begin{pmatrix}0&1\\-1&0\end{pmatrix}\otimes\one.
\]
Since \(\dCl=2m-1\), choose Hermitian generators
\(\eta_1,\ldots,\eta_{\dCl}\) of the irreducible complex
\(\Cl_{\dCl}(\C)\)-module on \(\C^{2^{m-1}}\), and set
\[
  \gamma^i=
  \begin{pmatrix}1&0\\0&-1\end{pmatrix}\otimes\eta_i.
\]
These operators are Hermitian involutions, commute with \(J_\chi\),
anticommute with \(\gamma^0\), and satisfy the required spatial Clifford
relations.  The resulting \(1+\dCl=2m\) generators give an irreducible
representation of the simple complex algebra
\(\Cl_{2m}(\C)\cong M_{2^m}(\C)\).  Unitary equivalence of any two full
completions follows from uniqueness of its irreducible representation.
\end{proof}


% =====================================================================
\section{The Lorentzian Dirac continuum}
\label{sec:lorentzian}

The preceding section has already produced the finite Clifford factor from
reversible predictive dynamics.  We now pass from that derived operator algebra
to its continuum Dirac symbol.  The deterministic product-order continuum and
its polyhedral obstruction are recorded separately in
\Cref{app:product-order}; the main line begins with the signed spectral cone,
then proves operator convergence, self-averaging, and the slowly varying
Levi--Civita limit.

\subsection{Signed Clifford symbols and Lorentzian frame universality}
\label{sec:isotropisation}

The continuum ladder has identified the obstruction: a deterministic finite
product order gives a polyhedral cone.  We now move from order alone to the
signed Dirac symbol.  The guiding idea is simple.  The signed sector supplies
the timelike Clifford generator, while a large family of reset directions
supplies the spatial second moment.  A nondegenerate second moment supplies a spatial vielbein; literal isotropy is only a convenient representative of the resulting flat isometry class.

There is, however, an important structural boundary.  Positive averaging by
itself cannot create Lorentzian signature.  The tensor
$\sum_a p_a\,v_a\otimes v_a$ formed from real reset directions and positive
weights is positive semidefinite.  Thus isotropisation supplies roundness, but
the fundamental symmetry $J_\chi$ supplies the indefinite sign.  We first record
this positivity obstruction, then introduce the Krein--Clifford reset datum,
then prove the model isotropic limit and its general form.

\subsubsection{Positive averaging is Riemannian}

\begin{proposition}[Symbol-level positivity obstruction]
\label{prop:positive-averaging}
Let $v_1,\dots,v_m\in\R^{n}$ be real vectors and $p_1,\dots,p_m\ge0$.  Then the
tensor
\[
  G=\sum_{a=1}^m p_a\,v_a\otimes v_a
\]
is positive semidefinite, and positive definite if $\spn\{v_a\}=\R^{n}$.  In
particular $G$ is never of Lorentzian signature, and the quadratic form
$\xi\mapsto G(\xi,\xi)$ has trivial null cone $\{0\}$ when $G$ is definite.
\end{proposition}

\begin{proof}
For every covector $\xi$, $\;G(\xi,\xi)=\sum_a p_a\,(v_a\cdot\xi)^2\ge0$, with
equality forcing $v_a\cdot\xi=0$ for all $a$ with $p_a>0$.  If the $v_a$ span,
this forces $\xi=0$.  A positive semidefinite form has no negative eigenvalue,
hence cannot have signature $(1,n-1)$.
\end{proof}

\Cref{prop:positive-averaging} is the principal-symbol shadow of
\Cref{cor:positivity-obstruction}: completely positive data, averaged with
positive weights, remain Riemannian.  The remedy is the same as in
\Cref{sec:signed}, and we now make it explicit at the level of the Clifford
symbol.



\subsubsection{The derived Krein--Clifford datum}

\begin{definition}[Derived Krein--Clifford reset datum]
\label{def:krein-clifford}
Let \(J=J^*=J^{-1}\) be the fundamental symmetry of the signed cover, and let
\(\alpha_0\) be a revision class with \(\omega(\alpha_0,t)=1\).  With the
Hermitian involutive normalization of
\Cref{prop:forced-timelike-generator}, define
\[
  \gamma^0:=J R_{\alpha_0}.
\]
Choose the compatible Clifford completion of
\Cref{lem:temporal-clifford-completion}, consisting of Hermitian spatial
generators \(\gamma^1,\ldots,\gamma^{\dCl}\) satisfying
\[
  J\gamma^iJ=\gamma^i,
  \qquad
  \{\gamma^0,\gamma^i\}=0,
  \qquad
  \{\gamma^i,\gamma^j\}=2\delta^{ij}\one.
\]
Then
\[
  (\gamma^0)^*=-\gamma^0,
  \qquad
  (\gamma^0)^2=-\one,
  \qquad
  J\gamma^0J=-\gamma^0.
\]
For \(\hat\theta\in\Sphere^{\dCl-1}\), write
\(\gamma(\hat\theta)=\sum_i\hat\theta^i\gamma^i\), so that
\(\gamma(\hat\theta)^2=\one\).
\end{definition}

\begin{definition}[Direction measure and coarse-grained Dirac symbol]
\label{def:dirac-symbol}
A direction system is a family
\(\{(\hat\theta_a,p_a)\}_{a=1}^{m}\) with
\(\hat\theta_a\in\Sphere^{\dCl-1}\), \(p_a>0\), and
\(\sum_a p_a=1\).  Its empirical measure and second moment are
\[
  \nu=\sum_a p_a\delta_{\hat\theta_a},
  \qquad
  M(\nu)^{ij}=\int_{\Sphere^{\dCl-1}}
  \hat\theta^i\hat\theta^j\,d\nu(\hat\theta).
\]
For \(\kappa>0\), define
\[
  \sym_\nu(\xi)
  =i\left(
     \gamma^0\xi_0+
     \kappa\int_{\Sphere^{\dCl-1}}
     \gamma(\hat\theta)(\hat\theta\cdot\vec\xi)
     \,d\nu(\hat\theta)
   \right).
\]
Equivalently, the spatial term is
\(\kappa\gamma^iM(\nu)^{ij}\xi_j\).
\end{definition}

\begin{lemma}[Forced Lorentzian symbol square]
\label{lem:symbol-square}
For every direction measure \(\nu\),
\[
  \sym_\nu(\xi)^2
  =\left(
      \xi_0^2-\kappa^2\vec\xi^{\,\top}M(\nu)^2\vec\xi
    \right)\one
  =g_\nu^{\mu\nu}\xi_\mu\xi_\nu\one,
\]
where
\[
  g_\nu^{-1}=\operatorname{diag}
  \bigl(1,-\kappa^2M(\nu)^2\bigr).
\]
\end{lemma}

\begin{proof}
The temporal square contributes
\(-i^2\xi_0^2=+\xi_0^2\), the spatial square contributes
\(-\kappa^2\vec\xi^{\,\top}M(\nu)^2\vec\xi\), and the mixed terms vanish by
\(\{\gamma^0,\gamma^i\}=0\).
\end{proof}

\begin{theorem}[Forced single-time Lorentzian signature]
\label{thm:signature-krein}
If the reset directions span the spatial module, equivalently
\(M(\nu)\succ0\), then \(g_\nu^{-1}\) has signature \((1,\dCl)\).  Its unique
positive eigenvalue is the temporal slot generated by the canonical hyperbolic
pair \((J,R_{\alpha_0})\); the positive spatial averaging contributes only the
negative-definite block \(-\kappa^2M(\nu)^2\).  In particular, no Euclidean
branch or multiple-time branch remains after the normalized derived temporal
generator is used.
\end{theorem}

\begin{proof}
Positive definiteness of \(M(\nu)\) makes
\(-\kappa^2M(\nu)^2\) negative definite, while the temporal coefficient is
\(+1\) by \Cref{prop:forced-timelike-generator,lem:symbol-square}.  Hence the
inertia is exactly \((1,\dCl)\).
\end{proof}

\begin{remark}[Phase independence of the temporal sign]
The metric sign is not the square of an arbitrarily phased projective lift.
It is the square of the composite \(J R_{\alpha_0}\) after the canonical
Hermitian-involution normalization of the order-two factors.  The residual
\(\pm1\) choices leave \((\gamma^0)^2=-\one\) unchanged.
\end{remark}

\subsubsection{Full generality}

The isotropic theorem is the round special case of a more general statement.
A non-isotropic but nondegenerate second moment still gives a Lorentzian
quadric; it is simply an anisotropic one until the spatial frame is rescaled.

\begin{theorem}[General Clifford rounding]
\label{thm:general-rounding}
Let $\{\nu_R\}$ be any sequence of direction measures with weak-$*$ limit
$\nu$ on $\Sphere^{\dCl-1}$ whose second moment $M(\nu)\succ0$ (equivalently,
$\supp\nu$ is not contained in the sphere of any proper subspace), and let any
Krein--Clifford reset datum of signature $(1,\dCl)$ be given.  Then:
\begin{enumerate}[label=\textup{(\roman*)},leftmargin=2.2em]
\item $g_{\nu_R}\to g_\nu=\operatorname{diag}(1,-\kappa^2M(\nu)^2)$, a
      nondegenerate Lorentzian inverse metric of signature $(1,\dCl)$;
\item the characteristic variety of $\sym_{\nu}$ is the smooth quadric Lorentz
      cone of $g_\nu$, hence is linearly equivalent to the standard round
      Lorentz cone;
\item if, in addition, the reset rays used for the order cone become dense in
      the null directions of this limiting quadric, then the corresponding
      order cones converge in Hausdorff distance on compact ray sections to the
      solid causal cone dual to $g_\nu$;
\item in the isotropic normalisation $M(\nu)=\dCl^{-1}I_{\dCl}$ and
      $\kappa=\dCl$, the preceding item reduces to the round-cone convergence of
      \Cref{thm:model-rounding};
\item the signature of $g_\nu$ is fixed by the Krein--Clifford datum and is
      independent of the direction measure, as long as $M(\nu)\succ0$.
\end{enumerate}
\end{theorem}

\begin{proof}
\emph{(i)} As in \Cref{thm:model-rounding}\,(i), $M(\nu_R)\to M(\nu)$ by
weak-$*$ convergence against the bounded continuous integrands
$\hat\theta^i\hat\theta^j$.  Since $M(\nu)\succ0$, the spatial block
$-\kappa^2M(\nu)^2$ is negative definite, so $g_\nu$ has signature $(1,\dCl)$ by
\Cref{thm:signature-krein}.

\emph{(ii)} By Lemma~\ref{lem:symbol-square} the characteristic variety is
$\{g_\nu^{\mu\nu}\xi_\mu\xi_\nu=0\}$, the zero set of a nondegenerate Lorentzian
quadratic form.  Any two nondegenerate Lorentzian quadratic forms of the same
signature are linearly equivalent, so a spatial frame change diagonalising
$M(\nu)$ followed by rescaling carries the quadric to the standard round cone.

\emph{(iii)} The order cone is determined by the conical hull of the reset rays,
not merely by the second moment.  Therefore density of the reset rays in the
null directions of the limiting quadratic cone is the required hypothesis.  Under
this hypothesis, convergence of closed convex conical hulls on compact ray
sections is the standard Hausdorff convergence of convex hulls of dense subsets.

\emph{(iv)} In the isotropic normalisation, the null directions of the limiting
quadric are exactly $(1,\hat\theta)$ with $\hat\theta\in\Sphere^{\dCl-1}$, so
(iii) is precisely \Cref{thm:model-rounding}\,(iii).

\emph{(v)} The form $g_\nu=\diag(1,-\kappa^2M(\nu)^2)$ has one positive and
$\dCl$ negative directions whenever $M(\nu)\succ0$.  The single positive slot is
the one created by the Krein-odd timelike Clifford generator, as in
\Cref{thm:signature-krein}; changing $\nu$ changes the spatial scale and shape,
not the inertia.
\end{proof}

\subsubsection{Model case: isotropic rounding}

The cleanest case is the isotropic one.  Here the spatial second moment becomes
a scalar multiple of the identity, so the Lorentzian quadric is round after the
normalisation $\kappa=\dCl$.

\begin{theorem}[Isotropic Clifford rounding]
\label{thm:model-rounding}
Let $\nu_R\to\nu_{\mathrm{unif}}$ weak-$*$ on $\Sphere^{\dCl-1}$, where
$\nu_{\mathrm{unif}}$ is the rotation-invariant probability measure; for
instance let $\nu_R$ be normalised counting measure on a net that
equidistributes on $\Sphere^{\dCl-1}$.  Take $\kappa=\dCl$.  Then:
\begin{enumerate}[label=\textup{(\roman*)},leftmargin=2.2em]
\item $M(\nu_R)\to \tfrac1{\dCl}\,I_{\dCl}$, hence $g_{\nu_R}\to
      \operatorname{diag}(1,-I_{\dCl})$, the standard Minkowski inverse metric of
      signature $(1,\dCl)$;
\item the spectral null cones $\{g_{\nu_R}^{\mu\nu}\xi_\mu\xi_\nu=0\}$ converge,
      in Hausdorff distance on the unit sphere of covectors, to the round null
      cone $\{\xi_0^2=|\vec\xi|^2\}$;
\item the order cones $\cone\{v_a^{(R)}\}$ converge, in Hausdorff distance on
      the unit sphere of directions, to the solid round future cone
      $\{x_0\ge|\vec x|\}$.
\end{enumerate}
Consequently both the spectral null cone and the causal order cone round to the
same standard Lorentz cone, and the polyhedral obstruction of
\Cref{thm:obstruction} is removed in the limit.
\end{theorem}

\begin{proof}
\emph{(i)} By rotation invariance $\int\hat\theta^i\hat\theta^j\,
d\nu_{\mathrm{unif}}=\lambda\,\delta^{ij}$ for some $\lambda$, and taking the
trace gives $\dCl\lambda=\int|\hat\theta|^2\,d\nu_{\mathrm{unif}}=1$, so
$\lambda=1/\dCl$ and $M(\nu_{\mathrm{unif}})=\tfrac1{\dCl} I_{\dCl}$.  The entries
$\hat\theta\mapsto\hat\theta^i\hat\theta^j$ are bounded and continuous on the
compact sphere, so weak-$*$ convergence $\nu_R\to\nu_{\mathrm{unif}}$ gives
$M(\nu_R)\to\tfrac1{\dCl} I_{\dCl}$.  With $\kappa=\dCl$, $\kappa^2M(\nu_R)^2\to
\dCl^2\cdot\tfrac1{\dCl^2}I_{\dCl}=I_{\dCl}$, so $g_{\nu_R}\to\operatorname{diag}(1,-I_{\dCl})$.

\emph{(ii)} The null cone of a nondegenerate quadratic form depends
continuously, in Hausdorff distance on the covector sphere, on the coefficients
of the form; apply this to $g_{\nu_R}\to\operatorname{diag}(1,-I_{\dCl})$.

\emph{(iii)} The future order cone is the closed convex conical hull of the
null rays $v=(1,\hat\theta)$.  As the directions $\hat\theta_a^{(R)}$ become
dense in $\Sphere^{\dCl-1}$, the conical hull of $\{(1,\hat\theta_a^{(R)})\}$
converges in Hausdorff distance (on rays) to the conical hull of
$\{(1,\hat\theta):\hat\theta\in\Sphere^{\dCl-1}\}$, which is exactly the solid cone
$\{x_0\ge|\vec x|\}$.  This is the standard convergence of convex hulls of
$\chi$-nets.
\end{proof}

\begin{example}[Spectral roundness with finitely many directions, $\dCl=2$]
\label{ex:tight-frame}
The two cones round under different conditions.  Take $\dCl=2$ and the three
equally spaced directions $\hat\theta_a=(\cos\tfrac{2\pi a}{3},
\sin\tfrac{2\pi a}{3})$, $a=0,1,2$, with $p_a=\tfrac13$.  Then
$\sum_a\hat\theta_a\hat\theta_a^{\top}=\tfrac32 I_2$, so
$M=\tfrac12 I_2$ is already isotropic, and with $\kappa=2$ the spectral null
cone is the round circular cone $\{\xi_0^2=|\vec\xi|^2\}$ \emph{exactly}, with
only three directions.  Yet the order cone $\cone\{(1,\hat\theta_a)\}$ is a
three-faced pyramid: polyhedral, not round.  Thus a finite tight frame rounds
the Dirac symbol but not the causal hull; rounding the order cone genuinely
requires the equidistributing limit $m_R\to\infty$ of
\Cref{thm:model-rounding}\,(iii).
\end{example}


% =====================================================================

\begin{example}[The octahedral \(3+1\) isotropic symbol]
\label{ex:3plus1}
Take \(\dCl=3\), \(n_{\mathrm{Cl}}=1+3\), and let the reset directions be the
six vertices of the spatial cross-polytope,
\[
  \hat\theta_a\in\{\pm e_1,\pm e_2,\pm e_3\}\subset\Sphere^{2},
  \qquad p_a=\frac16 .
\]
Then the second moment is exactly isotropic:
\[
  M(\nu)
  =\sum_a p_a\,\hat\theta_a\hat\theta_a^{\top}
  =\frac16\Bigl(2\sum_{i=1}^3 e_ie_i^{\top}\Bigr)
  =\frac13 I_3 .
\]
With the normalisation \(\kappa=\dCl=3\), this gives
\[
  \kappa^2M(\nu)^2
  =9\Bigl(\frac13 I_3\Bigr)^2
  =I_3,
  \qquad
  g_\nu=\operatorname{diag}(1,-I_3).
\]
Hence, by Lemma~\ref{lem:symbol-square}, the characteristic variety of
\(\sym_\nu\) is
\[
  \xi_0^2=\xi_1^2+\xi_2^2+\xi_3^2,
\]
the standard \(3+1\)-dimensional Minkowski light cone. Thus, once the primitive
endpoint has \(\dCl=3\), the round Lorentzian principal symbol is realised
exactly by a finite isotropic reset frame. The single timelike slot is still the
Krein-odd generator \(\gamma^0\); the octahedral averaging supplies only the
round spatial symbol.
\end{example}

\begin{remark}[Symbol roundness versus order roundness]
\label{rem:order-vs-symbol}
The preceding example is exact at the level of the Dirac principal symbol, but
not at the level of the order cone. The order cone generated by the six rays
\[
  v_a=(1,\hat\theta_a)\in\R^{1+3}
\]
has cross-section
\[
  \operatorname{conv}\{\pm e_1,\pm e_2,\pm e_3\},
\]
the octahedron. It is therefore polyhedral, with six extreme null rays, not a
round solid cone. Rounding the order cone requires the equidistributing limit
\(\nu_R\to\nu_{\mathrm{unif}}\) of Theorem~\ref{thm:model-rounding}; only then do
the causal cones converge in Hausdorff distance to
\[
  \{x_0\ge |\vec x|\}.
\]
Thus finite tight frames can make the spectral cone round exactly, while
rounding the causal hull is a genuine continuum-limit statement.
\end{remark}



\subsubsection{Scope of the theorem}

\begin{remark}[Consistency with the obstruction]
\label{rem:no-contradiction}
\Cref{thm:model-rounding,thm:general-rounding} do not contradict
\Cref{thm:obstruction,cor:lorentz-violation}.  Those concern the deterministic
finite $c$-generator product \emph{order} cone, which is polyhedral with $c$
extreme rays.  Here the number of directions tends to infinity and the cone is
read from the signed Dirac symbol.  The hypotheses are disjoint, and the escape
has a definite cost: one leaves the finite deterministic model and commits to
the spectral cone of Definition~\ref{def:dirac-symbol}.
\end{remark}

\begin{remark}[What is and is not proved]
\label{rem:isotropisation-scope}
The theorems in this section are statements about the principal symbol and its
characteristic cone.  The flat constant-coefficient operator-level upgrade is
proved separately in \Cref{thm:operator-limit}.  The following are still
\emph{not} proved at this stage: \textup{(a)} a position-dependent curved metric
$g(x)$, which would require the direction measure $\nu=\nu_x$ to vary slowly in
$x$ and a variable-coefficient/adiabatic analysis ; \textup{(b)} a
fully discrete spacetime evolution with a discretised timelike direction, which
requires a separate stability/CFL hypothesis; \textup{(c)} field
equations, curvature dynamics, or matter coupling.  The Lorentzian signature is carried by the minimal signed cover and the
forced temporal Clifford generator; it is not produced by positive averaging.
\end{remark}

\begin{remark}[Dimension count]
\label{rem:isotropisation-dimension}
With $n_{\mathrm{Cl}}=1+\dCl$, the case $\dCl=3$ gives a $(3+1)$-dimensional round Lorentz cone,
consistent with \Cref{rem:dimension-count}: the emergent geometry is
$1+(c-1)$-dimensional with $c=n$, the single timelike slot being the Krein-odd
generator $\gamma^0$ rather than an added coordinate.
\end{remark}

% =====================================================================
\subsection{Operator-level convergence of the rounded signed Dirac}
\label{sec:operator-convergence}

The preceding subsection was a statement about principal symbols and cones.  We
now upgrade it to operators.  The goal is to show that the rescaled signed
renewal Dirac operators converge to a constant-coefficient Lorentzian Dirac
evolution.  This requires one care point: the covariant Lorentzian Dirac
operator itself is hyperbolic, not elliptic, and has empty resolvent set in the
flat $L^2$ spacetime picture.  The operator limit therefore has two layers:
graph convergence for the full spacetime Dirac operator, and strong resolvent
convergence for the associated self-adjoint Dirac Hamiltonian.

Throughout this subsection the coefficients are constant.  Curved and
slowly-varying reset statistics are treated in \Cref{sec:curved-stability}.

\subsubsection{Why the Lorentzian Dirac operator forces the evolution picture}

Let $g=g_\nu$ be a normalised emergent inverse metric of Lemma~\ref{lem:symbol-square}
in the isotropic case $\kappa^2M(\nu)^2=I_{\dCl}$, so $g=\operatorname{diag}(1,-I_{\dCl})$,
and let $\Dirac_g$ be the constant-coefficient operator on
$L^2(\R^{1+\dCl},V)$ with Fourier multiplier $\sym_g(\xi)=i(\gamma^0\xi_0
+\gamma^i\xi_i)$, i.e.\ $\Dirac_g=\gamma^0\partial_0+\gamma^i\partial_i$.

\begin{proposition}[The Lorentzian Dirac operator has empty resolvent set]
\label{prop:empty-resolvent}
The operator $\Dirac_g$ is closed and densely defined on $L^2(\R^{1+\dCl},V)$, but
$\spec(\Dirac_g)=\C$; its resolvent set is empty.  Consequently strong (and
\emph{a fortiori} norm) resolvent convergence \emph{to} $\Dirac_g$ is not the
right notion, and the operator-level limit cannot be phrased as resolvent
convergence of the covariant Lorentzian Dirac operator itself.
\end{proposition}

\begin{proof}
The operator is the Fourier multiplier by the matrix $\sym_g(\xi)$ with domain
$\{u:\sym_g\hat u\in L^2\}$, hence is closed and densely defined.  We show that no
$z\in\C$ lies in the resolvent set.

First take $z=0$.  The multiplier $\sym_g(\xi)$ fails to be bounded below in
any neighbourhood of the nonzero null cone.  Indeed, after writing
$Q(\xi)=g^{\mu\nu}\xi_\mu\xi_\nu$, choose a nonzero null covector
$\eta$ and a transverse covector $\zeta$ with $dQ_\eta(\zeta)\ne0$.  For large
$R$ and $|s|<c/R$, the covectors $\xi=R\eta+s\zeta$ form positive-measure tubes
on which $|Q(\xi)|$ remains bounded while $|\xi|\simeq R$.  Since
$\det\sym_g(\xi)$ vanishes on $Q=0$ and the smallest singular value of
$\sym_g(\xi)$ tends to zero along these tubes as they approach the null cone,
$\sym_g(\xi)^{-1}$ is not essentially bounded.  Hence $0$ cannot belong to the
resolvent set.

Now let $z\ne0$.  Since
\[
  \sym_g(\xi)^2=-Q(\xi)\one,
\]
where the sign depends only on the convention for the factor $i$ in
$\sym_g$, the inverse, wherever it exists, has the form
\[
  (\sym_g(\xi)-z)^{-1}
  =\frac{\sym_g(\xi)+z\one}{-Q(\xi)-z^2} .
\]
Choose a constant $C>0$ such that the denominator stays bounded away from zero
whenever $|Q(\xi)|\le C$; this is possible after taking $C<|z|^2/2$.  The same
positive-measure tubular neighbourhoods of the null cone contain covectors with
$|Q(\xi)|\le C$ and $|\xi|\to\infty$.  On these sets the denominator is uniformly
bounded above and below, whereas
$\norm{\sym_g(\xi)+z\one}$ grows linearly in $|\xi|$ on a subregion of positive
measure.  Thus the inverse multiplier is not essentially bounded.  A Fourier
multiplier is bounded on $L^2$ only when its multiplier is essentially bounded.
Therefore $(\Dirac_g-z)^{-1}$ is unbounded for every $z\ne0$.  The resolvent set
is empty, i.e.\ $\spec(\Dirac_g)=\C$.
\end{proof}

\begin{remark}[The fix is the Dirac Hamiltonian]
\label{rem:hamiltonian-fix}
\Cref{prop:empty-resolvent} reflects the hyperbolicity of the wave operator
$\Dirac_g^2=g^{\mu\nu}\partial_\mu\partial_\nu$, which is not elliptic and not
sign-definite.  The standard remedy is to read $\Dirac_g\psi=0$ as an evolution
equation in the timelike direction.  Writing $\Dirac_g\psi=0$ as
$\gamma^0\partial_0\psi=-\gamma^i\partial_i\psi$ and multiplying by
$-\gamma^0=(\gamma^0)^{-1}$ gives $i\partial_0\psi=\Ham_g\psi$ with the Dirac
\emph{Hamiltonian}
\[
  \Ham_g=A^i(-i\partial_i),\qquad A^i:=\gamma^0\gamma^i .
\]
The matrices $A^i$ are the analogues of the Dirac $\alpha$-matrices; the
Krein-odd timelike generator $\gamma^0$ is exactly what makes them
self-adjoint.  This passage to the self-adjoint generator is where resolvent and
group convergence legitimately live.
\end{remark}

\begin{lemma}[The Dirac Hamiltonian is self-adjoint]
\label{lem:alpha-selfadjoint}
With the Hilbert adjoints induced by the fundamental symmetry $J$ of
Definition~\ref{def:J-eps} (so $\gamma^{0\dagger}=-\gamma^0$, $\gamma^{i\dagger}=\gamma^i$),
the matrices $A^i=\gamma^0\gamma^i$ are Hermitian and satisfy
$\{A^i,A^j\}=2\delta^{ij}\one$.  Hence for any positive-definite symmetric
$M$ the multiplier $h_g(\vec\xi)=\kappa\,A^iM^{ij}\xi_j$ is a Hermitian matrix
with $h_g(\vec\xi)^2=\kappa^2\,\vec\xi^{\top}M^2\vec\xi\,\one$, and the operator
$\Ham_g=\kappa A^iM^{ij}(-i\partial_j)$ is self-adjoint on the Sobolev domain
$H^1(\R^{\dCl},V)$, essentially self-adjoint on $\Sch(\R^{\dCl},V)$.
\end{lemma}

\begin{proof}
$(A^i)^\dagger=(\gamma^i)^\dagger(\gamma^0)^\dagger=\gamma^i(-\gamma^0)
=-\gamma^i\gamma^0=\gamma^0\gamma^i=A^i$, using $\{\gamma^0,\gamma^i\}=0$.  Next
$A^iA^j=\gamma^0\gamma^i\gamma^0\gamma^j$; since $\gamma^0\gamma^i\gamma^0
=-(\gamma^0)^2\gamma^i=\gamma^i$, this equals $\gamma^i\gamma^j$, whence
$\{A^i,A^j\}=\{\gamma^i,\gamma^j\}=2\delta^{ij}\one$.  Contracting the symmetric
tensor $(M\vec\xi)_i(M\vec\xi)_j$ against $A^iA^j$ replaces $A^iA^j$ by
$\tfrac12\{A^i,A^j\}$, giving $h_g(\vec\xi)^2=\kappa^2|M\vec\xi|^2\one$.  A
constant-coefficient first-order operator with Hermitian symbol of linear growth
is self-adjoint on $H^1$ and essentially self-adjoint on $\Sch$ by the standard
Fourier-multiplier criterion.
\end{proof}

\subsubsection{The discrete renewal Dirac and its finite-difference multipliers}

We now define the discrete operators whose limit will be taken.  The symmetric
difference is chosen so that the spatial momentum multipliers are real and the
Hamiltonians are self-adjoint.

\begin{definition}[Symmetric reset differences]
\label{def:reset-differences}
Fix mesh scales $h_R\downarrow0$.  For a spatial direction
$\hat\theta\in\Sphere^{\dCl-1}$ let
\[
  (\delta_{\hat\theta,R}f)(x)=\frac{f(x+h_R\hat\theta)-f(x-h_R\hat\theta)}{2h_R},
\]
whose Fourier multiplier is $i\,s_{\hat\theta,R}(\vec\xi)$, with
$\;s_{\hat\theta,R}(\vec\xi)=h_R^{-1}\sin(h_R\,\hat\theta\cdot\vec\xi)$.  The
self-adjoint lattice momentum is
\[
  P_{\hat\theta,R}:=-i\delta_{\hat\theta,R},
\]
with real multiplier $s_{\hat\theta,R}(\vec\xi)$.  One has
$s_{\hat\theta,R}(\vec\xi)\to\hat\theta\cdot\vec\xi$ as $R\to\infty$, uniformly on
compact $\vec\xi$-sets, with the uniform bound
$|s_{\hat\theta,R}(\vec\xi)|\le|\hat\theta\cdot\vec\xi|\le|\vec\xi|$.
\end{definition}

For a direction system $\nu_R=\sum_a p_a^{(R)}\delta_{\hat\theta_a^{(R)}}$ as in
Definition~\ref{def:dirac-symbol}, define the discrete renewal Dirac Hamiltonian
\[
  \Ham_R=\kappa\sum_a p_a^{(R)}\,A(\hat\theta_a^{(R)})\,P_{\hat\theta_a^{(R)},R},
  \qquad A(\hat\theta)=\gamma^0\gamma(\hat\theta)=\textstyle\sum_i\hat\theta^i A^i,
\]
with spatial multiplier $h_R(\vec\xi)=\kappa\sum_a p_a^{(R)}
A(\hat\theta_a^{(R)})\,s_{\hat\theta_a^{(R)},R}(\vec\xi)$.  Each
$A(\hat\theta)$ is Hermitian (Lemma~\ref{lem:alpha-selfadjoint}), so $h_R(\vec\xi)$ is
Hermitian and $\Ham_R$ is self-adjoint by Fourier multiplication.  The full
spacetime operator $\Dirac_R$ discretises the timelike direction and the spatial
Dirac part by
\[
  m_R(\xi)=i\gamma^0s_{0,R}(\xi_0)
  +i\kappa\sum_a p_a^{(R)}\gamma(\hat\theta_a^{(R)})
    s_{\hat\theta_a^{(R)},R}(\vec\xi),
\]
which is the finite-difference analogue of the Lorentzian Dirac multiplier.

\subsubsection{Strong convergence of the full operator}

For the full spacetime Dirac operator the correct convergence statement is
convergence on a common core.  This is enough to justify the symbol limit at the
operator level, while respecting the empty-resolvent obstruction above.

\begin{theorem}[Graph convergence of the full signed Dirac]
\label{thm:full-operator}
Assume the symbol convergence hypotheses of \Cref{thm:general-rounding}
($\nu_R\to\nu$ weak-$*$, $M(\nu)\succ0$) and $h_R\downarrow0$.  Then for every
$u\in\Sch(\R^{1+\dCl},V)$,
\[
  \Dirac_R\,u\longrightarrow \Dirac_g\,u\quad\text{in }L^2(\R^{1+\dCl},V),
\]
where $\Dirac_g$ has multiplier $\sym_\nu$.  Since $\Sch$ is a core for both
operators, this is convergence on a common core (graph convergence).  By
\Cref{prop:empty-resolvent} this is the strongest operator-level statement
available for $\Dirac_g$ itself; resolvent convergence is not.
\end{theorem}

\begin{proof}
Let $m_R(\xi)$ be the multiplier of $\Dirac_R$ and $\sym_\nu(\xi)$ that of
$\Dirac_g$.  By Definition~\ref{def:reset-differences}, $s_{0,R}(\xi_0)\to\xi_0$ and
\[
  \kappa\sum_a p_a^{(R)}\gamma(\hat\theta_a^{(R)})
  s_{\hat\theta_a^{(R)},R}(\vec\xi)
  \longrightarrow
  \kappa\gamma^iM(\nu)^{ij}\xi_j
\]
pointwise.  Hence $m_R(\xi)\to\sym_\nu(\xi)$
pointwise, with the uniform domination
$\norm{m_R(\xi)}\le C(1+|\xi|)$, $C$ independent of $R$, because
$|s_{\hat\theta,R}|\le|\vec\xi|$, $|s_{0,R}|\le|\xi_0|$ and
$\sum_a p_a^{(R)}\norm{A(\hat\theta_a)}\le\max_{\hat\theta}\norm{A(\hat\theta)}
<\infty$.  By Plancherel,
\[
  \norm{\Dirac_R u-\Dirac_g u}_{L^2}^2
  =\int_{\R^{1+\dCl}}\norm{\bigl(m_R(\xi)-\sym_\nu(\xi)\bigr)\hat u(\xi)}^2\,d\xi .
\]
The integrand tends to $0$ pointwise and is dominated by
$\bigl(\norm{m_R(\xi)}+\norm{\sym_\nu(\xi)}\bigr)^2\norm{\hat u(\xi)}^2
\le C'(1+|\xi|)^2\norm{\hat u(\xi)}^2$, which is integrable for
$u\in\Sch$.  Dominated convergence gives the claim.
\end{proof}

\subsubsection{Strong resolvent and propagator convergence of the Hamiltonian}

After passing to the Hamiltonian picture, the limiting operator is
self-adjoint.  Standard functional calculus then gives the usual strong
resolvent and propagator convergence.

\begin{theorem}[Operator-level Clifford continuum limit]
\label{thm:operator-limit}
Under the hypotheses of \Cref{thm:full-operator}, the self-adjoint discrete
Hamiltonians $\Ham_R$ converge to the self-adjoint continuum Dirac Hamiltonian
$\Ham_g=\kappa A^iM(\nu)^{ij}(-i\partial_j)$ in the strong resolvent sense:
\[
  (\Ham_R-z)^{-1}\longrightarrow(\Ham_g-z)^{-1}\quad\text{strongly, for all }
  z\in\C\setminus\R .
\]
Consequently:
\begin{enumerate}[label=\textup{(\roman*)},leftmargin=2.2em]
\item the Dirac propagators converge, $e^{-it\Ham_R}\to e^{-it\Ham_g}$ strongly,
      uniformly for $t$ in compact intervals;
\item $\varphi(\Ham_R)\to\varphi(\Ham_g)$ strongly for every bounded continuous
      $\varphi$, so spectral projections onto intervals whose endpoints are not
      atoms of the limit converge strongly;
\item the limiting dispersion is $\spec(\Ham_g)=\R$, absolutely continuous, with
      energies $\pm\kappa\,|M(\nu)\vec\xi|$, recovering the Lorentzian light cone
      $\xi_0^2=\kappa^2\,\vec\xi^{\top}M(\nu)^2\vec\xi$ of
      \Cref{thm:general-rounding}.
\end{enumerate}
\end{theorem}

\begin{proof}
By Lemma~\ref{lem:alpha-selfadjoint} both $\Ham_R$ and $\Ham_g$ are self-adjoint with
common core $\Sch(\R^{\dCl},V)$.  The argument of \Cref{thm:full-operator}, applied to
the Hermitian spatial multipliers $h_R(\vec\xi)\to h_g(\vec\xi)$ with domination
$\norm{h_R(\vec\xi)}\le\kappa\norm{A}\,|\vec\xi|$, gives $\Ham_R u\to\Ham_g u$ in
$L^2$ for every $u\in\Sch$.  Strong convergence on a common core of self-adjoint
operators implies strong resolvent convergence by the Trotter--Kato criterion
\cite[Thm.~VIII.25(a)]{reedsimon1}.  Item (i) is the Trotter--Kato theorem
relating strong resolvent convergence to strong convergence of the unitary groups,
uniformly on compact time intervals \cite[Thm.~VIII.21]{reedsimon1}; item (ii)
is the bounded-continuous functional-calculus form of strong resolvent convergence.  For (iii), $h_g(\vec
\xi)$ is Hermitian with $h_g(\vec\xi)^2=\kappa^2\vec\xi^{\top}M^2\vec\xi\,\one$ by
Lemma~\ref{lem:alpha-selfadjoint}, so its eigenvalues are $\pm\kappa|M\vec\xi|$; as
$\vec\xi$ ranges over $\R^{\dCl}$ these fill $\R$, with purely absolutely continuous
spectrum by the standard Fourier-multiplier (limiting absorption) argument, and
the energy relation $\xi_0=\pm\kappa|M\vec\xi|$ is the null condition
$g^{\mu\nu}\xi_\mu\xi_\nu=0$.
\end{proof}

\begin{theorem}[Frame universality of the nondegenerate flat endpoint]
\label{thm:frame-universality}
Let \(M\in\operatorname{Sym}_{\dCl}(\R)\) be positive definite and set
\(V_M=\kappa M\).  Define
\[
  (U_M\psi)(y):=(\det V_M)^{1/2}\psi(V_My),
  \qquad
  \psi\in L^2(\R^{\dCl},V).
\]
Then \(U_M\) is unitary, acts trivially on the spinor factor, and
\[
  U_M\,\bigl[\kappa A^iM^{ij}(-i\partial_j)\bigr]\,U_M^{-1}
  =A^i(-i\partial_i)=:\Ham_* .
\]
Consequently every nondegenerate constant-coefficient endpoint
\[
  g_M^{-1}=\operatorname{diag}(1,-\kappa^2M^2)
\]
is spatially linearly isometric to standard Minkowski space, and its Hamiltonian
is unitarily equivalent to the standard isotropic massless Dirac Hamiltonian.
The spatial change of variables leaves the Krein-odd timelike generator,
fundamental symmetry, Krein form, and single-time-slot conclusion unchanged.
\end{theorem}

\begin{proof}
With \(x=V_My\), one has \(dx=(\det V_M)\,dy\), so \(U_M\) is unitary.
For \(\phi=U_M\psi\),
\[
 (-i\partial_{y_k}\phi)(y)
 =(\det V_M)^{1/2}(V_M)_{jk}
   (-i\partial_{x_j}\psi)(V_My).
\]
Since \(V_M=\kappa M\) is symmetric,
\[
 A^k(V_M)_{jk}
 =\kappa A^kM^{kj},
\]
and hence \(\Ham_*U_M=U_M\Ham_g\) on
\(C_c^\infty(\R^{\dCl},V)\).  Essential self-adjointness on this common core
extends the identity to the closures.  On covectors, the dual change
\(\eta=V_M^\top\xi\) sends
\[
 \xi_0^2=\kappa^2\xi^\top M^2\xi
 \quad\text{to}\quad
 \xi_0^2=|\eta|^2.
\]
Because \(U_M\) acts only on the spatial argument, it commutes with the constant
spinor-space operators carrying the Krein and temporal data.
\end{proof}

\begin{remark}[Scope of frame universality]
\label{rem:frame-universality-scope}
The theorem removes constant anisotropy from the isometry class of the
\emph{isolated, single-sector, flat, free Dirac endpoint}.  It does not assert
that every microscopic tensor becomes physically irrelevant once additional
structures are retained.  If two matter or propagation sectors carry
inequivalent positive tensors \(M_1\) and \(M_2\), one spatial frame can
normalise one of them but not generally both; the relative tensor
\(M_1^{-1}M_2\) can then be invariant data.  Likewise, interactions,
higher-order dispersion, boundaries, or a separately fixed microscopic frame
may retain relative anisotropy.  No such independent spatial tensor is part of
the flat kinematic endpoint proved here.
\end{remark}

\begin{corollary}[Explicit $(3+1)$ free Dirac endpoint]
\label{cor:31-endpoint}
Take $\dCl=3$, $\kappa=3$, and $\nu_R\to\nu_{\mathrm{unif}}$ on $\Sphere^{2}$.  Then
$M(\nu_R)\to\tfrac13 I_3$, the continuum Hamiltonian is the free massless
$(3+1)$-dimensional Dirac Hamiltonian
\[
  \Ham_g=\sum_{j=1}^{3}A^j(-i\partial_j),\qquad \{A^i,A^j\}=2\delta^{ij},
\]
and $e^{-it\Ham_R}$ converges strongly, uniformly on compact time intervals, to
the free Dirac propagator $e^{-it\Ham_g}$.  The emergent light cone is the round
cone $\xi_0^2=|\vec\xi|^2$, i.e.\ a $(3+1)$-Minkowski signature, consistent with the dimension count $1+\dCl=4$.
\end{corollary}

\begin{proof}
By \Cref{thm:model-rounding}\,(i), $M(\nu_R)\to\tfrac13 I_3$; with $\kappa=3$,
$\kappa M(\nu)=I_3$, so $h_g(\vec\xi)=A^j\xi_j$ and $\Ham_g=A^j(-i\partial_j)$.
The conclusions are \Cref{thm:operator-limit} specialised to this case.
\end{proof}

\subsubsection{Scope and sharpness}

The preceding theorem gives the correct global convergence statement.  The next
remarks explain why it should not be upgraded to global norm-resolvent
convergence for the symmetric finite-difference scheme, and why a norm estimate
does hold after imposing a fixed momentum cutoff.

\begin{remark}[Strong, not norm, resolvent convergence]
\label{rem:no-norm-resolvent}
The convergence in \Cref{thm:operator-limit} is strong resolvent and no better:
for symmetric finite differences, norm resolvent convergence of lattice Dirac
operators to their continuum counterpart \emph{fails} in spatial dimension
$\dCl\ge2$, even in the elliptic (Riemannian) case.  For Euclidean lattice Dirac operators this is established by
Cornean--Garde--Jensen and by Schmidt--Umeda \cite{corneangardejensen2022,schmidtumeda2023}; the
same obstruction applies here.  Norm resolvent convergence can be restored by a
non-standard modification of the mesh operator (a mass-type correction of Wilson type
\cite{wilson1977}), which we do not pursue.
\end{remark}

\begin{remark}[Why norm-resolvent convergence of the mesh limit is unavailable]
\label{rem:norm-resolvent-status}
It is natural to ask whether the strong-resolvent convergence of
\Cref{thm:operator-limit} can be upgraded to norm-resolvent convergence.  For
the symmetric finite-difference scheme used here this is \emph{provably false} in
spatial dimension $\dCl\ge2$: the obstruction of
\cite{corneangardejensen2022,schmidtumeda2023}, recalled in
\Cref{rem:no-norm-resolvent}, rules it out already in the flat Euclidean case,
and our signed scheme inherits it.  The only norm-resolvent statement the
construction supports is the \emph{stability} estimate of \Cref{thm:stability}:
it compares two continuum operators at fixed mesh-free limit, where both lie in a
common uniformly elliptic class, and is therefore not a discretisation limit at
all (\Cref{rem:stability-stronger}).  Restoring norm-resolvent convergence of the
\emph{mesh} approximation would require a non-standard mesh correction of
Wilson/mass type \cite{wilson1977}, which modifies the renewal difference operator away from the
signed reset dynamics; we do not adopt it.  This item is thus recorded as a
boundary of the method, not a gap to be filled.
\end{remark}

\begin{theorem}[Band-limited norm-resolvent convergence]
\label{thm:band-norm-resolvent}
Let $\Ham_R$ and $\Ham_g$ be the self-adjoint Dirac Hamiltonians of
\Cref{thm:operator-limit}, with Hermitian Fourier multipliers
$h_R(\vec\xi)=\kappa\sum_a p_a^{(R)}A(\hat\theta_a^{(R)})s_{\hat\theta_a^{(R)},R}
(\vec\xi)$ and $h_g(\vec\xi)=\kappa\,A^iM(\nu)^{ij}\xi_j$ of
\Cref{def:reset-differences,thm:general-rounding}.  For $\Lambda_0>0$ let
$P_{\Lambda_0}$ be the momentum band projection on $L^2(\R^{\dCl},V)$ with multiplier
$\mathbf 1_{\{|\vec\xi|\le\Lambda_0\}}\otimes I_V$; being a Fourier multiplier it
commutes with $\Ham_R$ and $\Ham_g$.  Set
\[
  \chi_R(\Lambda_0):=\sup_{|\vec\xi|\le\Lambda_0}\bigl\lVert h_R(\vec\xi)
  -h_g(\vec\xi)\bigr\rVert .
\]
For the symmetric reset differences one has the explicit second-order bound
\[
  \chi_R(\Lambda_0)\ \le
  \kappa\norm{A}\Bigl(\tfrac16\,h_R^{2}\,\Lambda_0^{3}
  \;+\;\norm{M(\nu_R)-M(\nu)}\,\Lambda_0\Bigr)
  \xrightarrow[R\to\infty]{}0 .
\]
Then for every $z\in\C\setminus\R$,
\[
  \boxed{\ \bigl\lVert\bigl[(\Ham_R-z)^{-1}-(\Ham_g-z)^{-1}\bigr]
  P_{\Lambda_0}\bigr\rVert\ \le\ \frac{\chi_R(\Lambda_0)}{|\operatorname{Im}z|^{2}}
  \ =\ O(h_R^{2})\ } ,
\]
the $O(h_R^2)$ rate holding whenever $M(\nu_R)=M(\nu)$ \textup{(}e.g.\ an exactly
isotropic frame at every $R$, as in \Cref{ex:3plus1}\textup{)}.  Consequently,
for every bounded continuous $\varphi$,
\[
  \bigl\lVert\bigl[\varphi(\Ham_R)-\varphi(\Ham_g)\bigr]P_{\Lambda_0}
  \bigr\rVert\xrightarrow[R\to\infty]{}0 .
\]
Thus on any fixed momentum band the convergence improves from strong resolvent
\textup{(}\Cref{thm:operator-limit}\textup{)} to \emph{norm} resolvent, at
second order in the mesh.
\end{theorem}

\begin{proof}
Since $P_{\Lambda_0}$, $\Ham_R$, $\Ham_g$ are all Fourier multipliers they
commute, and on the range of $P_{\Lambda_0}$ the resolvent difference is the
multiplication operator with symbol
$(h_R(\vec\xi)-z)^{-1}-(h_g(\vec\xi)-z)^{-1}$ restricted to
$|\vec\xi|\le\Lambda_0$.  The norm of a multiplier operator is the essential
supremum of the pointwise matrix norm, so
\[
  \bigl\lVert[(\Ham_R-z)^{-1}-(\Ham_g-z)^{-1}]P_{\Lambda_0}\bigr\rVert
  =\sup_{|\vec\xi|\le\Lambda_0}
   \bigl\lVert(h_R-z)^{-1}-(h_g-z)^{-1}\bigr\rVert .
\]
By the resolvent identity $(h_R-z)^{-1}-(h_g-z)^{-1}
=(h_R-z)^{-1}(h_g-h_R)(h_g-z)^{-1}$, and since $h_R,h_g$ are Hermitian,
$\norm{(h_\bullet-z)^{-1}}\le|\operatorname{Im}z|^{-1}$; hence the pointwise norm
is $\le|\operatorname{Im}z|^{-2}\norm{h_R-h_g}$, giving the boxed bound after
taking the supremum.

For $\chi_R$: each scalar component obeys $|s_{\hat\theta,R}(\vec\xi)
-\hat\theta\!\cdot\!\vec\xi|=h_R^{-1}|\sin(h_R\hat\theta\!\cdot\!\vec\xi)
-h_R\hat\theta\!\cdot\!\vec\xi|\le h_R^{-1}\tfrac16|h_R\hat\theta\!\cdot\!
\vec\xi|^{3}\le\tfrac16h_R^{2}|\vec\xi|^{3}\le\tfrac16h_R^2\Lambda_0^3$ on the
band, using $|\sin t-t|\le|t|^3/6$ and $|\hat\theta\!\cdot\!\vec\xi|\le|\vec\xi|$;
summing the Hermitian matrices $A(\hat\theta_a)$ against $p_a^{(R)}$ contributes
$\kappa\norm{A}$.  The replacement of $M(\nu_R)$ by $M(\nu)$ in the limit symbol
adds $\kappa\norm{A}\,\norm{M(\nu_R)-M(\nu)}\,|\vec\xi|\le\kappa\norm{A}
\norm{M(\nu_R)-M(\nu)}\Lambda_0$.  This proves the displayed bound on
$\chi_R(\Lambda_0)$, which tends to $0$ since $h_R\downarrow0$ and
$M(\nu_R)\to M(\nu)$ \textup{(}\Cref{thm:general-rounding}\textup{)}.

For the functional-calculus statement, $\varphi$ is uniformly continuous on the
compact set $\{h_g(\vec\xi):|\vec\xi|\le\Lambda_0\}$ enlarged by an $\chi_R$
neighbourhood; uniform convergence $h_R\to h_g$ on the band then gives
$\sup_{|\vec\xi|\le\Lambda_0}\norm{\varphi(h_R(\vec\xi))-\varphi(h_g(\vec\xi))}
\to0$, which is the multiplier norm of $[\varphi(\Ham_R)-\varphi(\Ham_g)]
P_{\Lambda_0}$.
\end{proof}

\begin{remark}[Momentum band, not energy window; consistency with the no-go]
\label{rem:band-not-energy}
Two points fix the scope of \Cref{thm:band-norm-resolvent} and reconcile it with
\Cref{rem:no-norm-resolvent,rem:norm-resolvent-status}.

\emph{(a) The band is in momentum, and this is essential.}  The norm-resolvent
failure in $\dCl\ge2$ is caused by lattice fermion doublers
\cite{nielsenninomiya1981a,nielsenninomiya1981b}: the discrete symbol
$h_R(\vec\xi)$ has spurious near-zeros at high momentum
$|\vec\xi|\sim\pi/h_R$, where $h_R(\vec\xi)\approx0$ while $h_g(\vec\xi)$ is
large.  At such a covector the resolvent symbols differ by
$\bigl\lVert(h_R-z)^{-1}-(h_g-z)^{-1}\bigr\rVert\to|z|^{-1}\neq0$, so the
full-line supremum does not vanish.  These doublers sit at \emph{low energy} but
\emph{high momentum}; hence an energy-window truncation $\varphi(\Ham_R)$ with
$\varphi$ supported near $0$ does \emph{not} excise them, whereas the fixed
momentum band $P_{\Lambda_0}$ does, since $\pi/h_R\to\infty$.  This is why the
theorem is stated with a momentum projection and why the analogous energy-window
statement is false.

\emph{(b) No contradiction with the no-go.}  Norm-resolvent convergence of the
full operator is exactly the statement that
$\sup_{\vec\xi\in\R^{\dCl}}\lVert(h_R-z)^{-1}-(h_g-z)^{-1}\rVert\to0$; it fails because
of the doubler region $|\vec\xi|\sim\pi/h_R$, precisely as in
\cite{corneangardejensen2022,schmidtumeda2023}.  \Cref{thm:band-norm-resolvent}
restores norm convergence by removing that region, and says nothing about it: the
symmetric scheme is second-order consistent \emph{uniformly on compact momentum
sets} but not \emph{uniformly in $\vec\xi$}.  The band-limited statement is thus
the sharp positive complement of \Cref{rem:norm-resolvent-status}: full
norm-resolvent convergence needs a Wilson/mass mesh correction
\cite{wilson1977} \textup{(}there discussed and not adopted\textup{)}, while on any fixed momentum
band the unmodified renewal scheme already converges in norm resolvent at order
$h_R^2$.
\end{remark}


\begin{remark}[Scope of the flat operator theorem]
\label{rem:operator-scope}
The following are proved here, in the flat constant-coefficient regime:
graph convergence of the full signed Dirac (\Cref{thm:full-operator}) and strong
resolvent, propagator, and spectral-projection convergence of the self-adjoint
Dirac Hamiltonian (\Cref{thm:operator-limit}).  \Cref{sec:curved-stability}
upgrades the Hamiltonian convergence to a variable-coefficient adiabatic setting
and proves stability under perturbation of the reset statistics.  The curved refinement now reconstructs the
Levi--Civita spin connection from the renewal-generated coframe by the discrete
torsion-free Cartan rule and proves global strong-resolvent convergence through
subprincipal order.  What is still not claimed is different:
\textup{(a)} derivation of the exchange-symmetric diamond phase from bare channel associativity \textup{(}closure itself is derived from that phase in \Cref{thm:stationary-exchange}\textup{)}; \textup{(b)} a fully discrete spacetime unitary with the timelike
direction also discretised, which would require a separate stability/CFL
hypothesis; \textup{(c)} full norm-resolvent convergence of the mesh
approximation, which fails by \Cref{rem:no-norm-resolvent}, although fixed
momentum bands converge in norm resolvent in the flat case by
\Cref{thm:band-norm-resolvent}; and \textup{(d)} curvature dynamics, field
equations, or matter coupling.  As in
\Cref{rem:isotropisation-scope}, the Lorentzian signature is carried by the signed/Krein sector and the derived temporal generator, not by positive averaging.
\end{remark}

\begin{remark}[Relation to the modular signed Dirac]
\label{rem:modular-relation}
The continuum operators above are the coarse-grained, spatially isotropic
counterparts of the modular signed Dirac $D_-=J_\chi e^{\beta N}+B$ of
\Cref{thm:signed-dirac}: the modular factor governs the timelike generator, while
the equidistributed reset directions of Definition~\ref{def:reset-differences} furnish the
spatial Dirac matrices $A^i$.  The present theorems show that, after rescaling,
the spatial sector of the signed renewal dynamics converges to a genuine
self-adjoint Lorentzian Dirac evolution.
\end{remark}

\subsubsection{Self-averaging from ergodic reset statistics}
\label{subsec:self-averaging}

The preceding operator theorem is already anisotropic: it requires only
convergence of the second-moment tensors to a positive-definite limit.  Ergodicity
therefore need not create a rotationally invariant stationary law.  It is enough
that the stationary directions span the spatial slot.  The limiting tensor then
defines a flat vielbein, and \Cref{thm:frame-universality} removes its constant
frame anisotropy exactly.

\begin{theorem}[Self-averaged flat Dirac endpoint under nondegeneracy]
\label{thm:self-averaging}
Let $(\hat\theta_a)_{a\ge1}$ be a stationary ergodic sequence of spatial reset
directions in $\Sphere^{\dCl-1}$ on a probability space
$(\Omega,\mathcal F,\mathbb P)$, with one-dimensional marginal law $\zeta_*$.
Let $m_R\to\infty$ and $h_R\downarrow0$, form
\[
 \nu_R=\frac1{m_R}\sum_{a=1}^{m_R}\delta_{\hat\theta_a},
 \qquad
 M_*:=\int_{\Sphere^{\dCl-1}}
       \hat\theta\hat\theta^\top\,d\zeta_*(\hat\theta),
\]
and let $\Ham_R$ be the discrete renewal Dirac Hamiltonian of
\Cref{def:reset-differences} with equal weights \(p_a^{(R)}=1/m_R\).
Assume only
\[
  M_*\succ0,
\]
equivalently that \(\supp\zeta_*\) spans \(\R^{\dCl}\), and fix
\(\kappa>0\).  Then for \(\mathbb P\)-almost every realisation:
\begin{enumerate}[label=\textup{(\roman*)},leftmargin=2.4em]
\item \(\nu_R\to\zeta_*\) weak-$*$ and \(M(\nu_R)\to M_*\);
\item the random signed principal symbols converge pointwise to the
      deterministic Lorentzian symbol with inverse metric
      \[
        g_*^{-1}=\operatorname{diag}(1,-\kappa^2M_*^2),
      \]
      of signature \((1,\dCl)\);
\item \(\Ham_R\) converges in strong resolvent sense to
      \[
        \Ham_g[M_*]=\kappa A^iM_*^{ij}(-i\partial_j),
      \]
      and \(e^{-it\Ham_R}\to e^{-it\Ham_g[M_*]}\) strongly, uniformly for
      \(t\) in compact intervals;
\item on every fixed momentum band \(P_{\Lambda_0}\), the resolvents converge in
      norm with the estimate
      \[
      \begin{split}
       &\bigl\|
       [(\Ham_R-z)^{-1}-(\Ham_g[M_*]-z)^{-1}]P_{\Lambda_0}
       \bigr\|\\
       &\qquad\le
       \frac{\kappa\norm{A}}{|\operatorname{Im}z|^2}
       \left(
       \frac16h_R^2\Lambda_0^3+
       \norm{M(\nu_R)-M_*}\Lambda_0
       \right)\longrightarrow0;
      \end{split}
      \]
\item the fixed unitary \(U_{M_*}\) of
      \Cref{thm:frame-universality} intertwines the limit with the standard
      massless Dirac Hamiltonian:
      \[
        U_{M_*}\Ham_g[M_*]U_{M_*}^{-1}
        =A^i(-i\partial_i).
      \]
\end{enumerate}
Thus stationary ergodicity and the nondegeneracy already required for
Lorentzian signature suffice almost surely to produce the standard flat
Minkowski Dirac isometry class; rotational invariance of \(\zeta_*\) is not an
additional hypothesis.
\end{theorem}

\begin{proof}
For each \(i,j\), the function
\(\hat\theta\mapsto\hat\theta_i\hat\theta_j\) is bounded and continuous on the
compact sphere.  Birkhoff's ergodic theorem gives
\[
 M(\nu_R)^{ij}
 =\frac1{m_R}\sum_{a=1}^{m_R}
   (\hat\theta_a)_i(\hat\theta_a)_j
 \longrightarrow
 \int\hat\theta_i\hat\theta_j\,d\zeta_*
 =M_*^{ij}
\]
almost surely.  Applying the same theorem to a countable
convergence-determining family of continuous functions gives
\(\nu_R\to\zeta_*\) weak-$*$ almost surely.  Fix a realisation in this
probability-one event.  Since \(M_*\succ0\), the deterministic anisotropic
operator limit \Cref{thm:operator-limit} applies pathwise and proves
(ii)--(iii).  The fixed-band estimate is exactly
\Cref{thm:band-norm-resolvent} with \(M(\nu)=M_*\), proving (iv).
Finally, (v) is \Cref{thm:frame-universality}.
\end{proof}

\begin{proposition}[Literal isotropy is not implied by ergodic renewal dynamics]
\label{prop:isotropy-not-derived}
For every \(\dCl\ge2\) there exists a stationary ergodic finite reset law with
\(M_*\succ0\) but \(M_*\notin\R I_{\dCl}\).
\end{proposition}

\begin{proof}
Choose the directions \(e_1,\ldots,e_{\dCl}\) independently with probabilities
\(p_1,\ldots,p_{\dCl}>0\), not all equal.  The resulting i.i.d.\ sequence is
stationary and ergodic, while
\[
  M_*=\sum_{i=1}^{\dCl}p_i e_ie_i^\top
     =\operatorname{diag}(p_1,\ldots,p_{\dCl})
\]
is positive definite and non-scalar.  Such a finite label process can be
coupled to any primitive finite family of unital completely positive reset
channels, so it lies within the renewal class considered here.
\end{proof}

\begin{remark}[What self-averaging removes, and what it does not]
\label{rem:self-averaging-honest}
Three qualifications remain.  \emph{(a)}~Self-averaging produces the
deterministic tensor \(M_*\); it does not force \(M_*\) to be scalar, and
\Cref{prop:isotropy-not-derived} shows that no such implication follows from
stationarity and ergodicity.  Scalarity is unnecessary because the isolated
flat endpoint is frame-universal.  \emph{(b)}~Global convergence is strong
resolvent only; the doubler obstruction of
\Cref{rem:no-norm-resolvent} persists, while fixed-band norm-resolvent
convergence holds pathwise by \Cref{thm:band-norm-resolvent}.
\emph{(c)}~The conclusion concerns the single-sector flat free endpoint.
Relative anisotropy may remain when independent sectors carry inequivalent
second-moment tensors, as explained in
\Cref{rem:frame-universality-scope}.  The balancing layer of
\Cref{subsec:downstream-interfaces} is not needed to impose isotropy; its role
is instead to determine which nondegenerate stationary tensor is dynamically
selected.
\end{remark}

% =====================================================================




% =====================================================================
\subsection{Channel plaquettes and Levi--Civita reconstruction}
\label{sec:cix}
\label{sec:cxii}

The renewal coframe determines the principal metric, while the composed-channel
plaquette separates translational torsion from rotational curvature.  Covariance
of transported second moments yields isometric comparison and hence metric
compatibility.  In the exchange-symmetric stationary diamond phase, the
translational first moment vanishes, and the discrete fundamental theorem
selects the unique Levi--Civita connection without suppressing rotational
holonomy.

\subsubsection{Setup and the two defects}

Fix a smooth nondegenerate coframe field $E(x)=(e^a_j(x))\in GL(\dCl)$, $C^4_b$,
with metric $g_{ij}=\delta_{ab}e^a_ie^b_j=(E^\top E)_{ij}$ and frame Gram
$P:=EE^\top\succ0$.  Let $h\downarrow0$ be the mesh and $\delta_i$ the symmetric
difference (so $\delta_i e^a_j=\partial_i e^a_j+O(h^2)$).  A discrete connection is
an assignment $\gamma_{i\,ab}(x)$; its near-identity frame transport is
$U_i(x)=\exp(-h\,\gamma_i(x))$ with $(\gamma_i)^a{}_b=\delta^{ac}\gamma_{i\,cb}$.

Two tensors measure the failure of the two structural requirements:
\[
  Q_{i\,ab}:=\gamma_{i\,ab}+\gamma_{i\,ba}
  \quad\text{(metricity defect)},
  \qquad
  T^a_{ij}:=\delta_i e^a_j-\delta_j e^a_i
            +\gamma_i{}^a{}_b e^b_j-\gamma_j{}^a{}_b e^b_i
  \quad\text{(discrete torsion)}.
\]
Metric compatibility is $Q\equiv0$ (the connection is $\so(\dCl)$-valued);
torsion-freeness is $T\equiv0$.

\subsubsection{The discrete fundamental theorem}

Write $C^a_{ij}:=\delta_i e^a_j-\delta_j e^a_i$ for the coframe curl and
$C_{abc}:=\delta_{ad}\,e_b{}^i e_c{}^j C^d_{ij}$ (antisymmetric in $b,c$), with
$e_a{}^j$ the inverse coframe.

\begin{theorem}[Discrete metric / torsion-free uniqueness]
\label{thm:fundamental}
For every nondegenerate coframe there is a \emph{unique} connection
with $Q\equiv0$ and $T\equiv0$.  In frame components $\gamma_{d;ae}:=
e_d{}^i\gamma_{i\,ae}$ it is the discrete Koszul formula
\[
  \gamma_{d;ae}=\tfrac12\bigl(C_{dea}-C_{ade}-C_{ead}\bigr),
\]
which satisfies $\gamma_{d;ae}=-\gamma_{d;ea}$.  Moreover $\gamma_{i\,ab}
=\omega^{\mathrm{LC}}_{i\,ab}+O(h^2)$ locally uniformly, where
$\omega^{\mathrm{LC}}$ is the Levi--Civita spin connection of $g$.
\end{theorem}

\begin{proof}
\emph{Existence.}  Contracting $T^a_{ij}=0$ with $e_d{}^ie_e{}^j$ gives the frame
equation $C_{ade}+\gamma_{d;ae}-\gamma_{e;ad}=0$.  Using $C_{abc}=-C_{acb}$,
\[
  \gamma_{d;ae}-\gamma_{e;ad}
  =\tfrac12(C_{dea}-C_{ade}-C_{ead})-\tfrac12(C_{eda}-C_{aed}-C_{dae})
  =-C_{ade},
\]
since $C_{dea}=-C_{dae}$ and $C_{ead}=-C_{eda}$ cancel pairwise.  Antisymmetry in
$a,e$ ($Q=0$) is immediate: swapping $a\leftrightarrow e$ negates each term.

\emph{Uniqueness.}  If $\gamma,\gamma'$ both satisfy $Q=0,\ T=0$, the difference
$\Delta_{d;ae}$ is metric, $\Delta_{d;ae}=-\Delta_{d;ea}$, and obeys the
homogeneous torsion equation $\Delta_{d;ae}=\Delta_{e;ad}$.  Then
\[
  \Delta_{d;ae}=\Delta_{e;ad}=-\Delta_{e;da}=-\Delta_{a;de}
  =\Delta_{a;ed}=\Delta_{d;ea}=-\Delta_{d;ae},
\]
so $\Delta=0$.  Nondegeneracy of the coframe carries this back to coordinate
indices.

\emph{Convergence.}  $\delta_i e^a_j=\partial_i e^a_j+O(h^2)$ for $C^4_b$
coframes, so $C_{abc}=C^{\mathrm{cont}}_{abc}+O(h^2)$; the formula is linear in
$C$ with coefficients bounded by the uniformly bounded frame and inverse frame,
and its continuum value is the Levi--Civita spin connection, whence
$\gamma=\omega^{\mathrm{LC}}+O(h^2)$.
\end{proof}
\subsubsection{Isometric comparison forces metric compatibility}

\begin{assumption}[Covariance of transported second moments]
\label{ass:second-moment-covariance}
Let \(M_x,M_y\succ0\) be the spatial second moments of neighbouring local
direction-record laws, and let
\(T_e:\R^{\dCl}\to\R^{\dCl}\) be the invertible transport of those records along
an edge \(e:x\to y\).  Their transported quadratic statistics agree:
\[
   T_e M_x T_e^{\mathsf T}=M_y.
\]
The stable-record determinant orientation is preserved along the selected
oriented branch.
\end{assumption}

\begin{proposition}[Second-moment covariance derives isometric comparison]
\label{prop:second-moment-isometry}
Under \Cref{ass:second-moment-covariance}, the comparison in the orthonormal
frames selected by the two moments,
\[
   U_e:=M_y^{-1/2}T_eM_x^{1/2},
\]
lies in \(SO(\dCl)\).  Thus isometric comparison is a consequence of covariance
of renewal second moments, rather than an independent metric postulate.
\end{proposition}

\begin{proof}
The covariance identity gives
\[
   U_eU_e^{\mathsf T}
   =M_y^{-1/2}T_eM_xT_e^{\mathsf T}M_y^{-1/2}
   =I.
\]
Hence \(U_e\in O(\dCl)\).  Preservation of the selected determinant orientation
gives \(\det U_e>0\), and therefore \(U_e\in SO(\dCl)\).
\end{proof}

\begin{proposition}[Isometric comparison $\Leftrightarrow$ metric compatibility]
\label{prop:iso-metric}
A near-identity transport $U_i=\exp(-h\gamma_i)$ lies in $SO(\dCl)$
to first order in $h$ iff $\gamma_i\in\so(\dCl)$, i.e.\ $Q\equiv0$.
Consequently \Cref{ass:second-moment-covariance,prop:second-moment-isometry}
imply metric compatibility of the induced connection.
\end{proposition}

\begin{proof}
$\exp(-h\gamma_i)\in SO(\dCl)$ for all small $h$ iff $\gamma_i^\top=-\gamma_i$,
i.e.\ $\gamma_{i\,ab}=-\gamma_{i\,ba}$, i.e.\ $Q_{i\,ab}=0$.  An $\so$-valued
connection annihilates the constant frame metric $\delta_{ab}$ under transport,
which is metric compatibility.
\end{proof}

The coframe is \(e_a{}^j=\kappa M^{aj}\), with \(M\) the renewal second moment.
The proposition shows that covariance of the transported record law canonically
produces its orthogonal frame comparison.  Among such comparisons, the
orthogonal Procrustes factor remains the distinguished least-distortion
representative.

\begin{remark}[The Procrustes minimiser and why it is not the whole story]
\label{rem:procrustes}
The unique distortion-minimising isometry is the orthogonal polar
factor: among $R\in SO(\dCl)$, the minimiser of $\|R\,E(x)-E(x+he_i)\|_{\HS}$ is
the orthogonal factor of $E(x+he_i)E(x)^\top$.  At first order, writing
$R=\one-h\omega+O(h^2)$ with $\omega\in\so$ and $E(x+he_i)=E+hE_{,i}+O(h^2)$,
\[
  \|R E-E'\|_{\HS}^2=h^2\|\omega E+E_{,i}\|_{\HS}^2+O(h^3),
\]
and the minimiser solves the Sylvester equation $P\omega+\omega P=2\,
\Asym(EE_{,i}^\top)$ with $P=EE^\top\succ0$, which has a unique antisymmetric
solution.  This selects a \emph{particular} metric connection -- but, as the next
example shows, it is generally \emph{not} torsion-free, hence not Levi--Civita.
Thus least distortion delivers $Q=0$ (metric compatibility), not $T=0$.  The
two requirements are genuinely independent and must both be imposed.
\end{remark}

\subsubsection{Reset channels as Euclidean motions of the developed frame}

A coarse-grained elementary reset transports the renewal frame and advances the
developed position.  The natural carrier of these two data is the Euclidean
motion group $\Eucl=\SO(\dCl)\ltimes\R^\dCl$, acting on the developed
position $q\in\R^\dCl$ (recorded in the orthonormal frame) by
$(R,t)\cdot q=Rq+t$, with composition $(R_2,t_2)(R_1,t_1)=(R_2R_1,\,R_2t_1+t_2)$.

\begin{definition}[Elementary reset motion and first moment]
\label{def:reset-motion}
The coarse-grained elementary reset channel in coordinate direction $i$ at base
$x$, $\Phi_i^{(h)}(x)$, transports the local direction records by
\(T_i(x)\).  Under \Cref{ass:second-moment-covariance}, its
\emph{developed-frame data} are the Euclidean motion
\[
  g_i^{(h)}(x)=\bigl(R_i(x),\,h\,e_i(x)\bigr)\in\Eucl,
  \qquad
  R_i(x)=M_{x+he_i}^{-1/2}T_i(x)M_x^{1/2}
        =\exp\!\bigl(-h\,\gamma_i(x)\bigr)\in\SO(\dCl),
\]
where $e_i(x)=(e^a_i(x))_a$ is the coframe column \textup{(}the mean reset
displacement in the frame, a renewal datum\textup{)} and $\gamma_i(x)$ is the
\textup{(}a priori unknown\textup{)} internal frame transport generated by the
channel.  The \emph{first spatial moment} of a channel with data $(R,t)$ is its
translation part,
\[
  \mone\bigl(\Phi\bigr):=t\in\R^\dCl,
\]
i.e.\ the mean developed displacement of a localised renewal state.
\end{definition}

The displacement \(e_i\) is fixed by the renewal second moment, while the raw
record transport \(T_i\) is supplied by the channel.  The rotation \(R_i\) and
hence the metric-compatible part of the connection are not adopted separately:
they are the orthonormalised transport derived in
\Cref{prop:second-moment-isometry}.  An arbitrary Stinespring dilation need not
satisfy the covariance hypothesis; the load-bearing requirement is agreement of
the transported quadratic record statistics, together with determinant
orientation preservation.

\subsubsection{Translational and rotational parts of the plaquette defect}

For $i\neq j$ form the two opposite plaquette histories and their channel defect
\[
  \Pcal_{ij}^{(h)}(x)
  :=\Phi_j^{(h)}(x+he_i)\,\Phi_i^{(h)}(x)
   -\Phi_i^{(h)}(x+he_j)\,\Phi_j^{(h)}(x),
\]
read through the developed-frame data, with translational projection
$\Pitr\Pcal_{ij}^{(h)}$ (difference of first moments) and rotational projection
$\Pisp\Pcal_{ij}^{(h)}$ (difference of frame rotations).

\begin{lemma}[Torsion is the translational defect, curvature the rotational defect]
\label{lem:plaquette-split}
With $g_i^{(h)}=(R_i,he_i)$ as in \Cref{def:reset-motion} and
$\gamma_i\in C^1$,
\[
  \boxed{\ \Pitr\Pcal_{ij}^{(h)}(x)=h^2\,T^a_{ij}(x)+O(h^3),\qquad
  \Pisp\Pcal_{ij}^{(h)}(x)=-h^2\,\Fcal_{ij}(x)+O(h^3),\ }
\]
where $T^a_{ij}=\delta_i e^a_j-\delta_j e^a_i+\gamma_i{}^a{}_b e^b_j
-\gamma_j{}^a{}_b e^b_i$ is the discrete torsion and
$\Fcal_{ij}=\partial_i\gamma_j-\partial_j\gamma_i+[\gamma_i,\gamma_j]$ the
discrete curvature.
\end{lemma}

\begin{proof}
\emph{Translation.}  In $\Eucl$, the composite $g_j(x+he_i)g_i(x)$ has translation
part $R_j(x+he_i)\,h e_i(x)+h e_j(x+he_i)$.  Expanding with
$R_j=\one-h\gamma_j+O(h^2)$ and $e_j(x+he_i)=e_j(x)+h\partial_i e_j(x)+O(h^2)$,
\[
  R_j(x+he_i)\,h e_i(x)=h e_i(x)-h^2\gamma_j(x)e_i(x)+O(h^3),
  \qquad
  h e_j(x+he_i)=h e_j(x)+h^2\partial_i e_j(x)+O(h^2),
\]
so the path-A translation is
$h e_i+h e_j+h^2\bigl(\partial_i e_j-\gamma_j e_i\bigr)+O(h^3)$.  The path-B
translation is the $i\leftrightarrow j$ image,
$h e_j+h e_i+h^2\bigl(\partial_j e_i-\gamma_i e_j\bigr)+O(h^3)$.  Subtracting, the
$O(h)$ terms cancel and the connection terms combine as
$-\gamma_j e_i-(-\gamma_i e_j)=\gamma_i e_j-\gamma_j e_i$, giving
\[
  \Pitr\Pcal^{(h)}_{ij}
  =h^2\Bigl[(\partial_i e^a_j-\partial_j e^a_i)
   +(\gamma_i{}^a{}_b e^b_j-\gamma_j{}^a{}_b e^b_i)\Bigr]+O(h^3)
  =h^2 T^a_{ij}+O(h^3),
\]
using $\delta_i e^a_j=\partial_i e^a_j+O(h^2)$.

\emph{Rotation.}  The composite rotations are $R_j(x+he_i)R_i(x)$ and
$R_i(x+he_j)R_j(x)$.  With $R_i=\one-h\gamma_i+O(h^2)$,
\[
  R_j(x+he_i)R_i(x)
  =\one-h\gamma_i-h\gamma_j-h^2\partial_i\gamma_j+h^2\gamma_j\gamma_i+O(h^3),
\]
and subtracting the $i\leftrightarrow j$ image gives
\[
  \begin{aligned}
  \Pisp\Pcal^{(h)}_{ij}
  &=h^2\bigl(-\partial_i\gamma_j+\partial_j\gamma_i+\gamma_j\gamma_i
  -\gamma_i\gamma_j\bigr)+O(h^3)\\
  &=-h^2\bigl(\partial_i\gamma_j-\partial_j\gamma_i+[\gamma_i,\gamma_j]\bigr)+O(h^3)
  =-h^2\Fcal_{ij}+O(h^3).
  \end{aligned}
\]
\end{proof}

\subsubsection{The torsion theorem}

\begin{assumption}[Exchange-covariant stationary diamond phase]
\label{ass:stationary-diamond}
For each \(i\neq j\), let \(\mathcal D_{ij}\) be a finite set of resolved
two-reset diamond records, equipped with an involution
\(\mathsf s_{ij}\) exchanging the histories \(ij\) and \(ji\).  Let
\(P_{ij}\) be a primitive Markov kernel satisfying
\[
   P_{ij}(\mathsf s_{ij}a,\mathsf s_{ij}b)=P_{ij}(a,b),
\]
and let \(d^{\rm tr}_{ij}:\mathcal D_{ij}\to\R^{\dCl}\) be the developed
translational defect, with
\[
   d^{\rm tr}_{ij}(\mathsf s_{ij}a)=-d^{\rm tr}_{ij}(a).
\]
Writing \(\pi_{ij}\) for the stationary law of \(P_{ij}\), the coarse-grained
translational plaquette defect has the expansion
\[
   \Pitr\Pcal_{ij}^{(h)}
   =h^2\sum_{a\in\mathcal D_{ij}}
        \pi_{ij}(a)d^{\rm tr}_{ij}(a)+O(h^3).
\]
No exchange constraint is imposed on the rotational record.
\end{assumption}

\begin{theorem}[Stationary exchange theorem]
\label{thm:stationary-exchange}
Under \Cref{ass:stationary-diamond}, the stationary law is exchange invariant,
\[
   \pi_{ij}(\mathsf s_{ij}a)=\pi_{ij}(a),
\]
and therefore
\[
   \sum_a\pi_{ij}(a)d^{\rm tr}_{ij}(a)=0,
   \qquad
   \boxed{\ \Pitr\Pcal_{ij}^{(h)}=O(h^3)\ }.
\]
The rotational part of the plaquette is not constrained and may retain
nonzero curvature.
\end{theorem}

\begin{proof}
Define \(\pi_{ij}^{\mathsf s}(a):=\pi_{ij}(\mathsf s_{ij}a)\).  Exchange
covariance of \(P_{ij}\) shows that \(\pi_{ij}^{\mathsf s}\) is stationary.
Primitivity gives a unique stationary law, hence
\(\pi_{ij}^{\mathsf s}=\pi_{ij}\).  Pairing every record with its exchanged
partner and using exchange oddness gives
\[
   \sum_a\pi_{ij}(a)d^{\rm tr}_{ij}(a)
   =-\sum_a\pi_{ij}(\mathsf s_{ij}a)d^{\rm tr}_{ij}(a)
   =-\sum_a\pi_{ij}(a)d^{\rm tr}_{ij}(a),
\]
so the sum vanishes.  The asserted closure follows from the expansion in
\Cref{ass:stationary-diamond}.  Since no exchange condition was placed on the
rotational record, no analogous cancellation is implied there.
\end{proof}

\begin{proposition}[One-step detailed balance does not imply torsion freedom]
\label{prop:detailed-balance-not-torsion}
Ordinary reversibility of the microscopic diamond chain is insufficient to
derive the conclusion of \Cref{thm:stationary-exchange}.
\end{proposition}

\begin{proof}
Take two records \(+,-\), exchanged by \(\mathsf s(+)= -\), and the primitive
kernel
\[
   P=\begin{pmatrix}3/4&1/4\\[2pt]1/2&1/2\end{pmatrix}.
\]
Its stationary law is \(\pi=(2/3,1/3)\), and it satisfies detailed balance
because \(\pi_+P_{+-}=\pi_-P_{-+}=1/6\).  For any nonzero vector \(v\), set
\(d^{\rm tr}(+)=v\) and \(d^{\rm tr}(-)=-v\).  Then
\[
   \mathbb E_\pi d^{\rm tr}=\frac13v\neq0.
\]
The kernel is not exchange covariant, since \(P_{++}\neq P_{--}\).  Thus
ordinary one-step detailed balance does not force the exchange-invariant
stationary law needed for translational cancellation.
\end{proof}

\begin{theorem}[Channel-composition torsion theorem]
\label{thm:channel-torsion}
Suppose the elementary reset channels admit the smooth Euclidean
first-moment data of \Cref{def:reset-motion} with \(\gamma_i\in C^1\), their
transported second moments satisfy \Cref{ass:second-moment-covariance}, and
their resolved two-reset records lie in the stationary diamond phase of
\Cref{ass:stationary-diamond}.  Then the renewal coframe satisfies
\[
  T^a_{ij}\equiv0.
\]
The induced connection is the unique discrete Levi--Civita connection of
\Cref{thm:fundamental} and converges to
\(\omega^{\mathrm{LC}}+O(h^2)\).  The curvature \(\Fcal_{ij}\) is
unconstrained and survives as rotational holonomy.
\end{theorem}

\begin{proof}
By \Cref{prop:second-moment-isometry,prop:iso-metric}, covariance of the
transported second moments gives \(Q\equiv0\).  By
\Cref{thm:stationary-exchange}, the translational plaquette satisfies
\(\Pitr\Pcal_{ij}^{(h)}=O(h^3)\).  Comparing this with
\Cref{lem:plaquette-split},
\[
   \Pitr\Pcal_{ij}^{(h)}=h^2T^a_{ij}+O(h^3),
\]
forces \(T^a_{ij}=0\).  The discrete fundamental theorem then gives the unique
metric-compatible torsion-free connection and its
\(\omega^{\mathrm{LC}}+O(h^2)\) limit.  The rotational projection remains
\(-h^2\Fcal_{ij}+O(h^3)\) and is untouched.
\end{proof}

\begin{remark}[Why only the translational projection, and why exact commutativity over-kills]
\label{rem:trans-only}
The decomposition of \Cref{lem:plaquette-split} is the whole point:
torsion is the translational component of the channel plaquette defect and
curvature is the rotational component.  Requiring the \emph{full} channel defect
to vanish, $\|\Pcal_{ij}^{(h)}\|_{\cb}=O(h^3)$ -- equivalently full predictive
equivalence of the two histories, or exact commutativity of the elementary
channels -- would force \emph{both} $T=0$ and $\Fcal=0$, collapsing to a flat
(teleparallel-trivial) geometry and discarding curvature and gauge information.
This is the case of the commuting-CP regime of the preceding construction, where opposite
plaquette histories have identical accumulated channels and the geometry is flat.
The curved reconstruction therefore needs exactly the weaker
\Cref{thm:stationary-exchange}: stationary predictive coincidence of
\emph{positions}, not full equivalence of \emph{channels}.
\end{remark}

\begin{remark}[Scope]
\label{rem:cxii-status}
Given transported second-moment covariance and the exchange-covariant
stationary diamond phase, the curved Cartan reconstruction is unconditional
within this class.  What is not claimed is that torsion freedom follows from
bare channel associativity or from ordinary one-step detailed balance:
\Cref{prop:detailed-balance-not-torsion} gives a counterexample.  Torsion-free
and torsional phases are both admissible; the exchange-symmetric translational
phase is the genuine selector, while the closure equation itself is derived by
\Cref{thm:stationary-exchange}.
\end{remark}

\subsection{Adiabatic curved Dirac limit, discrete Cartan reconstruction, and stability}
\label{sec:curved-stability}

The flat operator limit assumed that the second-moment tensor of the reset
statistics was constant in space.  We now allow the reset statistics to vary
slowly.  The signed principal symbol then determines the variable Lorentzian
inverse metric
\[
  g^{\mu\nu}(x)=\diag\bigl(1,-\kappa^2M(x)^2\bigr).
\]
A variable principal symbol alone does not determine the lower-order spin
transport.  The preceding plaquette theorem identifies torsion with the
translational channel defect.  Second-moment covariance gives metric
compatibility, while stationary exchange symmetry derives the vanishing
translational first moment.  The resulting connection is therefore the unique
discrete Levi--Civita connection.  We now prove convergence of the corresponding
covariant renewal Hamiltonians.

Throughout, the Krein--Clifford datum of
Definition~\ref{def:krein-clifford} is fixed.  Write
\(A^i=\gamma^0\gamma^i\) for the Hermitian Dirac matrices of
Lemma~\ref{lem:alpha-selfadjoint}, with
\(\{A^i,A^j\}=2\delta^{ij}\one\), and fix \(\kappa>0\).
For a mesh \(h_R\downarrow0\), translations and symmetric differences act on
the common Hilbert space \(L^2(\R^{\dCl},V)\), as in
Definition~\ref{def:reset-differences}.  This common-space realization is
sufficient for the strong-resolvent argument and avoids any additional
sampling or interpolation hypothesis.

\subsubsection{The emergent coframe and the discrete Cartan connection}

\begin{definition}[Smooth nondegenerate reset field]
\label{def:reset-field}
A \emph{smooth reset field} is a finite family of unit directions
\(\hat\theta_1,\ldots,\hat\theta_m\in\Sphere^{\dCl-1}\) and weights
\(p_a\in C_b^4(\R^{\dCl})\), \(p_a\ge0\), \(\sum_a p_a=1\), with second moment
\[
  M(x)=\sum_{a=1}^m p_a(x)\hat\theta_a\hat\theta_a^\top .
\]
We assume uniform nondegeneracy: \(M(x)\succeq m_0I_{\dCl}\) for some \(m_0>0\).
Set
\[
  e_a^{\,j}(x):=\kappa M^{aj}(x),\qquad
  g^{ij}(x)=\delta^{ab}e_a^{\,i}(x)e_b^{\,j}(x)
           =\kappa^2(M(x)^2)^{ij},
\]
and let \(e^a_j\) denote the inverse coframe.  Then
\(g^{\mu\nu}=\diag(1,-g^{ij})\) has Lorentzian signature
\((1,\dCl)\) at every point.
\end{definition}

Let
\[
  \delta_i^{(R)}f(x)
  =\frac{f(x+h_Re_i)-f(x-h_Re_i)}{2h_R}
\]
and define the discrete coframe structure coefficients
\[
  C_R{}^a{}_{bc}
  :=e_b^{\,i}e_c^{\,j}
    \bigl(\delta_i^{(R)}e^a_j-\delta_j^{(R)}e^a_i\bigr),
  \qquad C_{R\,abc}:=\delta_{ad}C_R{}^d{}_{bc}.
\]
Define the discrete Levi--Civita coefficients by the Koszul formula
\[
  \gamma_{d;ae}^{(R)}
  :=
  \frac12\bigl(
    C_{R\,dea}-C_{R\,ade}-C_{R\,ead}
  \bigr),
  \qquad
  \gamma_{i,R\,ae}
  :=
  e^d_i\,\gamma_{d;ae}^{(R)} .
\]
Their spin lift and the corresponding link transport are
\[
  \Omega_{i,R}
  :=
  \frac14\gamma_{i,R\,ab}A^aA^b,
  \qquad
  U_{i,R}(x)
  :=
  \exp\!\bigl(-h_R\Omega_{i,R}(x)\bigr).
\]
By \Cref{thm:fundamental},
\[
  \gamma_{i,R\,ab}
  =
  \omega^{\mathrm{LC}}_{i\,ab}+O(h_R^2)
\]
locally uniformly.

\subsubsection{Midpoint spin transport and the curved Hamiltonian}

Evaluate the link at edge midpoints and define
\[
  \nabla_{i,R}\psi(x)
  :=
  \frac{
   U_{i,R}(x+\tfrac12h_Re_i)^{-1}\psi(x+h_Re_i)
   -
   U_{i,R}(x-\tfrac12h_Re_i)\psi(x-h_Re_i)
  }{2h_R}.
\]

\begin{lemma}[Second-order covariant consistency]
\label{lem:covariant-consistency}
For every \(\psi\in C_c^\infty(\R^{\dCl},V)\),
\[
  \nabla_{i,R}\psi
  =
  \bigl(\partial_i+\Omega_i^{\mathrm{LC}}\bigr)\psi
  +O_{L^2}(h_R^2),
  \qquad
  \Omega_i^{\mathrm{LC}}
  =\frac14\omega^{\mathrm{LC}}_{i\,ab}A^aA^b .
\]
More precisely, on each compact set containing the support of \(\psi\),
\[
  \left\|
    \nabla_{i,R}\psi-
    (\partial_i+\Omega_i^{\mathrm{LC}})\psi
  \right\|_{L^2}
  \le C_\psi h_R^2 .
\]
\end{lemma}

\begin{proof}
Taylor-expand the forward and backward spinors and the midpoint link
exponentials around \(x\).  The order-\(h_R^2\) brackets in the two transported
terms agree and cancel in the central difference; the first surviving
truncation term is order \(h_R^3\), hence order \(h_R^2\) after division by
\(2h_R\).  Replacing \(\gamma_R\) by \(\omega^{\mathrm{LC}}\) costs another
\(O(h_R^2)\) by Theorem~\ref{thm:fundamental}.  The orientation convention
in the definition transports both neighbouring spinors back to the fibre over
\(x\), yielding the displayed sign of \(\Omega_i^{\mathrm{LC}}\).
\end{proof}

Set \(B^j(x)=A^ae_a^{\,j}(x)=\kappa A^aM^{aj}(x)\).  On flat measure define
\[
  \Ham_g^{\mathrm{LC}}
  :=
  \frac12\sum_j
  \left[
    B^j\bigl(-i(\partial_j+\Omega_j^{\mathrm{LC}})\bigr)
    +
    \bigl(-i(\partial_j+\Omega_j^{\mathrm{LC}})\bigr)^\dagger B^j
  \right].
\]
\label{def:curved-ham}
This is the half-density realization of the geometric spatial Dirac
Hamiltonian.  Its principal symbol is \(B^j\xi_j\), and
\[
  (B^j\xi_j)^2
  =\kappa^2\xi^\top M(x)^2\xi\,\one .
\]

\begin{lemma}[Self-adjointness of the reconstructed curved Hamiltonian]
\label{lem:curved-selfadjoint}
The operator \(\Ham_g^{\mathrm{LC}}\) is essentially self-adjoint on
\(C_c^\infty(\R^{\dCl},V)\).  Its closure is self-adjoint and satisfies
\[
  \|u\|_{H^1}
  \le C\bigl(\|\Ham_g^{\mathrm{LC}}u\|_{L^2}+\|u\|_{L^2}\bigr).
\]
The same conclusions hold for
\(\Ham_g^{\mathrm{LC}}+V_{\mathrm{ren}}\) whenever
\(V_{\mathrm{ren}}\in L^\infty(\R^{\dCl},\R)\).
\end{lemma}

\begin{proof}
The principal coefficients \(B^j\) are bounded \(C_b^1\) Hermitian matrices,
the propagation speed is uniformly bounded, and the Weyl symmetrisation makes
the operator symmetric.  Essential self-adjointness follows from the same
Chernoff argument used in the flat-ordering formulation.  Uniform
nondegeneracy of \(M\) gives first-order ellipticity and the graph estimate.
A bounded real scalar potential is a bounded self-adjoint perturbation.
\end{proof}

\begin{definition}[Covariant discrete renewal Hamiltonian]
\label{def:curved-discrete}
Let \(V_{\mathrm{ren},R}\) be real bounded multiplication operators satisfying
\[
  \sup_R\|V_{\mathrm{ren},R}\|_\infty<\infty,
  \qquad
  V_{\mathrm{ren},R}\longrightarrow V_{\mathrm{ren}}
  \quad\text{locally uniformly}.
\]
Define
\[
  \Ham_R^{\mathrm{LC}}
  =
  \frac12\sum_j
  \left[
    B^j(-i\nabla_{j,R})
    +
    (-i\nabla_{j,R})^\dagger B^j
  \right]
  +V_{\mathrm{ren},R}.
\]
Each \(\Ham_R^{\mathrm{LC}}\) is bounded and self-adjoint for fixed \(R\).
The purely geometric model is recovered by setting
\(V_{\mathrm{ren},R}=V_{\mathrm{ren}}=0\).
\end{definition}

\begin{proposition}[Common-core consistency]
\label{prop:curved-core-consistency}
For every \(u\in C_c^\infty(\R^{\dCl},V)\),
\[
  \left\|
    \Ham_R^{\mathrm{LC}}u
    -
    \bigl(\Ham_g^{\mathrm{LC}}+V_{\mathrm{ren}}\bigr)u
  \right\|_{L^2}
  \longrightarrow0 .
\]
If \(V_{\mathrm{ren},R}-V_{\mathrm{ren}}=O(h_R^2)\) locally uniformly, then
the left-hand side is \(O_u(h_R^2)\).
\end{proposition}

\begin{proof}
Lemma~\ref{lem:covariant-consistency} controls the covariant differences.
Multiplication by the \(C_b^4\) coefficients \(B^j\), their discrete adjoints,
and the Weyl symmetrisation preserve the estimate and reproduce the
half-density divergence term in \(\Ham_g^{\mathrm{LC}}\).  The potential
contribution converges by local uniform convergence on the compact support of
\(u\).  Summing over the finitely many coordinate directions gives the claim.
\end{proof}

\begin{theorem}[Curved Levi--Civita strong-resolvent limit]
\label{thm:curved-limit}
Under Definitions~\ref{def:reset-field} and~\ref{def:curved-discrete},
\[
  \bigl(\Ham_R^{\mathrm{LC}}-z\bigr)^{-1}
  \longrightarrow
  \bigl(\Ham_g^{\mathrm{LC}}+V_{\mathrm{ren}}-z\bigr)^{-1}
  \quad\text{strongly}
\]
for every \(z\in\C\setminus\R\).  Consequently,
\[
  e^{-it\Ham_R^{\mathrm{LC}}}
  \longrightarrow
  e^{-it(\Ham_g^{\mathrm{LC}}+V_{\mathrm{ren}})}
\]
strongly, uniformly for \(t\) in compact intervals, and
\[
  \varphi(\Ham_R^{\mathrm{LC}})
  \longrightarrow
  \varphi(\Ham_g^{\mathrm{LC}}+V_{\mathrm{ren}})
  \quad\text{strongly}
\]
for every \(\varphi\in C_0(\R)\).

The limiting operator is the full geometric Dirac Hamiltonian through
subprincipal order: its principal symbol determines
\[
  \xi_0^2=\kappa^2\xi^\top M(x)^2\xi ,
\]
its spin subprincipal term is the Levi--Civita connection, and the only
additional term retained here is the explicitly separated real scalar
\(V_{\mathrm{ren}}\).
\end{theorem}

\begin{proof}
Proposition~\ref{prop:curved-core-consistency} gives convergence on
\(C_c^\infty\), which is a core for the limiting self-adjoint operator by
Lemma~\ref{lem:curved-selfadjoint}.  Every
\(\Ham_R^{\mathrm{LC}}\) is self-adjoint.  By the strong graph-convergence
criterion for self-adjoint operators, equivalently the common-core
Trotter--Kato theorem, convergence
\[
  \Ham_R^{\mathrm{LC}}u\longrightarrow
  (\Ham_g^{\mathrm{LC}}+V_{\mathrm{ren}})u,
  \qquad u\in C_c^\infty(\R^{\dCl},V),
\]
on a core of the limiting self-adjoint operator implies strong-resolvent
convergence.  The unitary-group statement follows uniformly on compact time
intervals by Stone's theorem and the Trotter--Kato equivalence, and the
\(C_0(\R)\) functional-calculus statement follows from the bounded-continuous
functional-calculus form of strong-resolvent convergence.
\end{proof}

\begin{proposition}[Worked reversible reset-diffusion residue]
\label{prop:renewal-scalar-potential}
Let \(\rho_*\in C_b^2(\R^{\dCl})\) satisfy
\(0<\rho_-\le\rho_*\le\rho_+<\infty\), and consider the reversible reset-diffusion
generator
\[
  L_{\mathrm{ren}}
  =D\rho_*^{-1}\nabla\!\cdot(\rho_*\nabla)
  =D\bigl(\Delta+\nabla\log\rho_*\cdot\nabla\bigr).
\]
Under the ground-state transform
\(\Theta\psi=\rho_*^{1/2}\psi\),
\[
  \Theta L_{\mathrm{ren}}\Theta^{-1}
  =D\Delta-V_{\mathrm{ren}},
  \qquad
  V_{\mathrm{ren}}
  =D\rho_*^{-1/2}\Delta\rho_*^{1/2}.
\]
Thus the gradient drift is removed and the residual renewal correction is a
bounded real scalar, with
\[
  \|V_{\mathrm{ren}}\|_\infty
  \le D\left(
    \frac{\|\Delta\rho_*\|_\infty}{2\rho_-}
    +
    \frac{\|\nabla\rho_*\|_\infty^2}{4\rho_-^2}
  \right).
\]
A consistent discrete realization is
\[
  V_{\mathrm{ren},R}
  =D\rho_{*,R}^{-1/2}\Delta_R\rho_{*,R}^{1/2},
\]
for positive \(\rho_{*,R}\to\rho_*\) with the corresponding discrete
derivatives converging locally uniformly.
\end{proposition}

\begin{proof}
The product rule applied to
\(\rho_*^{1/2}L_{\mathrm{ren}}(\rho_*^{-1/2}\psi)\) cancels the first-order
terms and gives the displayed scalar potential.  Expanding
\(\Delta\rho_*^{1/2}\) yields
\[
  V_{\mathrm{ren}}
  =D\left(
    \frac12\frac{\Delta\rho_*}{\rho_*}
    -\frac14\frac{|\nabla\rho_*|^2}{\rho_*^2}
  \right),
\]
from which the bound follows.  This is a worked reversible diffusion model,
not a claim that every finite CP reset memory produces the same scalar.
\end{proof}

\begin{remark}[Scope of the curved closure]
\label{rem:ordering}
The curved result has the following precise scope: the renewal second moment
determines the coframe and Lorentzian principal symbol; covariance of transported
second moments derives metric compatibility; exchange symmetry of the
stationary diamond law cancels the translational plaquette defect while leaving
its rotational part free to retain curvature; the resulting discrete Cartan
system uniquely reconstructs the Levi--Civita spin connection; and the
corresponding self-adjoint covariant Hamiltonians converge globally in
strong-resolvent sense.  A smooth-cutoff,
variable-coefficient norm-resolvent estimate remains a separate microlocal
problem, while full norm-resolvent convergence is obstructed by lattice
doublers.
\end{remark}

\begin{remark}[Why uniform ellipticity is needed]
\label{rem:ellipticity-needed}
Uniform ellipticity $M(x)\succeq m_0 I$ is essential: it bounds the propagation
speed away from degeneration and yields the elliptic estimate of
Lemma~\ref{lem:curved-selfadjoint}.  Where $M(x)$ degenerates, the cone closes and
the geometry becomes horizon-like; such loci are excluded here and would require
a genuinely different (degenerate-hyperbolic) analysis .  As in
\Cref{rem:operator-scope}, the signature is carried by the signed/Krein datum and the normalized temporal
generator, convergence is strong (not norm) resolvent, and no curvature dynamics, field
equation, or matter coupling is asserted.
\end{remark}

\subsubsection{Stability under perturbation of the reset statistics}

We now show the construction is robust: a small change of the reset directions
produces a small change of the metric, the cone, and the Hamiltonian.  We treat
the flat (constant-$M$) case in full and record the curved variant.

\begin{theorem}[Stability of the rounded Lorentzian endpoint]
\label{thm:stability}
Fix the Krein--Clifford datum.  Let two direction measures $\nu,\tilde\nu$ on
$\Sphere^{\dCl-1}$ have constant second moments $M=M(\nu)$, $\tilde M=M(\tilde\nu)$
with
\[
  \norm{M-\tilde M}\le\chi,
  \qquad M\succeq m_0 I,\qquad \chi<m_0 .
\]
Then, with all constants explicit:
\begin{enumerate}[label=\textup{(\roman*)},leftmargin=2.4em]
\item \textup{(signature is rigid)} $\tilde M\succeq(m_0-\chi)I\succ0$, so
      $g_{\tilde\nu}$ has the same Lorentzian signature $(1,\dCl)$ as $g_\nu$;
\item \textup{(metric)} $\norm{g_\nu-g_{\tilde\nu}}=\kappa^2\norm{M^2-\tilde
      M^2}\le 2\kappa^2\chi$;
\item \textup{(cone)} on the section $\{|\vec\xi|=1\}$ the null cones
      $\xi_0=\pm\kappa|M\vec\xi|$ and $\xi_0=\pm\kappa|\tilde M\vec\xi|$ are
      uniformly $\kappa\chi$-close, so the cones are $O(\chi)$-close in Hausdorff
      distance on compact sections;
\item \textup{(Hamiltonian)} $\Ham_{g_\nu}-\Ham_{g_{\tilde\nu}}
      =\kappa A^i(M-\tilde M)^{ij}(-i\partial_j)$ is first order with
      \[
        \norm{(\Ham_{g_\nu}-\Ham_{g_{\tilde\nu}})u}_{L^2}
        \le\kappa\norm{A}\,\chi\,\norm{u}_{H^1}\qquad(u\in H^1);
      \]
\item \textup{(norm resolvent)} for every $z\in\C\setminus\R$, the two operators
      have the common graph domain $H^1$ and
      \[
        \norm{(\Ham_{g_\nu}-z)^{-1}-(\Ham_{g_{\tilde\nu}}-z)^{-1}}
        \le \frac{\kappa\norm{A}\,C(z,\kappa,m_0)}{|\operatorname{Im}z|}\,\chi
        =O(\chi);
      \]
\item \textup{(propagator)} for $u\in H^1$,
      $\norm{(e^{-it\Ham_{g_\nu}}-e^{-it\Ham_{g_{\tilde\nu}}})u}_{L^2}
      \le|t|\,\kappa\norm{A}\,\chi\,\norm{u}_{H^1}$.
\end{enumerate}
\end{theorem}

\begin{proof}
\emph{(i)} $\tilde M=M-(M-\tilde M)\succeq m_0 I-\chi I=(m_0-\chi)I\succ0$, so by
\Cref{thm:signature-krein} $g_{\tilde\nu}=\diag(1,-\kappa^2\tilde M^2)$ has
signature $(1,\dCl)$.

\emph{(ii)} $g_\nu-g_{\tilde\nu}=\diag(0,-\kappa^2(M^2-\tilde M^2))$ and
$M^2-\tilde M^2=M(M-\tilde M)+(M-\tilde M)\tilde M$, so
$\norm{M^2-\tilde M^2}\le(\norm{M}+\norm{\tilde M})\chi\le2\chi$ since second
moments on the unit sphere satisfy $\norm{M},\norm{\tilde M}\le1$.

\emph{(iii)} For $|\vec\xi|=1$,
$\bigl||M\vec\xi|-|\tilde M\vec\xi|\bigr|\le|(M-\tilde M)\vec\xi|
\le\norm{M-\tilde M}\le\chi$, so the two cone graphs over the spatial sphere are
within $\kappa\chi$.

\emph{(iv)} The principal parts differ by $\kappa A^i(M-\tilde M)^{ij}
(-i\partial_j)$ (the constant-$M$ symmetrisation term vanishes); for $u\in H^1$,
$\norm{\kappa A^i(M-\tilde M)^{ij}(-i\partial_j)u}\le\kappa\norm{A}\norm{M-\tilde
M}\norm{\nabla u}\le\kappa\norm{A}\chi\norm{u}_{H^1}$.

\emph{(v)} By Lemma~\ref{lem:curved-selfadjoint} (flat case) both Hamiltonians have
common graph domain $H^1$ and are uniformly elliptic, so $(\Ham-z)^{-1}$ maps
$L^2\to H^1$ with norm bounded by $C(z,\kappa,m_0)$: indeed
$\norm{\nabla v}\le(\kappa m_0)^{-1}\norm{\Ham v}$ and
$\norm{\Ham(\Ham-z)^{-1}}\le1+|z|/|\!\operatorname{Im}z|$.  The perturbation in
(iv) is therefore bounded from the common graph domain $H^1$ to $L^2$.  The second resolvent
identity
\[
  (\Ham_{g_\nu}-z)^{-1}-(\Ham_{g_{\tilde\nu}}-z)^{-1}
  =(\Ham_{g_\nu}-z)^{-1}\bigl(\Ham_{g_{\tilde\nu}}-\Ham_{g_\nu}\bigr)
   (\Ham_{g_{\tilde\nu}}-z)^{-1},
\]
together with $\norm{(\Ham_{g_\nu}-z)^{-1}}_{L^2\to L^2}\le|\!\operatorname{Im}
z|^{-1}$, the $H^1\!\to\!L^2$ bound \emph{(iv)} for the middle factor, and the
$L^2\!\to\!H^1$ bound for the right factor, gives the stated norm bound.

\emph{(vi)} Duhamel's formula
$e^{-it\Ham_{g_\nu}}-e^{-it\Ham_{g_{\tilde\nu}}}
=-i\int_0^t e^{-i(t-s)\Ham_{g_\nu}}(\Ham_{g_\nu}-\Ham_{g_{\tilde\nu}})
e^{-is\Ham_{g_{\tilde\nu}}}\,ds$, applied to $u\in H^1$ (which $e^{-is
\Ham_{g_{\tilde\nu}}}$ preserves, the propagator commuting with the elliptic
$\Ham_{g_{\tilde\nu}}$), and the bound \emph{(iv)}, give the estimate after
taking norms under the integral.
\end{proof}

\begin{corollary}[Curved stability]
\label{cor:curved-stability}
If instead $M,\tilde M\in C^1_b$ are space-dependent with
$\norm{M-\tilde M}_{C^1}\le\chi$ and both $\succeq m_0 I$, then conclusions
\textup{(i)--(iii)} hold pointwise in $x$, and \textup{(iv)--(vi)} hold with the
same $O(\chi)$ rate, the only change being that the difference of Hamiltonians
now carries the bounded zeroth-order term
$-\tfrac{i\kappa}{2}A^i\partial_j(M-\tilde M)^{ij}$, of norm $\le
\tfrac{\kappa}{2}\norm{A}\,\chi$, which is absorbed into the $H^1\to L^2$ bound.
\end{corollary}

\begin{proof}
Identical to \Cref{thm:stability}, with the extra bounded term estimated by
$\norm{\partial(M-\tilde M)}_\infty\le\chi$; ellipticity holds for both operators
by Lemma~\ref{lem:curved-selfadjoint}.
\end{proof}

\begin{remark}[Stability is stronger than the discretisation limit]
\label{rem:stability-stronger}
The discretisation limits of \Cref{thm:operator-limit,thm:curved-limit} are
\emph{strong} resolvent only, the symmetric-difference obstruction
\cite{corneangardejensen2022,schmidtumeda2023} ruling out norm resolvent
convergence in $\dCl\ge2$.  By contrast, \Cref{thm:stability}\,(v) is \emph{norm}
resolvent stability: perturbing $M$ at fixed continuum keeps both operators in
the same uniformly elliptic class with common $H^1$ graph domain, so the
difference is a genuinely first-order $H^1\to L^2$ perturbation with no mesh
pathology.  Robustness of the emergent geometry is
therefore quantitatively sharper than its emergence.
\end{remark}

\begin{remark}[Novelty]
\label{rem:novelty-curved}
The analytic ingredients are classical: Chernoff's essential self-adjointness for
first-order symmetric hyperbolic systems \cite{chernoff1973}, the
strong-resolvent core criterion and Trotter--Kato group convergence
\cite{reedsimon1}, and second-resolvent/Duhamel perturbation bounds.  The
discrete-to-continuum Dirac literature
\cite{corneangardejensen2022,schmidtumeda2023} treats a \emph{constant} principal
part (flat metric) with at most a bounded matrix potential; it does not address a
position-dependent principal symbol $M(x)$, which is exactly the curved case here.
The contribution of this section is therefore the \emph{assembly}: position-dependent
reset statistics yield a self-adjoint variable-coefficient Dirac Hamiltonian whose
principal symbol is an emergent curved Lorentzian metric field, together with a
quantitative norm-resolvent stability of that field under perturbation of the reset
data.  The present statement is the renewal/Krein counterpart of these stability principles,
with an explicit $O(\chi)$ rate for perturbations of the reset statistics.
\end{remark}

\section{Interpretation, scope, and conclusion}

% =====================================================================

\subsection{From predictive memory to Lorentzian geometry}
\label{subsec:interpretation}
\label{subsec:self-averaging-interpretation}

The results of the preceding sections can be understood as a single
conceptual progression.  The construction begins with no manifold, metric,
causal cone, preferred time direction, or prescribed spacetime dimension.
Its primitive data consist only of renewal events and the completely positive
channels generated by their composition.  A history is therefore not yet a
trajectory through spacetime, but a sequence of updates distinguished only by
the predictive action it retains.
%
The decisive step is to identify histories that are indistinguishable by
every future continuation:
\[
    w\sim_{\CP}w'
    \qquad\Longleftrightarrow\qquad
    \Phi_w=\Phi_{w'}.
\]
This predictive quotient transforms a collection of symbolic histories into a
positive spectral geometry.  Distinguishable predictive actions become its
points; minimal renewal depth becomes length; continuation becomes order; and
the growth of distinguishable classes becomes dimension.
%
Positivity nevertheless imposes a sharp boundary.  It can organise what
prediction retains, but it cannot by itself distinguish a negative sector or
produce a Lorentzian signature.  Crossing that boundary requires the minimal
signed enrichment, a principal double cover of the positive predictive
geometry.  Reversible predictive transformations then generate a projective
commutator structure, whose noncentral reduction yields a Clifford factor.
After nondegenerate coarse-graining, this signed Clifford geometry admits
controlled flat and slowly varying Lorentzian Dirac limits.  In this precise sense, the paper does not place renewal dynamics inside an
already existing spacetime.  It shows instead how the kinematic architecture
normally attributed to spacetime can arise as the stable large-scale
organisation of predictive renewal memory.


This interpretation should not be confused with the claim that every renewal
process produces a universe, or that the complete dynamics of gravity has
already been derived.  The theorem chain distinguishes the finite UCP renewal
starting object and its analytic domain conditions from the phase conditions
selecting nontrivial signed holonomy, reversible projective transport,
nondegenerate spatial scaling, and an exchange-symmetric stationary
translational regime.  What it establishes is more specific: once a renewal
process enters the stated predictive, signed, reversible, and coarse-grained
phase, the resulting operator geometry is not merely analogous to Lorentzian
geometry.  Its principal symbol, causal orientation, Clifford structure, and
continuum Hamiltonian have the corresponding Lorentzian form.  The
significance of the construction is therefore structural rather than
cosmological.  It identifies a route by which geometry can be secondary to
memory without becoming metaphorical.

\paragraph{Renewal, opposite orientation, and cancellation.}
Renewal is ordinarily understood as a restart: part of the past is discarded,
a new local state is formed, and evolution continues from the retained
information.  In the present setting, renewal is more abstract.  It is the
operation by which a history is extended while only its future predictive
content remains relevant.  The predictive quotient formalises this economy of
memory.  Two microscopically different histories become one state when no
future continuation can distinguish them.  Identity is therefore not defined
by sameness of origin, but by sameness of consequence.  A predictive state is
what remains invariant after all details that make no difference to the
future have been removed.

The signed cover introduces a second possibility that is invisible to the
positive quotient.  Every renewal continuation has a deck-related partner
carrying the opposite Krein orientation.  This partner is not an inverse
completely positive channel, a retrodictive dual, a literal reversal of
microscopic time, or the destruction of a renewal event.  It is the
oppositely oriented realisation of the same positive predictive content once
the minimal signed layer is retained.  The two sheets therefore agree on the
positive facts of prediction while differing in the sign with which those
facts are embedded into the indefinite geometry.

This distinction makes cancellation possible.  Positive probabilities do not
cancel: the completely positive sector only adds and averages positive data.
By contrast, signed amplitudes or signed observables can receive equal
contributions from the two sheets with opposite orientation.  When the two
renewal orientations are exactly balanced, their signed contribution
disappears,
even though their common positive predictive content remains.  The positive
quotient records what both orientations agree upon; the signed sector records
the difference that their agreement suppresses.  Cancellation is thus not
the erasure of reality but the removal of oriented excess.  What survives
perfect cancellation is the unsigned geometry of distinguishability,
continuation, and volume.  What survives imperfect cancellation is an
orientation-sensitive residue capable of supporting a temporal Clifford
direction.

This gives the signed construction a clear metaphysical reading.  The positive
sector describes being insofar as it is stably reproducible in prediction.
The signed sector describes a difference of orientation that cannot be reduced
to additional positive information.  In that sense, Lorentzian signature does
not arise because one more positive degree of freedom has been appended to
the system.  It arises because the same predictive content admits two
oppositely oriented realisations whose relation is irreducibly signed.  The
negative direction is not another object among the objects.  It is a
structural opposition within the way predictive content can be carried.

\paragraph{From balance to temporal orientation.}
At the level of the double cover, the two renewal orientations are formally
symmetric.  The deck transformation exchanges the two sheets, and the
positive quotient forgets which sheet was used.  If the full state and
dynamics were exactly invariant under this exchange, neither orientation
would be physically preferred.  The signed geometry would still exist, and
its algebraic relations would still distinguish opposite orientations, but
there would be no net arrow selecting one as the realised direction of
becoming.

A temporal orientation appears when this balance is no longer exact.  The
mathematical construction identifies two ingredients of such an imbalance.
First, a nontrivial signed holonomy makes the opposition between the two
orientations globally non-removable: no local choice of sheet can eliminate
it around every recurrent cycle.  Second, modular homogeneity ties elementary
renewal steps to a common scaling flow, making renewal depth into an oriented
modular clock.  Together with the canonical hyperbolic pair in the reversible
revision sector, this produces the distinguished timelike Clifford generator.
The resulting orientation is not inserted as an external coordinate.  It is
the operator-theoretic expression of a persistent mismatch between the two
signed ways of carrying the same predictive process.

One may therefore interpret temporal becoming as the failure of
cancellation between the two renewal orientations to remain exact across the
recurrent
structure.  Under perfect balance, the signed excess vanishes and no
orientation is selected at the level of states.  Under broken balance, the
difference accumulates coherently and becomes macroscopic.  The emergence of
time is then not the appearance of an additional substance called ``time,''
but the stabilization of an oriented asymmetry within renewal itself.  The
future is the direction in which predictive continuation is modularly
organised; the opposite sheet represents the formally available
anti-orientation against which this direction is defined.

This interpretation must be stated with care.  The manuscript proves a
distinguished modular orientation and an intrinsic timelike Clifford slot.  It
does not prove that a particular physical state spontaneously breaks the deck
symmetry, nor does it derive thermodynamic irreversibility, entropy production,
or the observed cosmological arrow of time.  Those would require additional
state-selection and dynamical results.  The present construction instead
isolates the exact structural location at which such an arrow could arise:
not in the positive renewal data alone, but in a persistent imbalance between
the two signed orientations of that data.

\paragraph{Spacetime as stabilised predictive relation.}
The broader metaphysical conclusion is relational.  Spacetime is not treated
as a container that exists before events and gives them somewhere to occur.
The primitive notion is continuation: one predictive state can be prolonged
into another, some histories remain distinguishable, and others collapse to
the same future action.  Order expresses possible continuation, distance
expresses the renewal depth required to realise a predictive distinction, and
dimension expresses the rate at which new distinctions become available.
The signed layer adds the irreducible opposition needed for Lorentzian
signature, while reversible projective dynamics supplies the spinorial
structure through which those distinctions propagate.

Smooth geometry appears only after coarse-graining.  At finite resolution the
order cone may be polyhedral, the reset frame anisotropic, and the empirical
second moments history-dependent.  Nothing at this scale requires the system
to resemble a manifold.  Stationary self-averaging explains how these
fluctuating local frames can nevertheless produce a deterministic macroscopic
coframe.  The continuum metric is therefore not a microscopic object that has
been sampled imperfectly.  It is a stable law of organisation that emerges
when the variable details of renewal cease to matter at the scale of
propagation.  Constant anisotropy is absorbed into the choice of spatial
frame, while slowly varying statistics appear as a curved coframe whose
plaquette defects separate torsion from curvature.

The same perspective clarifies the dimensional result.  Dimension is not
chosen by declaring how many coordinates reality possesses.  It is constrained
by the rank of the noncentral projective revision sector and by the requirement
that the spatial part support a genuine two-direction interference or
plaquette structure.  One temporal direction leaves an odd spatial rank; the
first nontrivial spatial bivector sector occurs in three dimensions.  Thus
\(3+1\) is the first primitive Lorentzian endpoint of the mechanism, although
higher odd spatial ranks remain algebraically possible without an additional
selection principle.

The resulting picture is neither substance metaphysics nor pure
instrumentalism.  Predictive equivalence is operational, but the structures
reconstructed from it are objective invariants of the renewal process:
cohomology classes, spectral growth, projective commutators, Clifford rank,
and continuum operator limits.  The proposal is therefore that spacetime is
real as a stable relational form, rather than fundamental as a prior arena.
Its causal, spectral, temporal, and spinorial features are what remain when
renewal histories are quotiented by predictive equivalence, their hidden
orientation difference is retained, and their fluctuations are organised
into a coherent large-scale limit.


\subsection{What lies upstream of renewal memory?}
\label{subsec:upstream-ontology}

The construction developed in this paper begins with renewal memory, but it
does not follow that renewal memory is the final ontological layer.  Even the
primitive renewal object already contains several structural commitments.  It
assumes that elementary updates can be distinguished, that they can be
composed into histories, that their effects act on an algebra of possible
observables, and that two histories can be compared through the future action
they induce.  Later stages assume that some recurrent transformations are
reversible, that a nontrivial signed sector is realised, that the renewal
grading is compatible with a modular flow, and that large-scale statistics
enter a nondegenerate coarse-grained phase.  These are not assumptions of an
ambient spacetime, but they are nevertheless assumptions about what the
pregeometric world is capable of doing.

It is therefore useful to distinguish the mathematical starting point of the
manuscript from the most upstream metaphysical possibility suggested by it.
The mathematical starting point is a finite family of completely positive
renewal channels and the histories generated by their composition.  The more
primitive possibility is not a set of histories already given, but a capacity
for distinctions to be made, retained, combined, and forgotten.  On this
view, a history is not a fundamental object.  It is the record produced when
elementary distinctions are composed in an order.  Renewal is the rule by
which some part of that record ceases to matter, while its consequences remain
available for further composition.

The most economical upstream structure would therefore contain three
ingredients: possible distinctions, composability, and retained consequence.
One may represent them schematically as
\[
    \text{distinction}
    ;\longrightarrow;
    \text{update}
    ;\longrightarrow;
    \text{history}
    ;\longrightarrow;
    \text{predictive action}.
\]
The predictive quotient then removes everything in a history that makes no
difference to its future action.  What becomes geometrically real is consequently not the full microscopic
record, but the structure of distinctions that remain effective under future
continuation.  The primitive ``points'' of the emergent geometry are not
events considered in isolation.  They are equivalence classes of pasts that
retain different powers over what can happen next.

This suggests that the most fundamental entity is neither a spacetime point
nor a history in the ordinary sense.  A spacetime point is already a geometric
notion, while a history presupposes a criterion for distinguishing successive
updates.  More fundamental than either is a relational capacity: the capacity
of one effective distinction to condition another.  The basic reality is then
not a collection of things situated in a container, but a structured field of
possible continuations.  Histories are paths through this field; predictive
states are the equivalence classes of paths that have the same remaining
consequences; and spacetime is the stable geometry of those equivalence
classes when viewed at large scale.

From this perspective, the observable universe is a particular phase of a
deeper compositional structure.  The positive predictive geometry expresses
the stable content shared by the two renewal orientations.  The signed cover
retains the opposition between their two orientations.  Reversible revision
dynamics converts that opposition into a projective Clifford structure, while
modular imbalance selects a temporal orientation.  Coarse-graining then turns
the finite and fluctuating renewal frame into the smooth Lorentzian geometry
within which macroscopic propagation occurs.  What we call spacetime is thus
not the primitive arena of the process; it is the regular form assumed by the
process once predictive distinctions, signed orientation, and recurrent
composition have stabilised together.

The hypotheses of the present construction do not all have the same logical
status.  The finite UCP renewal memory is the mathematical starting object;
finite fibres, bounded increments, and continuum regularity are analytic
domain conditions.  Nontrivial signed holonomy, a realised reversible
projective sector, pressure supercriticality, elliptic spatial scaling, and
the exchange-symmetric stationary regime are phase conditions selecting the
Lorentzian universality class.  The choice of one oriented extremal state is
a physical or boundary-condition selection.  Regular variation, exact
isotropy, metric coding, and higher-rank uniqueness conditions are optional
strengthenings used only for sharper spectral or selection statements.
Within the exchange-symmetric stationary phase, translational plaquette
closure is derived rather than independently imposed, and the torsion-free
branch follows from the resulting vanishing of the exchange-odd
translational defect.

Some of these conditions may themselves arise dynamically.  Complete
positivity may be the effective remnant of a more primitive reversible
evolution after inaccessible degrees of freedom have been discarded.
Renewal may arise when only a finite sufficient record can be retained.
Predictive equivalence may emerge because no internal observer can access
differences that leave all subsequent interactions unchanged.  The signed
sector may originate in a deeper duality whose two branches are identical at
the level of positive probabilities but distinct at the level of coherent
orientation.  Modular homogeneity may be the fixed-point behaviour of a
renormalisation process rather than a microscopic law.  The axioms would then
describe the universality class in which Lorentzian spacetime appears, not the
absolute beginning of the ontology.

The deepest candidate suggested by the construction is therefore not
``history'' but \emph{effective difference under continuation}.  Something is
real at the pregeometric level insofar as its presence changes the space of
possible futures.  Two records that never differ in their future action are
one predictive state; two orientations that agree positively but differ
coherently form the signed pair; and a persistent asymmetry between those
orientations becomes temporal direction.  Geometry emerges when the network
of such effective differences acquires stable growth, order, orientation, and
propagation laws.

Under this interpretation, the universe is fundamentally processual but not
merely sequential.  Sequence alone would give a list of events, whereas the
renewal construction requires equivalence, composition, recurrence, and
orientation.  The fundamental reality is a compositional structure in which
distinctions may be created, erased, or retained according to their future
effects.  Our Lorentzian spacetime is the macroscopic fixed form of that
structure: the geometry seen when predictive memory has become sufficiently
regular that its relations can be represented as distance, causal order,
spinorial propagation, curvature, and a single oriented time.





\subsection{Conclusion}
\label{sec:conclusion}
\label{subsec:scope}
\label{subsec:limitations}

We have shown that a finite renewal process of completely positive channels
can carry substantially more geometric structure than its microscopic
definition appears to contain.  Predictive quotienting turns update histories
into a positive spectral geometry and a divisibility order: distinguishable
predictive actions become points, minimal renewal depth becomes length,
continuation becomes causal order, and the growth of distinguishable classes
becomes dimension and residue volume.  The same construction also exposes a
sharp boundary.  Positivity can organise what prediction retains, but it
cannot by itself produce a Lorentzian fundamental symmetry.

The minimal passage beyond this boundary is a principal \(\Z/2\)-cover, with a
\(\Z/4\)-lift only when the signed holonomy is retained at amplitude level.
On the dynamically realised reversible subgroup, predictive units act by
automorphisms, their inner implementers compose projectively, and the
noncentral reduction of the resulting commutator form is a Clifford matrix
factor.  The signed cover fixes the temporal row, while a normalized
hyperbolic pair gives the intrinsic timelike generator
\[
  \gamma^0=J_\chi R_{\alpha_0},
  \qquad
  (\gamma^0)^2=-\one .
\]
Under subexponential shell growth, the associated modular signed operator has
compact resolvent and regularized all-orders summability.  The dimensionless
modular exponent is not left arbitrary: on a supercritical recurrent weighted
component it is the unique positive root of the associated Perron--Frobenius
pressure equation.  After
nondegenerate coarse-graining, the derived Clifford symbols converge to
Lorentzian Dirac symbols with one distinguished timelike direction, and the
corresponding Hamiltonians converge globally in strong resolvent sense and in
norm resolvent sense on every fixed momentum band.

The slowly varying regime extends this reconstruction from flat propagation
to curved transport.  Covariance of neighbouring second moments makes the
normalised frame comparison orthogonal and, after orientation selection,
special orthogonal, so metric compatibility is derived rather than separately
postulated.  Channel composition separates translational torsion from
rotational curvature.  In the primitive exchange-covariant stationary
diamond phase, exchange invariance forces the exchange-odd translational
defect to vanish, while leaving the rotational record free.  The resulting
translational closure selects the torsion-free branch, and the discrete
fundamental theorem yields the Levi--Civita spin connection without
suppressing curvature.  The covariant renewal Hamiltonians then converge
through subprincipal order to the corresponding variable-coefficient
Lorentzian Dirac Hamiltonian.  Primitive rank parity, together with the
requirement of a nontrivial spatial bivector sector, makes \(3+1\) dimensions
the first nondegenerate Lorentzian endpoint of the mechanism, although higher
odd spatial ranks remain algebraically admissible without an additional
selection principle.

These results are kinematic and operator-theoretic.  They do not yet provide
a field equation governing the reconstructed curvature, and they do not
derive Einstein's equations, a gravitational variational principle, matter
multiplets, gauge interactions, black-hole thermodynamics, or
phenomenological predictions.  Such structures lie downstream of the present
construction.  It is equally important to distinguish theorem from
interpretation.  Predictive order, the positivity obstruction, the minimal
signed-cover classification, modular affineness under its stated covariance
hypothesis, the projective Clifford factor, the intrinsic temporal generator,
and the continuum convergence results are mathematical statements.  The
description of imbalance between the two renewal orientations as a physical
arrow of time, or
of finite-scale renewal fluctuations as quantum spacetime fluctuations, is a
metaphysical and physical interpretation of where further dynamics may enter;
it is not required by the theorem chain.

The central conclusion is therefore not that a complete theory of the
universe has already been obtained.  It is that the basic kinematic
architecture of Lorentzian spacetime need not be placed beneath predictive
dynamics as a prior arena.  Under the hypotheses isolated here, causal order,
spectral dimension, temporal orientation, spinorial propagation, and curved
Levi--Civita transport arise as different aspects of one renewal structure.
Spacetime can thus be understood as the stable large-scale form of predictive
memory: not the container in which renewal occurs, but the geometry that
renewal acquires when distinguishability, signed orientation, reversibility,
and coarse-grained propagation become coherent.




% =====================================================================
\appendix


\section{Logical status and reduction of the hypotheses}
\label{app:hypothesis-status}

The mathematical starting object of this paper is the finite UCP renewal memory
of \Cref{def:renewal-memory}; no claim is made here that bare operational
composition forces every renewal system into the geometric regime studied
below.  To keep the logical status of the assumptions explicit, four distinct
roles should not be conflated.
\begin{enumerate}[label=\textup{(\roman*)},leftmargin=2.4em]
\item \emph{Operational antecedents} concern composability, testing,
      predictive identification, and stability under ancillary extension.  They
      lie upstream of the present paper and are not used as hidden hypotheses in
      its theorem proofs.
\item \emph{Microscopic model specifications} are the concrete alphabet,
      channels, record transports, branch weights, and local propagation laws
      selecting one renewal model inside the admissible class.
\item \emph{Phase conditions} select the signed Lorentzian universality class:
      a recurrent component with nontrivial signed holonomy, a realised
      reversible projective sector of the required rank, a supercritical
      pressure transfer, a positive spatial second moment, and an
      exchange-symmetric translational diamond phase.
\item \emph{Physical selections} choose a realised component or extremal
      orientation when several stationary phases coexist.  Such a selection is
      needed to speak of a realised arrow or state, but not to establish the
      kinematic operator identities common to the deck-related phases.
\end{enumerate}

The remaining hypotheses have a separate analytic status.  Finite shift fibres,
bounded length increments, bounded propagation, and the stated smoothness
bounds are domain conditions for the particular spectral or convergence
theorem in which they occur.  Regular variation is required for the sharp Weyl
coefficient and residue volume, not for the existence of the predictive triple
or the Lorentzian continuum.  Exact isotropy is only a convenient representative
of the flat isometry class.  Lipschitz metric coding, uniform separation, and
the dimension-uniqueness criteria of the appendices are optional strengthenings,
not premises of the core Lorentzian result.

\begin{proposition}[Canonical reduction of global standing conditions]
\label{prop:canonical-condition-reductions}
Let \(G\) be a finite resolved predictive graph, let a finite group \(K\) act by
automorphisms on a finite-dimensional complex \(C^*\)-algebra \(\mathcal A\),
and let \(\omega\) be the alternating commutator form of any scalar projective
implementation on a selected simple block.  Then:
\begin{enumerate}[label=\textup{(\alph*)},leftmargin=2.4em]
\item every stationary recurrent branch law is supported on terminal strongly
      connected components, so all recurrent constructions may be performed on
      one selected terminal component;
\item after selecting one invariant central orbit of \(\mathcal A\), choose
      one simple block and restrict to its stabiliser in \(K\); quotienting the
      stabiliser's action kernel makes the action faithful, and every resulting
      automorphism of the selected matrix block is inner;
\item quotienting the projective label space by \(\rad\omega\) gives a
      nondegenerate alternating space and hence the canonical primitive matrix
      factor.
\end{enumerate}
Thus global strong connectivity, global faithfulness, global irreducibility,
and nondegeneracy of the unreduced projective form are not independent phase
axioms.
\end{proposition}

\begin{proof}
The condensation graph of a finite directed graph is acyclic.  A stationary
law assigns zero mass to a component from which probability can leave without
return, and every recurrent path eventually remains in a terminal component,
which proves \textup{(a)}.  A finite-dimensional complex \(C^*\)-algebra is a
finite direct sum of matrix blocks.  The action permutes its minimal central
projections.  After selecting one orbit and one block in that orbit, its
stabiliser acts on the block; quotienting the stabiliser's action kernel makes
that action faithful, and every automorphism of the resulting matrix block is
inner.  This proves \textup{(b)}.  A homogeneous projective implementer is central exactly
when its label lies in \(\rad\omega\).  The induced form on the quotient has
zero radical, and finite Stone--von Neumann identifies its twisted group
algebra with a full matrix algebra, proving \textup{(c)}.
\end{proof}

\begin{remark}[Residual genuine phase content]
\label{rem:residual-phase-content}
After \Cref{prop:canonical-condition-reductions}, the genuinely selective inputs
are the existence of a terminal component carrying nontrivial stable-record
sign holonomy, a nontrivial realised reversible projective sector, pressure
supercriticality, selection of a deck-oriented extremal state when a realised
orientation is claimed, an elliptic spatial scaling phase, an
exchange-symmetric translational phase rather than a torsional phase, and the
continuum regularity required by the desired operator-convergence statement.
Constant anisotropy is removed by \Cref{thm:frame-universality}; regular
variation and the optional dimension selectors do not enter this list.
\end{remark}

\section{Regularity criteria and auxiliary positive invariants}
\label{app:regularity-positive}

The following results provide concrete classes satisfying the main-text
regularity criterion and record auxiliary notions of predictive dimension that
are not needed for the Lorentzian continuum theorem.

\subsection{Finite fibres and bounded increments}
\label{subsec:fibre-criteria}

The definitions above give a quotient set and descended creation maps, but not
yet a spectral geometry.  For the length operator to play the role of a Dirac
operator, the creation maps must be bounded algebra elements and their
commutators with the length Dirac must be bounded.  These two requirements have
an elementary form in the present setting: the shifts must have uniformly finite
fibres, and the quotient length must change by a bounded amount under each
reset.

We record several regular classes in which this condition holds with explicit
constants.  This is useful because \Cref{ass:fibres} is not intended as an
opaque analytic assumption; in the examples that behave like finite automata,
commuting cancellative monoids, or bounded quotients of such objects, it is a
checkable finite-combinatorial condition.  After these examples we prove a sharp
dichotomy: either the predictive-length triple exists, or its failure has a
visible operator-theoretic cause.

\begin{proposition}[Graded finite-type automata]
\label{prop:fibres-automata}
Suppose the predictive quotient is presented by a graded labelled graph
\[
  Q=\bigsqcup_{R\ge0}Q_R,\qquad \partial|_{Q_R}=R,
\]
whose transitions labelled by $\sigma\in E$ send $Q_R$ to $Q_{R+1}$, and
suppose $\Lam_{\min}$ is the depth $\partial$ of the class.  Assume the graph has
uniformly bounded backward branching: for every vertex $q$ and every label
$\sigma$, there are at most $B$ vertices $p$ with a $\sigma$-edge $p\to q$.
Then the descended shifts have fibres of size $\le B$ and
\[
  |\Lam_{\min}([\sigma w])-\Lam_{\min}([w])|=1 .
\]
Hence \Cref{ass:fibres} holds with $C_\sigma=1$.
\end{proposition}

\begin{proof}
The fibre of $\widehat S_\sigma$ over a class $[u]$ is its set of
$\sigma$-predecessors, of cardinality $\le B$ by hypothesis.  Since labelled
edges raise the grading by one, prepending $\sigma$ raises $\Lam_{\min}$ by
exactly one.
\end{proof}

\begin{proposition}[Finitely generated cancellative commuting monoids]
\label{prop:fibres-monoid}
Let the predictive monoid $\M$ be finitely generated, commutative, and
cancellative, with the additive grading $\Lam_{\min}=\deg$ induced by a strictly
positive linear functional on $\mathrm{Gr}(\M)\cong\Z^{\rmon}$ \textup{(}as in
\Cref{thm:ehrhart-order}\textup{)}.  Then the descended shift for the generator
$\sigma$ has fibres of size $\le1$, and
$\Lam_{\min}([\sigma w])-\Lam_{\min}([w])=\deg(a_\sigma)$ is constant, where
$a_\sigma\in\M$ is the image of $\sigma$.  Hence \Cref{ass:fibres} holds with
$C_\sigma=\deg(a_\sigma)$.
\end{proposition}

\begin{proof}
Cancellativity makes the equation $a_\sigma+m=u$ have at most one solution
$m=u-a_\sigma$ in $\M$, so each class has at most one $\sigma$-predecessor.
Additivity of $\deg$ gives $\deg(a_\sigma+m)-\deg(m)=\deg(a_\sigma)$.  In the
free case $\M=\N^{c}$ with unit weights this is the word-length grading, $C_\sigma=1$.
\end{proof}

\begin{proposition}[Commuting CP channels with finite relation ideal]
\label{prop:fibres-cp}
Suppose the channels $\{\Phi_\sigma\}$ commute and the channel monoid
$\M=\{\Phi_w\}$ is finitely presented \textup{(}the kernel congruence of
$E^{*}\to\M$ is generated by finitely many relations\textup{)} and cancellative.
Then $\M$ is a finitely generated cancellative commutative monoid, and
\Cref{ass:fibres} holds by \Cref{prop:fibres-monoid}.
\end{proposition}

\begin{proof}
A finitely presented commutative monoid is finitely generated; with
cancellativity it satisfies the hypotheses of \Cref{prop:fibres-monoid}.
\end{proof}

\begin{remark}[Cancellativity is the real hypothesis]
\label{rem:cancellativity-needed}
Composition of completely positive maps is not cancellative in general, so
\Cref{prop:fibres-cp} genuinely requires the cancellativity of the channel
monoid; it holds, for example, when the $\Phi_\sigma$ act faithfully and
invertibly on the closure of their joint range, or are conjugations by a
commuting family of unitaries.  Without cancellativity the predictive shifts may
have infinite fibres and the positive triple need not be finitely summable.
\end{remark}

\begin{proposition}[Finite graph quotients with bounded rank distortion]
\label{prop:fibres-quotient}
Let $\pi:G\to G'=G/{\sim}$ be a quotient of a predictive graph satisfying
\Cref{ass:fibres}.  Assume the equivalence is shift-compatible, the fibres of
$\pi$ have size $\le K$, and the rank functions satisfy
\[
  |\Lam'_{\min}(\pi(x))-\Lam_{\min}(x)|\le D
\]
\textup{(}bounded rank distortion\textup{)}.  Then $G'$ satisfies
\Cref{ass:fibres}, with shift fibres bounded by
$K\cdot\sup_\sigma|\text{fibre}_\sigma(G)|$ and increments bounded by
$C_\sigma+2D$.
\end{proposition}

\begin{proof}
Shift-compatibility ensures that every $\sigma$-predecessor in $G'$ lifts to a
$\sigma$-predecessor in $G$ after choosing a lift of the target.  There are at
most $K$ choices of lift and, for each lift, at most
$\sup_\sigma|\text{fibre}_\sigma(G)|$ predecessors, giving the stated fibre
bound.  For the increment,
\[
|\Lam'_{\min}(\pi(\sigma w))-\Lam'_{\min}(\pi(w))|
 \le |\Lam_{\min}(\sigma w)-\Lam_{\min}(w)|+2D
 \le C_\sigma+2D .
\]
\end{proof}





\begin{corollary}[Automatic finite-fibre spectral triples in the regular classes]
\label{cor:automatic-triples}
Let a finite completely positive renewal memory lie in any of the following
classes:
\begin{enumerate}[label=\textup{(\Alph*)},leftmargin=2.4em]
\item graded finite-state automata with bounded backward branching
      \textup{(}\Cref{prop:fibres-automata}\textup{)};
\item finitely generated cancellative commuting monoids
      \textup{(}\Cref{prop:fibres-monoid}\textup{)};
\item commuting CP channels with finite relation ideal and cancellative channel
      monoid \textup{(}\Cref{prop:fibres-cp}\textup{)};
\item finite graph quotients of \textup{(A)--(C)} with bounded rank distortion
      \textup{(}\Cref{prop:fibres-quotient}\textup{)}.
\end{enumerate}
Then \Cref{ass:fibres} holds, with the explicit constants of the cited
propositions, and consequently the predictive quotient carries the positive
predictive spectral triple of \Cref{thm:triple}: the descended shifts
$\widehat S_\sigma$ have bounded commutators $[D_{\CP},\widehat S_\sigma]$, the
length operator $D_{\CP}$ has compact resolvent, and the zeta abscissa is the
algebraic predictive dimension $\qalg$.  In the commuting classes
\textup{(B)--(C)} the multiplicities are regularly varying with integer
$\qalg=\rmon$, so \Cref{ass:regular-variation} holds and the renewal Weyl law of
\Cref{thm:weyl-law} applies, exhibiting $\qalg$ as a genuine Weyl exponent.
Thus, in classes \textup{(A)--(D)}, the spectral triple exists with no standing
finite-fibre hypothesis.
\end{corollary}

\begin{proof}
Each class verifies \Cref{ass:fibres} by the cited proposition:
\Cref{prop:fibres-automata} for (A) \textup{(}fibres $\le B$, $C_\sigma=1$\textup{)},
\Cref{prop:fibres-monoid} for (B) \textup{(}fibres $\le1$,
$C_\sigma=\deg a_\sigma$\textup{)},
\Cref{prop:fibres-cp} for (C) \textup{(}via
(B)\textup{)}, and \Cref{prop:fibres-quotient} for (D) \textup{(}transfer under
bounded-fibre, bounded-distortion quotients\textup{)}.  With \Cref{ass:fibres} in
force, \Cref{thm:triple} produces the spectral triple, its bounded commutators,
and the compact resolvent, and identifies the abscissa as $\qalg$.  For (B)--(C)
the predictive monoid is a finitely generated commutative cancellative monoid, so
by \Cref{thm:commuting} $\qalg=\rmon\in\N$ and the length-$R$ multiplicities count
lattice points in an $\rmon$-dimensional rational cone, hence are regularly varying of
index $\rmon-1$; \Cref{ass:regular-variation} holds and \Cref{thm:weyl-law} applies.
\end{proof}

\subsection{Metric dimension and renewal calibration}
\label{sec:metric}

The preceding construction measures how many predictive classes are
combinatorially distinguishable at length at most $R$.  There is a second,
independent notion of size: the metric size of the set of accumulated channels
inside the completely bounded norm geometry of maps on $\A_{\rm int}$.  The two
notions agree in manifold-like polynomial regimes, but they need not agree in
contracting or fractal renewal systems.  This distinction is useful, because it
separates algebraic growth of predictions from metric resolution of the channel
space.

\begin{definition}[Metric predictive dimension]
\label{def:qmet}
Let $\mathcal S=\{\Phi_w:w\in E^*\}$ and let $\overline{\mathcal S}$ be its closure in completely bounded distance.  The metric predictive dimension is the upper box dimension
\[
 \qmet=\limsup_{\delta\downarrow0}\frac{\log N_{\cb}(\delta)}{\log(1/\delta)},
\]
where $N_{\cb}(\delta)$ is the minimal number of $\cb$-balls of radius $\delta$ needed to cover $\overline{\mathcal S}$.
\end{definition}

\begin{remark}[Relation to dimensions of causal-state presentations]
\label{rem:qmet-computational-mechanics}
Computational mechanics also assigns geometric dimensions to mixed-state and
causal-state presentations.  In the natural mixed-state realizations studied for
discrete-time and continuous-time \(\varepsilon\)-machines, finite or suitably
countable causal-state presentations can have vanishing box-counting dimension,
whereas uncountable predictive-state attractors generically carry a nonzero
fractal dimension \cite{Marzen2018MWC,MarzenCrutchfield2017,
JurgensCrutchfield2021}.  This comparison is useful but not literally identical
to \(\qmet\).  Here the covering number is taken in completely bounded norm on
the \emph{closure} of accumulated channels.  Countability of the raw predictive
set alone therefore does not force \(\qmet=0\): accumulation can leave a
positive upper box dimension.  The metric-collapse conclusion in
\Cref{thm:qmet-qalg} accordingly uses finite-scale stabilization, not merely a
countability hypothesis.
\end{remark}

\begin{theorem}[Contracting reset systems]
\label{thm:fractal}
Suppose the accumulated channels are represented, in a finite affine coordinate chart, by similarities $f_\sigma$ of ratios $0<r_\sigma<1$ satisfying the open set condition, and suppose the cb-metric on channels is bi-Lipschitz equivalent to the Euclidean metric on the attractor.  Then
\[
 \qmet=s,
 \qquad
 \sum_{\sigma\in E} r_\sigma^s=1.
\]
Thus the metric predictive dimension is the Moran--Falconer dimension and is generically non-integer.
\end{theorem}

\begin{proof}
The accumulated channel set is bi-Lipschitz to the self-similar attractor $K=\bigcup_\sigma f_\sigma(K)$.  Under the open set condition, the Hausdorff and box dimensions of $K$ are the unique solution of the Moran equation \cite{hutchinson,falconer}.  Bi-Lipschitz equivalence preserves box dimension.
\end{proof}

\begin{remark}[Provenance of the metric dimension]
\label{rem:provenance-qmet}
The identification of a box-counting dimension with the abscissa of a
spectral-triple zeta function, and of the associated Dixmier-trace measure with
the Hausdorff measure, is not new: it is due to Pearson--Bellissard for
ultrametric Cantor sets \cite{pearson-bellissard} and was developed for
subshifts, tilings, and aperiodic media by Kellendonk--Savinien
\cite{kellendonksavinien2012}, building on the length-module spectral triples of
Connes \cite{connes1989} and Christensen--Ivan \cite{christensenivan2006}; the
Moran--Falconer equation itself is classical \cite{hutchinson,falconer,falconertechniques}.
What is specific to the present setting is only that the self-similar attractor
in question is the completely-bounded closure of accumulated renewal channels.
The contracting-reset case is thus the renewal instance of an established
mechanism, recorded here to fix the metric/algebraic dimension dictionary, not to
claim the dimension formula.
\end{remark}

\begin{remark}[Two dimensions]
\label{rem:two-dimensions}
The two dimensions need not agree.  The dimension $\qalg$ is a growth exponent for exact distinguishability by word length.  The dimension $\qmet$ is a resolution exponent for the cb-metric closure of accumulated channels.  In algebraic commuting regimes $\qalg$ is integer and $\qmet$ may collapse; in contracting free regimes $\qalg$ can be infinite while $\qmet$ remains finite and fractal.  A third quantity, the renewal exponent $\alpha$ of \Cref{def:renewal-exponent}, measures neither growth nor resolution but the tail of calibrated first returns, equivalently the critical singularity of the predictive-order resolvent; the three invariants are logically independent without additional hypotheses.
\end{remark}

\begin{remark}[Non-integer dimensions are metric, not Clifford ranks]
\label{rem:noninteger-metric-dimension}
The construction therefore contains two different routes to dimension.  The
algebraic dimension $\qalg$ counts exact predictive classes by quotient length;
in the commuting manifold-like regimes and in the primitive modular endpoint it
is an integer rank.  The metric dimension $\qmet$, by contrast, measures the
completely bounded metric closure of accumulated channels.  In contracting reset
systems it is the Moran--Falconer exponent $s$ of \Cref{thm:fractal}, determined
by
\[
  \sum_{\sigma\in E} r_\sigma^s=1,
\]
and it is generically non-integer.  Thus non-integer dimension emerges as a
resolution exponent of the CP-channel attractor, while the $3+1$ theorem concerns
the integer Clifford rank of the primitive signed modular endpoint.  The two
numbers may interact through coding or balancing hypotheses, but they should not
be identified.
\end{remark}

\begin{assumption}[Lipschitz coding]
\label{ass:lipschitz-coding}
Let $S_R=\{\Phi_v:\Lam_{\min}([v])\le R\}$ be the length filtration of the
channel set, and $\overline{\mathcal S}$ its cb-closure.  There is a constant
$C>0$ such that $S_R$ is a $(C/R)$-net of $\overline{\mathcal S}$ for every $R$,
i.e.\ the Hausdorff cb-distance satisfies $\mathrm{dist}_{\cb}(S_R,
\overline{\mathcal S})\le C/R$.
\end{assumption}

\begin{theorem}[Metric dimension is dominated by algebraic dimension]
\label{thm:qmet-qalg}
Under \Cref{ass:lipschitz-coding},
\[
  N_{\cb}(\delta)\ \le\ N_{\CP}(C/\delta)\qquad(\delta>0),
  \qquad\text{hence}\qquad
  \boxed{\ \qmet\ \le\ \qalg.\ }
\]
The inequality is strict in the two collapsing/contracting regimes:
\begin{enumerate}[label=\textup{(\roman*)},leftmargin=2.4em]
\item \textup{(metric collapse)} if the channels are exactly distinguishable by
      length \textup{(}so $\qalg=c>0$\textup{)} but the cb-closure is finite, or
      the tail converges to a finite attractor faster than any power of the
      resolution scale, then $\qmet=0<c$;
\item \textup{(contracting IFS)} in the self-similar setting of
      \Cref{thm:fractal} with the words remaining exactly distinguishable, one
      has $\qalg=\infty$ while $\qmet=s\in(0,\infty)$ is the finite Moran
      dimension, so $\qmet<\qalg$.
\end{enumerate}
\end{theorem}

\begin{proof}
By \Cref{ass:lipschitz-coding} with $R=C/\delta$, the finitely many channels of
$S_{C/\delta}$ form a $\delta$-net of $\overline{\mathcal S}$, so a $\delta$-cover
needs at most $|S_{C/\delta}|=N_{\CP}(C/\delta)$ balls:
$N_{\cb}(\delta)\le N_{\CP}(C/\delta)$.  Taking logarithms and dividing by
$\log(1/\delta)$, write $R=C/\delta$ so that $\log(1/\delta)=\log R-\log C$; then
\[
  \qmet=\limsup_{\delta\downarrow0}\frac{\log N_{\cb}(\delta)}{\log(1/\delta)}
  \le\limsup_{R\to\infty}\frac{\log N_{\CP}(R)}{\log R-\log C}
  =\limsup_{R\to\infty}\frac{\log N_{\CP}(R)}{\log R}=\qalg .
\]
For (i), finite cb-closure has box dimension $0$, and the same conclusion holds
when the tail approaches a finite attractor faster than any power of the
resolution scale; exact distinguishability gives $\qalg=c>0$.  This rate
assumption is essential, since a general countable compact set can have positive
upper box dimension.  For (ii), \Cref{thm:fractal}
gives $\qmet=s<\infty$, while exact distinguishability of words makes
$N_{\CP}(R)$ grow super-polynomially, $\qalg=\infty$.
\end{proof}

\begin{remark}[The coding hypothesis is necessary]
\label{rem:coding-necessary}
Without \Cref{ass:lipschitz-coding} the two dimensions are genuinely independent,
consistent with \Cref{rem:two-dimensions}: if a $\delta$-separated family of
channels is realised only at lengths $\gg1/\delta$, then $\qmet$ can exceed
$\qalg$.  The Lipschitz coding is exactly the condition that metric resolution
$\delta$ be achievable at length $O(1/\delta)$, which is what forces the metric
dimension below the algebraic one.  Thus $\qalg$ bounds $\qmet$ precisely when
short histories already resolve the channel attractor.
\end{remark}

\begin{assumption}[Uniform coding regularity]
\label{ass:uniform-separation}
In addition to the $(C/R)$-net property of \Cref{ass:lipschitz-coding}, the
length filtration is \emph{uniformly cb-separated at the coding rate}: there
exist $c>0$ and $R_0$ such that any two distinct accumulated channels
$\Phi_u\neq\Phi_v$ with $\Lam_{\min}([u]),\Lam_{\min}([v])\le R$ satisfy
\[
  \mathrm{dist}_{\cb}(\Phi_u,\Phi_v)\ \ge\ \frac{c}{R}\qquad(R\ge R_0).
\]
Equivalently, the $N_{\CP}(R)$ distinct length-$\le R$ channels are
$(c/R)$-separated: short histories not only resolve the attractor
\textup{(}\Cref{ass:lipschitz-coding}\textup{)} but fill it without collision.
\end{assumption}

\begin{theorem}[Metric and algebraic predictive dimensions coincide]
\label{thm:qmet-qalg-equality}
Under \Cref{ass:lipschitz-coding,ass:uniform-separation},
\[
  N_{\CP}(R)\ \le\ N_{\cb}\!\Bigl(\tfrac{c}{3R}\Bigr)
  \qquad\text{and}\qquad
  N_{\cb}(\delta)\ \le\ N_{\CP}(C/\delta),
\]
and consequently
\[
  \boxed{\ \qmet=\qalg.\ }
\]
\end{theorem}

\begin{proof}
The upper bound $N_{\cb}(\delta)\le N_{\CP}(C/\delta)$ and hence
$\qmet\le\qalg$ are \Cref{thm:qmet-qalg}.  For the reverse, set
$\delta_R=c/(3R)$.  A $\cb$-ball of radius $\delta_R$ has diameter
$2\delta_R=\tfrac{2c}{3R}<\tfrac{c}{R}$, strictly below the separation of
\Cref{ass:uniform-separation}, so it contains at most one of the
$N_{\CP}(R)$ distinct length-$\le R$ channels.  Hence any $\delta_R$-cover of
$\overline{\mathcal S}\supseteq S_R$ needs at least $N_{\CP}(R)$ balls:
$N_{\cb}(\delta_R)\ge N_{\CP}(R)$.  Therefore
\[
  \qmet=\limsup_{\delta\downarrow0}\frac{\log N_{\cb}(\delta)}{\log(1/\delta)}
  \ \ge\ \limsup_{R\to\infty}\frac{\log N_{\cb}(\delta_R)}{\log(3R/c)}
  \ \ge\ \limsup_{R\to\infty}\frac{\log N_{\CP}(R)}{\log R+\log(3/c)}
  =\qalg,
\]
the last equality because the additive constant $\log(3/c)$ does not affect the
$\limsup$.  Combining the two inequalities gives $\qmet=\qalg$.
\end{proof}

\begin{remark}[Separation is exactly what promotes the inequality to equality]
\label{rem:equality-sharp}
\Cref{ass:uniform-separation} is sharp in the following sense: it fails in the
two strict regimes of \Cref{thm:qmet-qalg}, showing exactly why the inequality
need not be an equality.  In the metric-collapse case (i) the channels
cb-converge, so no fixed $c$ keeps them $(c/R)$-separated, and indeed $\qmet=0<
\qalg$.  In the contracting-IFS case (ii) the words stay exactly distinguishable,
so $N_{\CP}(R)$ grows super-polynomially while a bounded attractor admits only
$O(R^{\qmet})$ points at mutual distance $c/R$; the separation is impossible and
$\qmet<\qalg=\infty$.  Thus the algebraic and metric predictive dimensions
coincide exactly when the predictive classes are \emph{optimally distributed} ---
simultaneously a net (\Cref{ass:lipschitz-coding}) and a packing
(\Cref{ass:uniform-separation}) at the common rate $1/R$ --- and the gap
$\qmet<\qalg$ measures the failure of that optimality.  The integer Clifford rank
$\dCl$ of the signed endpoint is unrelated to either
\textup{(}\Cref{rem:noninteger-metric-dimension}\textup{)}.
\end{remark}





\begin{proposition}[Renewal calibration]
\label{prop:renewal-calibration}
In the conformal contracting case, the equation defining $\qmet$ can be written as a pressure-zero condition.  If $r_\sigma=e^{-\ell_\sigma}$, then $s=\qmet$ is the unique value for which
\[
 \Phi(s):=\sum_{\sigma\in E}e^{-s\ell_\sigma}=1.
\]
The renewal resolvent $(1-\Phi(s))^{-1}$ has residue equal to the reciprocal of the mean calibrated length, whenever that mean is finite.
For $|E|\ge2$ this root is positive and interior; the one-letter alphabet
$|E|=1$ gives the boundary root $s=0$ directly.
\end{proposition}

\begin{proof}
The identity $r_\sigma^s=e^{-s\ell_\sigma}$ rewrites the Moran equation as the
Bowen pressure-zero condition $\Phi(s)=1$ \cite{bowen1975,falconertechniques}.
Existence and uniqueness of the root are elementary for $|E|\ge2$: $\Phi$ is
continuous and strictly decreasing with $\Phi(0)=|E|>1$ and
$\Phi(s)\to0$ as $s\to\infty$, so the intermediate value theorem gives a
unique positive root.  For $|E|=1$ the equation $e^{-s\ell}=1$ with
$\ell>0$ forces $s=0$, and uniqueness is immediate.
The residue statement is the elementary renewal theorem: if the calibrated
probabilities are $p_\sigma=e^{-s\ell_\sigma}$ and $\mu=\sum p_\sigma\ell_\sigma$,
then $1-\Phi(s+t)=\mu t+o(t)$.
\end{proof}

% =====================================================================


\section{Renewal-tail singularities}
\label{app:renewal-tail}

\subsection{Renewal tail and the analytic type of the resolvent}
\label{subsec:renewal-singularity}

The preceding remark can be made precise.  This is the point at which the
\emph{renewal} character of the data---rather than only its automaton or
transition-monoid structure---becomes load-bearing.  Fix a recurrent predictive
component and a base class $x_0$.  Using the calibration of
\Cref{prop:renewal-calibration} at the Bowen critical point, normalised to
$z_c=1$, let $(f_n)_{n\geq1}$ be the first-return-chain probability law and set
$F(z)=\sum_{n\geq1}f_nz^n$.  If $u_n$ denotes the corresponding calibrated
return-chain weight, the first-passage decomposition identifies the rank-graded
diagonal of the chain-counting resolvent of \Cref{prop:incidence-zeta} as
\[
  \Acc_{\Pp}(x_0,x_0;z)=\sum_{n\geq0}u_nz^n=(1-F(z))^{-1},
  \qquad
  u_0=1,\qquad
  u_n=\sum_{k=1}^n f_k u_{n-k}.
\]
Thus the diagonal spectral object constructed from the predictive order is a
renewal generating function, and the tail of the first-return law measures the
order of memory.

\begin{theorem}[Renewal singularity dichotomy]
\label{thm:renewal-dichotomy}
Let $\mu=\sum_{n\geq1}n f_n\in(0,\infty]$ be the calibrated mean return length.
\begin{enumerate}[label=\textup{(\roman*)},leftmargin=2.4em]
\item \textup{(short memory)} If $\mu<\infty$, then
      \[
        \Acc_{\Pp}(x_0,x_0;z)\sim\frac{1}{\mu(1-z)}
        \qquad (z\uparrow1).
      \]
      If, in addition, the first-return law has an exponential moment---as for
      a finite irreducible aperiodic recurrent graph---then $F$ is analytic in
      a neighbourhood of $z=1$ and $\Acc_{\Pp}(x_0,x_0;z)$ has a genuine simple
      pole there, with
      \[
        \Res_{z=1}\Acc_{\Pp}(x_0,x_0;z)=-\frac1\mu.
      \]
      The singularity exponent is the integer $1$.
\item \textup{(long memory)} Suppose instead that, for some $0<\alpha<1$,
      \[
        \overline F(n):=\sum_{k>n}f_k=n^{-\alpha}\ell(n),
      \]
      where $\ell$ is slowly varying.  Then $\mu=\infty$ and Karamata's
      Tauberian theorem gives
      \[
        1-F(z)\sim \Gamma(1-\alpha)(1-z)^\alpha
        \ell\!\left(\frac1{1-z}\right),
        \qquad z\uparrow1.
      \]
      Consequently
      \[
        \boxed{\quad
        \Acc_{\Pp}(x_0,x_0;z)\sim
        \frac{(1-z)^{-\alpha}}
        {\Gamma(1-\alpha)\,\ell\!\left(\frac1{1-z}\right)}
        \qquad(z\uparrow1).
        \quad}
      \]
      The critical singularity has non-integer exponent $-\alpha$.  In
      particular, whenever this regularly varying equivalent is realised in a
      slit complex neighbourhood of $1$, the critical point is a branch point
      and the resolvent cannot be meromorphic across it.
\end{enumerate}
\end{theorem}

\begin{proof}
For \textup{(i)},
\[
  \frac{1-F(z)}{1-z}=\sum_{n\geq1}f_n(1+z+\cdots+z^{n-1})
  \longrightarrow \sum_{n\geq1}n f_n=\mu
\]
by monotone convergence as $z\uparrow1$.  Inversion gives the stated
pole-type asymptotic.  Under an exponential moment, $F$ extends analytically
past the unit circle near $1$ and $F'(1)=\mu\neq0$, so $1-F$ has a simple zero
and the residue formula follows.  For \textup{(ii)}, the power-series form of
Karamata's Tauberian theorem for non-negative coefficients gives the displayed
equivalent for $1-F(z)$ from the regularly varying tail
\cite{binghamgoldieteugels1987,feller}; inversion gives the resolvent
asymptotic.  A non-integer algebraic singularity in a slit neighbourhood is not
meromorphic.
\end{proof}

\begin{corollary}[Renewal counting asymptotics]
\label{cor:renewal-counting}
Under the hypotheses of \Cref{thm:renewal-dichotomy}\textup{(ii)}, the
cumulative return weight satisfies the unconditional Tauberian asymptotic
\[
  \sum_{k\leq n}u_k
  \sim \frac{\sin(\pi\alpha)}{\pi\alpha}\,
  \frac{n^\alpha}{\ell(n)},
  \qquad n\to\infty.
\]
If, in addition, the strong renewal condition holds---automatically for
$\alpha\in(\tfrac12,1)$ by Garsia--Lamperti, and under the sharp
Caravenna--Doney criterion for $\alpha\in(0,\tfrac12]$---then the local
asymptotic
\[
  u_n\sim \frac{\sin(\pi\alpha)}{\pi}\,
  \frac{n^{\alpha-1}}{\ell(n)}
\]
also holds
\cite{garsialamperti1962,erickson1970,doney1997,caravennadoney2019}.
\end{corollary}

\begin{proof}
The cumulative statement follows from the Hardy--Littlewood--Karamata theorem
applied to the non-negative coefficients of
$\Acc_{\Pp}(x_0,x_0;z)$, using
$\Gamma(1+\alpha)\Gamma(1-\alpha)=\pi\alpha/\sin(\pi\alpha)$.
The local statement is precisely the strong renewal theorem under the cited
range and criterion.
\end{proof}

\begin{definition}[Renewal exponent]
\label{def:renewal-exponent}
The number $\alpha\in(0,1]$ of \Cref{thm:renewal-dichotomy} is the
\emph{renewal exponent} of the recurrent predictive component: $\alpha=1$ in
the finite-mean pole regime and $\alpha\in(0,1)$ in the regularly varying
branch regime.
\end{definition}

\begin{remark}[A genuinely renewal-theoretic invariant]
\label{rem:renewal-invariant}
The renewal exponent is not determined by either predictive dimension.  Two
recurrent components may have the same algebraic growth exponent $\qalg$
\textup{(}\Cref{thm:triple}\textup{)} and the same metric dimension $\qmet$
\textup{(}\Cref{def:qmet}\textup{)}, while having different first-return tails
and hence different values of $\alpha$.  The quantities $\qalg$ and $\qmet$
measure, respectively, exact predictive growth and metric resolution, whereas
$\alpha$ measures the order of contact of the predictive-order resolvent with
its critical singularity.  It is therefore a third invariant of specifically
renewal-theoretic origin, integer in the short-memory regime and fractional in
the heavy-tailed regime.  This is the precise content of
\Cref{rem:markov-renewal}.
\end{remark}

\begin{example}[Parabolic renewal component realising $\alpha\in(0,1)$]
\label{ex:parabolic-renewal}
A finite recurrent-state graph has geometric return tails, so the long-memory
regime requires a countable recurrent \emph{induced component}; this is
compatible with a finite reset alphabet.  Consider the renewal structure
obtained by inducing a Pomeau--Manneville intermittent source
\cite{pomeaumanneville1980} on a base away from its indifferent fixed point.
For an indifferent fixed point of order $1+1/\alpha$, the induced
first-return law has
\[
  f_n\asymp n^{-1-\alpha},\qquad 0<\alpha<1,
\]
in the standard parameterisation of the model
\cite{aaronson1997,sarig2001}.  The induced predictive component is countable
and locally finite; whenever its calibrated shifts satisfy
\Cref{ass:fibres}, \Cref{thm:renewal-dichotomy}\textup{(ii)} applies and the
chain-counting resolvent has critical exponent $-\alpha$.  Varying the
indifference order sweeps $\alpha$ through $(0,1)$, so the fractional regime is
realised rather than merely formally admissible.
\end{example}


\begin{proposition}[Commuting order and interval growth]
\label{prop:commuting-order}
If the predictive monoid is freely generated by $c$ commuting channels modulo no further relations, then $\Pp\cong\N^c$ with product order.  For balanced intervals of rank separation $\tau$,
\[
 |[x,y]|\asymp \tau^c.
\]
Consequently the interval-growth dimension of the predictive order agrees with the algebraic predictive dimension $\qalg=c$.
\end{proposition}

\begin{proof}
The divisibility order of the free commutative monoid on $c$ generators is the product order on $\N^c$.  A box with coordinate gaps $n_1,\ldots,n_c$ has cardinality $\prod_i(n_i+1)$.  Along balanced intervals $n_i\simeq\tau/c$, this grows like $\tau^c$.
\end{proof}

\begin{theorem}[Ehrhart interval growth for normal affine commuting monoids]
\label{thm:ehrhart-order}
Let the predictive monoid $\M$ be a finitely generated, commutative,
cancellative, torsion-free, reduced affine semigroup.  Write
$L=\mathrm{Gr}(\M)\cong\Z^{\rmon}$ for the lattice generated by $\M$, and let
\[
  C=\R_{\ge0}\!\cdot\!\M\subseteq L\otimes\R
\]
be the salient rational polyhedral cone generated by $\M$.  Assume, for this
Ehrhart statement, that $\M$ is normal in its lattice:
\[
  \M=C\cap L .
\]
Then:
\begin{enumerate}[label=\textup{(\roman*)},leftmargin=2.4em]
\item \textup{(geometric model)} the order interval of $x\preceq y$ with
      $w:=y-x\in\M$ is
      \[
        [x,y]\ \cong\ \M\cap(w-\M)
        =\bigl(C\cap(w-C)\bigr)\cap L,
      \]
      a finite set, since $C$ is salient;
\item \textup{(Ehrhart asymptotics)} fix an interior lattice direction
      $v_0\in\operatorname{int}C\cap L$ and consider balanced intervals
      $w=\tau v_0$, $\tau\in\N$.  Then
      \[
        |[x,y]|
        =\operatorname{Vol}_L(P_{v_0})\,\tau^{\rmon}
        +O(\tau^{\rmon-1}),
        \qquad
        P_{v_0}:=C\cap(v_0-C),
      \]
      where $\operatorname{Vol}_L$ is normalised so that a fundamental domain of
      $L$ has volume $1$.
\end{enumerate}
Consequently, in this normal affine regime, the interval-growth dimension equals
the algebraic predictive dimension $\qalg=\rmon$.
\end{theorem}

\begin{proof}
Because $\M=C\cap L$, divisibility $z\in[x,y]$ is equivalent to
$z-x\in C\cap L$ and $y-z\in C\cap L$.  Writing $w=y-x$, this gives
\[
  [x,y]-x=\bigl(C\cap(w-C)\bigr)\cap L .
\]
The cone is salient, so $C\cap(w-C)$ is bounded and the interval is finite.

For $w=\tau v_0$, homogeneity gives
\[
  C\cap(\tau v_0-C)=\tau\bigl(C\cap(v_0-C)\bigr)=\tau P_{v_0}.
\]
The set $P_{v_0}$ is a rational polytope.  Ehrhart's theorem
\cite{stanley,brunsgubeladze,beckrobins2007}, applied with respect to the
intrinsic lattice $L$, gives
\[
  \#(\tau P_{v_0}\cap L)
  =\operatorname{Vol}_L(P_{v_0})\tau^{\rmon}+O(\tau^{\rmon-1}),
\]
which is the stated interval asymptotic.  The final assertion follows because
$\rmon$ is the rank of the affine semigroup, hence the polynomial growth exponent
in the commuting affine regime.
\end{proof}

\begin{remark}[Non-normal affine monoids]
\label{rem:non-normal-ehrhart}
For a non-normal affine semigroup the same exponent $\rmon$ is still expected
from Hilbert--Serre theory, but the leading coefficient can depend on the
semigroup module structure and may not be exactly the lattice volume of
$C\cap(v_0-C)$ without correction.  This is why \Cref{thm:ehrhart-order} is
stated in the normal affine case.  The paper only needs the exponent and the
intrinsic-lattice normalisation; any non-normal refinement should be treated as
a separate Hilbert-series calculation.
\end{remark}

\begin{corollary}[Free case and consistency]
\label{cor:ehrhart-free}
For $\M\cong\N^{c}$ one has $L=\Z^{c}$, $C=\R_{\ge0}^{c}$,
$v_0=(1,\dots,1)$, $P_{v_0}=[0,1]^{c}$ and $\operatorname{Vol}_L(P_{v_0})=1$,
so $|[x,y]|\sim\tau^{c}$, recovering \Cref{prop:commuting-order}; and the
interval-growth dimension $\rmon=c$ agrees with $\qalg=c$ of \Cref{thm:commuting}.
\end{corollary}

\begin{proof}
Specialise \Cref{thm:ehrhart-order}: the unit cube has unit lattice volume, and
$\tau^{c}+O(\tau^{c-1})\asymp\tau^{c}$.
\end{proof}

\begin{remark}[Why the intrinsic lattice is essential]
\label{rem:intrinsic-lattice}
The volume in \Cref{thm:ehrhart-order}\,(ii) must be normalised against
$L=\mathrm{Gr}(\M)$, not against any ambient lattice in which $\M$ happens to be
written.  For instance $\M=\langle(1,0),(1,2)\rangle\subset\Z^2$ has
$\mathrm{Gr}(\M)=\Z\times2\Z$ of index $2$, and the balanced interval count grows
like $\tfrac12\operatorname{Area}(P_{v_0})\tau^{2}$, exactly half the ambient
Euclidean value: the leading coefficient sees only the predictive lattice the
monoid actually generates.  This is the order/volume analogue of measuring
predictive content intrinsically rather than through a presentation.
\end{remark}



\begin{remark}[Order is not signature]
\label{rem:order-not-signature}
The order gives a causal skeleton and a time separation.  It does not by itself give an indefinite inner product.  The positive CP sector remains Hilbert-positive even when the predictive quotient is partially ordered.  This distinction is essential: order supplies causal comparability, while a Lorentzian signature requires a fundamental symmetry.
\end{remark}

% =====================================================================


\section{The deterministic product-order continuum}
\label{app:product-order}

\subsection{The renewal continuum ladder}
\label{sec:continuum}

The assembly theorem identifies the finite kinematic data.  We next ask what
their large-scale order geometry becomes in the most regular case,
\[
  \W_{\CP}\cong\N^c.
\]
This subsection gives a deliberately sharp answer.  The deterministic
commuting model has a continuum limit in every dimension, and the limit is a
perfectly good measured Lorentzian pre-length space.  But its cone is the
product cone.  Hence, except in the two-dimensional null-coordinate case, the
deterministic product order gives a polyhedral Lorentzian continuum rather than
a smooth Minkowski cone.  This is not a failure of the construction; it is the
reason the later Clifford-isotropisation step is needed.

\subsubsection{The discrete object and its rescaling}

We first isolate the purely order-theoretic continuum limit, before introducing
any Clifford data.  Throughout this subsection the predictive monoid is freely
generated by $c$ commuting reset channels with no further relations, so that
$\W_{\CP}\cong\N^c$ carries the product predictive order of
\Cref{prop:commuting-order}.  Let $\ell_1,\dots,\ell_c>0$ be calibrated
generator lengths, so that the quotient length is the linear rank
\[
  \Lam_{\min}(x)=\sum_{i=1}^c \ell_i x_i,\qquad x\in\N^c,
\]
and let $\mu$ be counting measure.

\begin{definition}[Calibrated rescaling]
\label{def:rescaling}
For $R>0$ let
\[
  \iota_R:\N^c\longrightarrow[0,\infty)^c,
  \qquad
  \iota_R(x)=R^{-1}(\ell_1x_1,\dots,\ell_cx_c),
\]
and let
\[
  \mu_R=\Bigl(\prod_i\ell_i\Bigr)R^{-c}(\iota_R)_*\mu.
\]
On $[0,\infty)^c$ write $X\le Y$ for componentwise order and define the geometric-mean time separation
\[
  \taug(X,Y)=
  \begin{cases}
  \left(\prod_{i=1}^c (Y_i-X_i)\right)^{1/c}, & X\le Y,\\[2pt]
  0, & \text{otherwise.}
  \end{cases}
\]
This separation is positive exactly for chronologically related pairs $X_i<Y_i$ for all $i$; it vanishes on the null faces of the product cone.
\end{definition}

The rescaling keeps three pieces of data visible at once: the order, the
rank/time function, and counting measure.  The first elementary observation is
that these already form a measured causal set.

\begin{proposition}[Renewal measured causal set]
\label{prop:measured-causet}
The tuple $(\W_{\CP},\preceq,\Lam_{\min},\mu)$ is a measured causal set: the rank $\Lam_{\min}$ is strictly increasing along strict chains, $(\W_{\CP},\preceq)$ is locally finite, every order interval has finite $\mu$-mass, and level sets of $\Lam_{\min}$ are antichains.
\end{proposition}

\begin{proof}
If $x\prec y$ strictly then $x_i\le y_i$ for all $i$ and strict inequality holds for at least one $i$, so
\[
  \Lam_{\min}(y)-\Lam_{\min}(x)=\sum_i\ell_i(y_i-x_i)>0.
\]
An interval is the finite box
\[
  [x,y]=\prod_i\{x_i,x_i+1,\dots,y_i\},
\]
with $\mu([x,y])=\prod_i(y_i-x_i+1)<\infty$.  Thus intervals are finite and the order is locally finite.  If two comparable points have the same rank, the displayed strict-increase identity forces them to be equal; hence rank level sets are antichains.
\end{proof}

\begin{theorem}[Polyhedral Lorentzian continuum]
\label{thm:taxicab-limit}
The ordered space $([0,\infty)^c,\le,\taug)$ is a Lorentzian pre-length space with compact diamonds on bounded regions.  On chronologically related triples it satisfies the reverse triangle inequality.  Moreover, on every bounded diamond, the rescaled renewal geometries
\[
  (\iota_R(\N^c),\le,\taug,\mu_R)
\]
converge locally to $([0,\infty)^c,\le,\taug,\Leb)$ in the order/length/measure sense, and interval volumes converge:
\[
  \mu_R\bigl([\iota_R(x_R),\iota_R(y_R)]\bigr)
  \longrightarrow
  \Leb([X,Y])=
  \prod_{i=1}^c(Y_i-X_i)
\]
whenever $\iota_R(x_R)\to X$, $\iota_R(y_R)\to Y$, and $X\le Y$.
\end{theorem}

\begin{proof}
Let
\[
  m(v)=\left(\prod_i v_i\right)^{1/c}
\]
on $(0,\infty)^c$.  The function $m$ is concave and homogeneous of degree one; hence it is superadditive.  Therefore, for $X\le Y\le Z$ with all coordinate gaps positive,
\[
  \taug(X,Z)=m((Y-X)+(Z-Y))\ge m(Y-X)+m(Z-Y)
  =\taug(X,Y)+\taug(Y,Z).
\]
The same inequality extends to boundary cases by continuity, with possible zero time separation on null faces.  Diamonds in bounded regions are compact boxes.

For volume convergence, write $X_i=R^{-1}\ell_i(x_R)_i$ and $Y_i=R^{-1}\ell_i(y_R)_i$.  Then
\[
\begin{aligned}
  \mu_R\bigl([\iota_R(x_R),\iota_R(y_R)]\bigr)
  &=\Bigl(\prod_i\ell_i\Bigr)R^{-c}
    \prod_i\bigl((y_R)_i-(x_R)_i+1\bigr) \\
  &\longrightarrow \prod_i(Y_i-X_i),
\end{aligned}
\]
which is the Lebesgue volume of the limiting box.  Finally, the grids $\iota_R(\N^c)$ are $O(R^{-1})$-dense on bounded boxes, preserve order exactly, and distort $\taug$ by $o(1)$ on compact diamonds.  This gives local order/length convergence.
\end{proof}

\begin{remark}[Two times on one order]
\label{rem:two-times}
The product order carries two natural time functions.  The longest-chain or rank time is the linear function $\sum_i(Y_i-X_i)$, while $\taug$ is the order-plus-volume time separation selected by the identity
\[
  \Leb([X,Y])=\taug(X,Y)^c.
\]
They share the same causal cones but are not the same metric datum.  In the two-dimensional case $\taug$ becomes the ordinary Minkowski proper time in null coordinates.
\end{remark}

The product cone becomes a round Lorentz cone only in the exceptional case
$c=2$.  This is the familiar fact that two-dimensional Minkowski space is most
naturally described by two null coordinates.

\begin{theorem}[Two-dimensional Minkowski emergence]
\label{thm:minkowski-2d}
For $c=2$, the limit $([0,\infty)^2,\le,\taug,\Leb)$ is isometric, up to the constant normalization of volume, to a causally convex region of $(1+1)$-dimensional Minkowski space.  In null coordinates $u=X_1$, $v=X_2$ and inertial coordinates
\[
  t=\frac{u+v}{2},\qquad \xi=\frac{v-u}{2},
\]
the product order is the Minkowski causal order, and
\[
  \taug(X,Y)=\sqrt{\Delta u\,\Delta v}
  =\sqrt{\Delta t^2-\Delta\xi^2}.
\]
\end{theorem}

\begin{proof}
The change of variables gives
\[
  \Delta u\,\Delta v=\Delta t^2-\Delta\xi^2.
\]
The condition $\Delta u\ge0$ and $\Delta v\ge0$ is exactly $\Delta t\ge |\Delta\xi|$, the future cone of the flat metric $-dt^2+d\xi^2$.  The time separation $\taug$ is therefore the usual proper time.  Also $du\,dv=2\,dt\,d\xi$, so counting volume fixes the conformal factor up to the displayed constant normalization.
\end{proof}

\begin{theorem}[Polyhedral obstruction to smooth higher-dimensional emergence]
\label{thm:obstruction}
For $c\ge3$, the deterministic product-order cone $[0,\infty)^c$ is not linearly isomorphic to a round Lorentz cone
\[
  L_c=\{x_0\ge |(x_1,\dots,x_{c-1})|\}.
\]
Consequently the deterministic finite commuting renewal model does not produce a smooth higher-dimensional Minkowski causal cone from product order and counting volume alone.
\end{theorem}

\begin{proof}
The orthant $[0,\infty)^c$ is a salient closed convex cone with exactly $c$ extreme rays, the coordinate axes.  The Lorentz cone $L_c$ has as extreme rays its null generators, the cone over $S^{c-2}$.  For $c\ge3$ this set is infinite, indeed a continuum.  Linear isomorphisms preserve extreme rays.  Hence no linear isomorphism carries the orthant cone to the round Lorentz cone.  Since the tangent cones of a smooth Lorentzian manifold are round Lorentz cones, the deterministic product-order continuum is polyhedral rather than smooth Lorentzian in dimensions $c\ge3$.
\end{proof}

\begin{corollary}[Deterministic lattice Lorentz violation]
\label{cor:lorentz-violation}
For $c\ge3$, the regular commuting renewal continuum selects the $c$ generator directions as distinguished null-like axes.  Thus the deterministic finite-generator model has a preferred polyhedral cone structure and does not possess the full Lorentz symmetry of a smooth Minkowski cone.
\end{corollary}

\subsubsection{Routes beyond the deterministic product cone}

The obstruction above is only an obstruction to the deterministic finite
commuting product cone.  There are two standard ways to escape it.  One is to
replace the lattice by a Lorentz-invariant point process; the other, used in
the rest of this section, is to keep the renewal origin but round the signed
Dirac symbol by isotropic Clifford coarse-graining.

\begin{theorem}[Stochastic renewal sprinkling, established input]
\label{thm:sprinkling}
If the deterministic lattice is replaced by a Poisson sprinkling into a causally
convex region of $(1+k)$-dimensional Minkowski space, with intensity proportional
to Minkowski volume, then the construction is Lorentz covariant: applying a
Lorentz transformation to the region carries the law of the sprinkled causal set
to the law of the transformed sprinkled causal set.  In Minkowski diamonds, the
induced longest-chain function recovers Minkowski proper time in probability
after the standard dimension-dependent normalisation \cite{Myrheim1978,brightwellgregory1991}.  Hence the obstruction of
\Cref{thm:obstruction} is an obstruction to the deterministic commuting
realisation, not to renewal-type point processes in general.
\end{theorem}

\begin{remark}
The renewal input in \Cref{thm:sprinkling} is the point-process realisation and
its intensity.  Lorentz covariance and longest-chain convergence are classical
causal-set inputs, not new consequences of the finite commuting predictive
quotient.
\end{remark}

\begin{remark}[Clifford isotropisation]
\label{rem:rounding-pointer}
The Clifford/isotropic route is proved in \Cref{sec:isotropisation}.  The key point is that positive averaging alone cannot create Lorentzian signature: a positive second-moment tensor is positive semidefinite.  Signature is supplied by the signed/Krein sector, while isotropic coarse-graining rounds the principal-symbol cone.
\end{remark}

\begin{remark}[What the continuum ladder establishes]
\label{rem:ladder-summary}
In the regular commuting regime, renewal data yield a measured causal set and an unconditional polyhedral Lorentzian continuum with convergent interval volumes.  Smooth Minkowski space emerges from product order and counting volume alone only for $c=2$, where the product coordinates are null coordinates.  For $c\ge3$, the deterministic finite-generator model is polyhedral and Lorentz-violating.  Higher-dimensional smooth Lorentzian spacetime therefore requires either stochastic sprinkling or an additional isotropic/Clifford coarse-graining mechanism.  No claim is made here of Einstein dynamics, matter coupling, or curvature equations.
\end{remark}

\begin{remark}[Signature is not an extra coordinate]
\label{rem:dimension-count}
The product order on $\N^c$ is already a $c$-dimensional causal order: its rank is a single time function and its generic antichains are $(c-1)$-dimensional spacelike slices.  The signed sector converts the positive order into Lorentzian signature through the fundamental symmetry $J_\chi$; it does not add a coordinate.  Thus the regular model should be read as a $c=1+(c-1)$ dimensional order geometry equipped with a signed/Krein Lorentzian structure.
\end{remark}


% =====================================================================


\section{Internal holonomy and conditional dimension selectors}
\label{app:conditional-dimension}

\subsection{Internal holonomy realization of the spatial block}

\begin{definition}[Coherent recurrent internal holonomy]
\label{def:coherent-internal-holonomy}
Let $U_e$ be the coherent internal transport attached to an oriented recurrent
edge, acting on the same internal Hilbert space and recurrent component as the
renewal channel.  For a recurrent cycle $\alpha$, set
$U_\alpha=\prod_{e\in\alpha}U_e$ and define the internally dressed revision
operator
\[
 R_\alpha^{\mathrm{int}}=\widetilde S_\alpha\otimes U_\alpha.
\]
\end{definition}

\begin{theorem}[Nonabelian internal holonomy supplies a spatial Clifford entry]
\label{thm:nonabelian-holonomy-spatial-entry}
Suppose two recurrent spatial classes $\alpha,\beta\in H$ have coherent internal
holonomies satisfying the scalar group-commutator identity
\[
 U_\alpha U_\beta U_\alpha^{-1}U_\beta^{-1}
 =-\one.
\]
Then the internally dressed revisions satisfy
\[
 \omega(\alpha,\beta)=1.
\]
More generally, a family of recurrent holonomies whose scalar commutators realise
a nondegenerate alternating form on a spatial subspace produces exactly that
spatial block of $\omega$.  Combined with the forced temporal row, a
nondegenerate rank-four restriction yields the primitive
$\Cl_4(\C)\cong M_4(\C)$ revision factor appropriate to a $3+1$ endpoint.
\end{theorem}

\begin{proof}
The commutator of the dressed revisions factorises into the flat sheet/base
commutator and the internal group commutator.  The former is trivial by
\Cref{prop:flat-spatial-block-zero}; the latter is $-\one$ by hypothesis.  Thus
the total scalar commutator is $-\one$, which is the statement
$\omega(\alpha,\beta)=1$.  The final assertion follows by applying the same
factorisation entrywise and then the rank criterion to a nondegenerate
four-dimensional restriction.
\end{proof}

\begin{remark}[One coherent transport family upstream and downstream]
\label{rem:shared-coherent-transport}
The operators $U_e$ in \Cref{def:coherent-internal-holonomy} are the coherent
recurrent transports of the Lorentzian construction.  Any downstream matter
construction using colour holonomy must restrict this same family on the same
recurrent component, rather than introduce an analogous second family.  With
that literal identification, spatial Clifford dimension and complex colour can
be traced to one coherent internal transport system.
\end{remark}

\begin{lemma}[Gauge invariance of the spatial commutator form]
\label{lem:spatial-omega-gauge-invariant}
Suppose the coherent recurrent transports form a scalar projective
representation on the internal recurrent sector,
\[
 U_\alpha U_\beta=\tau_H(\alpha,\beta)U_{\alpha+\beta},
 \qquad \tau_H(\alpha,\beta)\in\mathrm U(1),
\]
whose scalar commutator takes values in $\{\pm1\}$.  Then
\[
 (-1)^{\omega_H(\alpha,\beta)}
 :=U_\alpha U_\beta U_\alpha^{-1}U_\beta^{-1}
 =\tau_H(\alpha,\beta)\tau_H(\beta,\alpha)^{-1}
\]
depends only on the projective cohomology class $[\tau_H]$.  In particular it
is unchanged under scalar regauging $U_\alpha\mapsto\mu(\alpha)U_\alpha$.
\end{lemma}

\begin{proof}
A scalar regauging replaces $\tau_H$ by $\tau_H\,\delta\mu$.  Since the
coboundary $\delta\mu(\alpha,\beta)=\mu(\alpha)\mu(\beta)
\mu(\alpha+\beta)^{-1}$ is symmetric on the elementary abelian group $H$, its
alternation is trivial.  Hence the commutator bicharacter, and therefore
$\omega_H$, is gauge invariant.
\end{proof}

\begin{theorem}[Nonzero spatial block forces nonabelian recurrent transport]
\label{thm:nonabelian-holonomy-necessary}
Work in the local, covariant, internal scalar-projective class in which the
spatial form is realised by the coherent recurrent transports $U_\alpha$ as in
\Cref{lem:spatial-omega-gauge-invariant}.  Then:
\begin{enumerate}[label=\textup{(\roman*)},leftmargin=2.4em]
\item if the recurrent internal transport is abelian, then
      $\omega|_{H\times H}=0$;
\item equivalently, $\omega(\alpha,\beta)=1$ forces
      $U_\alpha U_\beta\ne U_\beta U_\alpha$ in every scalar gauge;
\item if $\rank_{\F_2}\omega_H=2m$, every faithful irreducible internal
      realisation has dimension $2^m$, and every realisation has internal
      dimension at least $2^m$ on the nondegenerate sector.
\end{enumerate}
Thus a nonzero spatial modular block cannot be carried by the abelian deck
algebra or by one-dimensional internal transport.
\end{theorem}

\begin{proof}
If all $U_\alpha$ commute, every group commutator is the identity, hence
$\omega_H=0$.  The contrapositive gives (ii), and
\Cref{lem:spatial-omega-gauge-invariant} makes the conclusion gauge invariant.
For (iii), quotient $H$ by $\rad(\omega_H)$.  The resulting $2m$-dimensional
symplectic $\F_2$-space defines the finite Heisenberg, equivalently extraspecial
$2$-group, central extension.  Finite Stone--von Neumann gives a unique faithful
irreducible projective representation of dimension $2^m$; any representation
containing that nondegenerate commutator sector has dimension at least $2^m$
\cite{karpilovsky,mackey}.
\end{proof}

\begin{theorem}[Internal-holonomy normal form]
\label{thm:internal-holonomy-normal-form}
Let $\omega_H=\omega|_{H\times H}$ be any alternating $\F_2$-form of rank $2m$
occurring in a local, covariant, internal revision sector.  After factoring the
radical and applying a scalar gauge, its nondegenerate spatial revision sector is
unitarily equivalent to the Heisenberg--Clifford transport on $\C^{2^m}$:
\[
 U_\alpha U_\beta U_\alpha^{-1}U_\beta^{-1}
 =(-1)^{\omega_H(\alpha,\beta)}\one.
\]
The representation on the nondegenerate quotient is unique up to unitary
equivalence, while $\rad(\omega_H)$ contributes precisely commuting central
transport directions.  Consequently, within this scalar-projective class,
\[
 \omega|_{H\times H}\ne0
 \quad\Longleftrightarrow\quad
 \text{the recurrent internal transport is nonabelian}.
\]
\end{theorem}

\begin{proof}
Choose a symplectic basis on $H/\rad(\omega_H)$.  The associated finite
Heisenberg extension has commutator pairing $(-1)^{\omega_H}$ and its unique
faithful irreducible representation has dimension $2^m$ by finite
Stone--von Neumann.  Pulling this representation back to $H$ realises the given
form; the radical acts through commuting central characters.  Conversely,
\Cref{thm:nonabelian-holonomy-necessary} shows that a nonzero spatial form cannot
come from abelian transport.  The scalar gauge freedom changes the cocycle by a
coboundary but not its alternation.
\end{proof}

\begin{remark}[Precise scope of the equivalence]
\label{rem:holonomy-equivalence-scope}
The equivalence in \Cref{thm:internal-holonomy-normal-form} concerns the coherent
internal transport that realises the scalar projective spatial revision sector.
It does not assert that every arbitrary nonabelian unitary holonomy automatically
produces a nonzero $\F_2$-valued block: a nonabelian family may have noncentral
commutators and therefore lie outside the scalar-bicharacter regime.  What is
proved is the exact normal form relevant to the modular revision datum of this
paper.
\end{remark}

\subsection{The common cycle-rank invariant}
\label{subsec:dimension-coincidence}

The positive predictive sector, the recurrent crystal, and the signed revision
sector introduce dimension by different constructions.  In general the resulting
numbers need not agree.  The metric link is genuinely independent: contracting
UCP memories can have positive algebraic growth but zero channel box dimension.
The remaining algebraic, modular, and heat-kernel links become consequences of
one regular topological-crystal package.

\begin{definition}[Finite displacement semigroupoid and accessed lattice]
\label{def:displacement-semigroupoid}
Let \(Q=(V,A)\) be a finite strongly connected directed control graph with base
state \(s_0\).  Each oriented edge \(e\) carries a displacement
\(\xi(e)\in\Gamma\), where \(\Gamma\cong\Z^{\rmon}\), and an additive positive
degree \(\ell(\xi(e))>0\).  For states \(s,t\), write
\[
 \M_{st}:=\left\{\sum_{e\in p}\xi(e):p:s\to t\right\}.
\]
These sets form a semigroupoid:
\(\M_{st}+\M_{tu}\subseteq\M_{su}\).  The return lattice is
\[
 \Gamma_{\rm ret}:=\Gr(\M_{s_0s_0})\subseteq\Gamma ,
\]
and the spatial accessed lattice is
\[
 \Lat:=\ker\!\left(\ell|_{\Gamma_{\rm ret}}\right),
 \qquad
 \beff:=\rank_{\Z}\Lat.
\]
\end{definition}

\begin{assumption}[Full-dimensional affine reachability]
\label{ass:full-rank-reachability}
There is a finitely generated affine monoid \(\M\subseteq\Gamma\) such that:
\begin{enumerate}[label=\textup{(\alph*)},leftmargin=2.2em]
\item every element of \(\M\) is realised by a path starting at \(s_0\);
\item \(\Gr(\M)=\Gamma\cong\Z^{\rmon}\) and
      \(\spn_{\R}\M=\Gamma\otimes\R\);
\item the quotient length is comparable to a strictly positive additive degree
      \(\ell:\Gamma\to\Z\);
\item \(\Ncp(R)\asymp R^{\rmon}\).
\end{enumerate}
Normal affine monoids of \Cref{thm:ehrhart-order} satisfy the growth clause, but
normality is not required for the return-rank argument below.
\end{assumption}

\begin{assumption}[Finite primitive control]
\label{ass:finite-primitive-control}
The control graph \(Q\) is finite, strongly connected, and aperiodic.
\end{assumption}

\begin{theorem}[Full-rank return completion]
\label{thm:full-rank-returns}
Under \Cref{ass:full-rank-reachability,ass:finite-primitive-control},
\[
 \rank_{\Z}\Gamma_{\rm ret}=\rmon .
\]
Consequently
\[
 \boxed{\beff=\rmon-1,\qquad \qalg=1+\beff.}
\]
\end{theorem}

\begin{proof}
For each state \(s\), choose a fixed return path \(q_s:s\to s_0\), and let
\(R_s\in\Gamma\) be its displacement.  Since \(V\) is finite,
\(K:=\max_s\|R_s\|<\infty\).  If \(x\in\M\) is realised by a path
\(p:s_0\to s\), then \(pq_s\) is a closed walk at \(s_0\), so
\(x+R_s\in\Gamma_{\rm ret}\).  Hence every \(x\in\M\) lies within distance
\(K\) of the real span \(W:=\R\Gamma_{\rm ret}\).

Choose \(u_1,\ldots,u_{\rmon}\in\M\) spanning \(\Gamma\otimes\R\).
Because \(\M\) is a monoid, \(nu_i\in\M\) for every \(n\ge1\).  Therefore
\[
 n\,\operatorname{dist}(u_i,W)
 =\operatorname{dist}(nu_i,W)\le K
\]
for all \(n\), which forces \(u_i\in W\).  Thus
\(W=\Gamma\otimes\R\) and \(\rank\Gamma_{\rm ret}=\rmon\).
Since \(\ell\) is nonzero on \(\Gamma_{\rm ret}\), its kernel \(\Lat\) has
rank \(\rmon-1\).  Finally \(\qalg=\rmon\) by
\Cref{ass:full-rank-reachability}\,(d).
\end{proof}

\begin{remark}[What the theorem does and does not use]
\label{rem:return-rank-scope}
The proof uses bounded return completion and full-dimensional scalable
reachability, not saturation of the affine monoid.  Normality remains useful for
the Ehrhart formula and for eliminating lattice holes, but it is not the
central hypothesis in \Cref{thm:full-rank-returns}.  No claim is made that
every vector in \(\ker\ell\cap\Gamma\) is represented by a single prescribed
difference of equal-degree closed walks; only the finite-index rank statement is
needed.
\end{remark}

For \(R\ge1\), let
\[
 S_R:=\{\Phi_v:\Lambda_{\min}([v])\le R\},\qquad
 \Ncp(R):=|S_R|,
\]
and let \(\Ncb(\delta)\) be the least number of cb-metric balls of radius
\(\delta\) covering the channel closure.

\begin{assumption}[Two-sided channel coding]
\label{ass:two-sided-channel-coding}
There are constants \(C\ge c>0\) such that \(S_R\) is a \(C/R\)-net of the
channel closure and distinct elements of \(S_R\) are separated by at least
\(c/R\) in cb-distance.
\end{assumption}

\begin{theorem}[Metric--predictive coincidence]
\label{thm:metric-predictive-coincidence}
Under \Cref{ass:two-sided-channel-coding},
\[
 \Ncb(C/R)\le\Ncp(R)\le\Ncb(c/(3R)),
 \qquad
 \boxed{\qmet=\qalg.}
\]
\end{theorem}

\begin{proof}
The net property gives the first inequality.  A ball of radius \(c/(3R)\)
contains at most one point of the \(c/R\)-separated set \(S_R\), giving the
second.  Taking logarithms and limsups along reciprocal integer scales proves
the equality.
\end{proof}

\begin{proposition}[Metric collapse without non-collapse]
\label{prop:metric-collapse}
There are genuine UCP renewal memories with \(\qalg>0\) and \(\qmet=0\).
Specifically, on \(M_d(\C)\) let \(\Psi(x)=d^{-1}\tr(x)\one\) and
\[
 \Phi_a=\tfrac12\Psi+\tfrac12\id,\qquad
 \Phi_b=\tfrac13\Psi+\tfrac23\id.
\]
Then the predictive quotient is \(\N^2\), so \(\qalg=2\), while all degree-\(n\)
channels lie within \(O(2^{-n})\) of \(\Psi\), whence
\(\Ncb(\delta)=O((\log(1/\delta))^2)\) and \(\qmet=0\).
\end{proposition}

\begin{proof}
A word with contents \((i,j)\) gives
\(\Phi_{i,j}=(1-2^{-i}3^{-j})\Psi+2^{-i}3^{-j}\id\).
Unique factorisation makes the classes distinct, while exponential contraction
to \(\Psi\) gives the stated covering bound.
\end{proof}

The modular comparison must use the actual accessed lattice, not a larger
finite-index lattice before reduction modulo \(2\).

\begin{definition}[Accessible revision block]
\label{def:accessible-revision-block}
Let
\[
 \bar\rho:H_1(\Go;\F_2)\twoheadrightarrow \Lat/2\Lat
\]
be the mod-\(2\) reduction of recurrent holonomy.  Universal coefficients give
the canonical subspace
\[
 H_{\rm acc}:=\Ann(\ker\bar\rho)
 \cong\Hom(\Lat/2\Lat,\F_2),
 \qquad
 \dim_{\F_2}H_{\rm acc}=\beff .
\]
Set
\[
 H_{\rm full,acc}:=H_{\rm acc}\oplus\langle t\rangle .
\]
\end{definition}

\begin{theorem}[Parity-correct accessible modular rank]
\label{thm:accessible-modular-rank}
Assume the full modular form is nondegenerate on \(H_{\rm full,acc}\).  Then
\[
 \beff\ \text{is odd},\qquad
 \dim\rad\!\left(\omega|_{H_{\rm acc}}\right)=1,
\]
the temporal generator pairs nontrivially with that radical, and
\[
 \boxed{\rank_{\F_2}\!\left(\omega|_{H_{\rm full,acc}}\right)=1+\beff.}
\]
If the inaccessible complement is modularly invisible and has no mixed block
with \(H_{\rm full,acc}\), the same equality holds for the full form.
\end{theorem}

\begin{proof}
A nondegenerate alternating form exists only in even dimension, so
\(1+\beff=\dim H_{\rm full,acc}\) is even.  The restriction to the odd-dimensional
space \(H_{\rm acc}\) is therefore degenerate.  If a vector in its radical also
paired trivially with \(t\), it would lie in the radical of the full accessible
form.  Thus the linear functional \(\alpha\mapsto\omega(\alpha,t)\) is injective
on the spatial radical, forcing that radical to be one-dimensional.  Adding
\(t\) supplies the missing hyperbolic partner and raises the rank from
\(\beff-1\) to \(1+\beff\).
\end{proof}

\begin{assumption}[Topological-crystal realisation]
\label{ass:topological-crystal}
The recurrent cover carries a free cocompact action of \(\Lat\) with finite
quotient motif, and the renewal time coordinate has bounded local distortion.
Thus the recurrent cover is quasi-isometric to \(\Lat\cong\Z^{\beff}\), and the
space--time crystal is quasi-isometric to \(\Z^{\beff}\times\N\).
\end{assumption}

\begin{assumption}[Crystal heat-kernel regularity]
\label{ass:crystal-heat-kernel}
The induced walk is lazy, reversible, uniformly elliptic, and bounded range.
\end{assumption}

\begin{corollary}[Return-probability spectral dimension]
\label{cor:return-spectral-dimension}
Under \Cref{ass:topological-crystal,ass:crystal-heat-kernel},
\[
 p_n(x,x)\asymp n^{-(1+\beff)/2},
 \qquad
 \boxed{\ds=1+\beff.}
\]
\end{corollary}

\begin{proof}
The free cocompact lattice action gives volume growth
\(|B(x,r)|\asymp r^{1+\beff}\).  Volume doubling and the scale-invariant
Poincar\'e inequality are stable for such bounded-geometry crystals; Delmotte's
theorem \cite{delmotte1999} gives the two-sided Gaussian estimate and hence the
diagonal exponent.
\end{proof}

Here \(d_s\) denotes the return-probability, or heat-kernel, spectral dimension
of the regular crystal.  It is not the zeta/Weyl dimension of
\Cref{thm:weyl-law}, nor a claim about dimensional flow in causal dynamical
triangulations or other quantum-gravity ensembles; its equality with
\(\qalg\) below is a theorem only under the stated crystal and channel
non-collapse hypotheses.

\begin{theorem}[Dimension coincidence in the regular Lorentzian crystal]
\label{thm:dimension-coincidence}
Assume
\Cref{ass:full-rank-reachability,ass:finite-primitive-control,ass:two-sided-channel-coding,ass:topological-crystal,ass:crystal-heat-kernel}, and suppose \(\omega\) is nondegenerate on
\(H_{\rm full,acc}\), with inaccessible classes modularly invisible.  Then
\[
 \boxed{
 \qmet=\qalg=\ds=\rank_{\F_2}\omega=1+\beff,
 \qquad
 \dCl=\beff=\rmon-1.
 }
\]
Without two-sided channel coding, the algebraic, accessible modular, and
heat-kernel dimensions still coincide, but the channel-metric equality is not
universal.
\end{theorem}

\begin{proof}
\Cref{thm:full-rank-returns} gives
\(\qalg=1+\beff\).
\Cref{thm:metric-predictive-coincidence} gives
\(\qmet=\qalg\).
\Cref{thm:accessible-modular-rank} gives
\(\rank\omega=1+\beff\) on the accessible block, and on the full form under the
stated invisibility condition.
\Cref{cor:return-spectral-dimension} gives
\(\ds=1+\beff\).
Finally \(\dCl=\dim H_{\rm acc}=\beff\).
\end{proof}

\begin{remark}[Logical scope]
\label{rem:dimension-coincidence-scope}
The growth--cycle equality is derived from full-dimensional affine reachability
and finite primitive return completion.  The mod-\(2\) comparison is made with
the actual accessed lattice \(\Lat\), so no finite-index substitution is used.
The crystal quasi-isometry and the cb-metric non-collapse remain distinct phase
conditions: rank equality alone does not imply either one.
\end{remark}

\begin{proposition}[Cycle covariance has full accessed rank]
\label{prop:cycle-covariance-rank}
Let the finite motif chain be irreducible, aperiodic, and reversible, and let
\(\xi:E(\Go)\to\Lat\otimes\R\) be an antisymmetric displacement cocycle.  If every
nonzero \(a\in\Hom(\Lat,\R)\) has a nonzero period on some closed directed cycle,
then the centred winding process has covariance \(Q\) with
\[
 \boxed{\rank Q=\beff.}
\]
\end{proposition}

\begin{proof}
The finite-state Markov-additive central limit theorem applies.  A scalar
projection has zero asymptotic variance precisely when its edge cocycle is
exact, equivalently when all closed-cycle periods vanish.  The hypothesis
excludes this for every nonzero projection.
\end{proof}

\begin{corollary}[Full regular-phase chain]
\label{cor:full-dimension-chain}
Under the hypotheses of
\Cref{thm:dimension-coincidence,prop:cycle-covariance-rank},
\[
 \boxed{
 \qmet=\qalg=\ds=\rank_{\F_2}\omega=1+\beff,
 \qquad
 \dCl=\rank Q=\beff.
 }
\]
\end{corollary}


\subsection{Survival of the volume-dual response as a primitive first-order observable}
\label{sec:cxi}

This selection statement cannot be made fully unconditional without an
explicit coarse-graining operator.  The following gives the
sharpest rigorous reduction: the algebraic degree statement (unconditional), a
\emph{survival criterion} that converts the vague ``survives coarse-graining'' into
a concrete parity/chirality condition on reset statistics (unconditional), and two
transparent closure routes -- renormalisation-group relevance and minimal field
content -- each replacing the abstract closure principle by a recognisable
hypothesis.

\subsubsection{The algebraic core, restated}

Let $\Vsp\cong\R^{\dCl}$ be the spatial revision module, $\Gamma:\Vsp\to\Cl(\Vsp)$
the Clifford embedding, and
\[
  \bren:\textstyle\bigwedge^2\Vsp\to\Cl(\Vsp),\qquad
  \bren(u\wedge v)=\tfrac12[\Gamma(u),\Gamma(v)]=\Gamma(u\wedge v),
\]
the canonical two-reset response, an $SO(\dCl)$-equivariant isomorphism onto the
bivector sector.  With the oriented volume element $\Volr=\Gamma(e_1)\cdots
\Gamma(e_{\dCl})$, the volume-dual response is
$\Hren=\Volr^{-1}\bren:\bigwedge^2\Vsp\xrightarrow{\sim}\bigwedge^{\dCl-2}\Vsp$.
The preceding analysis proves $\Img\Hren\subseteq\Vsp$ \textup{(}degree one\textup{)} iff
$\dCl=3$.  We take this as given and ask what survives coarse-graining.

\subsubsection{Survival is oriented circulation}

\begin{definition}[Oriented reset circulation]
\label{def:circulation}
Let $\nu$ be the coarse-grained reset-pair law with joint weights $p_{ab}\ge0$ on
ordered direction pairs $(\hat\theta_a,\hat\theta_b)$.  Its \emph{oriented
circulation} is the bivector
\[
  \Xi(\nu):=\sum_{a,b}p_{ab}\,\hat\theta_a\wedge\hat\theta_b
          =\sum_{a,b}p^{[ab]}\,\hat\theta_a\wedge\hat\theta_b
          \in\textstyle\bigwedge^2\Vsp,
  \qquad p^{[ab]}:=\tfrac12(p_{ab}-p_{ba}).
\]
The law is \emph{parity- (time-reversal-) symmetric} if $p_{ab}=p_{ba}$, i.e.\
$\Xi(\nu)=0$.
\end{definition}

\begin{theorem}[Survival criterion]
\label{thm:survival-criterion}
The coarse-grained two-reset response equals
\[
  \bren[\nu]:=\sum_{a,b}p_{ab}\,\tfrac12[\Gamma(\hat\theta_a),\Gamma(\hat\theta_b)]
  =\Gamma\bigl(\Xi(\nu)\bigr),
\]
and its volume dual is $\Hren[\nu]=\Volr^{-1}\Gamma(\Xi(\nu))$.  Consequently the
coarse-grained volume-dual response is a \emph{nonzero one-reset
\textup{(}degree-one\textup{)} observable} if and only if
\[
  \boxed{\ \dCl=3\quad\text{and}\quad \Xi(\nu)\neq0.\ }
\]
In particular it vanishes identically for every parity-symmetric reset law,
regardless of dimension.
\end{theorem}

\begin{proof}
Since $\hat\theta_a\wedge\hat\theta_b$ is antisymmetric in $(a,b)$ and
$\tfrac12[\Gamma(\hat\theta_a),\Gamma(\hat\theta_b)]=\Gamma(\hat\theta_a\wedge
\hat\theta_b)$, summing against $p_{ab}$ keeps only the antisymmetric part:
$\bren[\nu]=\sum_{ab}p^{[ab]}\Gamma(\hat\theta_a\wedge\hat\theta_b)
=\Gamma(\Xi(\nu))$ by linearity of $\Gamma$ on bivectors.  As $\Gamma$ and
$\Volr^{-1}$ are injective, $\Hren[\nu]=\Volr^{-1}\Gamma(\Xi(\nu))\neq0$ iff
$\Xi(\nu)\neq0$, and its image lies in $\bigwedge^{\dCl-2}\Vsp$, which is the
degree-one sector $\Vsp$ iff $\dCl-2=1$.  If $p_{ab}=p_{ba}$ then $p^{[ab]}=0$ and
$\Xi(\nu)=0$, so $\bren[\nu]=0$.
\end{proof}

\begin{remark}[Why survival was conditional]
\label{rem:survival-meaning}
\Cref{thm:survival-criterion} explains the conditionality precisely:
the two-reset response is built from the \emph{antisymmetric} part of the reset-pair
statistics, so it is killed by parity/time-reversal symmetry.  ``Survival under
coarse-graining'' is therefore equivalent to the reset law possessing nonzero
oriented circulation -- a handedness/chirality of the renewal dynamics.  This is a
concrete, measurable condition on microscopic reset statistics, replacing the
opaque survival postulate.  Existence of laws with $\Xi\neq0$ is immediate
(\Cref{ex:chiral-law}).
\end{remark}

\begin{example}[A chiral reset law with $\Xi\neq0$, and a symmetric one with $\Xi=0$]
\label{ex:chiral-law}
Take $\dCl=3$ and three coplanar directions $\hat\theta_1,
\hat\theta_2,\hat\theta_3$ at $120^\circ$ in the $e_1e_2$-plane, with a cyclic
\textup{(}chiral\textup{)} ordered pair law $p_{12}=p_{23}=p_{31}=\tfrac13$,
$p_{ba}=0$.  Then $\Xi=\tfrac13(\hat\theta_1\wedge\hat\theta_2+\hat\theta_2\wedge
\hat\theta_3+\hat\theta_3\wedge\hat\theta_1)=\tfrac12\,e_1\wedge e_2\neq0$, so the
volume-dual response is the nonzero one-reset observable $\Hren[\nu]=
\Volr^{-1}\Gamma(\Xi)\propto\Gamma(e_3)$.  Symmetrising the pair law,
$p_{ab}\mapsto\tfrac12(p_{ab}+p_{ba})$, gives $\Xi=0$ and kills the response.  Thus
chirality of the reset ordering is exactly what survives.
\end{example}

\subsubsection{Route B: Clifford-grade closure and the RG obstruction}

The oriented volume-dual response has a simple algebraic dimension count.  The
band-bottom curvature is a two-form.  Its Hodge dual belongs to
\(\Lambda^{\dCl-2}\Vsp\), and therefore closes into the elementary vector
generator sector precisely when
\[
  \dCl-2=1,
  \qquad\text{that is,}\qquad \dCl=3.
\]
Thus, if the oriented response is required to survive in the principal
grade-one Clifford sector, the minimal nondegenerate odd endpoint is selected.

This should not be confused with a canonical renormalisation-group relevance
argument.  A tempting strengthening would be to assign the degree-\(k\)
Clifford sector an engineering eigenvalue relative to a one-reset generator and
to infer a factor \(b^{3-\dCl}\).  That conclusion is not obtained from ordinary
field-theoretic engineering dimensions alone.  Once the response is coupled to
a fermion bilinear, the bilinear has engineering dimension \(\dCl\) in
\((1+\dCl)\)-dimensional spacetime, independently of the inserted Clifford
grade.  A gauge-potential-type coefficient must therefore have dimension one in
every spatial dimension.  Canonical power counting is grade-blind.

\begin{theorem}[Engineering power counting does not select the dimension]
\label{thm:rg-no-dimension-selection}
The coupling of the volume-dual oriented response to a fermion bilinear is a
dimension-one, gauge-potential-type coupling and is marginal by canonical
engineering power counting in every spatial dimension.  Hence ordinary power
counting does not produce a universal eigenvalue proportional to
\(b^{3-\dCl}\), and it does not by itself suppress the higher-grade currents
present for \(\dCl>3\).  The special role of \(\dCl=3\) is therefore an
exterior-algebra and Clifford-grade matching statement, not a consequence of
canonical renormalisation-group relevance.
\end{theorem}

\begin{proof}
A first-order Dirac coupling has the form of a Clifford-valued one-particle
operator inserted between fermions.  In \((1+\dCl)\)-dimensional spacetime, the
fermion field has engineering dimension \(\dCl/2\), so a fermion bilinear has
dimension \(\dCl\), independently of which Clifford grade is inserted in the
matrix acting on spinors.  Since the action density has dimension
\(\dCl+1\), the coefficient of a first-order gauge-potential-type coupling has
dimension one in every \(\dCl\).  Thus canonical engineering dimension does not
distinguish the Clifford grades.

The algebraic distinction is separate.  The curvature response is a two-form,
and
\[
  \star:\Lambda^2\Vsp\longrightarrow \Lambda^{\dCl-2}\Vsp .
\]
This lands in the vector sector \(\Lambda^1\Vsp\) only when \(\dCl-2=1\).
Therefore \(\dCl=3\) is selected by Hodge degree and Clifford-grade closure,
not by canonical power counting.
\end{proof}

\begin{corollary}[Dimensional status after the RG audit]
\label{cor:dimension-status-rg-audit}
Primitive modular nondegeneracy forces odd spatial rank and excludes the lower
degenerate endpoints.  Therefore \(\dCl=3\) remains the unconditional minimal
nondegenerate endpoint.  If, in addition, the oriented signed-cluster response
is required to close into and remain within the principal grade-one generator
sector, then \(\dCl=3\) is selected among the odd endpoints.  The present paper
does not derive this principal-grade survival condition from canonical power
counting: higher-grade couplings are not removed by engineering dimension
alone.  A genuine uniqueness theorem among higher odd endpoints would require
an interacting anomalous-dimension separation or an independent microscopic
operator-closure theorem.
\end{corollary}

\begin{remark}[What remains of Route B]
\label{rem:route-B-after-rg-audit}
The robust conclusion is the grade-closure statement:  the oriented response has Hodge degree
\(\dCl-2\), and hence can be identified with a primitive grade-one Dirac
generator only for \(\dCl=3\).  A special block-renewal model may still generate
additional anomalous scaling that suppresses higher grades, but that would be an
extra dynamical theorem about the interacting renewal flow, not a consequence of
canonical engineering power counting.
\end{remark}

\subsubsection{Route C: minimal continuum field content}

\begin{theorem}[Minimal-field-content closure]
\label{thm:minimal-field}
Suppose the continuum endpoint's one-reset observable space is
exactly $\Vsp$ -- the metric coframe and a first-order Dirac spinor, with no
independently propagating higher-form sector.  If the oriented volume-dual
response is nonzero \textup{(}equivalently $\Xi\neq0$, \Cref{thm:survival-criterion}\textup{)},
then it can be a one-reset observable only if its Clifford degree is one, i.e.
\[
  \dCl-2=1,\qquad \boxed{\dCl=3}.
\]
\end{theorem}

\begin{proof}
By the preceding construction's degree computation, $\Img\Hren=\bigwedge^{\dCl-2}
\Vsp$.  Under the minimal-field-content hypothesis the only one-reset observables
are degree one, so a nonzero $\Hren[\nu]$ is a one-reset observable iff $\dCl-2=1$.
\end{proof}

\subsubsection{Net status of the survival problem}

\begin{remark}[Exactly what is and is not discharged]
\label{rem:cxi-status}
The degree equivalence ($\Img\Hren=\bigwedge^{\dCl-2}\Vsp$, elementary
iff $\dCl=3$) and the survival criterion (survival $\Leftrightarrow$ $\dCl=3$ and
$\Xi\neq0$, \Cref{thm:survival-criterion}) are unconditional.  Together they reduce
the four-part principal-symbol-closure principle of the preceding construction to two concrete
questions: (i) does the reset law have nonzero oriented circulation $\Xi$? --
answered by a parity test, with chiral examples; (ii) does the response enter the
principal symbol? -- answered by \emph{either} RG relevance
(\Cref{thm:rg-eigenvalue}, on $\mathscr R$) \emph{or} minimal field
content (\Cref{thm:minimal-field}, ).  The irreducible
residue is exactly the construction of the block-renewal/OPE map needed to settle
(ii) unconditionally; absent it, $\dCl=3$ uniqueness from this route remains
conditional, consistent with the stated hypotheses.  What \emph{has} changed is
that the conditional content is now a sharp, checkable nonvanishing condition plus
one transparent closure route, rather than an opaque ``survives, is principal, is
primitive, is not an independent field'' bundle.
\end{remark}

% =====================================================================


\subsection{An explicit block-renewal coarse-graining map}
\label{sec:cxiii}

the survival analysis made the renormalisation-group selection of $\dCl=3$ conditional on the
existence of a block-renewal map with engineering eigenvalue $b^{3-\dCl}$.  We now
construct one such map explicitly in the graded-Clifford renewal class and prove
the eigenvalue with a controlled band error.  This is not a universal RG theorem;
it is one mathematically controlled model, which is what is needed to upgrade the
engineering argument to a theorem.

\subsubsection{Joint-degree reset-interference observables}

Work in the graded-Clifford renewal memory: $\Arev=\Cl_n(\C)$, $n=1+\dCl$, with
Hermitian generators $\Gamma_0,\dots,\Gamma_{\dCl}$, spatial Clifford embedding
$\Gamma:\Vsp\to\Cl(\Vsp)$, and the volume element
$\Volr=\Gamma(e_1)\cdots\Gamma(e_{\dCl})$.  Fix a reset-direction system
$\{\hat\theta_a\}\subset\Sphere^{\dCl-1}$ and the symmetric reset differences
$\delta_{\hat\theta_a,h}$ with Fourier multiplier
$s_{\hat\theta_a,h}(\xi)=h^{-1}\sin(h\,\hat\theta_a\cdot\xi)$.

\begin{definition}[Degree-$k$ reset-interference observable]
\label{def:degree-k-observable}
For $1\le k\le\dCl$, the canonical degree-$k$ reset-interference observable is the
$k$-fold antisymmetrised reset difference
\[
  \Ocal_k^{(h)}
  =\sum_{a_1<\dots<a_k}
   \Gamma(\hat\theta_{a_1}\wedge\dots\wedge\hat\theta_{a_k})\,
   \delta_{\hat\theta_{a_1},h}\cdots\delta_{\hat\theta_{a_k},h},
\]
with Fourier symbol, valued in $\Cl(\Vsp)$,
\[
  \sigma_k^{(h)}(\xi)
  =\sum_{a_1<\dots<a_k}
   \Gamma(\hat\theta_{a_1}\wedge\dots\wedge\hat\theta_{a_k})\,
   i^k\,s_{\hat\theta_{a_1},h}(\xi)\cdots s_{\hat\theta_{a_k},h}(\xi).
\]
Its \emph{joint degree} is $k$: it carries Clifford degree $k$ and momentum
homogeneity $k$.  The continuum symbol $\sigma_k(\xi)=\lim_{h\to0}\sigma_k^{(h)}$
is homogeneous of degree $k$ in $\xi$, $\sigma_k(\lambda\xi)=\lambda^k\sigma_k(\xi)$.
\end{definition}

\begin{remark}[The volume-dual response is a degree-$(\dCl-2)$ observable]
\label{rem:volume-dual-is-degree-d-2}
By Hodge duality in $\Cl(\Vsp)$, $\Volr^{-1}\Gamma(u\wedge v)
=\pm\Gamma\bigl(\star(u\wedge v)\bigr)$, a Clifford element of degree $\dCl-2$.
Hence the oriented volume-dual response $\Hren=\Volr^{-1}\bren$ lands in the
degree-$(\dCl-2)$ sector and its natural renewal realisation is the degree-$(\dCl-2)$
observable $\Ocal_{\dCl-2}^{(h)}$ of \Cref{def:degree-k-observable}: it is
$\Gamma$ of a $(\dCl-2)$-form, realised by $\dCl-2$ antisymmetrised reset
differences, of joint degree $\dCl-2$.  This is what makes its engineering
dimension $\dCl-2$, not $2$.
\end{remark}

\subsubsection{The block-renewal map and its eigenvalue}

\begin{construction}[Block-renewal coarse-graining]
\label{constr:block-rg}
Fix an integer block size $b\ge2$.  The block-renewal map $\Rb$ acts on
$\Cl(\Vsp)$-valued Fourier symbols by mesh dilation $h\mapsto bh$ followed by the
momentum rescaling $\xi\mapsto\xi/b$ and the normalisation prefactor $b$ chosen to
fix the degree-one Dirac symbol:
\[
  (\Rb\sigma)(\xi):=b\,\sigma^{(bh)}\!\bigl(\xi/b\bigr),
\]
where $\sigma^{(bh)}$ is the symbol recomputed at the block mesh $bh$
\textup{(}block reset differences $\delta_{\hat\theta_a,bh}$\textup{)}.  The
prefactor is fixed by requiring $\Rb$ to fix the degree-one continuum symbol
$\sigma_1(\xi)=\sum_a\Gamma(\hat\theta_a)\,i\,(\hat\theta_a\cdot\xi)$: indeed
$\Rb\sigma_1(\xi)=b\,\sigma_1(\xi/b)=b\cdot b^{-1}\sigma_1(\xi)=\sigma_1(\xi)$.
\end{construction}

\begin{theorem}[Engineering eigenvalue of the block-renewal map]
\label{thm:rg-eigenvalue}
On the continuum symbols, $\Rb$ acts diagonally on joint degree with
exact eigenvalue
\[
  \boxed{\ \Rb(\sigma_k)=b^{\,1-k}\,\sigma_k\ }\qquad(1\le k\le\dCl).
\]
On the lattice, for every momentum band $|\xi|\le\Lambda$ there is $C_k=C_k(\Lambda)$
with
\[
  \bigl\|\,\Rb(\sigma_k^{(h)})-b^{\,1-k}\sigma_k^{(h)}\,\bigr\|_{\,|\xi|\le\Lambda}
  \ \le\ C_k\,b^{\,1-k}\,(bh)^2\Lambda^2
  \ \xrightarrow[h\to0]{}\ 0.
\]
Consequently the volume-dual response, of joint degree $k=\dCl-2$
\textup{(}\Cref{rem:volume-dual-is-degree-d-2}\textup{)}, has eigenvalue
\[
  \boxed{\ \lambda_{\dCl}=b^{\,1-(\dCl-2)}=b^{\,3-\dCl}.\ }
\]
\end{theorem}

\begin{proof}
\emph{Continuum.}  $\sigma_k$ is homogeneous of degree $k$, so
$\sigma_k(\xi/b)=b^{-k}\sigma_k(\xi)$ and $\Rb\sigma_k(\xi)=b\,b^{-k}\sigma_k(\xi)
=b^{1-k}\sigma_k(\xi)$.

\emph{Lattice error.}  Each scalar factor obeys
$|s_{\hat\theta,bh}(\eta)-\hat\theta\cdot\eta|\le\tfrac16(bh)^2|\eta|^3$.  Writing
$\eta=\xi/b$ and using $|\hat\theta\cdot\eta|\le|\eta|\le\Lambda/b$ on the
rescaled band, the $k$-fold product symbol $\sigma_k^{(bh)}(\xi/b)$ differs from
its leading homogeneous term $\prod_l(\hat\theta_{a_l}\cdot\xi/b)
=b^{-k}\prod_l(\hat\theta_{a_l}\cdot\xi)$ by a relative error
$O\bigl((bh)^2(\Lambda/b)^2\bigr)$ per factor, hence
$O\bigl((bh)^2\Lambda^2 b^{-2}\bigr)$ overall after the $k$ bounded Clifford
factors.  Multiplying by the prefactor $b$ and the homogeneity factor $b^{-k}$
gives, on the band,
\[
  \bigl\|\Rb(\sigma_k^{(h)})-b^{1-k}\sigma_k^{(h)}\bigr\|
  \le C_k\,b\cdot b^{-k}\cdot(bh)^2\Lambda^2 b^{-2}\cdot b^{2}
  = C_k\,b^{1-k}(bh)^2\Lambda^2,
\]
the final $b^{2}$ restoring the comparison to the unrescaled lattice symbol
$\sigma_k^{(h)}$ on $|\xi|\le\Lambda$; the Clifford prefactors are unit norm and
contribute the constant $C_k$.  Letting $h\to0$ at fixed $b,\Lambda$ sends the
error to zero.  The volume-dual statement is the case $k=\dCl-2$.
\end{proof}

\subsubsection{Dimensional selection from the eigenvalue}

\begin{corollary}[Relevance trichotomy]
\label{cor:relevance}
\textup{} With $\lambda_{\dCl}=b^{3-\dCl}$ and $b\ge2$:
\[
  \begin{aligned}
  &\dCl=3:\ \lambda=1\ \text{(marginal/principal)};\qquad
   \dCl\ge5:\ \lambda<1\ \text{(irrelevant)};\\
  &\dCl=1:\ \text{no bivector domain}.
  \end{aligned}
\]
Hence, in the graded-Clifford renewal class with the explicit block-renewal map
$\Rb$, the oriented volume-dual response is a principal first-order observable
only at $\dCl=3$, is RG-irrelevant for every higher odd spatial rank, and has no
two-reset domain at $\dCl=1$.
\end{corollary}

\begin{proof}
$b^{3-\dCl}=1$ iff $\dCl=3$; $b^{3-\dCl}<1$ iff $\dCl>3$ (in particular for odd
$\dCl\ge5$); for $\dCl=1$, $\bigwedge^2\Vsp=0$ so $\Ocal_{\dCl-2}=\Ocal_{-1}$ is
empty.  Marginality of a $\lambda=1$ operator is precisely its persistence in the
principal symbol.
\end{proof}

\begin{remark}[Precise scope of CXIII]
\label{rem:cxiii-status}
The eigenvalue $b^{1-k}$ and the band error are theorems in the
graded-Clifford renewal model with the explicit map $\Rb$ of
\Cref{constr:block-rg}.  The engineering dimension is assigned by joint degree,
which is justified because the natural renewal realisation of a degree-$k$
Clifford observable is the $k$-fold reset difference $\Ocal_k$, of momentum
homogeneity $k$ (\Cref{rem:volume-dual-is-degree-d-2}).  \textup{} This is
\emph{one} controlled model, not a universal RG theorem; the selection of $\dCl=3$
by RG relevance therefore remains a model-class statement.  Within that
class, the eigenvalue and quantitative band estimate are proved rather than
inferred from engineering bookkeeping.
\end{remark}

% =====================================================================


\subsection{From interference postulate to the canonical volume-dual response}
\label{subsec:interference-closure}

The preceding results determine the common spatial Clifford rank
\[
  \dCl=\beff=q_{\mathrm{alg}}-1
\]
in the regular spacetime-crystal phase and prove that primitivity makes
\(\dCl\) odd. They do not, by themselves, exclude the higher odd values
\(5,7,\ldots\). We now sharpen the earlier interference-closure discussion in
three directions. First, we show by explicit models that the current renewal
axioms do not imply closure. Second, we identify the two-reset commutator as a
canonical bivector-valued renewal response in every odd dimension. Third, after
choosing an orientation, the Clifford volume element gives a canonical Hodge
dual of that response; this dual is a nonzero elementary revision exactly in
three spatial dimensions.

The relevant endpoint is not merely the abstract matrix algebra
\(\Arev\cong\Cl_{1+\dCl}(\C)\). It is the \emph{marked renewal--Clifford datum}
\[
  (\Arev,V_{\mathrm{full}},q,\Gamma,V_{\mathrm{sp}}),
\]
where \(\Gamma:V_{\mathrm{full}}\to\Arev\) is the Clifford embedding supplied by
the primitive revision construction and \(V_{\mathrm{sp}}\subset
V_{\mathrm{full}}\) is the spatial subspace selected by the modular time
direction. The abstract matrix algebra alone does not canonically remember a
degree-one subspace; all statements below are relative to this marked endpoint.

Set
\[
  \Vsp:=V_{\mathrm{sp}}\cong\R^{\dCl},
  \qquad
  \Iint:=\operatorname{span}_{\R}
  \{\Gamma(e_i)\Gamma(e_j):i<j\}.
\]

\begin{lemma}[The canonical two-reset response and the Lorentz symmetric pair]
\label{lem:interference-lorentz-pair}
The antisymmetric two-reset response
\[
  \bren:\bigwedge\nolimits^2\Vsp\longrightarrow\Arev,
  \qquad
  \bren(x\wedge y):=\frac12[\Gamma(x),\Gamma(y)]
\]
takes values in \(\Iint\), is \(O(\Vsp)\)-equivariant, and is a linear
isomorphism
\[
  \bren:\bigwedge\nolimits^2\Vsp\xrightarrow{\ \sim\ }\Iint.
\]
Under the associative commutator,
\[
  \Iint\cong\mathfrak{spin}(\dCl)\cong\mathfrak{so}(\dCl),
  \qquad
  \dim_{\R}\Iint=\binom{\dCl}{2},
\]
and
\[
  [\Iint,\Iint]\subseteq\Iint,\qquad
  [\Iint,\Vsp]\subseteq\Vsp,\qquad
  [\Vsp,\Vsp]\subseteq\Iint.
\]
With \(K_i=\frac12\Gamma(e_i)\) and
\(J_{ij}=-\frac12\Gamma(e_i)\Gamma(e_j)\), the real Lie algebra
\(\Iint\oplus\Vsp\) is isomorphic to \(\mathfrak{so}(\dCl,1)\).
\end{lemma}

\begin{proof}
For orthonormal \(e_i\ne e_j\), the Clifford relations give
\[
  \frac12[\Gamma(e_i),\Gamma(e_j)]
  =\Gamma(e_i)\Gamma(e_j).
\]
The ordered degree-two Clifford monomials are linearly independent in the
faithful full spinor representation, so \(\bren\) identifies
\(\bigwedge^2\Vsp\) with \(\Iint\). Equivariance follows from naturality of the
Clifford action. Direct computation gives
\[
 \left[\frac12\Gamma(e_i)\Gamma(e_j),\Gamma(e_k)\right]
 =\delta_{jk}\Gamma(e_i)-\delta_{ik}\Gamma(e_j),
\]
together with the standard \(\mathfrak{so}(\dCl)\) relations among bivectors.
Moreover
\[
  [K_i,K_j]=-J_{ij},\qquad
  [J_{ij},K_k]=\delta_{ik}K_j-\delta_{jk}K_i,
\]
which are the Cartan-decomposition relations of \(\mathfrak{so}(\dCl,1)\).
Jacobi is inherited from the associative commutator.
\end{proof}

The automatic Lie algebra therefore lives on
\(\Iint\oplus\Vsp\), not on \(\Vsp\) alone. A bracket on the spatial revision
module appears only after a degree-lowering identification of interferences
with revisions.

\begin{theorem}[Closure is satisfiable and falsifiable within the present model class]
\label{thm:interference-independence}
Let \(\mathsf A\) denote the standing package of finite fibres and bounded
increments, aperiodic local finiteness, primitivity, double recurrent
ergodicity, weak isotropy, and bounded geometry. Then:
\begin{enumerate}[label=\textup{(\roman*)},leftmargin=2.4em]
\item there is a model of \(\mathsf A\) with \(\dCl=3\) for which oriented
      interference closure holds;
\item for every odd \(\dCl\ge5\) there is a model of \(\mathsf A\) for which it
      fails.
\end{enumerate}
Hence closure is satisfiable and falsifiable within the present model class; in
this model-theoretic sense it is independent of, and in particular not implied
by, \(\mathsf A\).
\end{theorem}

\begin{proof}
Use the bouquet graded-Clifford renewal memory with the cross-polytope reset
frame. The construction has a finite reset alphabet and unit length increments,
is locally finite and aperiodic, has bounded degree and unit propagation, and
its cross-polytope second moment is \(M=\dCl^{-1}I\). For every odd \(\dCl\), the
full rank \(1+\dCl\) is even and the off-diagonal symplectic form used in the
primitive examples is nonsingular over \(\F_2\); the associated Clifford
representation is irreducible, giving the required primitive and doubly
ergodic endpoint. Thus the same axiom package is realized simultaneously in
every odd spatial rank. At \(\dCl=3\), the Hodge map constructed below is an
equivariant isomorphism \(\bigwedge^2\Vsp\cong\Vsp\). For odd
\(\dCl\ge5\),
\[
  \dim\bigwedge\nolimits^2\Vsp=\binom{\dCl}{2}\ne\dCl=\dim\Vsp,
\]
so no closure isomorphism exists.
\end{proof}

\begin{remark}[Why the current axioms do not see closure]
The standing axioms constrain the graph and length grading, irreducibility and
nondegeneracy of the primitive revision form, and the symmetric second moment
\(M\in\operatorname{Sym}^2\Vsp\). None is itself a condition on a map
\(\bigwedge^2\Vsp\to\Vsp\). This grading separation explains the preceding
models, but the explicit models, rather than this observation alone, establish
the non-implication.
\end{remark}

\begin{assumption}[Oriented spatial phase]
\label{hyp:oriented-spatial-phase}
The determinant line
\[
  \det\Vsp=\bigwedge\nolimits^{\dCl}\Vsp
\]
is trivialised compatibly with recurrent holonomy, so the acting structure group
reduces to
\[
  \Holp:=\Hol\cap SO(\Vsp).
\]
A time orientation alone does not provide this datum.
\end{assumption}

\begin{proposition}[Local controllability implies full oriented holonomy]
\label{prop:interference-controllability}
Suppose the recurrent transports extend to one-parameter oriented holonomy
families with infinitesimal generators \(X_e\in\mathfrak{so}(\Vsp)\), and
\[
  \Lie\langle X_e\rangle=\mathfrak{so}(\Vsp).
\]
Then
\[
  \overline{\Holp}=SO(\Vsp).
\]
\end{proposition}

\begin{proof}
The closed subgroup \(\overline{\Holp}\subseteq SO(\Vsp)\) contains every
\(\exp(tX_e)\). By the closed-subgroup theorem it is a Lie subgroup whose Lie
algebra contains \(\Lie\langle X_e\rangle=\mathfrak{so}(\Vsp)\). Since
\(SO(\Vsp)\) is connected, a closed subgroup with the full Lie algebra is the
whole group.
\end{proof}

\begin{remark}[Scope of controllability]
The proposition is established under the stated one-parameter-family
hypothesis. The microscopic renewal transports are discrete, so the existence
of such families is not automatic. For a purely discrete transport system one
must instead prove density of the generated subgroup. Thus controllability
replaces a global maximal-isotropy postulate by a local rank condition, but does
not derive continuous holonomy from the renewal axioms.
\end{remark}

\begin{proposition}[The oriented volume element of the marked endpoint]
\label{prop:renewal-volume-element}
Fix the marked Clifford embedding \(\Gamma\) and an orientation of \(\Vsp\). For
any oriented orthonormal basis \((e_1,\ldots,e_{\dCl})\), define
\[
  \Volr:=\Gamma(e_1)\cdots\Gamma(e_{\dCl}).
\]
Then \(\Volr\) is independent of the chosen oriented orthonormal basis, reversing
orientation sends \(\Volr\mapsto-\Volr\), and
\[
  g\Volr g^{-1}=\Volr\qquad(g\in\Holp).
\]
Moreover, for odd \(\dCl\),
\[
  \Volr^2=(-1)^{\dCl(\dCl-1)/2}\one,
\]
and \(\Volr\) is central in the spatial Clifford algebra. It need not be a
scalar on the full \(\Cl_{1+\dCl}(\C)\)-spinor module.
\end{proposition}

\begin{proof}
Two oriented orthonormal bases differ by an element of \(SO(\Vsp)\). The ordered
Clifford product represents the top exterior power, on which the change of basis
acts by its determinant, hence by \(1\); an orientation reversal contributes
\(-1\). The same determinant argument proves oriented-holonomy invariance.
The square and centrality are the standard odd-dimensional Clifford volume
identities. Since the temporal Clifford generator anticommutes with each spatial
generator and \(\dCl\) is odd, it anticommutes with \(\Volr\), so \(\Volr\) is
not central in the full spacetime Clifford algebra and therefore need not act as
one scalar on the full spinor module.
\end{proof}

\begin{definition}[Oriented volume-dual two-reset response]
\label{def:volume-dual-response}
Define
\[
  \Hren:=\Volr^{-1}\bren:
  \bigwedge\nolimits^2\Vsp\longrightarrow\Cl(\Vsp).
\]
This is the canonical \(SO(\Vsp)\)-equivariant degree-complementing map supplied
by the marked Clifford metric and orientation. Reversing orientation changes
\(\Hren\) by an overall sign and therefore does not change its image or its
dimension-selection property.
\end{definition}

\begin{proposition}[Degree and equivariance of the volume-dual response]
\label{prop:volume-dual-degree}
The map \(\Hren\) is \(\Holp\)-equivariant and is a linear isomorphism
\[
  \Hren:
  \bigwedge\nolimits^2\Vsp
  \xrightarrow{\ \sim\ }
  \bigwedge\nolimits^{\dCl-2}\Vsp.
\]
\end{proposition}

\begin{proof}
Equivariance follows from equivariance of \(\bren\) and invariance of \(\Volr\).
For basis bivectors,
\[
  \Volr^{-1}\Gamma(e_i)\Gamma(e_j)
\]
cancels the two factors \(e_i,e_j\) from the oriented top product, leaving, up
to sign, the ordered Clifford product of the remaining \(\dCl-2\) basis
vectors. This is the Clifford realization of Hodge duality
\(\bigwedge^2\Vsp\to\bigwedge^{\dCl-2}\Vsp\). It is invertible, with inverse the
opposite Hodge-degree map; equivalently both spaces have dimension
\(\binom{\dCl}{2}=\binom{\dCl}{\dCl-2}\).
\end{proof}

\begin{principle}[Oriented interference closure]
\label{prin:interference-closure}
The oriented volume-dual two-reset response is a nonzero elementary spatial
revision:
\[
  \Hren\ne0,
  \qquad
  \Img\Hren\subseteq\Vsp.
\]
Equivalently, pairwise interference, after the canonical metric-and-orientation
duality, introduces neither missing nor excess primitive revision modes.
\end{principle}

\begin{theorem}[Oriented volume-dual closure selects three spatial dimensions]
\label{thm:interference-closure-selects-three}
Assume a primitive marked renewal--Clifford endpoint and the oriented phase
\Cref{hyp:oriented-spatial-phase}. Then
\[
  \Hren\ne0\ \text{and}\ \Img\Hren\subseteq\Vsp
  \quad\Longleftrightarrow\quad
  \dCl=3.
\]
In particular, \Cref{prin:interference-closure} implies
\[
  \boxed{\dCl=3}.
\]
For \(\dCl=1\), \(\bigwedge^2\Vsp=0\) and \(\Hren=0\); for every odd
\(\dCl\ge5\), \(\Hren\ne0\) but its image has Clifford degree
\(\dCl-2>1\).
\end{theorem}

\begin{proof}
By \Cref{prop:volume-dual-degree},
\[
  \Img\Hren=\bigwedge\nolimits^{\dCl-2}\Vsp.
\]
At \(\dCl=1\), the domain \(\bigwedge^2\Vsp\) vanishes. For odd
\(\dCl\ge3\), \(\bren\) is an isomorphism and \(\Volr\) is invertible, hence
\(\Hren\ne0\). Its image is elementary, i.e. degree one, exactly when
\[
  \dCl-2=1,
\]
which is equivalent to \(\dCl=3\).
\end{proof}

\begin{proposition}[Realisation and induced bracket at \(\dCl=3\)]
\label{prop:hodge-interference-closure}
At \(\dCl=3\),
\[
  \Hren(x\wedge y)=\star(x\wedge y),
\]
up to the orientation-dependent overall sign. In an oriented orthonormal basis,
\[
  \Hren(e_i\wedge e_j)=\epsilon_{ijk}e_k.
\]
The induced bracket
\[
  [x,y]_\star:=\Hren(x\wedge y)
\]
satisfies Jacobi and gives
\[
  (\Vsp,[\,\cdot,\cdot\,]_\star)
  \cong\mathfrak{so}(3)\cong\mathfrak{su}(2).
\]
If the oriented holonomy acts irreducibly on \(\Vsp\), every
\(\Holp\)-equivariant closure map is a nonzero real multiple of \(\Hren\).
\end{proposition}

\begin{proof}
For \(\Volr=\Gamma(e_1)\Gamma(e_2)\Gamma(e_3)\),
\[
  \Volr^{-1}\Gamma(e_i)\Gamma(e_j)=\pm\epsilon_{ijk}\Gamma(e_k),
\]
which is precisely the Hodge cross product. The vector triple-product identity
proves Jacobi. The resulting structure constants are those of
\(\mathfrak{so}(3)\), and the compact real form is isomorphic to
\(\mathfrak{su}(2)\). For uniqueness, a real irreducible three-dimensional
representation cannot be of complex or quaternionic type, so its commutant is
\(\R\); Schur's lemma then gives uniqueness up to a real scalar.
\end{proof}

\begin{proposition}[Clifford-internal form]
\label{prop:clifford-internal-interference}
At \(\dCl=3\), the spatial volume multiplier satisfies
\[
  \Volr^2=-\one,
  \qquad
  \Gamma(e_i)\Gamma(e_j)
  =\epsilon_{ijk}\Volr\Gamma(e_k).
\]
It is central and invertible in the spatial algebra
\(\Cl_3(\C)\), but anticommutes with the temporal generator. On the full
spacetime spinor module it acts with the two opposite square roots of
\(-1\) on the two irreducible spatial summands and is not a scalar.
\end{proposition}

\begin{proof}
These are the odd-dimensional Clifford volume identities. The temporal
generator anticommutes with each of the three spatial generators, hence with
their product. Restricting the full even-dimensional spinor representation to
\(\Cl_3(\C)\cong M_2(\C)\oplus M_2(\C)\) yields the two simple summands, on which
the central spatial volume element acts by opposite square roots of \(-1\).
\end{proof}

\begin{remark}[Why the Hilbert--Schmidt projection does not work]
On the full spacetime spinor module, every spatial bivector is
Hilbert--Schmidt orthogonal to every spatial vector:
\[
  \Tr\!\left(\Gamma(e_k)^\dagger
  \Gamma(e_i)\Gamma(e_j)\right)=0.
\]
Thus orthogonal projection of \(\bren\) onto \(\Vsp\) vanishes in every odd
dimension, including \(\dCl=3\). The correct degree-lowering map is
multiplication by \(\Volr^{-1}\), not projection. The identity at
\(\dCl=3\) is ``bivector equals volume times vector,'' not ``bivector equals
vector'' on the full module.
\end{remark}

\begin{remark}[Why seven dimensions do not survive]
\label{rem:no-seven-dimensional-exception}
Although a norm-compatible cross product exists on \(\R^7\), its linearisation
is only a surjection
\[
  \bigwedge\nolimits^2\R^7\longrightarrow\R^7
\]
with a \(14\)-dimensional kernel:
\[
  \bigwedge\nolimits^2\R^7\cong\R^7\oplus\mathfrak g_2.
\]
It is not an isomorphism and is only \(G_2\)-equivariant, not
\(SO(7)\)-equivariant.
\end{remark}

\begin{theorem}[Full-isotropy route]
\label{thm:full-isotropy-interference}
If \(\overline{\Holp}=SO(\Vsp)\), then a nonzero
\(\Holp\)-equivariant map
\[
  \bigwedge\nolimits^2\Vsp\longrightarrow\Vsp
\]
exists only for \(\dCl=3\). At \(\dCl=3\) the space of such maps is
one-dimensional.
\end{theorem}

\begin{proof}
Dense-subgroup equivariance extends by continuity to \(SO(\Vsp)\). For
\(\dCl\ne3\), the defining representation does not occur in the adjoint module
\(\bigwedge^2\Vsp\); equivalently
\[
  \Hom_{SO(\dCl)}(\bigwedge\nolimits^2\Vsp,\Vsp)=0.
\]
For \(\dCl=3\), the Hodge star gives a nonzero map, and irreducibility together
with Schur's lemma makes the intertwiner space one-dimensional. The
seven-dimensional cross product does not contradict this because it is only
\(G_2\)-equivariant.
\end{proof}

\begin{principle}[Primitive first-order symbol closure]
\label{prin:first-order-symbol-closure}
The oriented volume-dual response survives the continuum scaling limit with
nonzero coefficient, contributes to the principal rather than subprincipal
symbol, lies in the primitive one-reset observable sector, and does not become
an independent higher-form or internal field.
\end{principle}

\begin{corollary}[Dimension consequence of primitive-symbol closure]
\label{cor:first-order-symbol-selects-three}
Under \Cref{prin:first-order-symbol-closure}, with \(\dCl\ge3\),
\[
  \boxed{\dCl=3}.
\]
\end{corollary}

\begin{proof}
The primitive first-order Dirac symbol has Clifford degree one, whereas
\Cref{prop:volume-dual-degree} places the nonzero volume-dual response in degree
\(\dCl-2\). The principle requires these degrees to agree, hence
\(\dCl-2=1\).
\end{proof}

\begin{remark}[Logical scope]
\label{rem:interference-closure-status}
The canonical response \(\bren\), its Hodge degree
\(\bigwedge^{\dCl-2}\Vsp\), the model-theoretic non-implication, and the
equivalence in \Cref{thm:interference-closure-selects-three} are established.
What remains conditional is the physical requirement that the volume-dual
response be an elementary revision, or equivalently the stronger continuum
claim \Cref{prin:first-order-symbol-closure}. The latter isolates four open
requirements: survival under scaling, appearance in the principal symbol,
membership in the primitive one-reset sector, and exclusion of a separate
higher-form field.
\end{remark}

\begin{corollary}[Dimension statement with hard conditional uniqueness]
\label{cor:dimension-final-status}
Within the primitive class,
\[
  \boxed{\dCl=3\ \text{is minimal unconditionally and unique when the nonzero
  oriented volume-dual response closes into the elementary revision sector}.}
\]
In the regular spacetime-crystal phase this is equivalently
\[
  \beff=3,\qquad
  q_{\mathrm{alg}}=q_{\mathrm{met}}=\rank\omega=4,
\]
and, under the optional heat-kernel assumptions, \(d_s=4\).
\end{corollary}


\subsection{Alternative conditional selection mechanisms}
\label{subsec:native-selection-principles}

\Cref{prop:primitive-all-odd} places a sharp ceiling on any unconditional
dimension theorem: every odd spatial rank is realised by a primitive revision
datum.  Hence no consequence of primitivity, locality, exact covariance, the derived
sheet/internal ergodicity, or isotropy alone can exclude
$\dCl=5,7,\ldots$.  Uniqueness therefore requires an additional structural
principle formulated inside the revision sector.

The structural interference-closure criterion has now been proved in full in
\Cref{subsec:interference-closure}.  The remaining criteria in this subsection
are retained only as independent comparisons: the real-even division condition
and the access-based principles provide different conditional routes, but they
are not needed for \Cref{thm:interference-closure-selects-three}.

\begin{theorem}[Real-even division condition]
\label{thm:real-even-division-selects-three}
Fix the Euclidean real Clifford convention carried by the signed spatial sector.
If one requires the real even revision algebra to be a division algebra, then
among odd ranks it permits only $\dCl=1$ and $\dCl=3$:
\[
 \Cl_{1}^{0}(\R)\cong\R,
 \qquad
 \Cl_{3}^{0}(\R)\cong\mathbb H,
\]
whereas for odd $\dCl\ge5$ the even algebra is a nontrivial matrix algebra and
contains zero divisors.  Spatial nondegeneracy excludes $\dCl=1$, leaving
$\dCl=3$.
\end{theorem}

\begin{proof}
The standard even-algebra identity reduces the classification to the real
Clifford table.  Here the spatial convention is \((\gamma^i)^2=+\one\), so
\(\Cl_1^0(\R)\cong\R\) and
\(\Cl_3^0(\R)\cong\Cl_{0,2}(\R)\cong\mathbb H\).  With this convention the
higher odd cases are matrix algebras of size at least two over
$\R$, $\C$, or $\mathbb H$; such matrix algebras contain nonzero zero divisors.
See \cite{atiyah-bott-shapiro,lawson-michelsohn}.
\end{proof}

\begin{corollary}[Concordant native selection]
\label{cor:concordant-native-selection}
Interference closure and the real-even division condition are logically distinct
conditions on the same signed revision sector, and both select
\[
 \boxed{\dCl=3}
\]
uniquely among spatially nondegenerate odd ranks.  They strengthen the dimension
argument by replacing a monotone cost preference with two exact, intrinsic and
concordant criteria.  They remain selection principles rather than consequences
of the renewal axioms.
\end{corollary}

\subsection{Why isotropy cannot make the selection unconditional}
\label{subsec:isotropy-dimension-no-go}

\begin{theorem}[Isotropy is dimension-blind]
\label{thm:isotropy-dimension-blind}
For every $\dCl\ge1$, the cross-polytope frame
$\{\pm e_1,\ldots,\pm e_{\dCl}\}$ with equal weights $1/(2\dCl)$ has second
moment
\[
 M=\frac1{\dCl}I_{\dCl}.
\]
Therefore isotropic self-averaging is available in every spatial dimension and
cannot imply interference closure or select $\dCl=3$.
\end{theorem}

\begin{proof}
Directly,
\[
 \sum_{i=1}^{\dCl}\frac1{2\dCl}
 \bigl(e_ie_i^{\top}+(-e_i)(-e_i)^{\top}\bigr)
 =\frac1{\dCl}I_{\dCl}.
\]
Since the construction works in every rank, isotropy excludes none of the higher
odd primitive endpoints.
\end{proof}

\begin{proposition}[Symmetric and alternating sectors are independent]
\label{prop:symmetric-alternating-independence}
The isotropic second moment belongs to $\operatorname{Sym}^2V_{\rm sp}$, whereas
revision interference belongs to $\bigwedge^2V_{\rm sp}$.  These are inequivalent
natural tensor sectors.  The Dirac-square contraction between metric and spin
connection reduces to the dimension-independent Lichnerowicz scalar term and does
not transfer a rank-selection condition from one sector to the other.
\end{proposition}

\begin{remark}[The dimension frontier]
\label{rem:honest-dimension-frontier}
The failure of isotropy to select the rank is a no-go, not merely an unfinished
proof.  Together with \Cref{prop:primitive-all-odd}, it shows that unconditional
uniqueness of $3+1$ cannot follow from the present renewal, primitive and isotropic
axioms.  The rigorous headline remains unconditional \emph{minimality}; uniqueness
is supported by several concordant native principles, but remains conditional
unless a new representation-level axiom is derived from a deeper layer.
\end{remark}

\subsection{Alternative conditional uniqueness: the access fixed point}
\label{subsec:access-selection-main}

Before stating the selection criterion, we emphasise what the paper actually
relies on.  The central dimension result is the \emph{unconditional}
minimality theorem \Cref{thm:minimal-nondegenerate-3plus1}: in the primitive
class \(2+1\) is impossible, \(1+1\) is spatially degenerate, and \(3+1\) is the
minimal spatially nondegenerate Lorentzian Dirac endpoint.  Nothing downstream of
this subsection feeds back into that theorem.  The access principle introduced
here, and the access-efficiency functional of the next subsection, are
\emph{optional} selection layers: they upgrade minimality to uniqueness among the
higher odd primitive ranks, but they are extra modelling choices, not theorems
about renewal memory, and a reader may stop at
\Cref{thm:minimal-nondegenerate-3plus1} with the headline intact.  We keep them
explicit precisely so that the conditional content is never silently absorbed
into the word ``emergence.''

The preceding subsection leaves a genuine gap: primitivity permits every higher
odd spatial rank.  To turn minimality into uniqueness, an additional selection
criterion is required.  The first such criterion is an access
self-consistency principle.  It asks that the dimension realised by the signed
renewal component agree with the integer dimension preferred by the same signed
budget.

Once the primitive modular revision algebra supplies the full Clifford grading,
the signed budget is the number of noncentral holonomy degrees.  Write
\(B_{\rm sign}=b_1\log2\).

\begin{proposition}[Additive access functional]
\label{prop:additive-access-main}
Among the three elementary ways of spending a signed budget \(e^{B_{\rm sign}}\)
over \(m\) candidate axes -- factored, replicated, and additively partitioned --
only the additive partition has a finite interior optimum.  Its maximal access is
\[
  \Acc_m(B_{\rm sign})=m(B_{\rm sign}-\log m),
\]
attained at the isotropic allocation.
\end{proposition}

\begin{proof}
For a factored allocation \(\prod_i W_i=e^{B_{\rm sign}}\), the logarithmic
access is constant in \(m\).  For a replicated allocation \(W_i=e^{B_{\rm sign}}\),
it is \(mB_{\rm sign}\) and hence runs away with \(m\).  For an additive
allocation \(\sum_iW_i=e^{B_{\rm sign}}\), the arithmetic--geometric mean
inequality gives \(\prod_iW_i\le(e^{B_{\rm sign}}/m)^m\), with equality at equal
allocation.  This gives the displayed functional.
\end{proof}

\begin{principle}[Access self-consistency]
\label{prin:access-self-consistency}
Among primitive signed renewal components, the realised spatial rank is an
access fixed point: the dimension present equals the integer dimension that the
additive access functional prefers for the same signed budget generated by the
component,
\[
  \dCl=\argmax_{m\ge1}\Acc_{m}(\dCl\log2).
\]
This is a selection principle, not a consequence of primitivity alone.
\end{principle}

The resulting fixed-point problem is elementary but useful.  Among odd primitive
ranks it leaves only the one-axis endpoint and the \(3\)-axis endpoint.

\begin{lemma}[Odd access fixed points]
\label{lem:odd-access-fixed-points}
For \(\Acc_m(B)=m(B-\log m)\), the access self-consistency equation
\[
  \dCl=\argmax_{m\ge1}\Acc_{m}(\dCl\log2)
\]
has, among odd positive integers, exactly the solutions \(\dCl=1\) and \(\dCl=3\).
\end{lemma}

\begin{proof}
For fixed \(B\), the increment is
\[
  \Acc_{n+1}(B)-\Acc_n(B)
  =B-\bigl((n+1)\log(n+1)-n\log n\bigr),
\]
and the bracketed term is strictly increasing in \(n\).  Equivalently, the real
stationary point satisfies \(d_*=e^{B-1}\).  At the self-consistent budget
\(B=\dCl\log2\), this is \(m_*=2^{\dCl}/e\).  For \(\dCl=1\) the integer maximizer is
\(1\); for \(\dCl=3\), \(8/e\) lies in the integer cell of \(3\).  For every odd
\(\dCl\ge5\), \(2^{\dCl}/e>\dCl+1/2\), so the integer maximizer is strictly larger than
\(\dCl\).  Thus the only odd fixed points are \(1\) and \(3\).
\end{proof}

\begin{theorem}[Conditional access selection of the nondegenerate endpoint]
\label{thm:conditional-access-3plus1}
Assume the signed renewal component is primitive, the finite-shell modular form
stabilises, and the access self-consistency principle \Cref{prin:access-self-consistency}
holds.  Then the only primitive access-self-consistent endpoints are \(1+1\)
and \(3+1\).  The first is spatially degenerate by
\Cref{def:spatially-nondegenerate}.  Therefore the unique nondegenerate
access-self-consistent primitive endpoint is \(3+1\).
\end{theorem}

\begin{proof}
By \Cref{thm:primitivity-canonical}, primitivity gives a Clifford revision algebra
on \(\Hf\).  The rank parity argument in
\Cref{thm:minimal-nondegenerate-3plus1} forces \(\dCl\) odd.  By
\Cref{lem:odd-access-fixed-points}, access self-consistency leaves only
\(\dCl=1\) and \(\dCl=3\).  Excluding the spatially degenerate one-axis endpoint leaves
\(\dCl=3\).  The corresponding \(3+1\) Lorentzian Dirac endpoint is the \(\dCl=3\)
specialisation of the isotropic Clifford rounding and operator-convergence
theorems.
\end{proof}


\subsection{A secondary optional functional: access efficiency}
\label{subsec:access-efficiency}

The access fixed-point principle is one way to rank the odd primitive
endpoints.  We now record a second, more transparent functional built directly
from two quantities already present in the construction: the operator dimension
of the primitive Clifford revision algebra and the smallest isotropic reset
frame needed for the spatial symbol.  This criterion is useful because it gives
the same conditional conclusion by a simpler monotonicity argument.

The functional ranks the odd primitive endpoints as
\[
  \dCl=1 > \dCl=3 > \dCl=5 > \dCl=7 > \cdots,
\]
and therefore selects \(\dCl=3\) once the independently established
spatial-nondegeneracy condition removes \(\dCl=1\).  As above, this is a
selection criterion rather than an algebraic exclusion: the higher odd endpoints
of \Cref{prop:primitive-all-odd} remain admissible.

\begin{lemma}[Primitive complexity gap]
\label{lem:complexity-gap}
For odd candidate spatial Clifford rank \(m\), a primitive revision datum has
\(\Arev\cong M_{2^{(m+1)/2}}(\C)\), equivalently
\(\Arev\cong\Cl_{m+1}(\C)\), by \Cref{thm:primitivity-canonical}; hence operator
dimension
\[
  \dim_\C\Arev(m)=\dim_\C\Cl_{m+1}=2^{\,m+1}.
\]
Consequently
\[
  \frac{\dim_\C\Arev(m+2)}{\dim_\C\Arev(m)}
  =\frac{2^{\,m+3}}{2^{\,m+1}}=4 .
\]
Higher odd endpoints are not forbidden, but each step \(m\mapsto m+2\) multiplies
the primitive revision-algebra dimension by four.
\end{lemma}

\begin{proof}
Primitivity is full \(\F_2\)-rank of the modular commutator form
(\Cref{thm:primitivity-canonical}), giving
\(\Arev\cong M_{2^{(m+1)/2}}(\C)\), equivalently the even complex Clifford
algebra on \(m+1\) generators.  Its complex vector-space dimension is
\(2^{m+1}\).  The ratio is
\(2^{(m+3)-(m+1)}=2^2=4\).
\end{proof}

\begin{lemma}[Minimal isotropic access frame]
\label{lem:minimal-frame}
Among finite, antipodal, equal-weight direction systems \(\nu\) on
\(\Sphere^{\dCl-1}\) whose second moment is isotropic,
\(M(\nu)=\tfrac1{\dCl} I_{\dCl}\), the minimal cardinality is
\[
  N_{\min}(\dCl)=2\dCl,
\]
attained, uniquely up to an orthogonal change of spatial frame, by the
cross-polytope frame \(\{\pm e_1,\dots,\pm e_{\dCl}\}\).
\end{lemma}

\begin{proof}
An antipodal equal-weight system is a union of \(k\) antipodal pairs
\(\{\pm u_1,\dots,\pm u_k\}\) of unit vectors, of cardinality \(2k\) and equal
weight \(1/(2k)\).  Its second moment is
\[
  M(\nu)=\sum_{a=1}^{k}\tfrac1{2k}\bigl(u_au_a^{\top}+(-u_a)(-u_a)^{\top}\bigr)
        =\frac1k\sum_{a=1}^{k}u_au_a^{\top}.
\]
Isotropy \(M(\nu)\propto I_{\dCl}\) forces \(\sum_a u_au_a^{\top}\propto I_{\dCl}\), i.e.\
the \(u_a\) form a tight frame for \(\R^{\dCl}\), which requires \(k\ge \dCl\).  With
\(k=\dCl\) unit vectors, a tight frame must be an orthonormal basis, giving
\(\sum_a e_ae_a^{\top}=I_{\dCl}\) and \(M(\nu)=\tfrac1{\dCl} I_{\dCl}\) (trace \(1\) as
required).  Hence the minimum is \(k=\dCl\), i.e.\ \(N_{\min}(\dCl)=2\dCl\), realised by
\(\{\pm e_i\}\).  This is the \(\dCl\)-dimensional cross-polytope frame; for \(\dCl=3\)
it is the octahedral frame of \Cref{ex:3plus1}.
\end{proof}

The two preceding lemmas isolate the numerator and denominator of the
efficiency ratio: isotropic spatial access divided by primitive revision cost.

\begin{definition}[Primitive access efficiency]
\label{def:access-efficiency}
For a candidate spatial Clifford rank \(m\), let
\[
  C_{\rm rev}(m):=\dim_\C\Arev(m)=2^{m+1}
\]
be the primitive revision cost, and let
\[
  A_{\rm iso}(m):=N_{\min}(m)=2m
\]
be the minimal isotropic access count.  The primitive access efficiency is
\[
  \eta(m):=\frac{A_{\rm iso}(m)}{C_{\rm rev}(m)}
          =\frac{2m}{2^{\,m+1}}
          =\frac{m}{2^{\,m}} .
\]
Thus \(\eta(m)\) measures isotropically accessible spatial directions per
primitive Clifford revision mode for rank \(m\); the denominator is the
operator-theoretic cost of the primitive revision algebra.
\end{definition}

The monotonicity of this efficiency is immediate, but it provides a useful
second reading of why the first nondegenerate odd endpoint is preferred.

\begin{theorem}[Access-efficiency selection of the nondegenerate endpoint]
\label{thm:access-efficiency-selection}
With \(\eta(m)=m/2^{m}\) of \Cref{def:access-efficiency}:
\begin{enumerate}[label=\textup{(\roman*)},leftmargin=2.4em]
\item \(\eta(1)=\eta(2)=\tfrac12\), and \(\eta\) is strictly decreasing
      for integer \(\dCl\ge2\) \textup{(}since
      \(\eta(m{+}1)/\eta(m)=(m{+}1)/(2m)<1\) iff \(m>1\)\textup{)}; in
      particular \(\eta\) is strictly decreasing along the odd integers
      \(m=1,3,5,\dots\), with
      \[
        \eta(1)=\tfrac12>\eta(3)=\tfrac38>\eta(5)=\tfrac5{32}>\cdots ;
      \]
\item among \emph{spatially nondegenerate} odd primitive endpoints
      \textup{(}\(\dCl\ge3\) odd, \Cref{def:spatially-nondegenerate}\textup{)},
      \(\eta\) is uniquely maximised at \(\dCl=3\);
\item read as a selection functional, the unique spatially
      nondegenerate primitive endpoint of maximal access efficiency is
      \[
        \boxed{\,\dCl=3\,},
      \]
      while \(\dCl=5,7,\dots\) remain algebraically admissible
      \textup{(}\Cref{prop:primitive-all-odd}\textup{)} but
      strictly less efficient.
\end{enumerate}
\end{theorem}

\begin{proof}
(i) \(\eta(1)=1/2\), \(\eta(2)=2/4=1/2\), and
\[
  \eta(m{+}1)/\eta(m)
  =\bigl((m{+}1)/2^{m+1}\bigr)\bigl(2^{m}/m\bigr)
  =\frac{m+1}{2m},
\]
which is \(<1\) exactly when \(m+1<2m\), i.e.\ \(m>1\).  Hence \(\eta\) is
strictly decreasing for \(\dCl\ge2\), and the displayed odd chain follows.  (ii)
Spatial nondegeneracy excludes \(\dCl=1\) (\Cref{def:spatially-nondegenerate}); on
the remaining odd \(\dCl\ge3\), (i) gives a strictly decreasing sequence, so the
maximum is the smallest admissible value \(\dCl=3\).  (iii) is the selection reading
of (ii), with the explicit values \(\eta(5)=5/32\), \(\eta(7)=7/128\) exhibiting
the strict loss; admissibility of higher ranks is \Cref{prop:primitive-all-odd}.
\end{proof}

\begin{remark}[Exactly what is, and is not, selected]
\label{rem:efficiency-honest}
Four points are needed for the statement to be mathematically precise.
\begin{enumerate}[label=\textup{(\alph*)},leftmargin=2.2em]
\item \emph{\(\dCl=1\) is the unconstrained maximum.}  Since \(\eta(1)=\tfrac12>
      \eta(3)=\tfrac38\), efficiency \emph{by itself} is maximised at \(\dCl=1\),
      not \(\dCl=3\).  The restriction to \(\dCl\ge3\) is the spatial-nondegeneracy
      condition of \Cref{def:spatially-nondegenerate} (the \(\dCl=1\) exclusion
      already used in \Cref{thm:minimal-nondegenerate-3plus1}), established
      independently of any cost.  Thus \(\eta\) selects \(\dCl=3\) only \emph{after}
      nondegeneracy removes \(\dCl=1\); it does not make \(\dCl=3\) the globally most
      efficient rank.
\item \emph{The selection is cost-dependent.}  It uses the full Clifford module
      dimension \(2^{\dCl+1}\) as the revision cost.  With the logarithmic cost
      \(\log_2\dim\Cl_{\dCl+1}=\dCl+1\), the efficiency \(2\dCl/(\dCl+1)\) is strictly
      increasing toward \(2\) and selects nothing.  The exponential cost is
      justified operator-theoretically: the spectral package acts on the entire
      Clifford module, so its dimension is the natural operator-level cost.  This
      is nevertheless a modelling choice, not an algebraic theorem.  Hence
      \(\eta\) is an internal selection criterion, tagged , exactly
      like \Cref{prin:access-self-consistency}.
\item \emph{\(\dCl=5\) is not defeated.}  \Cref{prop:primitive-all-odd} shows that higher odd ranks remain a valid
      nondegenerate primitive endpoint; \(\eta\) only ranks it lower.  This is
      consistent with \Cref{rem:minimal-not-unique}: no algebraic exclusion of
      higher odd Clifford ranks is claimed or implied.
\item \emph{Corroboration, not independence.}  Both \(\eta\) and the access
      self-consistency principle place the first two odd primitive endpoints
      \(\dCl=1\) and \(\dCl=3\) ahead of all higher odd Clifford ranks; after the independent
      spatial-nondegeneracy exclusion of \(\dCl=1\), both criteria select \(\dCl=3\).
      They are two readings of one exponential-cost backbone, so they reinforce
      the conditional \(\dCl=3\) selection rather than confirming it from logically
      independent premises.
\end{enumerate}
\end{remark}

\begin{remark}[Net effect on the dimension headline]
\label{rem:efficiency-headline}
The unconditional headline is unchanged: by
\Cref{thm:minimal-nondegenerate-3plus1}, \(3+1\) is the minimal nondegenerate
primitive endpoint and \(2+1\) is impossible.  What the efficiency functional adds
is a second, transparent \emph{conditional} selection that agrees with the access
self-consistency selection of \Cref{thm:conditional-access-3plus1}: \(3+1\) is the
unique spatially nondegenerate primitive endpoint of maximal access efficiency
\(\dCl/2^{\dCl}\).  The strengthened conditional statement is therefore
\[
  3+1=\text{minimal nondegenerate}\ \textbf{and}\ \text{uniquely access-efficient
  primitive endpoint},
\]
with the \(\dCl=3\) \emph{uniqueness} remaining, irreducibly, conditional on an
intrinsic cost choice --- two such choices now agreeing.
\end{remark}


\subsection{Why decimation does not discharge the access principle}
\label{subsec:decimation-limits}

One might hope that the access principle is not really an additional principle,
but merely the fixed point of an ordinary renewal coarse-graining.  This
subsection explains why that hope fails for quotient-type decimations.  Such
decimations identify predictive classes, contract or identify edges, delete
transient structure, or pass to a coarser shift congruence.  They do not create
new independent cycles.

\begin{lemma}[Quotient decimation is non-increasing in \(b_1\)]
\label{lem:quotient-decimation}
For any quotient-type decimation \(\mathcal D\),
\[
  b_1(\mathcal DG)\le b_1(G).
\]
\end{lemma}

\begin{proof}
A quotient-type decimation identifies vertices and/or edges and kills or
identifies cycles.  The induced map on first homology over \(\Z/2\) is therefore
surjective from the original cycle space onto the decimated cycle space, so the
rank cannot increase.
\end{proof}

\begin{proposition}[Quotient decimation cannot derive access selection]
\label{prop:decimation-cannot-access}
Quotient-type decimation cannot implement the access fixed-point map
\[
  F(b)=\argmax_{d'\ge1}\Acc_{d'}(b\log2).
\]
Indeed quotient decimation is non-increasing in \(b_1\), whereas \(F(b)\) grows
like \(2^b/e\) and satisfies \(F(b)>b\) for all sufficiently large \(b\) (in
particular for \(b\ge5\)).  Thus a rule sending \(b\) to \(F(b)\) is not an
ordinary renewal coarse-graining; it is precisely an additional access
optimisation principle.
\end{proposition}

\begin{proof}
The first assertion is \Cref{lem:quotient-decimation}.  For the access map, the
continuous maximizer of \(\Acc_{d'}(b\log2)\) is \(2^b/e\), and integer rounding
therefore grows exponentially in \(b\).  For \(b\ge5\), \(2^b/e>b+1/2\), hence
\(F(b)>b\).  A non-increasing quotient operation cannot produce this increasing
map.  Defining a ``decimation'' to output \(F(b)\) would simply restate the
access principle in different language.
\end{proof}

\begin{remark}[Scope of the dimension statement]
\label{rem:dimension-status}
The unconditional headline is: \emph{primitivity forbids \(2+1\) and makes
\(3+1\) the minimal nondegenerate Lorentzian Dirac endpoint}.  The preferred hard
conditional uniqueness theorem is
\Cref{thm:interference-closure-selects-three}: oriented interference closure
excludes every higher odd endpoint.  The access fixed point and efficiency
functional below remain alternative, softer conditional criteria and are not
consequences of ordinary coarse-graining.
\end{remark}

\subsection{Interfaces with downstream constructions}
\label{subsec:downstream-interfaces}

The primitive \(3+1\) result remains kinematic and operator-theoretic.  It does
not produce matter multiplets, field equations, or phenomenology.  It says that
the primitive signed renewal sector supplies the Clifford revision algebra and
that \(3+1\) is the minimal nondegenerate endpoint, with hard conditional
uniqueness under oriented interference closure and optional alternative access
criteria.  There are four natural downstream
interfaces.

\paragraph{Balanced reset statistics.}
The reset statistics used in the continuum limit may be treated as
self-consistent rather than externally prescribed.  In a balancing extension one
models reset ages by a piecewise-deterministic Markov process and asks for a
stationary law whose induced reset frequencies and channel weights reproduce the
renewal data.  Under covering and feasibility hypotheses such a balanced law
exists; when feasibility fails, a separating semidefinite witness gives an
obstruction; and under contraction or determined monotonicity hypotheses the
balanced law is unique and stable up to the evident predictive gauge freedom.
This layer is not used in the proofs of the present article.  Its role is to
explain how the second-moment fields and predictive measures used in the
continuum limit could arise self-consistently; it is not needed to enforce
rotational isotropy, since the flat endpoint depends only on the nondegenerate
metric defined by the selected tensor.

After predictive collapse, a balanced law pushes forward to a canonical measure
$\mu_{\CP}$ on predictive classes.  In conformal regular regimes its information
dimension is
\[
  \operatorname{Dim}_1(\mu_{\CP})=\frac{h}{\lambda},
\]
where $h$ is the entropy rate of the balanced reset law and $\lambda$ is the
corresponding contraction exponent.  Thus $\qmet$ is the unweighted capacity
dimension of the available predictive attractor, while $h/\lambda$ is the dimension
actually filled by the stationary renewal dynamics.  The gap is a free-energy
defect, equivalently a relative entropy from the dimension-maximising measure,
and should not be confused with entropy production.

\paragraph{Finite matter geometry.}
An almost-commutative or finite-real-spectral-triple matter sector can be
coupled to the signed renewal geometry by taking a product with the finite
internal renewal quotient.  Lorentzian versions of the noncommutative Standard
Model already show that signature \((1,3)\), KO-dimension, fermion reduction,
and the finite indefinite structure must be treated together
\cite{Barrett2007Lorentzian,BochniakSitarz2018,
BochniakSitarzZalecki2020}.  Twisted constructions provide another route from
the Euclidean almost-commutative model to Lorentzian fermionic data
\cite{DevastatoFarnsworthLizziMartinetti2018,
DevastatoFilaciMartinettiSingh2020}, and finite extensions such as
\(U(1)_{B-L}\) illustrate the downstream model-building interface
\cite{BesnardBrouder2019BL}.  The present paper does not reproduce those matter
models.  It supplies the external renewal-generated Lorentzian spectral datum;
the signed/Krein generator, real structure, and first-order conditions would
then constrain which finite matter sectors can be coupled consistently.

\paragraph{Amplitude and charge local systems.}
The finite \(\Zfour\) lift of \Cref{subsec:real-amplitude-lifts} is not used to
obtain Lorentzian signature; the \(\Ztwo\) signed cover already does that.  It is
instead the correct interface for downstream charged local systems.  Any matter
construction using phases, charged incidence, or square roots of projective
signs should therefore be formulated in the amplitude category and then reduced
to the real/Krein quotient when only signature is being read.  If a central
\(\Zsix\) screen is introduced in a finite matter sector, the renewal-compatible
form is a composite of the Krein parity with an independent three-cycle renewal
or generation grading,
\[
  \Ztwo^{\rm Krein}\times\Z/3\cong\Zsix,
\]
rather than a new primitive assumption.  Likewise, neutral or dark twisted
sectors constructed from real \(\Ztwo\) holonomy should be regarded as real
shadows until their amplitude, Pin/Majorana, or charged-local-system refinement
has been specified.  Finally, real-category protection statements for larger
ambient cells or exteriors do not by themselves rule out amplitude-level
splittings after a \(\Zfour\) lift; non-adoption of such ambient cells must rest
on independent minimality, projection, or dynamical-selection criteria.

\paragraph{Dynamical gravity.}
Dynamical gravity requires additional input not present here: a rule identifying
variations of reset statistics with a connection, curvature, and an entropy or
action principle.  The adiabatic curved Dirac theorem of
\Cref{sec:curved-stability} now supplies the complete geometric Dirac operator
through subprincipal order: the renewal-generated coframe determines the
discrete torsion-free Cartan connection, whose limit is the Levi--Civita spin
connection, and the covariant Hamiltonians converge in strong resolvent sense.
It does not derive Einstein equations, a curvature evolution law, or a
gravitational variational principle.

\paragraph{Modular and conformal completions.}
Modular or conformal extensions would require placing the renewal inclusion into
a genuine modular-theoretic framework.  The signed cover and the renewal
resolvent give the correct algebraic ingredients for such a comparison, but the
present paper does not construct a conformal net, a string worldsheet theory, or
a quantum-gravity model.

Thus the present article should be read as the foundation layer: it proves the
positive predictive geometry, the signed/Krein obstruction and classification,
the Lorentzian Dirac endpoint, universality and enhanced metric reconstruction,
and the primitive modular minimal-endpoint interface.  Matter, dynamics,
conformal extensions, and phenomenology are logically downstream.


\begin{remark}[Summary of the dimension statement]
\label{rem:primitive-dimension-summary}
The logical content of this section is therefore deliberately stratified.
Topology supplies signed holonomy but not Clifford anticommutation; modular
revision supplies the commutator form.  Under full-dimensional affine
reachability and finite primitive control, bounded return completion derives the
predictive growth--cycle identity.  The accessible modular rank then follows
from nondegeneracy, while topological-crystal and heat-kernel regularity supply
the return dimension.  The channel-metric equality remains a separate
non-collapse rider, as stated in \Cref{thm:dimension-coincidence}.  Primitivity then
forces even total Clifford rank and hence odd spatial rank; spatial
nondegeneracy removes the \(1+1\) endpoint; and \(3+1\) becomes the minimal
nondegenerate primitive endpoint.  Uniqueness among higher odd endpoints is conditional: the primary theorem uses
oriented interference closure, while the access fixed point and access-efficiency
functional are retained only as alternative optional criteria.  None is claimed
as a consequence of primitivity alone.
\end{remark}




\subsection{The canonicity theorem}
\label{sec:cxv}

We now consolidate the preceding characterisation results into a single
statement that complements the main assembly theorem.  The result gathers the
modular, symplectic, and Cartan reconstructions under one transparent hypothesis
package and isolates the sole remaining datum needed to upgrade dimensional
minimality to uniqueness.

\begin{axiom}[Canonical renewal hypotheses]
\label{ax:canonical}
A signed renewal memory is \emph{canonical} if:
\begin{enumerate}[label=\textup{(C\arabic*)},leftmargin=2.8em]
\item \textup{[connected recurrence]} the recurrent predictive component is
      finite, connected, graded with unit length increments, and its signed cover
      is connected;
\item \textup{[modular homogeneity]} the elementary resets are analytic modular
      eigenoperators of a faithful weight with a common eigenvalue $\beta$
      \textup{(}\Cref{ax:modular-homogeneity}\textup{)};
\item \textup{[reversible predictive revision phase]} the spatial revision
      labels are predictive units satisfying \Cref{ass:signed-revision-phase};
      hence their action on the finite recurrent factor is by automorphisms and
      has the automatic projective lift of
      \Cref{thm:spatial-multiplier-scalar-main}; the optional finite
      \(\Zfour\) square-root completion is retained only at the amplitude level;
\item \textup{[nondegenerate covariant second moment]} the
      renewal-generated coframe $e^a_j=\kappa M^{aj}$ is uniformly
      nondegenerate, $M\succeq m_0 I$, and $C^4_b$, and neighbouring
      second moments obey the transport covariance of
      \Cref{ass:second-moment-covariance};
\item \textup{[exchange-symmetric translational phase]} the resolved
      two-reset diamond records satisfy \Cref{ass:stationary-diamond}, so
      first-moment plaquette closure follows from
      \Cref{thm:stationary-exchange}.
\end{enumerate}
\end{axiom}

\begin{theorem}[Canonicity of the renewal Lorentzian datum]
\label{thm:canonicity}
For a canonical signed renewal memory:
\begin{enumerate}[label=\textup{(\arabic*)},leftmargin=2.6em]
\item \textup{[E/P]} the unbounded signed modular operator is uniquely
      \[
        D_-=\Jchi\,e^{\beta N}
      \]
      up to an overall scale and a bounded $\Jchi$-self-adjoint perturbation $B$;
      the modular Hamiltonian is affine in the clock, $\log\Delta_\varphi=\beta N
      +c\one$, with $\beta$ the common reset modular eigenvalue
      \textup{(}the modular characterisation and affine modular-flow results\textup{)};
\item the reversible predictive revision action canonically determines the
      projective multiplier, alternating form $\omega$, and reduced symplectic revision space
      \[
        \Hprim=\Hf/\rad\omega,\qquad \bar\omega\ \text{nondegenerate},
      \]
      with reduced revision algebra of Morita class $[M_{2^m}(\C)]=
      [\Cl_{2m}(\C)]$, factor representations parametrised by
      $\widehat{\rad\Omega}$ \textup{(}the canonical reduction and affine modular-flow results\textup{)};
\item the renewal coframe carries a unique metric-compatible
      torsion-free transport, with torsion the translational and curvature the
      rotational part of the reset-channel plaquette defect, converging to the
      Levi--Civita Dirac transport $\omega^{\mathrm{LC}}+O(h^2)$
      \textup{(}the Cartan and channel-composition results\textup{)};
\item \textup{[E/C]} the only remaining dimensional-selection datum is the
      nonzero projection of the oriented two-reset circulation $\Xi$ onto the
      principal first-order sector; this is nonzero iff $\dCl=3$ and $\Xi\neq0$,
      and is RG-marginal exactly at $\dCl=3$ in the explicit block-renewal model
      \textup{(}the survival criterion and block-renewal coarse-graining theorem\textup{)}.
\end{enumerate}
\end{theorem}

\begin{proof}
(1) combines the modular signed-Dirac characterisation with
\Cref{thm:affine-modular}, which proves that the modular Hamiltonian is affine in
the renewal clock.  (2) follows from the automatic projective lift and canonical symplectic
reduction, together with
\Cref{thm:factor-quotients-corrected,thm:radical-centre}.
(3): connected recurrence and the nondegenerate
$C^4_b$ coframe give the discrete fundamental theorem of the Cartan characterisation; isometric
comparison gives metric compatibility; second-order plaquette equivalence gives
$T=0$ via \Cref{thm:channel-torsion}; uniqueness then yields Levi--Civita with the
$O(h^2)$ rate.  (4) combines the survival criterion with the
block-renewal eigenvalue \Cref{thm:rg-eigenvalue} and \Cref{cor:relevance}.
\end{proof}

\begin{remark}[Reading of the canonicity theorem]
\label{rem:canonicity-reading}
Under \Cref{ax:canonical}, the modular weight, primitive revision sector, and
Cartan transport are forced by the displayed hypotheses.  Dimensional uniqueness
is reduced to the single sharp datum: whether oriented two-reset circulation projects nontrivially onto
the principal first-order sector.  Unconditionally this gives \emph{minimality}
($3+1$ is the minimal nondegenerate primitive endpoint); the \emph{uniqueness}
upgrade is the marginal/principal selection, a theorem in the explicit
block-renewal model and a principle in general.  This is the precise boundary of
the foundational layer.
\end{remark}

% =====================================================================


% =====================================================================


\section{Further continuum, propagation, and global-topology results}
\label{app:further-continuum}

\subsection{A least-distortion counterexample}

\begin{example}[Polar comparison is torsionful on a curved surface]
\label{ex:polar-not-LC}
Take $\dCl=2$ and the warped coframe $E(x)=\mathrm{diag}(1,f(x^1))$,
$f>0$ non-affine, i.e.\ metric $ds^2=(dx^1)^2+f(x^1)^2(dx^2)^2$ (a surface of
revolution, curved for $f''\neq0$).  The frame Gram is $P=\mathrm{diag}(1,f^2)$,
and $E E_{,1}^\top=\mathrm{diag}(0,ff')$ is symmetric, so $\Asym(EE_{,1}^\top)=0$;
the Sylvester equation of \Cref{rem:procrustes} gives the polar connection
$\omega^{\mathrm{polar}}\equiv0$.  Its torsion is
\[
  T^2_{12}=\partial_1 e^2_2-\partial_2 e^2_1=f'\neq0 ,
\]
so the polar (least-distortion) connection has nonzero torsion.  By contrast the
Levi--Civita spin connection in this frame is $\gamma^{\mathrm{LC}}$ with the only
nonzero component $\gamma_{2;12}=f'$, which is metric \emph{and} torsion-free.
Hence
\[
  \omega^{\mathrm{polar}}=0\ \neq\ \omega^{\mathrm{LC}} .
\]
Least distortion supplies metric compatibility but not Levi--Civita; the missing
ingredient is plaquette closure.
\end{example}

\subsection{Joint convergence of propagation geometry}
\label{sec:joint-propagation-geometry}

The preceding operator limit determines more than a principal symbol.  The same
second-moment field controls the characteristic cone, its optical distance, and
the macroscopic propagation of the renewal evolution.  This subsection records
the precise joint statement and separates it from the combinatorial
reset-ray order, which depends on a different statistic.

Write
\[
  \widehat g(x)=\kappa^{-2}M(x)^{-2},
  \qquad
  g_{\mu\nu}(x)=\diag\bigl(1,-\widehat g(x)\bigr).
\]
Uniform ellipticity gives
\[
  \kappa^{-2}|\vec v|^2
  \le \widehat g(x)[\vec v,\vec v]
  \le \kappa^{-2}m_0^{-2}|\vec v|^2,
\]
so \(\widehat g\) is complete and its distance \(d_{\widehat g}\) is
bi-Lipschitz equivalent to the Euclidean distance.  The future cone is
\[
  \cone_g(x)
  =\{(v^0,\vec v):v^0\ge0,\ (v^0)^2\ge
    \widehat g(x)[\vec v,\vec v]\},
\]
whose unit-time slice is the ellipsoid
\[
  \cone_g(x)\cap\{v^0=1\}=\kappa M(x)B^{\dCl}.
\]

\begin{lemma}[Stability of the optical metric and distance]
\label{lem:optical-distance-stability}
Let \(M,\widetilde M\) be symmetric fields with
\(m_0I\preceq M,\widetilde M\preceq I\).  Then
\[
 \norm{M^{-2}-\widetilde M^{-2}}
 \le 2m_0^{-4}\norm{M-\widetilde M}.
\]
Consequently, if \(M_R\to M\) locally uniformly in a common uniformly elliptic
class, then
\[
 d_{\kappa^{-2}M_R^{-2}}\longrightarrow d_{\widehat g}
\]
locally uniformly on \(\R^{\dCl}\times\R^{\dCl}\).
\end{lemma}

\begin{proof}
The algebraic identity
\[
 M^{-2}-\widetilde M^{-2}
 =M^{-2}(\widetilde M^2-M^2)\widetilde M^{-2}
\]
is valid without commutativity, and
\[
 \widetilde M^2-M^2
 =(\widetilde M-M)\widetilde M+M(\widetilde M-M).
\]
The displayed operator bound follows from
\(\norm{M^{-2}},\norm{\widetilde M^{-2}}\le m_0^{-2}\) and
\(\norm{M},\norm{\widetilde M}\le1\).

For the distance statement, fix a compact \(K_0\).  Uniform ellipticity confines
all minimizing geodesics between points of \(K_0\) to one larger compact set
\(K\), independent of \(R\).  For any absolutely continuous curve \(\gamma\)
in \(K\),
\[
 \left|\sqrt{\widehat g_R(\dot\gamma,\dot\gamma)}
       -\sqrt{\widehat g(\dot\gamma,\dot\gamma)}\right|
 \le C(m_0,\kappa)\,
      \norm{M_R-M}_{C^0(K)}\,|\dot\gamma|.
\]
Integrating along a minimizing curve for either metric and reversing the roles
of the two metrics gives a uniform estimate on \(K_0\times K_0\), which tends
to zero.
\end{proof}

\begin{proposition}[Frozen-symbol and cone convergence]
\label{prop:joint-symbol-cone}
Let \(M_R\to M\) locally uniformly and \(h_R\to0\).  On every compact
\(K\subset\R^{\dCl}\) and momentum ball \(|\xi|\le\Lambda\),
\[
 \sup_{x\in K,\,|\xi|\le\Lambda}
 \norm{\sigma_{\mathrm{fr}}(\Ham_R)(x,\xi)-\kappa A(M(x)\xi)}
 \longrightarrow0.
\]
Moreover
\[
 \sup_{x\in K}
 d_{\mathrm H}\!\left(\kappa M_R(x)B^{\dCl},
                       \kappa M(x)B^{\dCl}\right)
 \le \kappa\sup_{x\in K}\norm{M_R(x)-M(x)}
 \longrightarrow0.
\]
\end{proposition}

\begin{proof}
The frozen symbol is
\[
 \sigma_{\mathrm{fr}}(\Ham_R)(x,\xi)
 =\kappa\sum_a p_a^{(R)}(x)A(\hat\theta_a)
   \frac{\sin(h_R\hat\theta_a\cdot\xi)}{h_R}.
\]
Using \(|\sin s-s|\le |s|^3/6\), its difference from
\(\kappa A(M_R(x)\xi)\) is uniformly \(O(h_R^2\Lambda^3)\); the remaining
difference is bounded by
\(C\Lambda\norm{M_R-M}\).  The Hausdorff estimate follows by coupling the points
\(\kappa M_Ru\) and \(\kappa Mu\), \(|u|\le1\).
\end{proof}

\begin{proposition}[Energy-estimate cone and finite propagation]
\label{prop:joint-finite-propagation}
Let \(u(t)=e^{-it\Ham_g}u_0\).  If \(u_0\) is supported in a closed set \(S\),
then
\[
 \supp u(t)\subseteq
 J_t^+(S):=\{\vec q:d_{\widehat g}(\vec q,S)\le t\}.
\]
The frozen spatial propagation body is exactly
\(\kappa M(x)B^{\dCl}\): its support function in direction \(n\) is
\(\kappa|M(x)n|\), which is both the maximal characteristic speed and the
support function of the convex hull of the group velocities.
\end{proposition}

\begin{proof}
The equation is a symmetric hyperbolic system with Hermitian principal
coefficients
\(B^j(x)=\kappa A^iM^{ij}(x)\).  In a unit conormal direction \(n\),
\[
 \rho\!\left(\sum_j B^jn_j\right)
 =\rho\bigl(\kappa A(Mn)\bigr)=\kappa|Mn|.
\]
The standard energy estimate therefore propagates support inside the control
cone whose unit-time slice has support function \(n\mapsto\kappa|Mn|\), namely
\(\kappa MB^{\dCl}\).  Integrating this cone field gives precisely the optical
arrival-time condition \(d_{\widehat g}(\vec q,S)\le t\).

For the two dispersion branches
\(\lambda_\pm(\xi)=\pm\kappa|M\xi|\),
\[
 \nabla_\xi\lambda_\pm
 =\pm\kappa\frac{M^2\xi}{|M\xi|}.
\]
Their convex hull has the same support function
\[
 \sup_{\xi\ne0}
 n\cdot\kappa\frac{M^2\xi}{|M\xi|}
 =\kappa|Mn|,
\]
and hence equals the same ellipsoid.
\end{proof}

\begin{theorem}[Macroscopic domain-of-dependence convergence]
\label{thm:macroscopic-domain}
Under the hypotheses of \Cref{thm:curved-limit}, fix
\(u_0\in L^2\) with compact support \(S\).
\begin{enumerate}[label=\textup{(\roman*)},leftmargin=2.4em]
\item If \(F\) is closed and \(F\cap J_t^+(S)=\varnothing\), then
\[
 \norm{\one_F e^{-it\Ham_R}u_0}_{L^2}\longrightarrow0.
\]
The same conclusion holds for any \(L^2\)-precompact family
\(u_{0,R}\to u_0\); in particular it excludes lattice-scale doubler families.
\item For open \(U,V\subset\R^{\dCl}\), if
\[
 \one_Ue^{-it\Ham_g}\one_V\ne0,
\]
then there exists a fixed \(u_0\in L^2(V)\) such that
\[
 \liminf_{R\to\infty}
 \norm{\one_Ue^{-it\Ham_R}u_0}_{L^2}>0.
\]
\end{enumerate}
Thus the outer Lorentzian domain of dependence is unconditional, while inner
filling for variable \(M(x)\) is exactly conditional on local nonvanishing of
the continuum propagator.
\end{theorem}

\begin{proof}
Strong resolvent convergence in \Cref{thm:curved-limit} gives
\(e^{-it\Ham_R}u_0\to e^{-it\Ham_g}u_0\) strongly, uniformly for \(t\) in compact
sets.  By \Cref{prop:joint-finite-propagation},
\(\one_Fe^{-it\Ham_g}u_0=0\); hence
\[
 \norm{\one_Fe^{-it\Ham_R}u_0}
 \le \norm{(e^{-it\Ham_R}-e^{-it\Ham_g})u_0}\to0.
\]
For a precompact family, add and subtract \(e^{-it\Ham_R}u_0\).

For (ii), choose \(v\in L^2(V)\) with
\(c:=\norm{\one_Ue^{-it\Ham_g}v}>0\).  Then
\[
 \norm{\one_Ue^{-it\Ham_R}v}
 \ge c-\norm{(e^{-it\Ham_R}-e^{-it\Ham_g})v},
\]
and the second term tends to zero.
\end{proof}

\begin{remark}[No exact finite-mesh causal order]
\label{rem:no-finite-mesh-order}
Each \(\Ham_R\) is bounded and finite range, but
\(e^{-it\Ham_R}=\sum_{n\ge0}(-it\Ham_R)^n/n!\) contains arbitrarily long paths
at every \(t>0\).  The discrete evolution therefore has rapidly decaying
instantaneous tails rather than an exact light cone.  The statement of
\Cref{thm:macroscopic-domain} concerns macroscopic mass in the continuum limit,
not nonzero amplitude at finite mesh.
\end{remark}

\subsubsection{Exact constant-coefficient support}

Assume now \(M(x)\equiv M_0\), with \(M_0\) symmetric positive definite, and set
\[
 \Ham_0=\kappa A(M_0D),\qquad
 K_t=e^{-it\Ham_0}.
\]
After the linear change \(y=M_0^{-1}x\), this is the standard massless Dirac
evolution of speed \(\kappa\).

\begin{lemma}[Forward Dirac kernel and its support]
\label{lem:constant-dirac-kernel}
Let \(\mathcal K_t\) be the distribution kernel of \(K_t\).  Then
\[
 \mathcal K_t
 =\partial_tW_t\,\one-i\Ham_0W_t,
 \qquad
 \widehat W_t(\xi)=\frac{\sin(t\kappa|M_0\xi|)}
                        {\kappa|M_0\xi|},
\]
and
\[
 \supp\mathcal K_t=
 \begin{cases}
 \{x:|M_0^{-1}x|\le\kappa|t|\},
      & \dCl=1\ \text{or }\dCl\text{ even},\\[2mm]
 \{x:|M_0^{-1}x|=\kappa|t|\},
      & \dCl\ge3\text{ odd}.
 \end{cases}
\]
In particular, every neighbourhood of a point in the displayed support carries
a test spinor on which \(\mathcal K_t\) is nonzero.
\end{lemma}

\begin{proof}
Since
\(\Ham_0^2=\kappa^2|M_0D|^2\one\), functional calculus gives
\[
 K_t=\cos(t\kappa|M_0D|)\one
     -i\frac{\sin(t\kappa|M_0D|)}
              {\kappa|M_0D|}\Ham_0
     =\partial_tW_t\,\one-i\Ham_0W_t.
\]
The inverse Fourier transform of \(W_t\) is the standard wave sine kernel after
the invertible linear change \(y=M_0^{-1}x\).  Its distributional support is the
closed ball for \(\dCl=1\) and for even \(\dCl\), and the sphere for odd
\(\dCl\ge3\).  Constant-coefficient differentiation cannot enlarge support.
It also does not remove an open component of the interior support in the even
case: taking the matrix trace eliminates the Clifford-odd term
\(\Ham_0W_t\), so
\[
 \operatorname{tr}\mathcal K_t
 =(\dim V)\,\partial_tW_t.
\]
The standard wave kernel has distributional support equal to the same closed
ball, and \(\partial_tW_t\) is nonzero on every nonempty open subset of its
interior.  In odd dimension both terms are supported on the characteristic
sphere, and the nonzero wave front amplitude gives the full sphere as support.
The last assertion is precisely the definition of distributional support.
\end{proof}

\begin{theorem}[Constant-coefficient inner characterization]
\label{thm:constant-inner-characterization}
Let \(U,V\subset\R^{\dCl}\) be nonempty open sets and \(t>0\).
If
\[
 (U-V)\cap\supp\mathcal K_t\ne\varnothing,
\]
then
\[
 \one_Ue^{-it\Ham_0}\one_V\ne0,
\]
and the corresponding propagation transfers to the renewal approximants by
\Cref{thm:macroscopic-domain}(ii).  Consequently:
\begin{enumerate}[label=\textup{(\roman*)},leftmargin=2.4em]
\item for \(\dCl=1\) and even \(\dCl\), every open pair satisfying
\[
 \inf_{u\in U,v\in V}|M_0^{-1}(u-v)|<\kappa t
\]
has nonzero forward transfer;
\item for odd \(\dCl\ge3\), forward transfer at the fixed time \(t\) occurs
whenever the difference set meets the Huygens shell
\(|M_0^{-1}(u-v)|=\kappa t\);
\item the convex hull of \(\supp\mathcal K_t\) is the full causal ellipsoid
\(\kappa tM_0B^{\dCl}\), so the characteristic cone is attained in every
direction.
\end{enumerate}
\end{theorem}

\begin{proof}
Choose \((u_*,v_*)\in U\times V\) with
\(u_*-v_*\in\supp\mathcal K_t\).  By the definition of distributional support,
there is a matrix-valued test function \(\phi\), supported in a sufficiently
small neighbourhood of \(u_*-v_*\), with
\(\langle\mathcal K_t,\phi\rangle\ne0\).  Shrinking the neighbourhood, write
\(\phi\) as a finite sum of convolutions
\(\sum_j f_j*\widetilde g_j\) with \(f_j\in C_c^\infty(U,V)\) and
\(g_j\in C_c^\infty(V,V)\), where
\(\widetilde g(x)=g(-x)^*\).  Then for some \(j\),
\[
 \langle f_j,e^{-it\Ham_0}g_j\rangle\ne0,
\]
which proves the operator nonvanishing.  The parity statements follow from
\Cref{lem:constant-dirac-kernel}.  The final convex-hull statement is immediate:
the closed ball is already convex, while the convex hull of the characteristic
sphere is the ball.
\end{proof}

\begin{remark}[Variable coefficients]
\label{rem:variable-inner-status}
For general \(M(x)\), Green hyperbolicity yields only
\(\supp E^+\subseteq J^+\); it does not imply equality of supports or local
kernel nonvanishing at every causally related pair.  Conjugate points,
polarization transport, and subprincipal terms may create zeros or tails, and
Huygens behaviour is exceptional.  Accordingly the general inner statement is
the exact transfer theorem \Cref{thm:macroscopic-domain}(ii), not universal
filling of \(J^+\).
\end{remark}

\subsubsection{Ballistic reachability versus dispersive causality}

For a constant finite reset frame, put
\[
 \Theta=\operatorname{conv}\{\hat\theta_a:p_a>0\}.
\]
If one allows free concatenation of reset labels and an \(O(h)\) endpoint
tolerance, the reachable set after \(N\) steps is
\(Nh\,\Theta+O(h)\).  Hence the ballistic travel time converges to the
Minkowski gauge
\[
 d_{\mathrm{bal}}(p,q)=\gamma_\Theta(q-p).
\]
This depends only on the support of the reset law, whereas the Dirac cone depends
on its second moment \(M\).

\begin{proposition}[Drift hull versus dispersion ellipsoid]
\label{prop:drift-dispersion}
For constant reset data and free concatenation, the approximate ballistic
travel times converge locally uniformly to the Finsler distance with indicatrix
\(\Theta\).  They agree with the Lorentzian optical distance exactly when
\[
 \Theta=\kappa MB^{\dCl}.
\]
For a refining equidistributed family,
\(\Theta\to B^{\dCl}\) and \(M\to \dCl^{-1}I\); agreement therefore fixes the
clock normalization to \(\kappa=\dCl\).  For varying direction fields the
analogous Finsler convergence requires, in addition, Hausdorff continuity of
\(x\mapsto\Theta(x)\), a uniform interior-ball condition, and a separate
compactness/recovery-sequence argument; it is not used elsewhere in this paper.
\end{proposition}

\begin{proof}
For constant data, the \(N\)-step displacement set is the Minkowski sum of
\(N\) copies of the finite step set.  Its convexification is \(Nh\,\Theta\), and
the \(O(h)\) endpoint tolerance removes the arithmetic lattice obstruction.
The minimal \(Nh\) therefore converges to the gauge
\(\gamma_\Theta(q-p)\).

The optical distance is the length distance whose infinitesimal unit ball is
\(\kappa MB^{\dCl}\).  Equality of the two length distances near the diagonal
forces equality of their metric derivatives and hence equality of the
indicatrices.  Conversely equal indicatrices give equal length functionals.
The isotropic normalization follows from
\(M=\int\theta\theta^\top d\nu=\dCl^{-1}I\).
\end{proof}

\begin{theorem}[Joint propagation-geometry convergence]
\label{thm:joint-propagation}
Under the curved-limit hypotheses, the frozen symbols, ellipsoidal cone fields,
optical distances, macroscopic outer domains of dependence, and self-adjoint
Hamiltonians converge simultaneously to the data of the single ultrastatic
metric
\[
 g_{\mu\nu}(x)=\diag\bigl(1,-\kappa^{-2}M(x)^{-2}\bigr).
\]
Inner propagation channels are inherited wherever the limiting propagator is
locally nonzero; for constant \(M_0\) this nonvanishing is characterized exactly
by \Cref{lem:constant-dirac-kernel,thm:constant-inner-characterization}.
Ballistic reset-ray reachability is a distinct drift-hull geometry and coincides
with this Lorentzian propagation geometry only under the indicatrix equality of
\Cref{prop:drift-dispersion}.
\end{theorem}

\begin{proof}
Combine \Cref{lem:optical-distance-stability,prop:joint-symbol-cone,prop:joint-finite-propagation,thm:curved-limit,thm:macroscopic-domain,thm:constant-inner-characterization,prop:drift-dispersion}.
\end{proof}


\subsection{Universality and metric reconstruction}
\label{sec:universality-reconstruction}

We finish the continuum section by recording what the limiting construction
forgets and what it remembers.  It forgets much of the microscopic renewal
presentation: after scaling, only the second-moment field of the reset
directions enters the principal symbol.  This is the universality statement.
But if the second-moment field is included as part of the scaling data, then the
signed spectral package reconstructs the limiting Lorentzian principal symbol
up to the unavoidable gauge and spatial-frame choices.  This is the
reconstruction statement.

\begin{definition}[Enhanced signed spectral package]
\label{def:enhanced-package}
An \emph{enhanced signed renewal spectral package} consists of the signed
predictive package
\[
  (\B_{\CP,0},\Hh,D_{\CP},J_\chi,D_-,\Arev)
\]
together with the reset-direction scaling datum
\[
  \Theta:x\longmapsto \nu_x,
  \qquad
  M_\Theta(x)=\int_{\Sphere^{\dCl-1}}\hat\theta\hat\theta^{\top}
  \,d\nu_x(\hat\theta),
\]
up to the evident flat gauge on the signed cover and orthogonal changes of the
spatial Clifford frame.  The adjective \emph{enhanced} is important: the bare
package reconstructs the ordered signed predictive monoid and the holonomy
class, while the continuum metric also uses the scaling statistics encoded by
\(\Theta\).
\end{definition}

\begin{lemma}[Continuum factorisation through the second moment]
\label{lem:continuum-factorisation}
The continuum Hamiltonian \(\Ham_g[M]\) of \Cref{def:curved-ham} depends on the
microscopic renewal data only through the normalisation \(\kappa\), the fixed
Krein--Clifford datum, and the limiting second-moment field \(M(x)\) together
with its first derivatives.  In particular, two reset fields with the same
\(M(x)\) and the same clock normalisation give the same continuum operator.
\end{lemma}

\begin{proof}
By \Cref{def:curved-ham},
\[
  \Ham_g[M]
  =\kappa A^iM^{ij}(x)(-i\partial_j)
        -\tfrac{i\kappa}{2}A^i(\partial_jM^{ij})(x).
\]
The direction measure \(\nu_x\) appears only through
\(M(x)=\int \hat\theta\hat\theta^{\top}\,d\nu_x\); no higher moment or
microscopic presentation of the reset alphabet enters the limiting operator.
\end{proof}

\begin{theorem}[Universality of Lorentzian Dirac continua]
\label{thm:universality}
Let two primitive signed renewal memories satisfy the adiabatic hypotheses of
\Cref{thm:curved-limit}.  Suppose that, after the standard continuum scaling,
they have the same clock normalisation \(\kappa\), the same fixed
Krein--Clifford datum, and the same uniformly elliptic limiting second-moment
field
\[
  M(x)=M'(x)\succeq m_0I.
\]
Then their rescaled discrete renewal Hamiltonians converge in strong resolvent
sense to the same variable-coefficient Lorentzian Dirac Hamiltonian
\(\Ham_g[M]\).  Consequently
\[
  e^{-it\Ham_R}\longrightarrow e^{-it\Ham_g[M]}
  \longleftarrow e^{-it\Ham'_R}
\]
strongly, uniformly for \(t\) in compact intervals.  Thus the Lorentzian
continuum endpoint depends only on the renormalised second-moment field, not on
the microscopic renewal presentation.
\end{theorem}

\begin{proof}
Apply \Cref{thm:curved-limit} to each family.  The first family converges to
\(\Ham_g[M]\), and the second to \(\Ham_g[M']\).  By
\Cref{lem:continuum-factorisation}, the two limits coincide when
\(M=M'\), with common \(\kappa\) and common Clifford datum.  The propagator
statement follows from strong resolvent convergence by the Trotter--Kato
theorem.
\end{proof}

\begin{theorem}[Inverse theorem for the Lorentzian principal symbol]
\label{thm:metric-reconstruction}
Let two enhanced signed renewal spectral packages have adiabatic scaling limits
with uniformly elliptic second-moment fields \(M\) and \(M'\).  Suppose a unitary
intertwines the diagonal predictive algebra, descended shifts, length operator,
fundamental symmetry, modular revision algebra, and reset-direction scaling
datum.  Then the unitary induces a graph isomorphism \(\varphi:G\to G'\) and a
spatial frame change \(O(x)\in O(\dCl)\) such that
\[
  M'(\varphi(x))=O(x)M(x)O(x)^{\top},
\]
and hence the limiting Lorentzian inverse metrics are isometric:
\[
  g'(\varphi(x))
  =\bigl(1\oplus O(x)\bigr)g(x)\bigl(1\oplus O(x)\bigr)^{\top}.
\]
Thus the Lorentzian principal symbol is an invariant of the enhanced signed
renewal spectral package, up to flat gauge and spatial-frame equivalence.
\end{theorem}

\begin{proof}
The signed reconstruction theorem \Cref{thm:reconstruction} gives an
isomorphism \(\varphi\) of graded ordered predictive monoids and identifies the
sign class.  Intertwining the reset-direction datum means that the direction
measure at \(x\) is carried to the direction measure at \(\varphi(x)\) up to the
orthogonal relabelling \(O(x)\) of the spatial Clifford frame.  Taking second
moments gives
\[
  M'(\varphi(x))
  =\int \hat\theta\hat\theta^{\top}\,d\nu'_{\varphi(x)}
  =O(x)\Bigl(\int \hat\theta\hat\theta^{\top}\,d\nu_x\Bigr)O(x)^{\top}.
\]
The metric formula follows from
\(g=\diag(1,-\kappa^2M^2)\) and the common clock normalisation.
\end{proof}

\begin{corollary}[Lorentzian principal-symbol reconstruction]
\label{cor:symbol-reconstruction}
The enhanced signed renewal spectral package of \Cref{def:enhanced-package}
determines the limiting Lorentzian inverse metric
\[
  g^{\mu\nu}(x)=\operatorname{diag}\bigl(1,\,-\kappa^2 M_\Theta(x)^2\bigr)
\]
as an invariant, up to signed-cover flat gauge and orthogonal spatial
Clifford-frame equivalence $O(x)\in O(\dCl)$.  That is, the assignment
\[
  \bigl(\B_{\CP,0},\Hh,D_{\CP},J_\chi,D_-,\Arev,\Theta\bigr)
  \ \longmapsto\ \bigl[g^{\mu\nu}(x)\bigr]
\]
is well defined on unitary-equivalence classes of enhanced packages with the
common clock normalisation $\kappa$ and Krein--Clifford datum, and is complete
for the principal symbol: two enhanced packages have isometric Lorentzian
principal symbols if and only if they are intertwined as in
\Cref{thm:metric-reconstruction}.
\end{corollary}

\begin{proof}
By \Cref{thm:metric-reconstruction} an intertwining unitary yields
$\varphi:G\to G'$ and $O(x)\in O(\dCl)$ with
$M'_\Theta(\varphi(x))=O(x)M_\Theta(x)O(x)^{\top}$.  Substituting into
$g^{\mu\nu}=\diag(1,-\kappa^2M_\Theta^2)$ and using
$O(x)M_\Theta^2O(x)^{\top}=(O M_\Theta O^{\top})^2$ gives
$g'^{\mu\nu}(\varphi(x))=(1\oplus O(x))\,g^{\mu\nu}(x)\,(1\oplus O(x))^{\top}$,
an isometry of inverse metrics fixing the time slot.  Conversely, equal
$M_\Theta$ up to $O(\dCl)$ frame and equal $\kappa$ give equal $g^{\mu\nu}$ up to
the same frame, and the package data determine $M_\Theta$ via its definition in
\Cref{def:enhanced-package}.  The timelike slot is frame-independent because it
is the image of the Krein-odd generator $\gamma^0$ (\Cref{thm:signature-krein}).
\end{proof}



\begin{remark}[Why the enhancement is necessary]
\label{rem:enhancement-necessary}
The unenhanced signed package reconstructs the discrete predictive monoid and
its \(\Z/2\)-holonomy class.  It does not by itself determine the continuum
second-moment field, since different microscopic direction measures can have the
same signed quotient but different scaling limits.  The reconstruction theorem
therefore uses the enhanced package of \Cref{def:enhanced-package}; with this
datum included, the Lorentzian principal symbol is recovered exactly up to the
unavoidable frame and gauge choices.
\end{remark}


\begin{remark}[Summary of the continuum passage]
\label{rem:continuum-summary-final}
The logical path of this section is therefore layered.  Product renewal order
gives a measured polyhedral continuum; signed Clifford isotropisation rounds the
principal symbol to a Lorentzian cone; the Hamiltonian formulation gives the
operator-level limit; slow spatial variation gives the curved adiabatic
extension; and the enhanced package records exactly the second-moment data
needed to reconstruct the limiting metric.  None of these steps adds field
equations or matter couplings.  They identify the Lorentzian Dirac continuum
which the renewal geometry supplies kinematically.
\end{remark}



% =====================================================================
\subsection{A marked flat-torus correspondence carrying spin structure}
\label{subsec:marked-torus-correspondence}

The local continuum results above identify the Lorentzian principal symbol but,
on a contractible chart, cannot retain the cohomology class of the signed cover.
A genuine two-sided statement therefore requires a global topology on which the
sign class survives.  We give such a statement on a marked flat torus.  The
result is deliberately band-limited: naive symmetric differences have
high-momentum lattice doublers, so no global norm-resolvent or global kernel
statement is asserted.

Let
\[
  \T^{\dCl}=\R^{\dCl}/\Z^{\dCl}
\]
carry its standard marking and flat parallelisation, and fix a reference spin
structure \(\mathfrak s_0\).  Relative to \(\mathfrak s_0\), the marked spin
structures are indexed by
\[
  \rho\in H^1(\T^{\dCl},\Ztwo)\cong(\Ztwo)^{\dCl},
  \qquad
  \mathfrak s_\rho=\mathfrak s_0\otimes L_\rho ,
\]
where the flat real line bundle \(L_\rho\) has holonomy
\((-1)^{\rho_k}\) around the \(k\)-th marked generator.  We realise its spinors
as equivariant functions on the universal cover,
\[
  L^2_\rho(\T^{\dCl},V)
  =
  \left\{\psi\in L^2_{\rm loc}(\R^{\dCl},V):
  \psi(x+e_k)=(-1)^{\rho_k}\psi(x)\right\}.
\]
Their physical Fourier momenta are
\[
  \xi\in2\pi\bigl(\Z+\rho/2\bigr)^{\dCl}.
\]

For constant \(M\succ0\), \(\tr M=1\), define
\[
  \Ham^\rho_g[M]
  =
  \kappa A^iM^{ij}(-i\partial_j)
  \quad\hbox{on }L^2_\rho(\T^{\dCl},V).
\]
The marked deck transformations
\[
  (\widetilde T_k\psi)(x)=\psi(x+e_k)
\]
act as the scalars
\[
  \widetilde T_k=(-1)^{\rho_k}\id .
\]
Hence the full marked class \(\rho\) is retained by the deck representation.
The bare spectrum always distinguishes the trivial class through its zero modes,
\[
  0\in\spec\Ham^\rho_g[M]\quad\Longleftrightarrow\quad\rho=0,
\]
but need not distinguish all nonzero classes; the complete class belongs to the
operator together with its marked deck holonomies.

\begin{lemma}[Exact integer-stencil decomposition]
\label{lem:integer-stencil-decomposition}
Every positive definite \(M\in\operatorname{Sym}_{\dCl}(\R)\) with \(\tr M=1\) admits a finite
decomposition
\[
  M=\sum_{a=1}^{m}p_a\,\hat\theta_a\hat\theta_a^\top,
  \qquad
  \hat\theta_a=\frac{v_a}{|v_a|},
  \quad
  v_a\in\Z^{\dCl}\setminus\{0\},
  \quad
  p_a>0,\quad \sum_ap_a=1 .
\]
\end{lemma}

\begin{proof}
The rank-one projections \(uu^\top\), \(u\in\Sphere^{\dCl-1}\), are the extreme
points of the compact convex set
\[
  \mathcal Q_{\dCl}
  =
  \{S\in\operatorname{Sym}_{\dCl}(\R):S\succeq0,\ \tr S=1\}.
\]
Because \(M\succ0\), it lies in the relative interior of \(\mathcal Q_{\dCl}\).
Choose finitely many rank-one projections whose convex hull contains \(M\) in
its relative interior.  The set of directions \(v/|v|\), with
\(v\in\Z^{\dCl}\setminus\{0\}\), is dense in the sphere, so the chosen vertices
may be perturbed to integer-direction projections.  Since \(M\) has positive
distance from every facet of the original finite polytope, sufficiently small
vertex perturbations preserve the strict containment of \(M\).  The resulting
positive barycentric coordinates give the stated decomposition.
\end{proof}

For \(h_R=N_R^{-1}\), define the integer-stencil symmetric difference
\[
  (\delta_{v_a,R}u)(x)
  =
  \frac{u(x+h_Rv_a)-u(x-h_Rv_a)}
       {2h_R|v_a|}.
\]
The twist is imposed intrinsically through the boundary condition
\[
  u(n+N_Re_k)=(-1)^{\rho_k}u(n),
  \qquad n\in\Z^{\dCl},
\]
equivalently by a flat cellular \(\Ztwo\)-cocycle on the cubical torus whose
plaquette holonomies vanish and whose marked winding holonomies are
\((-1)^{\rho_k}\).  This convention also fixes the sign of a long stencil shift
that crosses one or more fundamental boundaries.

With an exact decomposition from
\Cref{lem:integer-stencil-decomposition}, set
\[
  \Ham^\rho_R
  =
  \kappa\sum_a p_aA(\hat\theta_a)(-i\delta^\rho_{v_a,R}).
\]
Let
\[
  V_*=\max_a|v_a|,
  \qquad
  \mathcal B_\Lambda(\rho)
  =
  \left\{
  \xi\in2\pi(\Z+\rho/2)^{\dCl}:|\xi|\le\Lambda
  \right\},
\]
and let \(P_\Lambda\) be the Fourier projection onto this finite set.  The
smallest twisted momentum is
\[
  \min_{\xi\in2\pi(\Z+\rho/2)^{\dCl}}|\xi|
  =
  \pi\sqrt{|\rho|_0},
\]
where \(|\rho|_0\) is the number of nonzero components of \(\rho\).

\begin{theorem}[Band-limited twisted convergence]
\label{thm:marked-torus-band-limit}
Fix \(\Lambda>0\).  For all sufficiently large \(R\), the discrete and continuum
physical bands are canonically identified by twisted Fourier modes, and
\[
 \left\|
 P_\Lambda
 \bigl(J_R\Ham^\rho_RJ_R^*-\Ham^\rho_g[M]\bigr)
 P_\Lambda
 \right\|
 \le
 \frac{\kappa}{6}V_*^2h_R^2\Lambda^3
 \longrightarrow0.
\]
Consequently the resolvents converge in norm on the fixed band, and the
band eigenvalues converge with multiplicities.  If
\[
  \delta_\Lambda(\rho)
  :=
  \kappa
  \min_{\substack{\xi\in\mathcal B_\Lambda(\rho)\\\xi\ne0}}
  |M\xi|,
  \qquad \min\varnothing:=+\infty,
\]
then, whenever
\[
  \frac{\kappa}{6}V_*^2h_R^2\Lambda^3<\delta_\Lambda(\rho),
\]
the band contains no additional zero modes beyond the genuine continuum zero
mode.  In particular, for fixed
\(0<\varepsilon<\delta_\Lambda(\rho)\) and large \(R\),
\[
  \dim\mathbf1_{[-\varepsilon,\varepsilon]}
  \bigl(P_\Lambda\Ham^\rho_RP_\Lambda\bigr)
  =
  \dim V\;\mathbf1[\rho=0].
\]
\end{theorem}

\begin{proof}
On a physical mode \(\xi\), the multiplier of the \(a\)-th stencil is
\[
  s_{v_a,R}(\xi)
  =
  \frac{\sin(h_Rv_a\cdot\xi)}{h_R|v_a|}.
\]
Using \(|\sin t-t|\le |t|^3/6\),
\[
 \left|
 s_{v_a,R}(\xi)-\hat\theta_a\cdot\xi
 \right|
 \le
 \frac{h_R^2|v_a|^2|\xi|^3}{6}.
\]
Since \(\|A(\hat\theta_a)\|=1\) and \(\sum_ap_a=1\), taking the supremum over
\(\mathcal B_\Lambda(\rho)\) gives the displayed operator-norm estimate.  The
second resolvent identity gives norm-resolvent convergence on the finite band,
and finite-dimensional spectral stability gives convergence with
multiplicities.

The condition \(\pi N_R/V_*>\Lambda\) excludes the first nonzero sine zeros of
every individual stencil from the physical band.  This alone does not exclude
a cancellation in the total matrix symbol.  The latter is ruled out by the
uniform estimate: away from the genuine zero mode the continuum symbol has
spectral gap \(\delta_\Lambda(\rho)\), so any perturbation of norm strictly
smaller than that gap remains invertible.  Hence no additional band zero can
occur for sufficiently large \(R\).
\end{proof}

\begin{definition}[Strictly marked Dirac datum]
\label{def:strictly-marked-dirac-datum}
A strictly marked Dirac datum is
\[
  \mathfrak D=(\pi,\{\widetilde T_k\},c,\Ham),
\]
consisting of the multiplication representation, the marked deck holonomies,
the Clifford representation, and the Dirac Hamiltonian.  A strict equivalence
intertwines all four structures and fixes the marking.  A marking-preserving
geometric equivalence may additionally act by
\[
  O\in O(\dCl,\Z):=O(\dCl)\cap GL_{\dCl}(\Z),
\]
the signed-permutation isometry group of the standard marked torus.
\end{definition}

\begin{theorem}[Classification of marked flat-torus data]
\label{thm:marked-torus-classification}
For constant \(M,M'\succ0\), \(\tr M=\tr M'=1\), and
\(\rho,\rho'\in(\Ztwo)^{\dCl}\):
\begin{enumerate}
\item the strictly marked data of \(\Ham^\rho_g[M]\) and
      \(\Ham^{\rho'}_g[M']\) are equivalent if and only if
      \[
        M'=M,\qquad \rho'=\rho;
      \]
\item they are equivalent through a marking-preserving torus isometry
      \(O\in O(\dCl,\Z)\) if and only if
      \[
        M'=OMO^\top,\qquad \rho'\equiv O\rho\pmod2.
      \]
\end{enumerate}
No transformation law for the symmetric-gauge matrix \(M\) under a general
\(GL_{\dCl}(\Z)\) mapping class is asserted: the geometric co-metric transforms
tensorially, whereas its symmetric square root does not, and the fixed
normalisation \(\tr h^{1/2}=\kappa\) is not preserved by a general nonorthogonal
mapping class.
\end{theorem}

\begin{proof}
A strict intertwiner of the multiplication representation is a matrix-valued
multiplier \(U(x)\).  Intertwining the irreducible constant Clifford
representation forces \(U(x)=e^{i\alpha(x)}\one\).  Conjugation of the Dirac
operator then adds the multiplication term
\[
  \kappa A^iM^{ij}\partial_j\alpha.
\]
Equality with another pure constant-coefficient operator forces
\(d\alpha=0\), hence \(U\) is constant and equality of principal symbols gives
\(M'=M\).  Intertwining the deck scalars gives \(\rho'=\rho\).
For a marking-preserving isometry, the simultaneous coordinate and Clifford
rotation is restricted to \(O(\dCl,\Z)\), yielding the displayed transformations.
\end{proof}

\begin{theorem}[Marked torus realization and retention]
\label{thm:marked-torus-correspondence}
For every constant trace-normalised flat co-metric
\[
  h\succ0,\qquad \tr h^{1/2}=\kappa,
\]
and every marked spin structure \(\rho\in(\Ztwo)^{\dCl}\), there is an
integer-stencil signed renewal torus whose winding class is \(\rho\) and whose
Hamiltonians converge in norm resolvent on every fixed physical momentum band
to
\[
  \Ham^\rho_g[M],
  \qquad M=\kappa^{-1}h^{1/2}.
\]
The continuum strictly marked datum retains
\[
  h=\kappa^2M^2
\]
through its principal symbol and retains the full spin class \(\rho\) through
its deck holonomies.  The bare spectrum always detects the trivial class by its
harmonic spinors but does not determine the full nonzero class in general.
Thus
\[
  (M,[\chi])
  \longmapsto
  \bigl(h=\kappa^2M^2,\mathfrak s_0\otimes L_{[\chi]}\bigr)
\]
is a bijection of strictly marked enriched moduli, while the microscopic
integer stencil realising \(M\) remains non-unique.
\end{theorem}

\begin{proof}
The integer-stencil decomposition is
\Cref{lem:integer-stencil-decomposition}; the prescribed spin class is imposed by
the twisted periodicity, whose winding signs equal the deck holonomies; the
band-limited convergence is \Cref{thm:marked-torus-band-limit}; metric and
holonomy retention are the principal-symbol and deck-representation statements
above; and the strict classification is
\Cref{thm:marked-torus-classification}.  Non-uniqueness of the microscopic frame
is the continuum factorisation of \Cref{lem:continuum-factorisation}.
\end{proof}

\begin{remark}[Scope]
\label{rem:marked-torus-scope}
The theorem is global but flat.  A periodic variable field \(M(x)\) would combine
the local adiabatic limit with twisted boundary conditions, while a general
closed spin manifold would require a separate graph-Dirac-to-manifold-Dirac
convergence theorem.  Neither extension is used here.
\end{remark}


% =====================================================================
\section{Examples and sharp model cases}
\label{sec:examples}

This appendix collects elementary models realising the main mechanisms of the
paper.  The examples are organised by role.  The first group illustrates the
positive predictive geometry; the second isolates signed holonomy and its
separation from spectral deficiency; the third separates finite polyhedral order
cones from Lorentzian Clifford symbols; the last group records the primitive
modular patterns behind the \(3+1\) dimension statement.

\[
\begin{array}{p{0.25\linewidth}p{0.53\linewidth}p{0.13\linewidth}}
\toprule
\textbf{Example type} & \textbf{Purpose} & \textbf{Reference} \\
\midrule
Positive predictive examples &
Separate algebraic distinguishability, cb-metric dimension, and order-volume
growth. &
\Cref{subsec:examples-positive} \\[2pt]

Signed holonomy examples &
Show when non-removable Krein sectors exist, and why holonomy is not the same as
spectral-radius deficiency. &
\Cref{subsec:examples-signed} \\[2pt]

Continuum examples &
Separate finite polyhedral order cones from Lorentzian Clifford symbols, and
show how isotropy is only a special representative of the flat endpoint. &
\Cref{subsec:examples-continuum} \\[2pt]

Primitive modular examples &
Realise the \(3+1\) primitive shell and the higher-rank obstruction to
unconditional uniqueness. &
\Cref{subsec:examples-primitive} \\
\bottomrule
\end{array}
\]

% =====================================================================
\subsection{Positive predictive geometry}
\label{subsec:examples-positive}

\begin{example}[Separation of algebraic and metric predictive dimension]
\label{ex:alg-metric-separation}
Let the reset channels act, in a finite affine chart, by similarities of
contraction ratios \(r_1,\ldots,r_m\) satisfying the open set condition.  If the
products remain exactly distinguishable as words, then the algebraic predictive
quotient has exponential word growth, so the additive degree triple is not
finitely summable.  Nevertheless the cb-metric closure is a self-similar
attractor, and the metric predictive dimension is the unique solution \(s\) of
\[
  \sum_{i=1}^{m} r_i^s=1.
\]
Thus \(q_{\mathrm{alg}}\) and \(q_{\mathrm{met}}\) measure genuinely different
features: exact predictive distinguishability and metric resolution of the
channel attractor.
\end{example}

\begin{example}[Free commuting \(\N^{c}\): dimension and order volume]
\label{ex:free-nc}
For \(\M=\N^{c}\), the channels are indexed by exponent vectors.  The number of
classes of rank at most \(R\) satisfies
\[
  N_{\CP}(R)\asymp R^{c},
\]
so \(\qalg=c\) by \Cref{thm:commuting}.  The predictive order is the product
order, and balanced intervals have volume growth
\[
  |[x,y]|\sim \tau^{c}
\]
by \Cref{prop:commuting-order}.  This is the basic manifold-like regime:
algebraic dimension, order-growth dimension, and interval-volume exponent all
equal \(c\).
\end{example}

% =====================================================================
\subsection{Signed covers and holonomy}
\label{subsec:examples-signed}

\begin{example}[Tree graph: no non-removable signed sector]
\label{ex:tree}
If the recurrent predictive graph \(G\) is a tree up to orientation, then
\[
  b_1(G)=0,\qquad H^1(G,\Z/2)=0 .
\]
By \Cref{thm:classification,cor:sector-count}, there is one gauge class of
compatible signed covers and no non-removable Krein sector.  A tree therefore
carries no intrinsic signed holonomy datum; a cycle is needed.
\end{example}

\begin{example}[One-cycle graph: one non-removable Krein class]
\label{ex:one-cycle}
If \(b_1(G)=1\), then
\[
  H^1(G,\Z/2)\cong \Z/2 .
\]
Hence there are \(2\) signed gauge classes and exactly one non-removable class.
The nontrivial class is the loop-holonomy cocycle.  This is the minimal renewal
carrier of a signed/Krein sector.
\end{example}

\begin{example}[Half-period sign: nontrivial holonomy without spectral deficiency]
\label{ex:half-period}
On the two-vertex cycle, take a sign cocycle \(\chi\) with holonomy \(-1\) around
the loop.  The unsigned adjacency matrix
\[
  A=\begin{pmatrix}0&1\\1&0\end{pmatrix}
\]
has spectral radius \(\rho(A)=1\), while the signed transfer operator
\[
  A_\chi=\begin{pmatrix}0&1\\-1&0\end{pmatrix}
\]
has eigenvalues \(\pm i\), hence \(\rho(A_\chi)=1\).  Thus
\[
  [\chi]\neq0\in H^1(G,\Z/2),
  \qquad
  r_\chi:=\frac{\rho(A_\chi)}{\rho(A)}=1 .
\]
The class is topologically nontrivial but spectrally silent.  This illustrates
why the signed sector is classified by holonomy, not by spectral-radius
deficiency.
\end{example}

% =====================================================================
\begin{example}[The smallest nontrivial amplitude lift]
\label{ex:z4-smallest-lift}
Let \(G\) be a single cycle.  Then
\[
  H^1(G,\Ztwo)\cong\Ztwo,
  \qquad
  H^1(G,\Zfour)\cong\Zfour .
\]
The nontrivial signed cover sends the generator to \(-1\).  Its amplitude lifts
send the generator to \(i\) or \(-i\), and both square to \(-1\).  Thus the two
amplitude lifts have the same real/Krein shadow but distinct fourth-root phase
data.  This is the local model for the general distinction between the signed
cover and its amplitude completion.
\end{example}

\subsection{Polyhedral order cones and Lorentzian Clifford symbols}
\label{subsec:examples-continuum}

\begin{example}[Finite \(c\ge3\) product cone: the polyhedral obstruction]
\label{ex:finite-cone}
For \(\M=\N^{c}\) with \(c\ge3\), the product-order cone
\[
  [0,\infty)^c
\]
has exactly \(c\) extreme rays, the coordinate axes.  A round Lorentz cone
\[
  L_c=\{x_0\ge |(x_1,\dots,x_{c-1})|\}
\]
has the continuum of null generators over \(S^{c-2}\).  Linear isomorphisms
preserve extreme rays.  Hence no deterministic finite \(c\)-generator product
order can be linearly equivalent to a smooth Lorentz cone in dimension
\(c\ge3\).  This is the obstruction of \Cref{thm:obstruction}: finite
deterministic order cones remain polyhedral, even when a separate signed
Clifford symbol has a Lorentzian continuum limit.
\end{example}

\begin{example}[Finite tight frames round the symbol in a chosen frame]
\label{ex:finite-frame-symbol}
A finite isotropic tight frame can round the Dirac \emph{symbol} exactly without
rounding the order cone.  In spatial rank \(\dCl=3\), take the six octahedral
directions
\[
  \{\pm e_1,\pm e_2,\pm e_3\}\subset S^2
\]
with equal weights.  Then
\[
  M(\nu)=\frac13 I_3 ,
\]
and with \(\kappa=3\) the signed Clifford symbol has the exact Minkowski
characteristic cone
\[
  \xi_0^2=\xi_1^2+\xi_2^2+\xi_3^2 .
\]
However, the order cone generated by the six null rays \((1,\pm e_i)\) has
octahedral cross-section, hence remains polyhedral.  Thus finite isotropic
frames round the spectral cone in the chosen renewal frame, while order-cone
rounding requires an equidistributing limit.  This is a special isotropic
representative of the general nondegenerate symbol theorem: non-isotropic
positive-definite second moments give linearly isometric flat endpoints after
choosing the corresponding vielbein.
\end{example}

\begin{example}[Non-isotropic second moment and flat isometry]
\label{ex:anisotropic-frame-equivalence}
Let the signed reset directions in spatial rank \(3\) have positive-definite
second moment
\[
  M_*=\operatorname{diag}(a_1,a_2,a_3),
  \qquad a_i>0,
\]
with not all \(a_i\) equal.  In the original renewal frame the characteristic
quadric is anisotropic:
\[
  \xi_0^2=\kappa^2\bigl(a_1^2\xi_1^2+a_2^2\xi_2^2+a_3^2\xi_3^2\bigr).
\]
Equivalently, with the fixed spatial vielbein
\[
  V=\kappa M_*,
\]
the spatial covector \(\eta=V\xi\) brings the symbol to the standard flat form
\[
  \xi_0^2=|\eta|^2 .
\]
Thus isotropy of the reset law is not required for the flat Lorentzian Dirac
isometry class.  Isotropy merely makes the standard Euclidean frame coincide
with the original renewal frame.
\end{example}

\begin{example}[Equidistributed directions: order-cone rounding]
\label{ex:equidistributed}
Let \(\nu_R\to\nu_{\mathrm{unif}}\) weak-\(*\) on
\(\Sphere^{\dCl-1}\), with directions becoming dense.  Then
\[
  M(\nu_R)\longrightarrow \frac1{\dCl}I_{\dCl},
\]
and with \(\kappa=\dCl\), the spectral cones converge to the round Lorentz cone
in the original renewal frame.  Moreover, the conical hulls of the null rays
\((1,\hat\theta_a^{(R)})\) converge in Hausdorff distance on rays to the solid
future cone
\[
  \{x_0\ge|\vec x|\}.
\]
This realises order-cone rounding.  It is stronger than what is needed for the
flat Dirac endpoint, where nondegenerate second moments already determine a
Lorentzian symbol up to linear isometry.
\end{example}

% =====================================================================
\subsection{Primitive modular dimension patterns}
\label{subsec:examples-primitive}

\begin{example}[Primitive \(3+1\) modular revision pattern]
\label{ex:primitive-d3}
Let
\[
  \Hf=H^1(G,\Z/2)\oplus\langle t\rangle\cong(\Z/2)^4
\]
with basis \((t,e_1,e_2,e_3)\).  Define the modular commutator form by
\[
  \Omega=J_4-I_4
  =
  \begin{pmatrix}
  0&1&1&1\\
  1&0&1&1\\
  1&1&0&1\\
  1&1&1&0
  \end{pmatrix}
  \quad\text{over }\F_2 .
\]
Since \(4\) is even, this matrix is nonsingular over \(\F_2\).  Hence
\[
  \rank_{\F_2}\Omega=4,
\]
and the modular revision algebra is a full matrix factor,
equivalently the even complex Clifford algebra
\[
  \Arev\cong \Cl_4(\C)\cong M_4(\C).
\]
The spatial \(3\times3\) block has a radical generated by
\(e_1+e_2+e_3\), and the temporal generator \(t\) pairs with this class to make
the full \(4\times4\) form nondegenerate.  This is the minimal primitive pattern
realising the \(3+1\) endpoint.
\end{example}

\begin{example}[Primitive \(\dCl=5\) shell: no unconditional uniqueness]
\label{ex:primitive-d5}
Let
\[
  \Hf\cong(\Z/2)^6
\]
and take the same off-diagonal form
\[
  \Omega=J_6-I_6 .
\]
Again \(6\) is even, so \(\Omega\) is nonsingular over \(\F_2\).  Thus
\[
  \rank_{\F_2}\Omega=6,
\]
and the corresponding modular revision algebra is the full matrix factor
\[
  \Arev\cong \Cl_6(\C)\cong M_8(\C).
\]
This gives a primitive revision datum with spatial Clifford rank \(\dCl=5\).
Therefore primitivity alone does not exclude the higher odd ranks
\(\dCl=5,7,\ldots\).  The selection of \(\dCl=3\) is consequently a minimality
statement unconditionally, and a uniqueness statement only after imposing the
stated oriented-circulation or volume-response closure; access and efficiency
principles are optional alternative comparisons.
\end{example}

\begin{remark}[Why these two primitive examples are useful]
\label{rem:primitive-examples-use}
The \(3+1\) example realises the endpoint selected by the minimality theorem.
The \(\dCl=5\) example prevents overclaiming: it shows concretely why primitive
modular nondegeneracy cannot by itself prove uniqueness of \(3+1\).  Together
they make the status of the dimension theorem transparent.
\end{remark}


\begin{table}[t]
\centering
\small
\renewcommand{\arraystretch}{1.18}
\begin{tabularx}{\textwidth}{%
  >{\raggedright\arraybackslash}p{0.22\textwidth}
  >{\raggedright\arraybackslash}X
  >{\raggedright\arraybackslash}p{0.22\textwidth}}
\toprule
\textbf{Layer} & \textbf{Result in this paper} & \textbf{Status} \\
\midrule

Positive predictive geometry &
The predictive quotient carries a positive spectral triple, zeta dimension,
residue volume, metric predictive dimension, predictive order, and renewal
resolvent. &
Proved under finite-fibre, bounded-increment, and regular-growth hypotheses. \\[3pt]

Positivity boundary &
The completely positive sector is intrinsically Hilbert-positive and cannot
functorially supply a nontrivial Lorentzian fundamental symmetry. &
Proved as a structural obstruction. \\[3pt]

Signed/Krein enrichment &
The minimal renewal-positive indefinite enrichment over a connected recurrent
component is a principal \(\mathbb Z/2\)-cover; its gauge classes lie in
\(H^1(G,\mathbb Z/2)\), and the sheet label gives the Krein fundamental
symmetry. &
Proved in the renewal-positive edge-local class; label-local variants may
restrict to \(H^1_{\mathrm{loc}}\). \\[3pt]

Modular signed Dirac &
The signed cover and renewal clock determine the modular signed operator
\(D_-=J_\chi e^{\beta N}+B\), giving a twisted Krein spectral triple under the
bounded twisted-commutator hypotheses. &
Proved; nontrivial differential content is carried by the bounded renewal
correction \(B\). \\[3pt]

Flat continuum and self-averaging &
Nondegenerate Clifford coarse-graining gives a Lorentzian characteristic
quadric.  Under stationary ergodic reset statistics with \(M_*\succ0\), the
flat endpoint is deterministic and unitarily equivalent to the standard
Minkowski massless Dirac Hamiltonian. &
Proved at symbol level and for the flat Hamiltonian; no rotational-isotropy
hypothesis is required. \\[3pt]

Adiabatic and curved endpoint &
In the slowly varying regime, channel plaquette composition separates
translational torsion from rotational curvature; second-order predictive
equivalence gives the discrete Levi--Civita transport, and the covariant
Hamiltonians converge through subprincipal order, up to a separated bounded real
scalar renewal potential. &
Proved under uniform nondegeneracy and regularity assumptions. \\[3pt]

Propagation and marked tori &
The same second-moment field controls frozen symbols, optical distance, and
macroscopic domains of dependence.  Flat marked tori with marked spin structure
admit integer-stencil signed-renewal realisations with fixed-band norm-resolvent
convergence. &
Proved in the stated flat/adiabatic regimes; curved periodic tori and general
spin manifolds remain open. \\[3pt]

Dimension and primitive interface &
Full-dimensional affine reachability with finite primitive return completion
gives \(q_{\mathrm{alg}}=1+b_1^{\mathrm{eff}}\).  Primitive stabilised
\(\mathbb F_2\)-commutator data generate the Clifford revision factor and make
\(3+1\) the minimal spatially nondegenerate Lorentzian Dirac endpoint. &
Minimality proved in the primitive class; uniqueness among higher odd endpoints
requires the stated oriented-circulation or volume-response closure. \\

\bottomrule
\end{tabularx}
\caption{Summary of the theorem chain and logical status of each layer.}
\label{tab:theorem-chain}
\end{table}


\bibliographystyle{unsrt} 
\bibliography{biblio}



\end{document}